% One-file collected edition. Compile with pdfLaTeX twice (three times for a new build).
% All manuscript text, numerical certificates and supplemental data are inline.
\documentclass[11pt,a4paper]{article}
\usepackage[T1]{fontenc}
\usepackage{lmodern,amsmath,amssymb,amsthm,microtype,booktabs,fancyhdr,fvextra,xcolor,enumitem,multicol}
\usepackage[a4paper,margin=25mm,headheight=15pt]{geometry}
\usepackage{hyperref}
\usepackage{bookmark}
\definecolor{volblue}{RGB}{23,66,84}
\definecolor{volgray}{RGB}{61,72,80}
\definecolor{ink}{RGB}{23,66,84}\colorlet{teal}{volblue}\colorlet{inklink}{volblue}
\hypersetup{colorlinks=true,linkcolor=volblue,citecolor=volblue,urlcolor=volblue,
  unicode=true,bookmarksnumbered=false,pdfstartview=FitH,
  pdftitle={Collected research papers on the Randic index, graph distances, and rational Lambert series},pdfauthor={Keio and Shaitetsu}}
\DeclareUnicodeCharacter{27E8}{\ensuremath{\langle}}\DeclareUnicodeCharacter{27E9}{\ensuremath{\rangle}}
\DeclareUnicodeCharacter{00B2}{\ensuremath{{}^2}}
\newtheorem{theorem}{Theorem}[section]
\newtheorem{lemma}[theorem]{Lemma}\newtheorem{corollary}[theorem]{Corollary}
\newtheorem{proposition}[theorem]{Proposition}
\newcommand{\rad}{\operatorname{rad}}\newcommand{\diam}{\operatorname{diam}}
\newcommand{\dist}{\operatorname{dist}}\newcommand{\ecc}{\operatorname{ecc}}
\newcommand{\dd}{\overline d}\newcommand{\pos}[1]{(#1)_+}
\pagestyle{fancy}\fancyhf{}\newcommand{\runningname}{Contents}
\fancyhead[L]{\small\itshape\runningname}\fancyhead[R]{\small Collected research papers}\fancyfoot[C]{\thepage}
\renewcommand{\headrulewidth}{0.25pt}\setlength{\emergencystretch}{2em}\allowdisplaybreaks[1]
\frenchspacing\raggedbottom
\fvset{breaklines=true,breakanywhere=true}
\DefineVerbatimEnvironment{verbatim}{Verbatim}{fontsize=\small,breaklines=true,breakanywhere=true}
\newcounter{volumeunit}\newif\ifvolumeappendix
\newcommand{\volappendix}{\setcounter{section}{0}\setcounter{subsection}{0}\volumeappendixtrue\renewcommand{\thesection}{\Alph{section}}}
\renewcommand{\theHsection}{u\arabic{volumeunit}.\ifvolumeappendix a\else s\fi.\arabic{section}}
\renewcommand{\theHsubsection}{\theHsection.\arabic{subsection}}
\renewcommand{\theHsubsubsection}{\theHsubsection.\arabic{subsubsection}}
\renewcommand{\theHequation}{u\arabic{volumeunit}.e\arabic{equation}}
\renewcommand{\theHtheorem}{\theHsection.t\arabic{theorem}}
\renewcommand{\theHtable}{u\arabic{volumeunit}.tbl\arabic{table}}
\renewcommand{\theHfigure}{u\arabic{volumeunit}.fig\arabic{figure}}
\newcommand{\unitstart}[5]{\clearpage\stepcounter{volumeunit}\volumeappendixfalse
\setcounter{section}{0}\setcounter{subsection}{0}\setcounter{theorem}{0}\setcounter{equation}{0}
\setcounter{figure}{0}\setcounter{table}{0}\setcounter{footnote}{0}
\renewcommand{\thesection}{\arabic{section}}\renewcommand{\runningname}{#3}
\bookmarksetup{startatroot}\pdfbookmark[0]{#1. #4}{vol-unit-#2}\phantomsection\label{vol-#2}\thispagestyle{plain}
{\small\sffamily\bfseries\color{volblue}#1}\par\vspace{8pt}
{\raggedright\LARGE\bfseries #4\par}#5\vspace{10pt}}
\newcommand{\subtitlevolume}[1]{{\vspace{3pt}\noindent\large\itshape #1\par}}
\newcommand{\tocrow}[3]{\par\noindent\begin{tabular*}{\linewidth}{@{}p{10mm}@{\extracolsep{\fill}}p{.80\linewidth}r@{}}{\small\sffamily #1}&\hyperref[vol-#2]{#3}&\pageref{vol-#2}\end{tabular*}\par\vspace{3pt}}
\newcommand{\tocgroup}[1]{\vspace{9pt}\noindent{\small\bfseries\sffamily\color{volblue}#1}\par\vspace{3pt}}
\begin{document}
\pdfbookmark[0]{Title and results}{vol-title}\thispagestyle{empty}
\vspace*{19mm}\noindent{\large\sffamily\color{volblue}Keio and Shaitetsu}\par\vspace{13mm}
{\noindent\fontsize{33}{39}\selectfont\bfseries Collected research papers\par}
\vspace{5mm}
{\noindent\LARGE The Randi\'c index, graph distances,\\[3pt]and rational Lambert series\par}
\vspace{10mm}\noindent{\color{volblue}\rule{0.25\linewidth}{1.1pt}}\par\vspace{9mm}
\noindent This edition brings together fourteen manuscript texts, the original ten-problem research notes, and the supplementary proofs, programs and data distributed with them. Two successive versions of the radius paper appear in their original order. Each paper keeps its own section, theorem and reference numbering.\par
\vspace{6mm}
\noindent\textbf{Radius.} The Caporossi--Hansen paper characterizes the connected graphs with $R(G)<\rad(G)$ as the even-order paths on at least four vertices. Paper 05 gives the sharp fixed-maximum-degree bound and characterizes its extremizers.\par\vspace{4mm}
\noindent\textbf{Mean distance.} The last paper proves $R(G)\ge\operatorname{ad}(G)$ for whole-graph diameter at most 53, and for arbitrary assemblies of cyclic blocks satisfying its diameter or coloring condition. An unrestricted proof for all connected graphs remains open in these manuscripts.\par\vspace{4mm}
\noindent\textbf{Lambert series.} Paper 04 proves irrationality in its stated region of rational bases. The appendix records the fixed upstream revision, the verification programs and the 860 rational intervals used in its certificate.\par
\vfill\noindent{\small\sffamily Collected edition\hfill November 2026}\clearpage
\pdfbookmark[0]{Contents}{vol-contents}\thispagestyle{plain}
{\noindent\huge\bfseries Contents\par}\vspace{3mm}

\tocgroup{Radius, rational Lambert series, and degree indices}
\tocrow{01}{p01}{Degree cuts and the Randi\'c--radius inequality}
\tocrow{02}{p02}{The Randi\'c index and the vertex radius of a graph}
\tocrow{03}{p03}{A proof of the Caporossi--Hansen conjecture on the Randi\'c index and graph radius}
\tocrow{04}{p04}{A larger irrationality region for rational Lambert bases}
\tocrow{05}{p05}{Degree indices and distance in connected graphs}
\tocgroup{Fajtlowicz's average-distance conjecture}
\tocrow{06}{p06}{Spanning trees and Fajtlowicz's average-distance conjecture}
\tocrow{07}{p07}{Threshold forests and Fajtlowicz's average-distance conjecture}
\tocrow{08}{p08}{Leaf attachment and bridge gluing in Fajtlowicz's conjecture}
\tocrow{09}{p09}{Subdividing a bridge and joining weighted graph cores}
\tocrow{10}{p10}{Vertex deletion and leaf covers for the Randi\'c index and average distance}
\tocrow{11}{p11}{A weighted cut-vertex reduction for the Randi\'c index and average distance}
\tocrow{12}{p12}{Clipped distances and the Randi\'c average-distance conjecture}
\tocrow{13}{p13}{An iterated distance cutoff for the Randi\'c index}
\tocrow{14}{p14}{The Randi\'c index, large degrees, and mean distance}
\tocgroup{Research notes and supplements}
\tocrow{A}{notes}{Ten problems from the Verifiable Mathematical Reasoning Corpus}
\tocrow{B}{supp}{Lambert determinant estimates, programs, and interval data}

\unitstart{Paper 01}{p01}{Degree cuts}{Degree cuts and the Randi\'c--radius inequality}{}
\makeatletter\@removefromreset{theorem}{section}\makeatother\renewcommand{\thetheorem}{\arabic{theorem}}
% BEGIN SOURCE BODY Randic_Radius_Degree_Cuts
\begin{abstract}
For a finite connected simple graph, we prove that the sum of the inverse
maximum endpoint degrees over the edges is at least the absolute radius,
where the demand points are the original vertices. As a consequence the
classical Randi\'c index is at least the vertex radius minus one half.
The proof reduces a minimal counterexample to a two-connected graph
with attached leaves. For each threshold on full vertex degree, deleting
the leaves on the low side alone pays for the edges crossing the cut.
Summation over thresholds bounds the degree imbalance by the drop from
deleting all leaves. A spanning tree through a vertex of maximum degree
then gives the required metric comparison. The argument uses no graph
enumeration or numerical certificate.
\end{abstract}

\section{Statement and scope}
Write $d_v$ for the degree of $v$ in the graph under discussion. Define
\[
 M(G)=\sum_{uv\in E(G)}\frac1{\max(d_u,d_v)},\qquad
 R(G)=\sum_{uv\in E(G)}\frac1{\sqrt{d_ud_v}}.
\]
Give each edge length one and let $|G|$ denote the resulting metric graph.
Write $\rho(G)=\min_{x\in|G|}\max_{v\in V(G)}d(x,v)$ and let $\rad(G)$
be the same minimum restricted to vertices. Sums on the one-vertex graph
are empty. A center exists by compactness. Rounding it to a nearest vertex
gives $\rho(G)\ge\rad(G)-1/2$.

\begin{theorem}\label{p01-main}
Every finite connected simple graph satisfies
\[
 R(G)\ \ge\ M(G)\ \ge\ \rho(G)\ \ge\ \rad(G)-\tfrac12.
\]
\end{theorem}

This implies the 1988 conjecture $R(G)\ge\rad(G)-1$ of Fajtlowicz
\cite{p01-Faj,p01-Deng}. The separate conjecture of Caporossi and Hansen asks
for $R(G)\ge\rad(G)$ outside the even-order paths. The half-unit estimate
leaves cases of that conjecture outstanding. For the modified index $M$,
the half-unit loss from the vertex radius is attained on every even-order
path of order at least four. On such a path $M=\rho=(|V|-1)/2$.

The minimum-counterexample construction below follows the arguments of
VMRC [087] and [106]. The degree-flow identity used in Section~\ref{p01-cuts}
comes from VMRC [094] \cite{p01-VMRC}. The step used here is to charge each threshold
cut to the drop for the corresponding subset of leaves, then sum the
degree cuts and apply a spanning-tree bound. All ingredients used in the
proof are derived below. References to prior graph inequalities give
context and are unnecessary for the deductions.

\section{The structure of a minimal counterexample}
Suppose $M(G)<\rho(G)$. Choose such a graph minimizing its cyclomatic
number $|E|-|V|+1$, and then its number of edges. A singleton and a
single edge satisfy the assertion. A star with at least two leaves has
$M=\rho=1$. A path on at least three vertices has $M=\rho=(|V|-1)/2$.

Deleting a nonbridge edge cannot decrease $\rho$, since the remaining
center locations are a subset of the original ones and distances cannot
shorten. It reduces the cyclomatic number. Thus every nonbridge edge $uv$
satisfies $M(G-uv)>M(G)$; call such an edge \emph{supportive}.

The minimum defining $\rho$ is attained at a vertex or at an edge midpoint.
Indeed, at position $t\in[0,1]$ on $uv$, the distance to $w$ is the smaller
of $d(u,w)+t$ and $d(v,w)+1-t$. The endpoint distances are integers
differing by at most one. On each half-edge these distances are affine
with slopes $\pm1$ and integer intercepts. The maximum over $w$ can have
a change of slope only at a half-integer. In particular the radius and
the distances from a chosen vertex or midpoint center are half-integers.

\paragraph{Leaf deletion.}
If $z$ is a leaf with support $v$ of degree $d\ge2$, let $b$ count the
other neighbors of $v$ of degree less than $d$. Recomputing the edge
weights gives
\[
 M(G)-M(G-z)=\frac1d-\frac{b}{d(d-1)}\ge0.
\]
The radius stays the same or drops by at most $1/2$. To see the upper bound
on the drop, choose a vertex or midpoint center $y$ of $G-z$ with radius
$r$. If $d(y,v)\le r-1/2$, the old center covers $G$ within $r+1/2$.
Otherwise $d(y,v)=r$. Move $1/2$ toward $v$, which puts every vertex of
$G$ within $r+1/2$ of the new point. The case $r=0$ is a single edge.
Projection away from a deleted pendant edge shows that deleting a leaf
cannot increase the radius. Deleting leaves at several supports at once
lowers it by at most one, since each deleted leaf is one step from a
retained vertex.

If $d=2$ and $G$ is not a star, its other neighbor has degree at least
two. Thus the drop in $M$ is $1/2$, contrary to minimality. If a support
has two leaves and $G$ is not a star, deleting one preserves $\rho$:
a center of the smaller graph can be chosen outside the other pendant
edge. Indeed an interior point there can be projected to the support
without increasing eccentricity, since there is a vertex at distance
at least one from the support on the other side. Restoring the deleted
leaf gives it the same distance as its twin. This again contradicts
minimality. We have shown that every support has degree at least three
and carries at most one leaf.

\paragraph{Cut vertices.}
Suppose $v$ is a cut vertex, and a partition $A,B$ of $V(G)\setminus\{v\}$
has no edges across it, but each side contains an edge. Form $G_A$ from
$G[A\cup\{v\}]$ by adding at $v$ as many new leaves as $v$ has neighbors in
$B$. Form $G_B$ symmetrically. Each graph is connected, has no greater
cyclomatic number, and has fewer edges than $G$. Indeed the cyclomatic
numbers add up to that of $G$, and the internal edges on the omitted side
are lost. The degrees on retained edges are preserved and all the new
edges together have weight one. Thus
\[
 M(G)=M(G_A)+M(G_B)-1.
\]
Let $r_A,r_B$ denote their absolute radii. Choose centers $y_A,y_B$ in
the retained parts: each artificial pendant edge can be projected to
$v$ if necessary. Write $s_i=d(y_i,v)$. The artificial leaves imply
$s_i\le r_i-1$, and $r_i\ge1$. Along a geodesic through $v$ from $y_A$
to $y_B$, let $u$ denote signed distance from $v$, positive toward $B$.
Distances to vertices on the two sides are bounded by
$r_A+s_A+u$ and $r_B+s_B-u$, respectively. If these are equal at a
point of $[-s_A,s_B]$, the common value is at most $r_A+r_B-1$. Otherwise
the better endpoint gives a bound of $r_A$ or $r_B$. Thus
$\rho(G)\le r_A+r_B-1\le M(G_A)+M(G_B)-1=M(G)$, a contradiction.

Consequently deleting any cut vertex leaves at most one component containing
an edge. If $G$ is a tree of diameter at least four, deleting the middle
vertex of a path of length four gives two such components. A tree of
diameter at most three is a star or a double star; the bound on the number
of leaves per support leaves only paths or stars. If $G$ has a cycle,
take a maximal two-connected block $K$ of order at least three.
Every component outside $K$ attaches to a single vertex of $K$, by
maximality of the block. An edge wholly within such a component would
give two components with edges after deleting its attachment.
Consequently $G$ is $K$ with at most one leaf at each core vertex.

\section{Charging a degree cut to leaf deletion}\label{p01-cuts}
All degrees in this section are measured in $G$. Write $n=|V(K)|$ and
$D=\Delta(G)$. For $v\in V(K)$ let $p_v\in\{0,1\}$ indicate the
presence of a leaf, and set
\[
 \ell_v=\#\{w\sim v:d_w<d_v\},\quad h_v=d_v-\ell_v,\quad
 g_v=d_v(d_v-1),\quad a_v=\ell_v/g_v,\quad b_v=h_v/g_v .
\]
For a core edge $uv$ with $d_u<d_v$, direct edge deletion gives
\[
 M(G-uv)-M(G)=\frac{\ell_u}{g_u}
                 +\frac{\ell_v-1}{g_v}-\frac1{d_v}
               =a_u-b_v>0 .
\]
For equal degrees $d$ it gives $(\ell_u+\ell_v)/(d(d-1))-1/d>0$.
Every core vertex has a lower-degree neighbor. Otherwise it has no
leaf, and on an incident edge to a higher-degree neighbor the strict
inequality just proved is impossible. On an incident edge of equal
degree the other endpoint has at most $d-1$ lower neighbors, again
impossible. A full degree of two is therefore impossible: the vertex
would need a leaf, leaving only one neighbor in the two-connected
core. Thus $3\le d_v\le D$ and $\ell_v\ge1$.

Put $\delta=M(G)-M(K)$. Deleting leaves one by one gives $\delta\ge0$.
At least one leaf exists, since a core vertex of smallest full degree
needs a lower-degree neighbor in $G$. Hence $K$ is smaller, and
\[
 \delta<\rho(G)-\rho(K)\le1.
\]

For $3\le k<D$, let $S=\{v\in K:d_v\le k\}$, and let $N_k$ be the
number of core edges crossing from $S$ to its complement. Orient each
strictly unequal core edge from smaller degree to larger degree and
write $e_v$ for the number of same-degree core neighbors of $v$.
The identity $(d_v-\ell_v)a_v=\ell_v b_v$ gives
\begin{align*}
 \sum_{v\in S}p_v b_v
  ={}&\sum_{u\to v\text{ within } S}(a_u-b_v)
      +\sum_{u\to v\text{ across the cut}} a_u
      +\sum_{v\in S}e_v a_v .
\end{align*}
Indeed $\ell_v-p_v$ is the number of incoming core edges and
$h_v-e_v$ is the number of outgoing ones. No edge enters $S$ from
a smaller degree outside it.

Let $G_S$ result from deleting just the leaves attached in $S$.
A host of degree $d$ loses a leaf of weight $1/d$, and its
$\ell_v-1$ edges to smaller core neighbors each gain $1/g_v$.
These gains depend only on deletion at the upper endpoint, including
when the two degrees initially differ by one. An edge with equal
endpoint degrees gains $1/g_v$ exactly when both endpoints lose a leaf.
Consequently
\[
 M(G)-M(G_S)=\sum_{v\in S}p_v b_v-
 \sum_{\substack{uv\text{ equal in }S\\p_u=p_v=1}}\frac1{g_u}.
\]
In the flow identity, the first sum on the right is nonnegative. Each
equal host--host edge in the correction contributes $a_u+a_v\ge2/g_u$
to $\sum e_v a_v$. Thus the cut sum is at most $M(G)-M(G_S)\le\delta$,
the second inequality following by deleting the remaining leaves.
At the lower endpoint of a crossing edge, $a_u\ge1/[k(k-1)]$. We obtain
\begin{equation}\label{p01-cut}
 \boxed{N_k\le k(k-1)\delta\qquad(3\le k<D).}
\end{equation}

\section{Summation and the metric bound}
Summing $1/d_u+1/d_v$ over all edges counts all vertices of $G$ once.
Subtracting twice $M$ therefore gives the sum of $|1/d_u-1/d_v|$ over
the edges. Separating the pendant edges and telescoping the remaining
reciprocal-degree differences yields the exact identity
\[
 n-2M(G)=\sum_{k=3}^{D-1}\frac{N_k}{k(k+1)}
               -\sum_{v\in K}\frac{p_v}{d_v}.
\]
Deleting all leaves removes their weights and can only increase the
weights of retained edges, so $\sum_v p_v/d_v\ge\delta$.
For $D\ge4$, equation~\eqref{p01-cut} implies
\[
 n-2M(G)\le
 \left(\sum_{k=3}^{D-1}\frac{k-1}{k+1}-1\right)\delta
 \le (D-4)\delta\le D-4 .
\]
If $D=3$, all core vertices have full degree three and one lower
neighbor, necessarily a leaf. Every core edge would have deletion
change $2/6-1/3=0$, contradicting supportiveness.

Finally, choose a vertex of degree $D$. At least $D-1$ of its
neighbors lie in $K$. Its incident core edges form a star, which
extends to a spanning tree of $K$. A tree with a vertex of degree
at least $D-1$ has at least $D-1$ leaves, since each branch at that
vertex contains a leaf. A diametral path uses at most two leaves, so
at least $D-3$ core vertices lie outside it. The tree thus has diameter
at most $n-D+2$. Reattach the original leaves. The diameter of the
resulting spanning tree of $G$ is at most $n-D+4$. Its diametral
midpoint covers every vertex within half that distance, and distances
can only shorten in $G$. Hence
\[
 2\rho(G)\le n-D+4\le2M(G).
\]
This contradicts minimality. The inequalities on either side of
$M(G)\ge\rho(G)$ in Theorem~\ref{p01-main} follow from
$\sqrt{d_ud_v}\le\max(d_u,d_v)$ and from rounding the center.
\hfill$\square$

\begin{thebibliography}{9}
\bibitem{p01-Faj} S. Fajtlowicz, On conjectures of Graffiti,
\emph{Discrete Mathematics} \textbf{72} (1988), 113--118.
\bibitem{p01-Deng} H. Deng, Z. Tang and J. Zhang, On a conjecture of
Randi\'c index and graph radius, \emph{Filomat}.
\href{https://journal.pmf.ni.ac.rs/filomat/index.php/filomat/article/view/710}{Journal article}.
\bibitem{p01-CPS} M. Cygan, M. Pilipczuk and R. \v{S}krekovski,
On the inequality between radius and Randi\'c index for graphs,
\emph{MATCH Commun. Math. Comput. Chem.} \textbf{67} (2012), 451--466.
\bibitem{p01-VMRC} Keio, \emph{Verifiable Mathematical Reasoning Corpus},
research drafts [087], [094] and [106], commit
\href{https://github.com/QEEDDD/Verifiable-Mathematical-Reasoning-Corpus/tree/8a3f6895d0b9b9f125e43157e4a003a2a2b51704}{8a3f6895d0b9}.
\end{thebibliography}

% END SOURCE BODY Randic_Radius_Degree_Cuts
\unitstart{Paper 02}{p02}{Vertex radius / first version}{The Randi\'c index and the vertex radius of a graph}{\subtitlevolume{First version}}
\makeatletter\@removefromreset{theorem}{section}\makeatother\renewcommand{\thetheorem}{\arabic{theorem}}
% BEGIN SOURCE BODY Randic_Index_and_Vertex_Radius_Preliminary
\begin{abstract}
For a finite connected simple graph of maximum degree $D\ge2$, we obtain
$R(G)-\rad(G)\ge\sqrt D-3/2$. This proves the Caporossi--Hansen inequality
and determines exactly the exceptions and equality cases. First we show
that the sum of the reciprocal maximum endpoint degrees is at least the
absolute radius with demand at the original vertices. In a minimal
counterexample, deleting leaves on the low-degree side of each cut
bounds the number of crossing edges. Summing these cut bounds gives
the metric inequality. A pointwise identity for the squared differences
of reciprocal square roots of degrees then relates the Randi\'c index
to the vertex radius. The proof uses no graph enumeration.
\end{abstract}

\section{Statement and scope}
Write $d_v$ for the degree of $v$ in the graph under discussion. Define
\[
 M(G)=\sum_{uv\in E(G)}\frac1{\max(d_u,d_v)},\qquad
 R(G)=\sum_{uv\in E(G)}\frac1{\sqrt{d_ud_v}}.
\]
Give each edge length one and let $|G|$ denote the resulting metric graph.
Write $\rho(G)=\min_{x\in|G|}\max_{v\in V(G)}d(x,v)$ and let $\rad(G)$
be the same minimum restricted to vertices. Sums on the one-vertex graph
are empty. A center exists by compactness. Rounding it to a nearest vertex
gives $\rho(G)\ge\rad(G)-1/2$.

\begin{theorem}\label{p02-main}
Every finite connected simple graph satisfies
\[
 R(G)\ \ge\ M(G)\ \ge\ \rho(G)\ \ge\ \rad(G)-\tfrac12.
\]
\end{theorem}

Fajtlowicz conjectured $R(G)\ge\rad(G)-1$ in 1988 \cite{p02-Faj,p02-Deng}.
Caporossi and Hansen proposed $R(G)\ge\rad(G)$ for every connected
graph other than an even-order path \cite{p02-CH,p02-Deng,p02-CPS}.
Theorem~\ref{p02-main} proves the first inequality. Theorem~\ref{p02-ch}
deduces the second and its equality classification. For the modified index
$M$, the half-unit loss is attained on all even-order paths of order at
least four, which satisfy $M=\rho=(|V|-1)/2$.

The minimum-counterexample construction below follows the arguments of
VMRC [087] and [106]. The degree-flow identity used in Section~\ref{p02-cuts}
comes from VMRC [094] \cite{p02-VMRC}. The step used here is to charge each threshold
cut to the drop for the corresponding subset of leaves, then sum the
degree cuts and apply a spanning-tree bound. All ingredients used in the
proof are derived below. References to prior graph inequalities give
context and are unnecessary for the deductions.

\section{The structure of a minimal counterexample}
Suppose $M(G)<\rho(G)$. Choose such a graph minimizing its cyclomatic
number $|E|-|V|+1$, and then its number of edges. A singleton and a
single edge satisfy the assertion. A star with at least two leaves has
$M=\rho=1$. A path on at least three vertices has $M=\rho=(|V|-1)/2$.

Deleting a nonbridge edge cannot decrease $\rho$, since the remaining
center locations are a subset of the original ones and distances cannot
shorten. It reduces the cyclomatic number. Thus every nonbridge edge $uv$
satisfies $M(G-uv)>M(G)$; call such an edge \emph{supportive}.

The minimum defining $\rho$ is attained at a vertex or at an edge midpoint.
Indeed, at position $t\in[0,1]$ on $uv$, the distance to $w$ is the smaller
of $d(u,w)+t$ and $d(v,w)+1-t$. The endpoint distances are integers
differing by at most one. On each half-edge these distances are affine
with slopes $\pm1$ and integer intercepts. The maximum over $w$ can have
a change of slope only at a half-integer. In particular the radius and
the distances from a chosen vertex or midpoint center are half-integers.

\paragraph{Leaf deletion.}
If $z$ is a leaf with support $v$ of degree $d\ge2$, let $b$ count the
other neighbors of $v$ of degree less than $d$. Recomputing the edge
weights gives
\[
 M(G)-M(G-z)=\frac1d-\frac{b}{d(d-1)}\ge0.
\]
The radius stays the same or drops by at most $1/2$. To see the upper bound
on the drop, choose a vertex or midpoint center $y$ of $G-z$ with radius
$r$. If $d(y,v)\le r-1/2$, the old center covers $G$ within $r+1/2$.
Otherwise $d(y,v)=r$. Move $1/2$ toward $v$, which puts every vertex of
$G$ within $r+1/2$ of the new point. The case $r=0$ is a single edge.
Projection away from a deleted pendant edge shows that deleting a leaf
cannot increase the radius. Deleting leaves at several supports at once
lowers it by at most one, since each deleted leaf is one step from a
retained vertex.

If $d=2$ and $G$ is not a star, its other neighbor has degree at least
two. Thus the drop in $M$ is $1/2$, contrary to minimality. If a support
has two leaves and $G$ is not a star, deleting one preserves $\rho$:
a center of the smaller graph can be chosen outside the other pendant
edge. Indeed an interior point there can be projected to the support
without increasing eccentricity, since there is a vertex at distance
at least one from the support on the other side. Restoring the deleted
leaf gives it the same distance as its twin. This again contradicts
minimality. We have shown that every support has degree at least three
and carries at most one leaf.

\paragraph{Cut vertices.}
Suppose $v$ is a cut vertex, and a partition $A,B$ of $V(G)\setminus\{v\}$
has no edges across it, but each side contains an edge. Form $G_A$ from
$G[A\cup\{v\}]$ by adding at $v$ as many new leaves as $v$ has neighbors in
$B$. Form $G_B$ symmetrically. Each graph is connected, has no greater
cyclomatic number, and has fewer edges than $G$. Indeed the cyclomatic
numbers add up to that of $G$, and the internal edges on the omitted side
are lost. The degrees on retained edges are preserved and all the new
edges together have weight one. Thus
\[
 M(G)=M(G_A)+M(G_B)-1.
\]
Let $r_A,r_B$ denote their absolute radii. Choose centers $y_A,y_B$ in
the retained parts: each artificial pendant edge can be projected to
$v$ if necessary. Write $s_i=d(y_i,v)$. The artificial leaves imply
$s_i\le r_i-1$, and $r_i\ge1$. Along a geodesic through $v$ from $y_A$
to $y_B$, let $u$ denote signed distance from $v$, positive toward $B$.
Distances to vertices on the two sides are bounded by
$r_A+s_A+u$ and $r_B+s_B-u$, respectively. If these are equal at a
point of $[-s_A,s_B]$, the common value is at most $r_A+r_B-1$. Otherwise
the better endpoint gives a bound of $r_A$ or $r_B$. Thus
$\rho(G)\le r_A+r_B-1\le M(G_A)+M(G_B)-1=M(G)$, a contradiction.

Consequently deleting any cut vertex leaves at most one component containing
an edge. If $G$ is a tree of diameter at least four, deleting the middle
vertex of a path of length four gives two such components. A tree of
diameter at most three is a star or a double star; the bound on the number
of leaves per support leaves only paths or stars. If $G$ has a cycle,
take a maximal two-connected block $K$ of order at least three.
Every component outside $K$ attaches to a single vertex of $K$, by
maximality of the block. An edge wholly within such a component would
give two components with edges after deleting its attachment.
Consequently $G$ is $K$ with at most one leaf at each core vertex.

\section{Charging a degree cut to leaf deletion}\label{p02-cuts}
All degrees in this section are measured in $G$. Write $n=|V(K)|$ and
$D=\Delta(G)$. For $v\in V(K)$ let $p_v\in\{0,1\}$ indicate the
presence of a leaf, and set
\[
 \ell_v=\#\{w\sim v:d_w<d_v\},\quad h_v=d_v-\ell_v,\quad
 g_v=d_v(d_v-1),\quad a_v=\ell_v/g_v,\quad b_v=h_v/g_v .
\]
For a core edge $uv$ with $d_u<d_v$, direct edge deletion gives
\[
 M(G-uv)-M(G)=\frac{\ell_u}{g_u}
                 +\frac{\ell_v-1}{g_v}-\frac1{d_v}
               =a_u-b_v>0 .
\]
For equal degrees $d$ it gives $(\ell_u+\ell_v)/(d(d-1))-1/d>0$.
Every core vertex has a lower-degree neighbor. Otherwise it has no
leaf, and on an incident edge to a higher-degree neighbor the strict
inequality just proved is impossible. On an incident edge of equal
degree the other endpoint has at most $d-1$ lower neighbors, again
impossible. A full degree of two is therefore impossible: the vertex
would need a leaf, leaving only one neighbor in the two-connected
core. Thus $3\le d_v\le D$ and $\ell_v\ge1$.

Put $\delta=M(G)-M(K)$. Deleting leaves one by one gives $\delta\ge0$.
At least one leaf exists, since a core vertex of smallest full degree
needs a lower-degree neighbor in $G$. Hence $K$ is smaller, and
\[
 \delta<\rho(G)-\rho(K)\le1.
\]

For $3\le k<D$, let $S=\{v\in K:d_v\le k\}$, and let $N_k$ be the
number of core edges crossing from $S$ to its complement. Orient each
strictly unequal core edge from smaller degree to larger degree and
write $e_v$ for the number of same-degree core neighbors of $v$.
The identity $(d_v-\ell_v)a_v=\ell_v b_v$ gives
\begin{align*}
 \sum_{v\in S}p_v b_v
  ={}&\sum_{u\to v\text{ within } S}(a_u-b_v)
      +\sum_{u\to v\text{ across the cut}} a_u
      +\sum_{v\in S}e_v a_v .
\end{align*}
Indeed $\ell_v-p_v$ is the number of incoming core edges and
$h_v-e_v$ is the number of outgoing ones. No edge enters $S$ from
a smaller degree outside it.

Let $G_S$ result from deleting just the leaves attached in $S$.
A host of degree $d$ loses a leaf of weight $1/d$, and its
$\ell_v-1$ edges to smaller core neighbors each gain $1/g_v$.
These gains depend only on deletion at the upper endpoint, including
when the two degrees initially differ by one. An edge with equal
endpoint degrees gains $1/g_v$ exactly when both endpoints lose a leaf.
Consequently
\[
 M(G)-M(G_S)=\sum_{v\in S}p_v b_v-
 \sum_{\substack{uv\text{ equal in }S\\p_u=p_v=1}}\frac1{g_u}.
\]
In the flow identity, the first sum on the right is nonnegative. Each
equal host--host edge in the correction contributes $a_u+a_v\ge2/g_u$
to $\sum e_v a_v$. Thus the cut sum is at most $M(G)-M(G_S)\le\delta$,
the second inequality following by deleting the remaining leaves.
At the lower endpoint of a crossing edge, $a_u\ge1/[k(k-1)]$. We obtain
\begin{equation}\label{p02-cut}
 \boxed{N_k\le k(k-1)\delta\qquad(3\le k<D).}
\end{equation}

\section{Summation and the metric bound}
Summing $1/d_u+1/d_v$ over all edges counts all vertices of $G$ once.
Subtracting twice $M$ therefore gives the sum of $|1/d_u-1/d_v|$ over
the edges. Separating the pendant edges and telescoping the remaining
reciprocal-degree differences yields the exact identity
\[
 n-2M(G)=\sum_{k=3}^{D-1}\frac{N_k}{k(k+1)}
               -\sum_{v\in K}\frac{p_v}{d_v}.
\]
Deleting all leaves removes their weights and can only increase the
weights of retained edges, so $\sum_v p_v/d_v\ge\delta$.
For $D\ge4$, equation~\eqref{p02-cut} implies
\[
 n-2M(G)\le
 \left(\sum_{k=3}^{D-1}\frac{k-1}{k+1}-1\right)\delta
 \le (D-4)\delta\le D-4 .
\]
If $D=3$, all core vertices have full degree three and one lower
neighbor, necessarily a leaf. Every core edge would have deletion
change $2/6-1/3=0$, contradicting supportiveness.

Finally, choose a vertex of degree $D$. At least $D-1$ of its
neighbors lie in $K$. Its incident core edges form a star, which
extends to a spanning tree of $K$. A tree with a vertex of degree
at least $D-1$ has at least $D-1$ leaves, since each branch at that
vertex contains a leaf. A diametral path uses at most two leaves, so
at least $D-3$ core vertices lie outside it. The tree thus has diameter
at most $n-D+2$. Reattach the original leaves. The diameter of the
resulting spanning tree of $G$ is at most $n-D+4$. Its diametral
midpoint covers every vertex within half that distance, and distances
can only shorten in $G$. Hence
\[
 2\rho(G)\le n-D+4\le2M(G).
\]
This contradicts minimality. The inequalities on either side of
$M(G)\ge\rho(G)$ in Theorem~\ref{p02-main} follow from
$\sqrt{d_ud_v}\le\max(d_u,d_v)$ and from rounding the center.
\hfill$\square$

\section{The vertex radius and the maximum degree}
Write $N=|V(G)|$, $r=\rad(G)$ and $D=\Delta(G)$.
\begin{theorem}\label{p02-ch}
If $D\ge2$, then
\[
 R(G)-r\ \ge\
 \frac{N-2r-(\sqrt D-1)(r-\rho(G))}{\sqrt D+1}
 \ \ge\ \sqrt D-\tfrac32 .
\]
Consequently $R(G)>r$ whenever $D\ge3$. The only connected graphs
with $R(G)<r$ are the paths on an even number of vertices of order
at least four. Equality $R(G)=r$ holds precisely for the singleton,
the single edge and the even cycles.
\end{theorem}
\begin{proof}
For each edge order the endpoint degrees as $d\le e$. Its contribution
to $R-M$ is $1/\sqrt{de}-1/e$. The exact identity
\[
 (d^{-1/2}-e^{-1/2})^2
  =(\sqrt{e/d}-1)\left(\frac1{\sqrt{de}}-\frac1e\right)
\]
controls the sum of squared differences, because each vertex contributes
one to the sum of the reciprocal endpoint degrees. In particular
\[
 N-2R=\sum_{uv\in E(G)}(d_u^{-1/2}-d_v^{-1/2})^2
       \le (\sqrt D-1)(R-M).
\]
Using $M\ge\rho$ from Theorem~\ref{p02-main} and rearranging gives the first inequality.

The $D$ incident edges at a maximum-degree vertex form a star that extends
to a spanning tree $T$. Each of its $D$ branches contains a leaf of $T$.
At most two leaves lie on a diametral path. Hence
$\diam(T)\le N-D+1$, and the vertex radius satisfies
\[
 2r\le 2\rad(T)=2\lceil\diam(T)/2\rceil\le N-D+2 .
\]
It follows that $N-2r\ge D-2$. Since $0\le r-\rho\le1/2$, the stated
bound follows from
\[
 \frac{D-2-(\sqrt D-1)/2}{\sqrt D+1}=\sqrt D-\tfrac32 .
\]
For $D\ge3$ this is positive. The connected graphs of maximum degree
at most two are paths, cycles, and the singleton. For $n\ge3$,
\[
 R(P_n)=(n-3)/2+\sqrt2,\qquad R(C_n)=n/2,\qquad
 \rad(P_n)=\rad(C_n)=\lfloor n/2\rfloor .
\]
For $K_1$ and $K_2$, $R=r$ equals zero and one respectively.
These values give the classification.
\end{proof}

Combining the theorem with the path formula yields the sharp universal bound
$R(G)-\rad(G)\ge\sqrt2-3/2$, attained precisely on even-order paths
with at least four vertices. The coefficient of $\sqrt D$ in
Theorem~\ref{p02-ch} also has the correct asymptotic size. Subdividing one
edge of $K_{1,D}$ once, for $D\ge3$, gives a tree of radius two and
\[
 R-r=\sqrt D-2+\frac1{\sqrt2}
         -\frac{1-1/\sqrt2}{\sqrt D}.
\]

\begin{thebibliography}{9}
\bibitem{p02-Faj} S. Fajtlowicz, On conjectures of Graffiti,
\emph{Discrete Mathematics} \textbf{72} (1988), 113--118.
\bibitem{p02-Deng} H. Deng, Z. Tang and J. Zhang, On a conjecture of
Randi\'c index and graph radius, \emph{Filomat}.
\href{https://journal.pmf.ni.ac.rs/filomat/index.php/filomat/article/view/710}{Journal article}.
\bibitem{p02-CH} G. Caporossi and P. Hansen, Variable neighborhood search for
extremal graphs 1: the AutoGraphiX system, \emph{Discrete Mathematics}
\textbf{212} (2000), 29--44.
\bibitem{p02-CPS} M. Cygan, M. Pilipczuk and R. \v{S}krekovski,
On the inequality between radius and Randi\'c index for graphs,
\emph{MATCH Commun. Math. Comput. Chem.} \textbf{67} (2012), 451--466.
\bibitem{p02-VMRC} Keio, \emph{Verifiable Mathematical Reasoning Corpus},
research drafts [087], [094] and [106], commit
\href{https://github.com/QEEDDD/Verifiable-Mathematical-Reasoning-Corpus/tree/8a3f6895d0b9b9f125e43157e4a003a2a2b51704}{8a3f6895d0b9}.
\end{thebibliography}

% END SOURCE BODY Randic_Index_and_Vertex_Radius_Preliminary
\unitstart{Paper 03}{p03}{Caporossi--Hansen}{A proof of the Caporossi--Hansen conjecture on the Randi\'c index and graph radius}{\subtitlevolume{Expanded version}}
\makeatletter\@removefromreset{theorem}{section}\@addtoreset{theorem}{section}\makeatother\renewcommand{\thetheorem}{\thesection.\arabic{theorem}}
% BEGIN SOURCE BODY Randic_Index_and_Vertex_Radius
\begin{abstract}
For every finite connected simple graph of maximum degree $\Delta\ge2$,
we prove $R(G)-\rad(G)\ge\sqrt{\Delta}-3/2$, where $R$ is the classical
Randi\'c index. Thus the paths of even order at least four are exactly
the connected graphs whose Randi\'c index is smaller than their vertex
radius. Equality between the index and vertex radius holds exactly for the singleton, the single edge and
the even cycles. The main estimate bounds the absolute radius $\rho(G)$
by the sum $M(G)$ of reciprocal maximum endpoint degrees. In a minimal counterexample,
deleting the leaves attached to vertices below a degree threshold gives
an upper bound for the number of edges crossing that threshold. Summing
these bounds and constructing a spanning tree of the two-connected block
gives $2\rho(G)\le n-\Delta+4\le2M(G)$ in the proposed counterexample,
where $n$ is the order of that block.
\end{abstract}

\section{Introduction and principal results}
For a finite connected undirected simple graph $G$, let $d_v$ denote the
degree of a vertex $v$. Its Randi\'c index and its vertex radius are
\[
 R(G)=\sum_{uv\in E(G)}\frac{1}{\sqrt{d_ud_v}},\qquad
 \rad(G)=\min_{v\in V(G)}\max_{w\in V(G)}\dist_G(v,w).
\]
Here $uv$ denotes an unordered edge, counted once. The one-vertex graph
has radius zero and an empty edge sum. The notation $P_N$ and $C_N$
denotes a path and a cycle on $N$ vertices, respectively. Throughout,
$\Delta$ denotes maximum degree and $\diam$ denotes diameter. The usual
bound $\diam(G)\le2\rad(G)$ follows by joining two vertices through a
vertex center.

Fajtlowicz proposed $R(G)\ge\rad(G)-1$ \cite{p03-Faj}. Caporossi and Hansen
conjectured $R(G)\ge\rad(G)$ for every connected graph other than a
path of even order and proved the tree case \cite{p03-CH,p03-Zuo}. We obtain a
bound in terms of maximum degree, as well as a complete classification.
\begin{theorem}\label{p03-thm:main}
Let $G$ be a finite connected simple graph with maximum degree $\Delta\ge2$.
Then
\begin{equation}\label{p03-eq:main}
 R(G)-\rad(G)\ \ge\ \sqrt{\Delta}-\tfrac32.
\end{equation}
For the class of all finite connected simple graphs, the following alternatives hold.
\begin{enumerate}
\item $R(G)<\rad(G)$ precisely for $P_{2m}$ with $m\ge2$.
\item $R(G)=\rad(G)$ precisely for $K_1$, $K_2$, and $C_{2m}$ with $m\ge2$.
\item Every other graph satisfies $R(G)>\rad(G)$.
\end{enumerate}
The minimum of $R(G)-\rad(G)$ over all such graphs is $\sqrt2-3/2$,
attained precisely in the first case.
\end{theorem}

The modified Randi\'c index
\[
 M(G)=\sum_{uv\in E(G)}\frac1{\max(d_u,d_v)}
\]
occurs in the diameter analysis of Dvo\v{r}\'ak, Lidick\'y and
\v{S}krekovski \cite{p03-DLS} and in the radius argument of Cygan, Pilipczuk
and \v{S}krekovski \cite{p03-CPS}. The latter proved the Caporossi--Hansen
inequality for graphs of maximum degree at most four and established
$M(G)\ge\rad(G)-1/2$ in that class. Zuo proved the conjecture for
unicyclic and bicyclic graphs \cite{p03-Zuo}. Deng, Tang and Zhang gave a
bound depending on the cyclomatic number \cite{p03-Deng}. Doig studied the
inequality for cactus graphs \cite{p03-Doig}. The spanning-tree estimate used
in Section~\ref{p03-sec:metric} also appears as Lemma~2.1 of \cite{p03-CPS}.

Give every edge of $G$ length one and denote the resulting metric graph
by $|G|$. Write
\[
 \rho(G)=\min_{x\in|G|}\ \max_{w\in V(G)}\dist_{|G|}(x,w).
\]
Thus a center may lie inside an edge, and the vertices are the demand points.
\begin{theorem}[Absolute-radius inequality]\label{p03-thm:abs}
Every finite connected simple graph satisfies $M(G)\ge\rho(G)$.
\end{theorem}
The proof of Theorem~\ref{p03-thm:abs} builds on the minimal-counterexample
arguments in VMRC [087] and [106] and the local degree identity in
VMRC [094] \cite{p03-VMRC}. All reductions are proved in
Sections~\ref{p03-sec:metric}--\ref{p03-sec:sum}. The core with attached leaves
is a bag in the sense of Definition~4.1 of \cite{p03-CPS}. The step that
removes the degree
bound is Proposition~\ref{p03-prop:cut}, which compares each cut with the effect
of deleting just the leaves on its lower-degree side. Section~\ref{p03-sec:finish}
then deduces Theorem~\ref{p03-thm:main} by a comparison edge by edge.

\section{Metric centers and spanning trees}\label{p03-sec:metric}
Distances between vertices in $|G|$ agree with graph distances. With
finitely many vertices, the maximum of their distance functions is a
continuous function on the compact space $|G|$, so an absolute center
exists.
\begin{lemma}\label{p03-lem:half}
An absolute center can be chosen at a vertex or an edge midpoint.
In particular $\rad(G)-1/2\le\rho(G)\le\rad(G)$.
\end{lemma}
\begin{proof}
Consider an edge $uv$ with coordinate $t\in[0,1]$ measured from $u$.
For each vertex $w$ the distance from $t$ to $w$ is
\[
 \min\{\dist_G(u,w)+t,\ \dist_G(v,w)+1-t\}.
\]
Both affine functions have integer intercepts and slopes $\pm1$. Two
such functions of opposite slopes can meet only at a half-integer.
Consequently every individual distance, and their maximum, is affine
on each of the open half-edges. Its minimum on the closed edge is
attained at an endpoint or midpoint. Distances from a center so chosen
to vertices are multiples of $1/2$. Rounding any absolute center to a
nearest vertex increases all its distances by at most $1/2$. The vertex
centers are allowed when minimizing for $\rho$.
\end{proof}
\begin{lemma}\label{p03-lem:tree}
For a tree $T$ of diameter $L$, one has $\rho(T)=L/2$ and
$\rad(T)=\lceil L/2\rceil$. If a connected graph $H$ on $m$ vertices
has a vertex with at least $j\ge2$ incident edges, then it admits a
spanning tree of diameter at most $m-j+1$. In particular
$2\rad(H)\le m-j+2$.
\end{lemma}
\begin{proof}
Let a longest path in $T$ have endpoints $a,b$. For any vertex $w$, its
branch meets the path at a unique point at distance $s$ from $a$; write
$t$ for the distance from that point to $w$. The two paths to the
endpoints give $t\le\min(s,L-s)$. Thus the distance from the midpoint
to $w$ is at most $|s-L/2|+\min(s,L-s)=L/2$. The endpoints force a
radius of at least $L/2$ by the triangle inequality. Rounding the
midpoint yields the asserted integer radius.

The $j$ chosen incident edges form a tree. Extend this tree, one edge
to a new vertex at a time, to a spanning tree $T$ of $H$. At least $j$
components of $T-v$ meet the chosen vertex $v$ by distinct edges.
Each contains a leaf of $T$, as a farthest vertex of that component
from $v$ is a leaf. A longest path in $T$ uses at most two leaves. It
therefore uses at most $m-(j-2)$ vertices and has at most $m-j+1$ edges.
Distances in $H$ can only be shorter than those in its spanning tree.
The vertex-radius bound now follows from $2\lceil L/2\rceil\le L+1$.
\end{proof}

\section{Reductions of a minimal counterexample}\label{p03-sec:red}
Suppose Theorem~\ref{p03-thm:abs} fails. Among counterexamples $G$ choose
one minimizing the pair $(\beta(G),|E(G)|)$ lexicographically, where
$\beta(G)=|E(G)|-|V(G)|+1$ is the cyclomatic number. This order on
pairs of nonnegative integers is well founded. A graph smaller in this
order satisfies $M\ge\rho$. Stars with at least two leaves have
$M=\rho=1$. On a path of order $m\ge3$, each edge contributes $1/2$
to $M$ and Lemma~\ref{p03-lem:tree} gives $M=\rho=(m-1)/2$. Also
$(M,\rho)=(0,0)$ for $K_1$ and $(1,1/2)$ for $K_2$.

\begin{lemma}\label{p03-lem:delete}
If $e$ is a nonbridge edge of $G$, then $M(G-e)>M(G)$. If $z$ is
a leaf with neighbor $v$ of degree $d\ge2$, and $b$ of the remaining
neighbors of $v$ have degree less than $d$, then
\begin{equation}\label{p03-eq:oneleaf}
 M(G)-M(G-z)=\frac1d-\frac{b}{d(d-1)}\ge0,\qquad
 0\le\rho(G)-\rho(G-z)\le\tfrac12.
\end{equation}
The leaf formulas apply to any connected graph in the stated range.
\end{lemma}
\begin{proof}
Deleting a nonbridge preserves the demand vertices, removes potential
center locations, and cannot shorten distances. It reduces $\beta$ by
one. Hence $M(G-e)\ge\rho(G-e)\ge\rho(G)>M(G)$.

Deleting $z$ removes $1/d$ from $M$. If the other endpoint of an edge
at $v$ has degree less than $d$, its weight increases from $1/d$ to
$1/(d-1)$. The other weights stay fixed. As $b\le d-1$, this proves
the index formula.

Project a center lying on the removed pendant interval to $v$. All
distances to retained vertices weakly decrease, proving the first
radius inequality. Conversely, choose a vertex or midpoint center $y$
of $G-z$, and put $r_0=\rho(G-z)$. If $\dist(y,v)\le r_0-1/2$,
then $\dist(y,z)\le r_0+1/2$ after restoration. Otherwise
$\dist(y,v)=r_0$ by Lemma~\ref{p03-lem:half}. As $G-z$ has at least two
vertices, $r_0\ge1/2$. Move $y$ a distance $1/2$ along a shortest
path toward $v$. The distances to retained vertices increase by at
most $1/2$, and the distance to $z$ is $r_0+1/2$.
\end{proof}

A vertex adjacent to a leaf will be called its support. A support of
degree two in $G$ has a second neighbor of degree at least two, because
$G$ is not a star. Formula~\eqref{p03-eq:oneleaf} then gives a drop in
$M$ of $1/2$, at least the drop in $\rho$. Deleting this leaf would
produce a smaller counterexample with the same $\beta$. Thus every
support in $G$ has degree at least three.

If $z,z'$ are leaves with the same support $v$, delete $z$. In the
remaining graph a center can be chosen outside the pendant interval
leading to $z'$, except possibly at $v$. Indeed, projecting a point
of that interval to $v$ decreases distances to all other vertices,
and puts $z'$ at distance one. Any point on that interval has
eccentricity at least one, since another neighbor of $v$ remains.
Restoring $z$ now gives $z$ the same distance from the center as $z'$.
The radius is unchanged and $M$ weakly decreases upon deletion.
Minimality gives a contradiction. Each support therefore has exactly
one leaf.

\begin{lemma}[Splitting at a cut vertex]\label{p03-lem:split}
For every vertex $v$ of the minimal counterexample $G$, at most one
component of $G-v$ contains an edge.
\end{lemma}
\begin{proof}
Otherwise partition $V(G)\setminus\{v\}$ into unions $A,B$ of components
of $G-v$, each containing an edge. For $i\in\{A,B\}$ let $t_i$ be
the number of neighbors of $v$ in $i$, and $m_i$ the number of edges
within $i$. All four numbers are positive. Form $G_A$ by retaining
$G[A\cup\{v\}]$ and adding $t_B$ new leaves at $v$. Define $G_B$
symmetrically. These graphs are connected, $|E(G_A)|=|E(G)|-m_B$,
and $\beta(G_A)=m_A+t_A-|A|$. The analogous expression for $G_B$
gives $\beta(G_A)+\beta(G_B)=\beta(G)$. Both graphs are therefore
smaller. Each retained edge has its original endpoint degrees.
The $t_A+t_B=d_v$ artificial edges across the two graphs each have
weight $1/d_v$, which gives
\begin{equation}\label{p03-eq:splitindex}
 M(G_A)+M(G_B)=M(G)+1 .
\end{equation}

Set $r_i=\rho(G_i)$. In each $G_i$, an artificial leaf and a retained
neighbor of $v$ are at distance two, so $r_i\ge1$. Choose a center
$y_i$ in the retained part: if a center lies on an artificial pendant
edge, project it to $v$. This weakly decreases its distances to all
vertices other than the tip of that edge, and places that tip at
distance $1\le r_i$. If $s_i=\dist(y_i,v)$, the distance to an
artificial leaf is $s_i+1$, so $s_i\le r_i-1$.

In $|G|$ join $y_A$ to $v$ and $v$ to $y_B$ by shortest paths. Their
concatenation is a geodesic since the two retained parts meet only at
$v$. With signed distance $u$ from $v$ along this path, its ends have
coordinates $-s_A,s_B$. Via $y_A$ and $y_B$, respectively, the
distances to the original vertices on the two sides (including $v$)
are at most
\[
 L(u)=r_A+s_A+u,\qquad H(u)=r_B+s_B-u .
\]
If these functions meet on the interval, their common value is
$(r_A+s_A+r_B+s_B)/2\le r_A+r_B-1$. If $L$ dominates throughout,
at $u=-s_A$ their maximum is $r_A$. If $H$ dominates throughout,
at $u=s_B$ the maximum is $r_B$. Both are at most $r_A+r_B-1$.
Consequently $\rho(G)\le r_A+r_B-1\le M(G_A)+M(G_B)-1=M(G)$ by
minimality and \eqref{p03-eq:splitindex}, a contradiction.
\end{proof}

\begin{proposition}\label{p03-prop:core}
The minimal counterexample consists of a two-connected induced subgraph
$K$ of order at least three together with at most one leaf attached to
each vertex of $K$. Every edge of $K$ is a nonbridge in $G$.
\end{proposition}
\begin{proof}
A tree containing a path of length four violates Lemma~\ref{p03-lem:split}
at the middle of that path. A tree of diameter at most two is a
singleton, an edge or a star. In a tree of diameter three, every vertex
lies on the central edge of a diametral path or is adjacent to one of
its ends: the midpoint covers the tree within $3/2$ by
Lemma~\ref{p03-lem:tree}. Each central vertex supports at most one leaf,
so this tree is a path. These possibilities were excluded above.

Thus $G$ contains a cycle. Its vertices induce a two-connected graph,
where two-connected means having at least three vertices and remaining
connected after deletion of any one vertex. Extend that set to a
maximal vertex set inducing a two-connected subgraph $K$. If a component
outside $K$ has two distinct attachment vertices in $K$, it gives a
path between them with all interior vertices outside $K$. Adding the
path preserves two-connectivity: deletion of an interior vertex leaves
each remaining piece attached at an end, and deletion of a vertex of
$K$ leaves $K$ connected with at least one end of the path attached.
This contradicts maximality. Each outside component therefore has a
unique attachment $v$. If it contained an edge, deleting $v$ would
separate that edge from $K-v$, which is connected and has at least
two vertices. Lemma~\ref{p03-lem:split} excludes this. Each outside
component is a single leaf, and the previous reduction bounds their
number per support. Finally, every vertex of $K$ has at least two
neighbors in $K$. For a core edge $uv$, choose another core neighbor
$w$ of $u$; a path in $K-u$ from $w$ to $v$ completes an alternative
route from $u$ to $v$. Hence $uv$ is a nonbridge.
\end{proof}

\section{A bound for each degree cut}\label{p03-sec:cuts}
Fix the subgraph $K$ of Proposition~\ref{p03-prop:core} and write
$n=|V(K)|$. Degrees $d_v$ in this section are always measured in $G$.
At a core vertex $v$, let $p_v$ be $1$ if $v$ supports a leaf and
$0$ otherwise. Define
\[
 \ell_v=\#\{w\sim v:d_w<d_v\},\quad h_v=d_v-\ell_v,\quad
 g_v=d_v(d_v-1),\quad a_v=\ell_v/g_v,\quad b_v=h_v/g_v .
\]
In particular $\ell_v$ includes the leaf when $p_v=1$.

The next lemma includes the local-minimum and minimum-degree properties
of bags from Lemma~4.2 of \cite{p03-CPS}.
\begin{lemma}\label{p03-lem:local}
Every core vertex has $d_v\ge3$ and $\ell_v\ge1$. For a core edge $uv$
with $d_u<d_v$ we have $a_u>b_v$.
\end{lemma}
\begin{proof}
On deleting $uv$, each remaining edge at an endpoint $x$ gains $1/g_x$
precisely when its other endpoint has degree strictly below $d_x$.
No remaining edge occurs at both $u$ and $v$, since $G$ is simple.
Subtracting the weight of the deleted edge gives
\begin{align}
 M(G-uv)-M(G)
 &=a_u+\frac{\ell_v-1}{g_v}-\frac1{d_v}
   =a_u-b_v &&(d_u<d_v),\label{p03-eq:strict}\\
 &=\frac{\ell_u+\ell_v}{d(d-1)}-\frac1d
   &&(d_u=d_v=d).\label{p03-eq:equal}
\end{align}
Both changes must be positive by Lemma~\ref{p03-lem:delete}. If
$\ell_u=0$, there is no leaf at $u$. A core neighbor of higher degree
would contradict \eqref{p03-eq:strict}, since $b_v\ge0$. A neighbor of
equal degree would have $\ell_v\le d-1$, contradicting
\eqref{p03-eq:equal}. These exhaust the neighbors of $u$. Finally, a
vertex of full degree two already has two neighbors in $K$; all
vertices of $K$ have degree at least two. It could have no lower-degree
neighbor.
\end{proof}

A core vertex of minimum full degree must have a leaf, by
Lemma~\ref{p03-lem:local}. Thus $K$ is smaller than $G$ in the order of
Section~\ref{p03-sec:red}: removing the leaves leaves $\beta$ unchanged
and reduces the edge count. Let $\delta=M(G)-M(K)$. Repeated leaf
deletion and minimality give
\begin{equation}\label{p03-eq:delta}
 0\le\delta<\rho(G)-\rho(K)\le1 .
\end{equation}
The last bound holds because a center of $K$ reaches each added leaf
within $\rho(K)+1$. If $\Delta=3$, every core vertex has full degree
three. Its only possible lower-degree neighbor is its leaf, giving
$\ell_v=1$ at every core vertex. Equation~\eqref{p03-eq:equal} gives a
zero change at every core edge, a contradiction. Hence $\Delta\ge4$.

\begin{proposition}[Degree-cut estimate]\label{p03-prop:cut}
For an integer $3\le k<\Delta$ let $S_k=\{v\in V(K):d_v\le k\}$,
and let $N_k$ count edges of $K$ with exactly one end in $S_k$. Then
\begin{equation}\label{p03-eq:cut}
 N_k\le k(k-1)\delta .
\end{equation}
\end{proposition}
\begin{proof}
Write $S=S_k$. Orient each unequal-degree core edge from its lower
endpoint to its higher endpoint. Denote by $I$ the directed edges
inside $S$ and by $F$ the directed edges leaving $S$. Every crossing
edge leaves $S$. Let $e_v$ count the core neighbors with degree $d_v$.
At $v$ there are $\ell_v-p_v$ incoming and $h_v-e_v$ outgoing edges.
Since $h_v a_v=\ell_v b_v$, we have
\[
 p_v b_v=(h_v-e_v)a_v-(\ell_v-p_v)b_v+e_v a_v .
\]
Sum this identity over $v\in S$. An internal directed edge $u\to v$
contributes $a_u-b_v$, and an edge leaving $S$ contributes $a_u$:
\begin{equation}\label{p03-eq:flow}
 \sum_{v\in S}p_v b_v
 =\sum_{(u,v)\in I}(a_u-b_v)+\sum_{(u,v)\in F}a_u
        +\sum_{v\in S}e_v a_v .
\end{equation}
Each term of the first sum is positive by Lemma~\ref{p03-lem:local}.

Delete precisely the leaves whose support belongs to $S$ and call the
result $G_S$. To compute the change of $M$, first consider a deleted
leaf at $v$. It has weight $1/d_v$, and the $\ell_v-1$ edges to
strictly lower core neighbors each gain $1/g_v$. Indeed, on a strict
edge with original degrees $d<e$, the new maximum remains $e$ when
the higher endpoint keeps its leaf or has none. If the higher endpoint
loses a leaf, the new maximum is $e-1$, as $d\le e-1$, irrespective
of any deletion at the lower endpoint. Thus the loss from that leaf,
after these gains, is $1/d_v-(\ell_v-1)/g_v=b_v$.

On a core edge of equal degrees $d$, the maximum decreases precisely
when both endpoints lose a leaf. In that case the edge gains
$1/[d(d-1)]$. Let $Q$ consist of these unordered core edges. The
exact total change is therefore
\begin{equation}\label{p03-eq:subset}
 M(G)-M(G_S)=\sum_{v\in S}p_v b_v-\sum_{uv\in Q}\frac1{g_u}.
\end{equation}
For $uv\in Q$, the last sum of \eqref{p03-eq:flow} counts $a_u+a_v$
on that edge. This is at least $2/g_u$ since $g_u=g_v$ and
$\ell_u,\ell_v\ge1$. All other contributions to that last sum are
nonnegative. Subtracting the correction in \eqref{p03-eq:subset} from
\eqref{p03-eq:flow}, and then deleting the remaining leaves one at a time
by Lemma~\ref{p03-lem:delete}, gives
\[
 \sum_{(u,v)\in F}a_u\le M(G)-M(G_S)\le M(G)-M(K)=\delta .
\]
For every $(u,v)\in F$, one has $\ell_u\ge1$ and $d_u\le k$, so that
$a_u\ge1/[k(k-1)]$. There are $N_k$ such edges, giving \eqref{p03-eq:cut}.
\end{proof}

\section{Summing the cuts and bounding the absolute radius}\label{p03-sec:sum}
\begin{proof}[Completion of the proof of Theorem~\ref{p03-thm:abs}]
Continue with the hypothetical minimal counterexample. Write
$t=\sum_{v\in V(K)}p_v$, so $G$ has $n+t$ vertices. In a graph
without isolated vertices, the sum of $1/d_u+1/d_v$ over the edges
counts each vertex once. Subtracting $2M(G)$ gives
\[
 n+t-2M(G)=\sum_{uv\in E(G)}\left|\frac1{d_u}-\frac1{d_v}\right|.
\]
A leaf at $v$ contributes $1-1/d_v$. A core edge with degrees
$d<e$ contributes $\sum_{k=d}^{e-1}1/[k(k+1)]$, counting exactly
the cuts it crosses. Equal-degree edges contribute zero. We obtain
\begin{equation}\label{p03-eq:sum}
 n-2M(G)=\sum_{k=3}^{\Delta-1}\frac{N_k}{k(k+1)}
                     -\sum_{v\in V(K)}\frac{p_v}{d_v}.
\end{equation}
When all leaves are deleted, the retained core-edge weights can only
increase. Thus the sum of the removed weights $\sum_v p_v/d_v$ is at
least $\delta$. By Proposition~\ref{p03-prop:cut},
\begin{align}
 n-2M(G)
 &\le \left(\sum_{k=3}^{\Delta-1}\frac{k-1}{k+1}-1\right)\delta
 \le(\Delta-4)\delta
 \le\Delta-4 . \label{p03-eq:lowerM}
\end{align}
There are $\Delta-3$ summands, each at most one; the last two
inequalities use $\delta\ge0$, $\Delta\ge4$ and \eqref{p03-eq:delta}.

Choose a core vertex of full degree $\Delta$. It has at least
$\Delta-1$ neighbors in $K$, since at most one is a leaf. By
Lemma~\ref{p03-lem:tree}, $K$ has a spanning tree of diameter at most
$n-(\Delta-1)+1=n-\Delta+2$. Restore each original leaf and its
incident edge to that tree. Between any two vertices of the resulting
tree $T$, the path consists of a core-tree path and at most two leaf
edges. Thus $\diam(T)\le n-\Delta+4$. Distances in $G$ from a
diametral midpoint of $T$ are at most their distances in $T$. Applying
Lemma~\ref{p03-lem:tree} and \eqref{p03-eq:lowerM} gives
\begin{equation}\label{p03-eq:contradiction}
 2\rho(G)\le\diam(T)\le n-\Delta+4\le 2M(G).
\end{equation}
This contradicts $M(G)<\rho(G)$ and proves Theorem~\ref{p03-thm:abs}.
\end{proof}

\section{From absolute radius to vertex radius}\label{p03-sec:finish}
Let $N=|V(G)|$, $r=\rad(G)$, and assume $\Delta\ge2$. We first keep
the difference $r-\rho(G)$ in the estimate; on graphs with a vertex
absolute center this term vanishes.
\begin{proposition}\label{p03-prop:energy}
For each such graph,
\begin{equation}\label{p03-eq:strong}
 R(G)-r \ \ge\
 \frac{N-2r-(\sqrt{\Delta}-1)(r-\rho(G))}{\sqrt{\Delta}+1}.
\end{equation}
\end{proposition}
\begin{proof}
The degree-square identity of Caporossi, Gutman, Hansen and
Pavlovi\'c \cite{p03-CGHP} expresses $N-2R(G)$ as a sum of squares.
For clarity we derive it here and compare the summands, following the
bounded-degree argument in Section~6 of \cite{p03-CPS}. On an edge with
degrees $d\le e$, the contribution to $R(G)-M(G)$ is
$1/\sqrt{de}-1/e\ge0$, and factoring gives
\[
 (d^{-1/2}-e^{-1/2})^2
  =(\sqrt{e/d}-1)\left(\frac1{\sqrt{de}}-\frac1e\right).
\]
Summing the expanded squares counts each vertex once in the terms
$1/d_u+1/d_v$. Since $\sqrt{e/d}\le\sqrt{\Delta}$, we have
\[
 N-2R(G)=\sum_{uv\in E(G)}(d_u^{-1/2}-d_v^{-1/2})^2
 \le(\sqrt{\Delta}-1)(R(G)-M(G)).
\]
Put $s=R(G)-r$ and $\eta=r-\rho(G)$. Theorem~\ref{p03-thm:abs} gives
$R(G)-M(G)\le s+\eta$, whence
$N-2r-2s\le(\sqrt{\Delta}-1)(s+\eta)$. Solving for $s$
gives \eqref{p03-eq:strong}.
\end{proof}

\begin{proof}[Proof of Theorem~\ref{p03-thm:main}]
Apply Lemma~\ref{p03-lem:tree} with $j=\Delta$ to obtain $N-2r\ge\Delta-2$.
Lemma~\ref{p03-lem:half} gives $0\le r-\rho(G)\le1/2$.
Substitution in \eqref{p03-eq:strong} yields
\[
 R(G)-r\ge
 \frac{\Delta-2-(\sqrt{\Delta}-1)/2}{\sqrt{\Delta}+1}
 =\sqrt{\Delta}-\tfrac32 .
\]
This is strictly positive for $\Delta\ge3$. A connected graph of
maximum degree two is a path or a cycle: following successive edges
either closes a cycle, which leaves no available edge to another
vertex, or ends at the two ends of a spanning path. For $N\ge3$ a
path has two edges of degrees $(1,2)$ and $N-3$ of degrees $(2,2)$.
In a cycle all edge degrees are $(2,2)$. Their vertex radii equal
$\lfloor N/2\rfloor$, by Lemma~\ref{p03-lem:tree} on a path and by measuring
the two arc lengths on a cycle. The resulting values are
\[
\begin{array}{c|cc}
 & N\text{ even}&N\text{ odd}\\ \hline
 R(P_N)-\rad(P_N) & \sqrt2-3/2 & \sqrt2-1\\
 R(C_N)-\rad(C_N) & 0 & 1/2
\end{array}
\]
The singleton has $R=r=0$ and the single edge has $R=r=1$.
These cases exhaust all connected simple graphs. Since
$\sqrt2-3/2<0$, the first entry gives exactly the exceptional paths
and the global minimum.
\end{proof}

The coefficient of $\sqrt{\Delta}$ in \eqref{p03-eq:main} cannot be
increased while retaining a constant additive term. For $\Delta\ge3$,
subdivide one edge of the star $K_{1,\Delta}$ once. Its diameter is
three and its vertex radius is two. There are $\Delta-1$ edges of
degrees $(1,\Delta)$, one of degrees $(2,\Delta)$ and one of degrees
$(1,2)$. For this family,
\[
 R(G)-r=\sqrt{\Delta}-2+\frac1{\sqrt2}
            -\frac{1-1/\sqrt2}{\sqrt{\Delta}} .
\]

\begin{thebibliography}{9}
\bibitem{p03-CH}
G. Caporossi and P. Hansen.
Variable neighborhood search for extremal graphs 1: The AutoGraphiX system.
\emph{Discrete Mathematics} \textbf{212} (2000), 29--44.
\href{https://doi.org/10.1016/S0012-365X(99)00206-X}{doi:10.1016/S0012-365X(99)00206-X}.

\bibitem{p03-CGHP}
G. Caporossi, I. Gutman, P. Hansen and L. Pavlovi\'c.
Graphs with maximum connectivity index.
\emph{Computational Biology and Chemistry} \textbf{27} (2003), 85--90.
\href{https://doi.org/10.1016/S0097-8485(02)00016-5}{doi:10.1016/S0097-8485(02)00016-5}.

\bibitem{p03-CPS}
M. Cygan, M. Pilipczuk and R. \v{S}krekovski.
On the inequality between radius and Randi\'c index for graphs.
\emph{MATCH Communications in Mathematical and in Computer Chemistry}
\textbf{67} (2012), 451--466.
\href{https://match.pmf.kg.ac.rs/electronic_versions/Match67/n2/match67n2_451-466.pdf}{Journal PDF}.

\bibitem{p03-Deng}
H. Deng, Z. Tang and J. Zhang.
On a conjecture of Randi\'c index and graph radius.
\emph{Filomat} \textbf{29} (2015), 1369--1375.
\href{https://doi.org/10.2298/FIL1506369D}{doi:10.2298/FIL1506369D}.

\bibitem{p03-Doig}
M. I. Doig.
Block eccentricity, a radius bound, and an application to the Randi\'c index.
Preprint (2023). \href{https://arxiv.org/abs/2309.11613}{arXiv:2309.11613}.

\bibitem{p03-DLS}
Z. Dvo\v{r}\'ak, B. Lidick\'y and R. \v{S}krekovski.
Randi\'c index and the diameter of a graph.
\emph{European Journal of Combinatorics} \textbf{32} (2011), 434--442.
\href{https://doi.org/10.1016/j.ejc.2010.12.002}{doi:10.1016/j.ejc.2010.12.002}.

\bibitem{p03-Faj}
S. Fajtlowicz.
On conjectures of Graffiti.
\emph{Discrete Mathematics} \textbf{72} (1988), 113--118.
\href{https://doi.org/10.1016/0012-365X(88)90199-9}{doi:10.1016/0012-365X(88)90199-9}.

\bibitem{p03-VMRC}
Keio.
\emph{Verifiable Mathematical Reasoning Corpus}.
Research drafts [087], [094] and [106].
\href{https://github.com/QEEDDD/Verifiable-Mathematical-Reasoning-Corpus/tree/8a3f6895d0b9b9f125e43157e4a003a2a2b51704}{Repository snapshot 8a3f6895d0b9}.

\bibitem{p03-Zuo}
L. Zuo.
About a conjecture on the Randi\'c index of graphs.
\emph{Bulletin of the Malaysian Mathematical Sciences Society}
(2) \textbf{35} (2012), 411--424.
\href{https://math.usm.my/bulletin/pdf/v35n2/v35n2p16.pdf}{Journal PDF}.
\end{thebibliography}

% END SOURCE BODY Randic_Index_and_Vertex_Radius
\unitstart{Paper 04}{p04}{Rational Lambert bases}{A larger irrationality region for rational Lambert bases}{}
\makeatletter\@removefromreset{theorem}{section}\makeatother\renewcommand{\thetheorem}{\arabic{theorem}}
% BEGIN SOURCE BODY Rational_Lambert_Cutoff
\begin{abstract}
Changing the dimensions in Lambrinoudis's six-presentation determinant
gives an irrationality criterion for $F(a/b)=\sum_{n\ge1}((a/b)^n-1)^{-1}$
at $\log b/\log a<0.57329026631440$. This includes $20814/299$,
a base beyond the best possible integration of the six previously used
local estimates at their original dimensions and endpoint costs.
We give the dimension-dependent endpoint calculation, the exact conditions
checked on 860 rational intervals, and a separate upper bound for the old
dimensions. The proof uses the generic determinant identities and valuation
bounds in the author's pinned development. The new specialization is
accompanied by rational Python certificates.
\end{abstract}

\begin{theorem}
For coprime integers $a>b\ge1$, write $\theta=\log b/\log a$.
If $\theta<0.57329026631440$, then $F(a/b)$ is irrational.
In particular, $F(20814/299)$ and $F(21548/305)$ are irrational.
\end{theorem}
\begin{center}
\begin{tabular}{lrr}
\toprule
Base & $\theta$ (rounded) & Position relative to the cutoff\\
\midrule
$21548/305$ & 0.573290221375 & below\\
$20814/299$ & 0.573290262817 & below\\
$11/4$ & 0.578129652636 & above\\
\bottomrule
\end{tabular}
\end{center}

\section{Determinant and endpoint costs}
Use the functional and pole determinant in \cite{p04-lam}. Equivalently,
$\mu_X(t^i)=(U^{i+1}-1)^{-1}$ and
$\mu_X((U^j-t)^{-1})=X-\sum_{v=1}^j(U^v-1)^{-1}$; partial fractions
determine the entries. Fix
\[
 (h,k,L,s)=(3000,161,3097,245),\qquad d=L-h=97 .
\]
For scale $N$, let $\Delta_N$ be the Hankel determinant of order $hN$
with entries
\[
 \mu_X\left(\frac{t^{sN+i+j}}
   {\prod_{kN\le m<(k+L)N}(U^m-t)}\right),\qquad 0\le i,j<hN .
\]
For $q>1$ and $X=F(q)$ it is a positive Gram determinant. Normalize by
$Q_N=U^{-g_N}\Delta_N$, where
$g_N=(s-d)N\sum_{i=0}^{hN-1}((k+d)N+i)$ and
$g=\lim g_N/N^3=(s-d)h(k+d+h/2)$.
The generic infinity rearrangement in \cite{p04-lam} bounds the degree at
infinity of each coefficient of $Q_N$ by zero.
Comparison with the geometric Cauchy determinant gives
\[
 \lim_{N\to\infty}N^{-3}\log Q_N(q,F(q))=-S\log q,\qquad
 S=h^3/3+sh^2/2+h(Lk+L^2/2)+g=26\,766\,016\,500 .
\]

At $U=0$ use the finite-part expansion and its two slots per pole from
\cite{p04-lam}. With scaled row index $r\in[0,h]$ the scaled node has lower
bound $r$ before $r=k$ and $(r+k)/2$ thereafter, up to rounding by
$O(1/N)$. Before the pole window its cost is
$(2h+s-L-2r)y$. Inside the window $k\le y\le k+L$ it is
\[
 w(r,y)=\bigl[y^2+2(2h+s-k-L-2r)y+k^2\bigr]/2 .
\]
After the window the cost is nondecreasing. The affine and quadratic
pieces join with the same slope. Thus the minimizing arguments are
$r$ for $r\le k$, $(r+k)/2$ for $k\le r\le2045$, and
$2r-(2h+s-k-L)$ thereafter. The last two arguments meet at
$(4h+2s-k-2L)/3=2045$ and stay inside the pole window. The
finite costs differ by $O(N)$ per slot from the scaled continuous
costs after multiplication by $N^2$. Integrating these three
quadratics and accounting for $g_N$ gives a sufficient asymptotic
clearing cost at zero
\[
 A=g-\int_0^h\min_y w_{\rm piecewise}(r,y)\,dr
 =2\,054\,711\,745 .
\]
All costs here are cubic coefficients, with rounding of lower order.

\section{Cyclotomic costs and the contradiction}
Write $x=n/N$ for the scaled cyclotomic index and
$v_K(x)=\sum_{j\ge1}(K-jx)_+$. For $x<k+L$ the original local bound is
majorized by $B(x)=h+v_L(x)-v_L(2x)$; its clearing cost vanishes for
$x\ge k+L$. The bordered local filtration and residue corank in \cite{p04-lam}
apply to each numerator-weighted state. In its $(A_0,b,d',s')$ notation
the six states used here are
\[
\begin{gathered}
(0,k,d,s),\quad(d,s,0,k),\quad(s,k+d,s-d,0),\\
(0,s-d,d,k+d),\quad(d,k+d,0,s-d),\quad(s-d,k,s,d).
\end{gathered}
\]
The generic transposition and prefix-swap identities transport these
bounds back to the same determinant. Their applicability requires
$0\le d\le s\le k+d$, which the chosen dimensions satisfy.

Let $I$ be the sum of $\int_l^u x(\alpha+\beta x)\,dx$ over the
860 rows $(l,u,j,\alpha,\beta)$ in the accompanying \texttt{cells.tsv}.
Each row certifies a nonnegative saving $\alpha+\beta x$ compared with
$B(x)$ using the transported presentation $j$ on the entire interval.
The checker rebuilds its expression using \cite{p04-lam}, checks each floor
and maximum on both endpoints, and verifies that intervals have pairwise
disjoint interiors. Exact rational three-piece quadrature also recovers
$A$. Summatory totients yield projected weight $6x\,dx/\pi^2$.
The odd tents $v_L(x)-v_L(2x)$ integrate to $L^3/8$ with this
weight, using $\sum_{j\ {\rm odd}}j^{-2}=\pi^2/8$. The constant
term $h$ contributes $3h(k+L)^2/\pi^2$. Thus the degree of the
monic denominator clearer $J_N$ satisfies
\[
 \limsup N^{-3}\deg J_N\le
 C=A+L^3/8+\{3h(k+L)^2-6I\}/\pi^2.
\]
Indeed the only finite coefficient poles are at zero and roots of
unity. Monic division gives $J_NQ_N\in\mathbb Z[U,X]$ with
$U$-degree at most $\deg J_N$ and $X$-degree at most $hN$.
For fixed real $q>1$, $\log\Phi_n(q)=\varphi(n)\log q+O_q(1)$.
The total of the exponents is $O(N^2\log N)$, since only $n=O(N)$
occur and each exponent is $O(N+N^2/n)$. It follows that
$\log J_N(q)=(\deg J_N)\log q+o(N^3)$.
\par
Suppose $q=a/b$ and $F(q)=u/v$ with positive denominator $v$.
Consider the positive integer
\[ b^{\deg J_N}v^{hN}(J_NQ_N)(a/b,u/v). \]
Its logarithm has normalized upper limit at most
$C\log a-S\log(a/b)$. The expression is strictly negative when
$\theta<1-C/S$. The rational checker bounds $\pi$ and logarithms
by convergent series with explicit remainders and verifies
\[
 1-C/S>0.57329026631440>\log299/\log20814 .
\]
This proves the theorem. The old interval where $F(113590/791)$ was
certified is contained in the present one.

\section{Barrier for the old dimensions}
Hold $(h,k,L,s)=(500,27,516,41)$ and the six published local majorants
and endpoint bounds fixed. On $5\le x\le543$, exact comparison on 762
rational subintervals proves that the maximum available saving has the
same weighted integral as the 400 intervals of the upstream certificate.
Above $543$ its cost is already zero.
For $0<x\le5$ the filtration dominates the residue corank in all six
states. With $t=\{K/(cx)\}$,
\[
 v_K(cx)=K^2/(2cx)-K/2+(cx/2)t(1-t).
\]
In the six signed savings the inverse and constant terms cancel.
Bounding positive fractional terms by $cx/8$ times their multiplicities
gives the six upper bounds $0,0,5x/8,x/2,x/2,x/2$, including the
truncation of the pole order at zero. The integral missed below 5
is thus at most $(5/8)\int_0^5 x^2\,dx=625/24$. Adding this to the
exact unrounded old mass gives
\[
 \theta_{\rm old,\,six}<0.573290237682<\log299/\log20814.
\]
These inequalities localize the gain to a change of dimensions.
The new sufficient cutoff still falls short of $\log4/\log11$.
The irrationality of $F(11/4)$ remains unresolved in this work.

\begin{thebibliography}{9}
\bibitem{p04-lam} C. Lambrinoudis, \emph{Nonquadraticity of Lambert series
values}, manuscript and Lean sources, commit
\href{https://github.com/clambro/lambert-nonquadraticity/tree/008e3136a95e1e1ea605837ca9742b8613b2034a}{008e3136a95e}.
The parameter-independent inputs are in \texttt{Determinants/Decay},
\texttt{Refinement/Infinity}, \texttt{Cyclotomic},
\texttt{Transpose}, and \texttt{Symmetry}, and the zero-place moment
block and integrality arguments in the same development.
\end{thebibliography}

% END SOURCE BODY Rational_Lambert_Cutoff
\unitstart{Paper 05}{p05}{Degree indices and distance}{Degree indices and distance in connected graphs}{}
\makeatletter\@removefromreset{theorem}{section}\@addtoreset{theorem}{section}\makeatother\renewcommand{\thetheorem}{\thesection.\arabic{theorem}}
% BEGIN SOURCE BODY Graph_Distance_and_Degree_Indices
\begin{abstract}
We prove Liu's two sharp bounds on the harmonic index in terms of
diameter and order, and the stronger non-tree bounds proposed by
Jerline and Michaelraj. We characterize equality in the modified
Randi\'c index bound $M\ge\rho$ by a recursive construction on connected
graphs. For trees the construction reduces to a path with leaves attached
at pairwise nonadjacent internal vertices. We determine the exact minimum
of the classical Randi\'c index minus the vertex radius at each maximum
degree $\Delta\ge2$ and characterize all extremizers.
For the average-distance conjecture of Fajtlowicz, we prove a stronger
modified-index inequality under a cycle-rank condition on the vertices of
degree at least three. It applies to all cacti and all graphs obtained by
subdividing every edge at least once. A separate estimate proves
$R>\mu$ for every connected graph of order greater than two and maximum
degree at most nine. The unrestricted average-distance
problem remains open in this work.
\end{abstract}
\section{Notation and scope}
Throughout, $G$ is a finite connected undirected simple graph with $n$
vertices. Write $d_v$ for the degree of $v$, $\Delta$ for maximum degree,
$D$ for diameter, $r=\rad(G)$ for vertex radius, and $\mu$ for the mean
distance between two distinct vertices. The three edge sums are
\[
 R=\sum_{uv\in E(G)}\frac1{\sqrt{d_ud_v}},\qquad
 M=\sum_{uv\in E(G)}\frac1{\max(d_u,d_v)},\qquad
 H=\sum_{uv\in E(G)}\frac2{d_u+d_v}.
\]
The absolute radius $\rho$ permits centers inside edges and has demands
only at the vertices. Thus $r$ and $\rho$ can differ by half a unit.

Sections~\ref{p05-sec:harmonic} and~\ref{p05-sec:nontree} use the published
diameter theorem of Dvo\v{r}\'ak, Lidick\'y and \v{S}krekovski
\cite{p05-DLS}. Section~\ref{p05-sec:mean} gives the average-distance results.
Section~\ref{p05-sec:equality} proves the absolute-radius inequality and
classifies its equality cases. Its cut-vertex decomposition and degree
cuts refine the argument of the preceding research manuscript \cite{p05-KS}. The reduced two-connected core is related to the bags studied in
\cite{p05-CPS}. Section~\ref{p05-sec:broom} uses the resulting absolute-radius
bound to determine the sharp vertex-radius inequality.

\section{The harmonic index and diameter}\label{p05-sec:harmonic}
Let $D$ denote the diameter and let $\Delta$ denote the maximum degree.
For the harmonic index
$H(G)=\sum_{uv\in E(G)}2/(d_u+d_v)$,
Liu proved two sharp diameter inequalities for trees and proposed
them for all connected graphs of order $n\ge4$
\cite[Theorem~2.5 and Conjecture~2.6]{p05-Liu}. The following estimate
implies both inequalities.

\begin{proposition}\label{p05-prop:harmonic}
If $G$ is connected and $n\ge3$, then
\begin{equation}\label{p05-eq:harmonic}
 H(G)\ \ge\ \frac D2+\frac{n-D}{\Delta+1}
       \ \ge\ \frac D2+\frac{n-D}{n-D+2}.
\end{equation}
\end{proposition}
\begin{proof}
On an edge with endpoint degrees $a\le b$, let $h=2/(a+b)$ and
$m=1/b$. The identities
\[
 h-m=\frac{b-a}{b(a+b)},\qquad
 \frac1a+\frac1b-2h
   =\left(\frac ba-1\right)(h-m)
\]
imply the following exact formula, since summing reciprocal degrees
over all edge incidences gives $n$:
\begin{equation}\label{p05-eq:harmonic-slack}
 (\Delta+1)H
 =n+(\Delta-1)M+
   \sum_{uv\in E(G)}\left(\Delta-\frac ba\right)(h-m).
\end{equation}
Every summand is nonnegative.
Dvo\v{r}\'{a}k, Lidick\'{y}, and \v{S}krekovski proved
$M(G)\ge D/2$ for every connected graph \cite[Theorem~3]{p05-DLS}.
There is also an internal derivation from Theorem~\ref{p05-thm:absolute}:
the triangle inequality for two diametral endpoints gives
$\rho(G)\ge D/2$. Substitution gives the first inequality of~\eqref{p05-eq:harmonic}.

Take a shortest path joining vertices at distance $D$. Its $D+1$
vertices induce a path. A vertex on the path has at most two neighbors
on it. A vertex outside the path has at most three neighbors on it,
since the first and last such neighbors can be at most two path edges
apart. In either case its degree is at most $n-D+1$. Thus
$\Delta+1\le n-D+2$, proving the second inequality.
\end{proof}

\begin{theorem}[Liu's two diameter bounds]\label{p05-thm:liu}
For every connected graph of order $n\ge3$,
\[
 H(G)-D\ \ge\ \frac56-\frac n2,\qquad
 \frac{H(G)}{D}\ \ge\ \frac12+\frac{1}{3(n-1)}.
\]
Each equality holds if and only if $G$ is the path $P_n$.
\end{theorem}
\begin{proof}
Set $k=n-D\ge1$. By~\eqref{p05-eq:harmonic},
$H-D/2\ge k/(k+2)\ge1/3$. Subtracting $D/2$ and using
$D\le n-1$ gives the first assertion. Dividing by $D$ and using
the same bound gives the second. If $k>1$ then $k/(k+2)>1/3$
and both inequalities are strict. If $k=1$ the diametral path
contains every vertex and is induced, so $G=P_n$. The two end edges
of $P_n$ contribute $2/3$ each and the other $n-3$ edges
contribute $1/2$ each. Therefore $H(P_n)=(n-1)/2+1/3$.
\end{proof}

For the additive assertion alone there is an even shorter argument.
On every connected graph here other than a path one has $D\le n-2$ and hence
$H-D\ge M-D\ge-D/2\ge1-n/2$. The path computation above
supplies equality in Liu's bound.

\section{The harmonic index on graphs containing a cycle}\label{p05-sec:nontree}
Jerline and Michaelraj established stronger inequalities for unicyclic graphs
and proposed the corresponding extremal problem for connected graphs
containing a cycle \cite{p05-JM}. Deng, Vetr\'ik and Balachandran obtained
$H\ge D/2+2/3$ on unicyclic graphs of order at least seven \cite{p05-DVB}. Abdolghafourian and Iranmanesh treated
quasi-trees and recorded the non-tree conjecture for $n\ge7$
\cite[Conjecture~2]{p05-AI}.

\begin{theorem}\label{p05-thm:jm}
Let $G$ be a connected graph of order $n\ge7$ containing a cycle.
Then
\[
 H(G)-D\ \ge\ \frac53-\frac n2,\qquad
 \frac{H(G)}D\ \ge\ \frac12+\frac{2}{3(n-2)}.
\]
For either inequality, equality holds precisely on the graph obtained
from $C_4$ by attaching paths of lengths $1$ and $n-5$ at opposite
vertices.
\end{theorem}
\begin{proof}
Put $k=n-D$. Every diametral path is induced, so $k\ge2$.
If $k\ge3$, Proposition~\ref{p05-prop:harmonic} gives
\[
 H-D+\frac n2\ge\frac{k}{2}+\frac{k}{k+2}>\frac53.
\]
For the ratio we need $H-D/2\ge 2D/[3(n-2)]$.
If $k\ge4$, then $k/(k+2)\ge2/3$ and $D<n-2$, giving strict
inequality. If $k=3$, the bound $3/5$ is strictly sufficient
when $D<9$, since $n-2=D+1$. For $D\ge9$ we show $H-D/2\ge2/3$.
As $\Delta\le k+1=4$, the first bound of~\eqref{p05-eq:harmonic}
already gives $3/4$ when $\Delta\le3$.

Suppose $k=3$, $D\ge9$, and $\Delta=4$. Each of the two vertices
outside a fixed diametral path $P$ has at most three neighbors on $P$.
As $P$ has $D-1\ge8$ internal vertices, one of them has degree two in $G$.
The component containing it in the subgraph induced by vertices of
degree two is a path, possibly a single vertex. Indeed, a cycle of such
vertices would be a whole component of $G$, contrary to connectivity and
$\Delta=4$. This path has exactly two boundary edges. The possible
degree pairs on a boundary edge are $(1,2),(2,3),(2,4)$, whose
summands in~\eqref{p05-eq:harmonic-slack} are $1/3,1/6,1/6$, respectively.
Consequently the total slack is at least $1/3$, whence
$H-D/2\ge (3+1/3)/5=2/3$. The ratio bound is strict since $D<n-2$.

It remains to treat $k=2$, when $D=n-2\ge5$. Write
$P=v_0\ldots v_D$, let $x$ be the sole vertex outside it, and put
$S=N(x)$, $q=|S|$. A cycle forces $q\ge2$. The indices in $S$
span at most two consecutive path edges, so $q\in\{2,3\}$.
Here both proposed inequalities are equivalent to $H\ge D/2+2/3$.
Since $H(P)=D/2+1/3$, it suffices to show $H(G)-H(P)\ge1/3$.

For each $v_i\in S$ the added edge increases $d(v_i)$ by one.
For an original path edge of degree sum $s$, one increment lowers
its harmonic weight by $2/[s(s+1)]$. This is $1/6$ on an end edge
and $1/10$ on an internal edge. If both endpoints are incremented,
the second decrement is strictly smaller.
Suppose first that $S$ contains no endpoint of $P$. The $q$ new edges
have total weight $2q/(q+3)$. Among the $2q$ incidences incremented
on $P$, at most one is on an end edge, since $D\ge5$ and the indices
in $S$ span at most two. It follows that
\[
 H(G)-H(P)\ge\frac{2q}{q+3}-\frac16-\frac{2q-1}{10}
 =\frac13 \qquad(q=2,3).
\]
Equality requires exactly one end incidence and no two consecutive
indices in $S$. For $q=3$ two indices must be consecutive. For
$q=2$ equality is possible precisely when $S=\{v_1,v_3\}$ up
to reversing $P$.

Finally suppose that $S$ contains an endpoint. It contains only one,
since its indices span at most two. There are $2q-1$ path incidences
and at most two are on an end edge. Thus
\[
 H(G)-H(P)\ge
 \frac{2}{q+2}+\frac{2(q-1)}{q+3}
 -\frac{2}{6}-\frac{2q-3}{10}>\frac13
 \qquad(q=2,3).
\]
When $S=\{v_1,v_3\}$ the only cycle is the four-cycle through
$v_1,v_2,v_3,x$ and all the incidence losses above are exact.
Conversely, the indicated graph has diameter $n-2$ and this
neighborhood, so it attains both equalities.
\end{proof}

The order threshold is necessary. Attach one leaf to each of two opposite
vertices of $C_4$. Then $n=6$, $D=4$ and $H=13/5$, which violates both
bounds of Theorem~\ref{p05-thm:jm}. This example also appears in \cite{p05-AI}.

\section{Average distance and the modified Randi\'c index}\label{p05-sec:mean}
Write $\mu(G)=\binom n2^{-1}\sum_{\{u,v\}\subseteq V(G)}d_G(u,v)$.
The proposed sufficient inequality for Fajtlowicz's conjecture \cite{p05-Faj} is
\begin{equation}\label{p05-eq:strong-mean}
 M(G)\ \ge\ \mu(G)-1+2/n.
\end{equation}
Stars attain equality. We give a certificate for \eqref{p05-eq:strong-mean} in
terms of degrees and cycle rank, then prove the implication for $R$.
All statements in this section assume $n\ge2$.

\begin{lemma}\label{p05-lem:mean-leaves}
Suppose that $n\ge3$ and $G$ has a spanning tree $T$ with $I$ internal
vertices. If $2M(G)\ge I+1$, then
\[
 M(G)-\mu(G)+1-2/n\ \ge\
 (I-1)\left(\frac{2}{n}-\frac{1}{2(n-1)}\right).
\]
\end{lemma}
\begin{proof}
If deleting an edge of $T$ gives parts of orders $s$ and $n-s$, exactly
$s(n-s)$ unordered pairs use this edge. The $n-I$ pendant edges each
contribute $n-1$, and the other $I-1$ edges contribute at most $n^2/4$
each. Since distances in $G$ are at most those in $T$,
\[
 \mu(G)\le \mu(T)\le\frac{2(n-I)}n+\frac{(I-1)n}{2(n-1)}.
\]
Subtract this bound from $(I+1)/2+1-2/n$.
\end{proof}

\begin{theorem}\label{p05-thm:mean-core}
Let $J$ be the subgraph induced by the vertices of $G$ with degree at
least three, let $\beta_J=|E(J)|-|V(J)|+\kappa(J)$ be its cycle rank, and put
\[
 B=\sum_{uv\in E(J)}\left(1-\frac{2}{\min(d_u,d_v)}\right),
\]
where $d_u,d_v$ are degrees in $G$ and $\kappa(\varnothing)=0$.
If $B\ge\beta_J$, then \eqref{p05-eq:strong-mean} holds, with equality exactly
for stars. This applies whenever the cycles in $J$ are pairwise
edge-disjoint. In particular, it applies to all cacti, all trees,
and all graphs obtained by subdividing every edge of a simple connected
graph at least once.
\end{theorem}
\begin{proof}
The case $K_2$ follows directly, so take $n\ge3$.
Write $m=|E(G)|$, $c=m-n+1$, and $L=|\{v:d_v=1\}|$. By summing the
reciprocals of the endpoint degrees,
\[
 n-2M(G)=\sum_{uv\in E(G)}|1/d_u-1/d_v|.
\]
For each edge, split the difference at $1/2$ using the triangle
inequality. Its slack is zero when an endpoint has degree two or the two
reciprocals are on opposite sides of $1/2$. There is no edge joining two
leaves since $n\ge3$. When $d_u,d_v\ge3$, the slack equals
$1-2/\min(d_u,d_v)$. The sum of the split terms is
$\sum_v|1-d_v/2|=L+m-n=L+c-1$. Hence
\begin{equation}\label{p05-eq:core-identity}
 n-2M(G)=L+c-1-B.
\end{equation}
Extend a spanning forest of $J$ to a spanning tree $T$ of $G$.
Of the $c$ omitted edges, exactly $\beta_J$ have both ends in $J$,
since the initial forest already connects each component of $J$.
Every other omitted edge has an endpoint of degree two in $G$,
because all original pendant edges are bridges. Such an endpoint
becomes a leaf of $T$. Different omitted edges cannot use the same
degree-two endpoint, since $T$ is connected. Thus $T$ has at least
$L+c-\beta_J$ leaves. By \eqref{p05-eq:core-identity} and $B\ge\beta_J$, its number
$I$ of internal vertices satisfies $I+1\le2M(G)$.
Lemma~\ref{p05-lem:mean-leaves} proves \eqref{p05-eq:strong-mean}.
Its bound is strict for $I>1$. If $I=1$, then $T$ is a star.
When $G\ne T$ we have $M(G)\ge m/(n-1)>1$ and
$\mu(G)\le\mu(T)=2-2/n$, so equality requires $G=T$.
Conversely every star has $M=1$ and $\mu=2-2/n$.

If the cycles of $J$ are pairwise edge-disjoint, removing one edge
from each yields a spanning forest of $J$, so there are exactly $\beta_J$
cycles. Each cycle contributes at least $3(1-2/3)=1$ to $B$.
This gives $B\ge\beta_J$. A cactus has pairwise edge-disjoint cycles,
and subdividing every edge makes $J$ edgeless.
\end{proof}

\begin{proposition}\label{p05-prop:mean-degree}
Inequality~\eqref{p05-eq:strong-mean} also holds if
$\Delta\le2\delta$, or if $\Delta\le3$, where $\delta$ and $\Delta$
denote the minimum and maximum degrees.
\end{proposition}
\begin{proof}
First $\mu(G)\le(n+1)/3$. To see this, use a spanning tree, remove a
leaf, and bound its distance sum by $1+\cdots+(n-1)$: the $j$th
distance in nondecreasing order is at most $j$, as all intermediate
distances occur. Induction bounds the whole distance sum by
$\binom{n+1}{3}$. Next, for endpoint degrees $a\le b$,
\[
 \frac{1/b}{1/a+1/b}=\frac{a}{a+b}\ge\frac{\delta}{\delta+\Delta}.
\]
Summing gives $M(G)\ge n\delta/(\delta+\Delta)$. If $\Delta\le2\delta$,
this yields $M(G)\ge n/3\ge(n+1)/3-1+2/n$ for $n\ge3$.
If $\Delta\le3$, either $G$ is a tree and Theorem~\ref{p05-thm:mean-core}
applies, or $m\ge n$ and $M(G)\ge m/3\ge n/3$. Check $K_2$ directly.
\end{proof}

\begin{proposition}\label{p05-prop:mean-transfer}
If $G$ satisfies~\eqref{p05-eq:strong-mean}, then $R(G)\ge\mu(G)$,
with strict inequality for $n>2$.
\end{proposition}
\begin{proof}
For endpoint degrees $a\le b$, put $x=1/\sqrt a$ and $y=1/\sqrt b$.
Then
$(x-y)^2=(\sqrt{b/a}-1)y(x-y)\le(\sqrt\Delta-1)(xy-y^2)$.
Summing over edges gives $n-2R\le(\sqrt\Delta-1)(R-M)$, hence
\[
 (\sqrt\Delta+1)(R-\mu)\ge
 n-2\mu-(\sqrt\Delta-1)(1-2/n).
\]
Using $\mu\le(n+1)/3$ and $\Delta\le n-1$, the right-hand side is at least
\[
 (n-2)\left(\frac13-\frac{\sqrt{n-1}-1}{n}\right).
\]
The bracket is positive since $(n/3+1)^2-(n-1)=(n^2-3n+18)/9>0$.
For $K_2$, $R=\mu=1$.
\end{proof}

\begin{theorem}\label{p05-thm:mean-nine}
Every connected graph of order at least two and maximum degree at most nine satisfies
$R(G)\ge\mu(G)$, strictly for $n>2$.
\end{theorem}
\begin{proof}
The preceding results cover $\Delta\le3$, so suppose $4\le\Delta\le9$.
Let $L$ be the number of leaves and $j=\max\{\Delta,L\}$.
A spanning tree containing the full star at a vertex of degree
$\Delta$ has at least $\Delta$ leaves, one in each branch. It also
retains all $L$ original leaves.
For $a=(n-2)^2/[2n(n-1)]$, the tree distance estimate in
Lemma~\ref{p05-lem:mean-leaves} and the path bound give
\[
 \mu\le\min\{(n+1)/3, U(j)\},\qquad
 U(j)=\frac{2j}{n}+\frac{(n-j-1)n}{2(n-1)}.
\]
Here $U$ has slope $-a$ and equals $(n+1)/3$ at
$j_0=n(n-1)/[3(n-2)]$.

For a nonpendant edge with degrees $s\le t$, its contribution to $R$
divided by $1/s+1/t$ is $\sqrt{t/s}/(1+t/s)$, decreasing in $t/s\ge1$.
Thus this ratio is at least $\sqrt{2\Delta}/(\Delta+2)$.
At a pendant edge with support degree $d$, the deficit from
$k(1+1/d)$ is $k(1+1/d)-1/\sqrt d$. This increases with $d\ge2$
whenever $k\le1/2$. For all $\Delta\le9$ we may take
$k=140/363$ and $b=35/363$, since
$k\le3\sqrt2/11$ and $10k/9-1/3\le b$. When $\Delta\le6$ we may
instead take $k=43/100$, $b=1/10$, as
$k\le\sqrt{12}/8$ and $7k/6-1/\sqrt6\le b$.
Summing over edges gives $R\ge kn-bL\ge kn-bj$.

For $n\ge5$, $a\ge9/40>b$ for either choice (the expression for $a$
increases with $n$). Hence the function
$kn-bj-\min\{(n+1)/3,U(j)\}$, on $j\ge d_0$, attains its minimum at
$j_* = \max\{d_0,j_0\}$. For $4\le\Delta\le6$ use $d_0=4$ and $k=43/100$, $b=1/10$.
For $7\le\Delta\le9$ use $d_0=7$ and $k=140/363$, $b=35/363$.
It remains to bound the minimum as a function of $n$.

Since $j_0=(n+1+2/(n-2))/3$ increases for $n\ge5$, the regime
$j_0\le d_0$ consists of the integer intervals $5\le n\le10$ when
$d_0=4$, and $8\le n\le19$ when $d_0=7$. With $j_*=d_0$ the minimum
equals
\[
 (k-\tfrac12)n-bd_0+\tfrac{d_0}{2}
   +\frac{d_0}{2(n-1)}-\frac{2d_0}n.
\]
It is concave in $n\ge5$, since its second derivative is
$d_0[(n-1)^{-3}-4n^{-3}]<0$. The endpoint values on the first
interval are $3/20$ and $29/90$, and those on the second are
$959/1452$ and $9145/82764$. They are all positive.
In the complementary regime $j_*=j_0$, the minimum is
\[
 (k-b/3-1/3)n-1/3-\frac{bn}{3(n-2)}.
\]
This increases since $k-b/3-1/3>0$ for both choices of $k,b$.
At the first eligible orders $n=11$ and $n=20$ its values are,
respectively, $871/2700$ and $343/9801$, both positive.
\end{proof}

Previous results for $R\ge\mu$ include trees \cite{p05-CPSMean}, unicyclic
graphs \cite{p05-YL}, and a degree-based sufficient condition implying an
almost-all-graphs result \cite{p05-Dal}. The unrestricted cases of
\eqref{p05-eq:strong-mean} and $R\ge\mu$ remain unresolved here.

\section{Equality in the absolute-radius inequality}\label{p05-sec:equality}
Let $|G|$ be the metric graph in which every edge is a unit interval.
Demand points are the vertices of $G$. Define
\[
 \rho(G)=\min_{x\in|G|}\max_{w\in V(G)}\dist_{|G|}(x,w).
\]
Bisect every edge once to form $S(G)$. For a half-integer $t\ge0$ write
$W(G,t)$ for the following distance condition:
\[
 \text{for each }z\in V(S(G))\text{ some }w\in V(G)
       \text{ satisfies }\dist_{S(G)}(z,w)\ge 2t .
\]
To see that $W(G,t)$ is equivalent to $\rho(G)\ge t$, take an edge $uv$
with coordinate $s$ measured from $u$. Its distance to $w$ is the
minimum of $\dist_G(u,w)+s$ and $\dist_G(v,w)+1-s$. All the affine
functions appearing here have integer intercepts and slopes $\pm1$.
No two distinct slopes cross in the interior of either half-edge, so the
maximum distance to the demands is affine on each half-edge.
A vertex or midpoint therefore realizes the minimum. Rounding such a
center gives $\rad(G)-1/2\le\rho(G)\le\rad(G)$.

For a tree $T$ of diameter $D$, one has $\rho(T)=D/2$ and
$\rad(T)=\lceil D/2\rceil$. Indeed, a branch of height $h$ meeting a
diametral path at position $i$ has $h\le\min(i,D-i)$, so the midpoint
covers it within $D/2$. The path endpoints force the opposite inequality.
We use the following description as the base of a recursive construction.

\begin{lemma}[All equality trees]\label{p05-lem:eqtree}
Every tree satisfies $M(T)\ge\rho(T)$. Equality holds for $K_1$ and for
exactly the following trees of order at least three. Take a path of length
at least two and attach any number of additional leaves to each vertex of
an independent subset of its internal vertices.
\end{lemma}
\begin{proof}
For a tree of order at least three let $I$ be the number of nonleaves and
let $J$ be the subgraph induced by vertices of degree at least three.
Splitting each $|1/d_u-1/d_v|$ at $1/2$, as in
\eqref{p05-eq:core-identity}, gives
\[
 2(M(T)-\rho(T))=
 I-(D-1)+\sum_{uv\in E(J)}\left(1-\frac{2}{\min(d_u,d_v)}\right).
\]
A diametral path contains $D-1$ internal vertices, all nonleaves in $T$.
Both terms are nonnegative. Equality says that every nonleaf lies on
that path and that no two vertices of degree at least three are adjacent.
These are exactly the stated trees. Conversely the original path stays
diametral after the permitted attachments, since their supports are
internal. Directly, $(M,\rho)=(0,0)$ on $K_1$ and $(1,1/2)$ on $K_2$.
\end{proof}

We next give a construction for graphs of arbitrary cycle rank. A label
$t$ is carried along from the base trees. The tests in the construction
use only degrees and distances in a subdivision. The cut operation can
have two cyclic inputs.

\smallskip
\noindent\textbf{Base.} Use any tree of Lemma~\ref{p05-lem:eqtree} achieving
equality, with label $t=D/2$.

\smallskip
\noindent\textbf{L. Add a leaf.} In a labelled graph of order at least two,
choose a vertex $v$ of degree $d$ and write $b$ for the number of its
neighbors whose degrees are at most $d$. Set
$c=(d-b)/[d(d+1)]$. Attach a new leaf to $v$ if
$c\in\{0,1/2\}$ and the resulting graph $G$ satisfies $W(G,t+c)$.
Its label is $t+c$.

\smallskip
\noindent\textbf{E. Add a cycle edge.} Join two nonadjacent vertices $u,v$
of one labelled graph to form $G$. Use the degrees \emph{after} insertion.
At each endpoint $x$ write $g_x=d_x(d_x-1)$ and let $\ell_x$ count its
neighbors of strictly smaller degree. Relabel so $d_u\le d_v$.
Use the edge when $W(G,t)$ holds and the corresponding expression is zero:
\[
 \begin{cases}
 \ell_u/g_u-(d_v-\ell_v)/g_v, & d_u<d_v,\\
 (\ell_u+\ell_v)/g_u-1/d_u, & d_u=d_v .
 \end{cases}
\]
The label remains $t$.

\smallskip
\noindent\textbf{C. Join at a cut vertex.} Choose disjoint labelled graphs
$G_A,G_B$, labels $t_A,t_B$, with roots of the same degree $a+b$ where
$a,b\ge1$. Remove $b$ chosen leaf neighbors of the first root and $a$
of the second, then identify the roots to form $G$. Each retained side
must contain an edge not incident with the root. Use this join if
$W(G,t_A+t_B-1)$ holds. Its label is $t_A+t_B-1$.

\begin{theorem}[Recursive classification]\label{p05-thm:absolute}
Every finite connected simple graph satisfies $M(G)\ge\rho(G)$.
Equality holds if and only if $G$ is obtained from the base trees by a
finite sequence of operations L, E and C. For every such construction the
final label equals $M(G)=\rho(G)$.
\end{theorem}
The distance test in E is necessary. Start with $P_4=v_0v_1v_2v_3$
and add $v_1v_3$. This keeps $M=3/2$, but $v_1$ now covers every
vertex within one and $\rho$ drops to one.
Cyclic equality graphs do occur. For example, the triangle with one
pendant leaf at each vertex has $M=\rho=2$, as does the four-cycle with
a pendant leaf at each of two opposite vertices. For the triangle
example every vertex or midpoint has some leaf at distance at least two,
and a triangle vertex covers all demands within two. The four-cycle
example has diameter four and a vertex covering all demands within two.

\begin{proof}
We prove the lower bound and necessity together by induction on the
lexicographic pair $(|E|-|V|+1,|E|)$. Lemma~\ref{p05-lem:eqtree} handles
trees. Consider a cyclic graph $G$ with $M(G)\le\rho(G)$. We show it
has equality and admits a reverse move to smaller equality graphs.

\emph{Leaves and cycle edges.} On adding a leaf to $v$, its new edge
contributes $1/(d+1)$. Each of the $b$ old edges to a neighbor of degree
at most $d$ loses $1/[d(d+1)]$. Thus the change in $M$ is exactly the
nonnegative number $c$ in L. Deleting a leaf $z$ decreases $\rho$ by at
most $1/2$ and cannot increase it. For the latter statement project any
center lying on the deleted pendant interval to its support $v$. For the
former let $r_0=\rho(G-z)$ and take a vertex or midpoint center $x$.
If $\dist(x,v)\le r_0-1/2$, the same center works within $r_0+1/2$
after restoration. Otherwise $\dist(x,v)=r_0\ge1/2$ by half-integrality.
Move $x$ half a step toward $v$ and restore the leaf.

If $z$ has support of degree two, its other neighbor has degree at least
two since $G$ contains a cycle. Deletion drops $M$ by $1/2$. Induction
on $G-z$ and the radius bound force equality on both graphs, with
labels differing by $1/2$. If two leaves share their support, delete
one. A center of the remaining graph can be chosen off the other leaf
interval: projecting to its support decreases distances to the rest
and puts that leaf at distance one, within the radius since there is
another neighbor. Restoring the twin preserves $\rho$. Monotonicity of
$M$ and induction force equality on both graphs with $c=0$. Both
reductions reverse by L.

For a nonbridge $uv$, deleting it cannot decrease $\rho$. Each surviving
edge incident with an endpoint $x$ gains $1/g_x$ if its other end had
degree strictly below $d_x$. Subtracting the removed edge weight gives
exactly the expression in E for $M(G-uv)-M(G)$. If this quantity is
nonpositive, induction gives
\[
 M(G)\ge M(G-uv)\ge\rho(G-uv)\ge\rho(G)\ge M(G).
\]
Hence the operation reverses by E. We may now assume every nonbridge
deletion strictly increases $M$.

\emph{Cut vertices.} If some $v$ separates two parts each containing an
edge away from $v$, partition the components of $G-v$ into two such
groups $A,B$. Denote the numbers of neighbors of $v$ on the two sides
by $a,b$. In the $A$ side attach $b$ new leaves at $v$, and in the $B$
side attach $a$ there, giving $G_A,G_B$. Each graph has fewer edges
and their cycle ranks sum to that of $G$. Original edge weights are
preserved on their sides and the $a+b$ new edges each weigh $1/(a+b)$.
Thus $M(G)=M(G_A)+M(G_B)-1$.

Write $r_i=\rho(G_i)$. Each $r_i\ge1$, since an artificial leaf and a
retained neighbor have distance two. Choose centers $x_i$ in the retained
part by projecting off any artificial leaf edge. Their distances $s_i$
from $v$ satisfy $s_i+1\le r_i$. On the path $x_A v x_B$ use coordinate
$u\in[-s_A,s_B]$ with $v$ at zero. The two groups are covered from $u$
within $r_A+s_A+u$ and $r_B+s_B-u$, respectively, by going through
their centers. At the intersection of these affine bounds the common
value is at most $r_A+r_B-1$. If it lies outside the interval, one
endpoint has covering bound $r_A$ or $r_B$. Thus
$\rho(G)\le r_A+r_B-1$. Together with induction and $M(G)\le\rho(G)$
this forces equality in both inputs and in the bound on the joined
radius. This is the reverse of C.

\emph{Excluding a terminal cyclic graph.} It remains to show that the
preceding reductions cannot all fail. Choose a maximal induced
two-connected subgraph $K$ containing a cycle. Each outside component
attaches at a unique vertex of $K$, since an outside path between two
distinct attachments can be added while preserving two-connectivity.
If such a component contained an edge, its attachment would give the
cut reduction. Thus $G$ is $K$ with at most one pendant leaf at each
vertex, since twin leaves were dealt with above. Every core edge is a
nonbridge. Set $p_v=1$ when a core vertex $v$ bears a leaf and $0$
otherwise, with all degrees measured in $G$. At a core vertex set
\[
 \ell_v=\#\{w\sim v:d_w<d_v\},\quad h_v=d_v-\ell_v,\quad
 g_v=d_v(d_v-1),\quad a_v=\ell_v/g_v,\quad b_v=h_v/g_v .
\]
The deletion formula in E gives $a_u>b_v$ on each ascending core edge
$u\to v$. At an equal-degree edge it gives $\ell_u+\ell_v>d_u-1$.
If $\ell_u=0$, a larger neighbor violates the first inequality, and an
equal neighbor violates the second since $\ell_v\le d_v-1$ there.
Such a vertex has no smaller neighbor, which exhausts the possibilities.
Thus $\ell_v\ge1$ everywhere on $K$ and every full degree there is at
least three. A core vertex of minimum full degree has a pendant leaf.
Write $q=|V(K)|$ and $\delta_M=M(G)-M(K)$. Leaf monotonicity and induction
on the smaller graph $K$ give
\[
 0\le\delta_M\le\rho(G)-\rho(K)\le1 .
\]
The upper bound follows by using a center of $K$ in $G$. When
$\Delta=3$, all core vertices have degree three and exactly one smaller
neighbor. Deleting an equal-degree edge then costs zero, a contradiction.
Hence $\Delta\ge4$.

For $3\le k<\Delta$ put $S_k=\{v\in K:d_v\le k\}$ and let $N_k$ count
core edges crossing from $S_k$ to its complement. We claim
\begin{equation}\label{p05-eq:cutclassification}
 N_k\le k(k-1)\delta_M .
\end{equation}
Orient unequal-degree core edges upward. Write $I_k$ for the directed
edges internal to $S_k$ and $F_k$ for the ones leaving it. Write $e_v$
for the number of equal-degree neighbors of $v$. The incoming and
outgoing counts are $\ell_v-p_v$ and $h_v-e_v$. Summing the identity
$h_va_v=\ell_vb_v$ over $S_k$ in this form gives
\[
 \sum_{v\in S_k}p_v b_v
 =\sum_{(u,v)\in I_k}(a_u-b_v)
    +\sum_{(u,v)\in F_k}a_u+\sum_{v\in S_k}e_v a_v .
\]
Delete all leaves whose support is in $S_k$. The net loss of $M$ is
$\sum_{v\in S_k}p_vb_v-\sum_{uv\in Q}1/g_u$, where $Q$ is the set of
equal-degree edges both of whose ends lose a leaf. Indeed a removed
leaf has weight $1/d_v$, and its $\ell_v-1$ strictly lower core neighbors
give a gain of $1/g_v$ each. On a strict edge this gain depends only on
deletion at the higher end. On an equal edge a gain occurs exactly for
the edges in $Q$. Each correction $1/g_u$ is covered by $a_u+a_v$ in
the last sum of the display because $\ell_u,\ell_v\ge1$. The first sum
is nonnegative, so the net loss is at least $\sum_{F_k}a_u$. Deleting
the remaining leaves cannot increase $M$. Hence
$\delta_M\ge\sum_{F_k}a_u\ge N_k/[k(k-1)]$, proving the claim.

The reciprocal-degree identity telescopes along the core edges:
\[
 q-2M(G)=\sum_{k=3}^{\Delta-1}\frac{N_k}{k(k+1)}
                 -\sum_{v\in K}\frac{p_v}{d_v}.
\]
Indeed a strict edge of degrees $d<e$ contributes
$1/d-1/e=\sum_{k=d}^{e-1}1/[k(k+1)]$, and each pendant edge contributes
$1-1/d_v$ to $|V(G)|-2M(G)$. The removed leaf weights sum to at least
$\delta_M$, since core weights can only grow under leaf deletion. If
$\delta_M>0$, \eqref{p05-eq:cutclassification} gives
\[
 q-2M(G)\le
 \left(\sum_{k=3}^{\Delta-1}\frac{k-1}{k+1}-1\right)\delta_M
 <(\Delta-4)\delta_M\le\Delta-4 .
\]
If $\delta_M=0$, every $N_k$ vanishes and at least one $p_v$ is positive,
so $q-2M(G)<0\le\Delta-4$ as well.

Finally, a vertex of full degree $\Delta$ has at least $\Delta-1$
neighbors in $K$. Extend its core star to a spanning tree of $K$. Each
branch has a leaf, so this tree has at least $\Delta-1$ leaves. At most
two lie on a longest path, giving diameter at most $q-\Delta+2$.
Restoring the pendant edges gives a spanning tree of $G$ of diameter
at most $q-\Delta+4$. Its diametral midpoint covers all demands in $G$,
so $2\rho(G)\le q-\Delta+4<2M(G)$. This excludes the terminal case
and proves the inequality and necessity.

For sufficiency, proceed along a construction. The local change formulas
in L and E and the exact sum for C show that $M(G)$ equals the assigned
label $t$ when this holds for the inputs. By $W(G,t)$, $\rho(G)\ge t$.
The lower bound $M(G)\ge\rho(G)$ just proved forces equality. Base labels
are correct by Lemma~\ref{p05-lem:eqtree}.
\end{proof}

\section{The sharp radius bound at fixed maximum degree}\label{p05-sec:broom}
For an integer $\Delta\ge2$ write $c=1/\sqrt2$ and
\[
 B(\Delta)=\sqrt{\Delta}-2+c-(1-c)/\sqrt{\Delta}.
\]
\subsection{Trees and the equality cases}
\begin{theorem}\label{p05-thm:broomtree}
Every finite tree $T$ with maximum degree $\Delta\ge2$ satisfies
$R(T)-\operatorname{rad}(T)\ge B(\Delta)$.
For $\Delta\ge3$, equality holds exactly on a star with $\Delta-1$ arms of length one
and one arm of positive even length. For $\Delta=2$ the equality cases are
the even-order paths of order at least four.
\end{theorem}
\begin{proof}
If a leaf is removed from a vertex of degree $m\ge2$, direct subtraction
shows that $R$ decreases by at least
\[
 m^{-1/2}-(m-1)((m-1)^{-1/2}-m^{-1/2})
   =\sqrt m-\sqrt{m-1}>0,
\]
since the reciprocal square root of each other neighbor's degree is at
most one. In $T$ choose a vertex $v$ of degree $\Delta$ and a diametral path $P$.
Let $K$ be the smallest subtree containing $P$, $v$, and all neighbors of
$v$. Successively pruning leaves outside $K$ preserves its diameter (because
$P$ survives), and hence preserves its radius. It also preserves $\Delta$, since
all arms at $v$ survive. We have $R(T)\ge R(K)$, strictly unless $T=K$.

Suppose $v\in P$. The tree $K$ consists of $\Delta-2$ single-edge arms and two
other arms of lengths $a,b\ge1$. Put $t=\mathbf{1}_{a\ge2}+\mathbf{1}_{b\ge2}$.
Lengthening an arm from one to $j\ge2$ adds
$(j-1)/2+c-1/2-(1-c)/\sqrt{\Delta}$ to $R$. Since
$\operatorname{rad}(K)=\lceil(a+b)/2\rceil$, we obtain
\[
 R(K)-\operatorname{rad}(K)=
 \sqrt{\Delta}-1-\tfrac12\mathbf{1}_{a+b\ {\rm odd}}
  +t\left(c-\tfrac12-(1-c)/\sqrt{\Delta}\right).
\]
For $\Delta>2$ the coefficient of $t$ is positive, is less than $1/2$, and
$a+b$ odd implies $t\ge1$. The lowest value is $B(\Delta)$, exactly when
$a+b$ is odd and $t=1$. For $\Delta=2$ the coefficient vanishes, and an
odd-length path of at least three edges gives precisely equality.

If $v\notin P$, let $p\in P$ be the point at which the path from $v$
reaches $P$. Its distance to each endpoint exceeds or equals
$\operatorname{dist}(p,v)+1\ge2$, since otherwise an endpoint could
be replaced by a neighbor of $v$ not on the connecting path to yield
a longer path. In $K$, $d_p=3$ and its two neighbors on $P$ have
degree two. Relative to $R(P)$ the edges of $P$ thus lose
$2(1/2-1/\sqrt6)$. The $\Delta-1$ leaves at $v$ contribute
$(\Delta-1)/\sqrt{\Delta}$; the remaining edge at $v$ contributes at least
$1/\sqrt{3D}$. Since $R(P)-\operatorname{rad}(P)\ge\sqrt2-3/2$
and the radii are equal, we get
\[
 R(K)-\operatorname{rad}(K)-B(\Delta)
 \ge c-\tfrac12+2/\sqrt6+(1/\sqrt3-c)/\sqrt{\Delta}>0 .
\]
This proves both claims.
\end{proof}


\subsection{Graphs containing a cycle}\label{p05-sec:broomcyclic}
We use the absolute-radius estimate of Theorem~\ref{p05-thm:absolute},
together with $\rho(G)\ge\rad(G)-1/2$, obtained by rounding an absolute
center to a nearest vertex. The tree argument above is independent of
this estimate.

\begin{theorem}\label{p05-thm:broomglobal}
Every connected graph $G$ containing a cycle and having maximum degree
$\Delta\ge2$ satisfies $R(G)-\rad(G)>B(\Delta)$. Consequently the minimum over
all connected graphs of fixed maximum degree $\Delta$ is $B(\Delta)$, with precisely
the equality cases of Theorem~\ref{p05-thm:broomtree}.
\end{theorem}
\begin{proof}
Write $r=\rad(G)$, $n=|V(G)|$, and $L$ for the number of leaves. A cycle
has $R=n/2$ and $r=\lfloor n/2\rfloor$, so the claim holds for $\Delta=2$.
Suppose a cyclic counterexample to strictness exists, and choose one
of smallest order.

No vertex $v$ of degree $m$ supports two leaves. Otherwise delete one,
$w$, to form $G'$. Any vertex other than the surviving twin leaf has the
same distance to the two leaves. The twin itself cannot be a center of
$G'$ since moving to its support strictly decreases its eccentricity
in a connected graph containing a cycle. Thus $\rad(G')=r$.
Put $h=(m-1)^{-1/2}-m^{-1/2}>0$. Direct subtraction gives
\[
 R(G)-R(G')=\sqrt m-\sqrt{m-1}
  +h\sum_{u\in N(v)\setminus\{w\}}(1-d_u^{-1/2}).
\]
One of these neighbors is a nonleaf since $G$ contains a cycle, so
the sum is at least $1-c$, where $c=1/\sqrt2$. The gain is at least
$B(m)-B(m-1)>0$. If $\Delta$ falls upon deleting $w$, then $m=\Delta$ and the
new maximum degree is $\Delta-1$. Otherwise the maximum degree stays $\Delta$.
In both cases strictness for the smaller graph implies strictness for $G$.

Write $s=n-2r-\Delta+2$. We claim that
\begin{equation}\label{p05-eq:broom-count}
 s\ge\max\{0,L-3\}.
\end{equation}
Choose any vertex $v$ of degree $\Delta$, extend its full star to a spanning tree
$S$, and take a diametral path $Q$ of $S$. As $\rad(S)\ge r$ and
$\rad(S)=\lceil\diam(S)/2\rceil$, the path has at least $2r$ vertices.
At most two neighbors of $v$ in $S$ lie on $Q$; if $v\notin Q$, at most
one does. At least $\Delta-2$ neighbors of $v$ and at least $L-2$ original
leaves thus lie outside $Q$. The two sets overlap in at most one vertex
because $v$ supports at most one original leaf. Counting their union
among the at most $n-2r$ vertices outside $Q$ proves~\eqref{p05-eq:broom-count}.

Set $y=\sqrt{\Delta}$, $z=1/y$, $t=cy>1$, and $p=(1-c)(1-z)$.
For an edge write $a\ge b$ for the reciprocal square roots of its endpoint
degrees. Then $R-M=\sum_{E}b(a-b)$ and $n-2R=\sum_{E}(a-b)^2$, since
$\sum_{E}(a^2+b^2)=n$. Consequently
\[
 (t+1)(R-r)= n-2r+(t-1)(M-r)+\sum_{E}(a-b)(tb-a).
\]
For a nonleaf edge $a\le c$ and $b\ge z$, so its last summand is
nonnegative. At a leaf $a=1$; if $tb<1$, we have $1-b\le1-z$
and $1-tb\le1-c$. Thus every leaf edge contributes at least $-p$,
with equality only if $b=z$. From $M\ge r-1/2$ it follows that
\begin{align}
 (t+1)(R-r-B(\Delta))&\ge E(y)+s-(L-3)p,\label{p05-eq:broom-excess}\\
 E(y)&=\Delta-2-(t-1)/2-3p-(t+1)B(\Delta)\notag\\
 &= (1-c)(y-2)(y^2+y/2-2)/y.\notag
\end{align}
Here $0<p<1$, so the correction $s-(L-3)p$ is nonnegative by
\eqref{p05-eq:broom-count}, and is positive unless $L=3$ and $s=0$.
For $\Delta\ge5$, $E(y)>0$. For $\Delta=4$, $E(y)=0$. Equality at $\Delta=4$ would
therefore require exactly three leaves, every leaf adjacent to a
degree-four vertex, and every nonleaf edge satisfying $a=b$ or $a=tb$.
In the latter case the only unequal degree pair is $(2,4)$ since $t=\sqrt2$.
A degree-three vertex could only meet degree-three vertices; connectedness
and the existence of a leaf exclude such a vertex. This leaves exactly
three vertices of odd degree, contrary to the handshaking lemma.

It remains to consider $\Delta=3$. At $y=\sqrt3$ one has
$E(y)=-(1-c)/(2\sqrt3)$. If $L\le2$, the correction in
\eqref{p05-eq:broom-excess} is at least $p$ and
$E(y)+p=(1-c)(1-\sqrt3/2)>0$. If $L\ge4$, the correction is at least
$1-p$ and $E(y)+1-p=c+(1-c)/(2\sqrt3)>0$. When $L=3$ and $s\ge1$,
the latter estimate also suffices.

In the sole remaining case $L=3$ and $n-2r=1$. For each vertex $v$ of
degree three, take the spanning tree and path used in~\eqref{p05-eq:broom-count}.
At most one vertex lies outside $Q$. Both a neighbor of $v$ and an
original leaf lie outside $Q$, so they coincide. Hence every vertex of
degree three supports one leaf. If $b_3$ is the number of degree-three
vertices, these distinct supported leaves imply
$b_3\le3$. Since $G$ contains a cycle, $|E(G)|\ge n$, and handshaking
gives $2(|E(G)|-n)=b_3-3$. Hence $b_3=3$ and every leaf is adjacent to
one of these three vertices. After removing all leaves, the graph is
connected and every vertex has degree two, so it is a cycle of length
$q=n-3=2r-2$. Each support has at most one antipodal cycle vertex.
Since $q$ is even and at least four, some cycle vertex is antipodal to
none of the three supports. Its distance to every support is at most
$q/2-1$, so its distance to every vertex of $G$ is at most $q/2=r-1$.
This contradicts the definition of $r$.
\end{proof}

\begin{thebibliography}{12}
\bibitem{p05-AI} A. Abdolghafourian and M. A. Iranmanesh.
On harmonic index and diameter of quasi-tree graphs.
\emph{Journal of Mathematics} (2021), article 6650407.
\href{https://doi.org/10.1155/2021/6650407}{Publisher}.
\bibitem{p05-CPS} M. Cygan, M. Pilipczuk and R. \v{S}krekovski.
On the inequality between radius and Randi\'c index for graphs.
\emph{MATCH Commun. Math. Comput. Chem.} \textbf{67} (2012), 451--466.
\href{https://match.pmf.kg.ac.rs/electronic_versions/Match67/n2/match67n2_451-466.pdf}{Journal PDF}.
\bibitem{p05-CPSMean} M. Cygan, M. Pilipczuk and R. \v{S}krekovski.
Relation between Randi\'c index and average distance of trees.
\emph{MATCH Commun. Math. Comput. Chem.} \textbf{66} (2011), 605--612.
\bibitem{p05-Dal} C. Dalf\'o. On the Randi\'c index of graphs.
\emph{Discrete Mathematics} \textbf{342} (2019), 2792--2796.
\href{https://doi.org/10.1016/j.disc.2018.08.020}{Publisher}.
\bibitem{p05-DVB} H. Deng, T. Vetr\'ik and S. Balachandran.
On the harmonic index and diameter of unicyclic graphs.
\emph{Mathematical Reports} \textbf{22(72)} (2020), 11--18.
\href{https://imar.ro/journals/Mathematical_Reports/Pdfs/2020/1/2.pdf}{Journal PDF}.
\bibitem{p05-DLS} Z. Dvo\v{r}\'ak, B. Lidick\'y and R. \v{S}krekovski.
Randi\'c index and the diameter of a graph.
\emph{European Journal of Combinatorics} \textbf{32} (2011), 434--442.
\href{https://doi.org/10.1016/j.ejc.2010.12.002}{Publisher}.
\bibitem{p05-Faj} S. Fajtlowicz. On conjectures of Graffiti.
\emph{Discrete Mathematics} \textbf{72} (1988), 113--118.
\href{https://doi.org/10.1016/0012-365X(88)90199-9}{Publisher}.
\bibitem{p05-JM} J. A. Jerline and L. B. Michaelraj.
On a conjecture of harmonic index and diameter of graphs.
\emph{Kragujevac Journal of Mathematics} \textbf{40} (2016), 73--78.
\href{https://imi.pmf.kg.ac.rs/kjm/pub/kjom40/kjm_40_1-6.pdf}{Journal PDF}.
Also, On harmonic index and diameter of unicyclic graphs.
\emph{Iranian Journal of Mathematical Sciences and Informatics} \textbf{11}
(2016), 115--122.
\bibitem{p05-KS} Keio and Shaitetsu.
A proof of the Caporossi--Hansen conjecture on the Randi\'c index and
graph radius. Research manuscript, 13 November 2026.
\bibitem{p05-Liu} J. Liu. On harmonic index and diameter of graphs.
\emph{Journal of Applied Mathematics and Physics} \textbf{1}(3) (2013), 5--6.
\href{https://doi.org/10.4236/jamp.2013.13002}{Publisher}.
\bibitem{p05-YL} L. You and B. Liu.
On a conjecture between Randi\'c index and average distance of unicyclic graphs.
\emph{Filomat} \textbf{28} (2014), 767--773.
\href{https://doi.org/10.2298/FIL1404767Y}{Publisher}.
\end{thebibliography}

% END SOURCE BODY Graph_Distance_and_Degree_Indices
\unitstart{Paper 06}{p06}{Spanning trees}{Spanning trees and Fajtlowicz's average-distance conjecture}{}
\makeatletter\@removefromreset{theorem}{section}\@addtoreset{theorem}{section}\makeatother\renewcommand{\thetheorem}{\thesection.\arabic{theorem}}
% BEGIN SOURCE BODY Fajtlowicz_Average_Distance
\begin{abstract}
For the conjecture that the Randi\'c index of a connected graph bounds its
average distance, we prove strict inequality for graphs of order at least
three in each of the following classes: maximum degree at most fifteen,
circuit rank at most seventeen, and graphs admitting a shortest path
within distance four of every vertex. The circuit-rank argument also
gives an explicit order threshold for each fixed rank. The proofs compare
edge-degree contributions with the leaf count of a spanning tree, or
with the cost of projecting the vertex demands onto a shortest path.
All finite parameter checks are included in the text. The unrestricted
conjecture remains unresolved by these arguments.
\end{abstract}
\section{The conjecture and the estimates}
Let $G$ be a finite connected undirected simple graph of order $n\ge2$.
We write $d_v$ for the degree of $v$, $\Delta$ for maximum degree and $D$
for diameter. Fajtlowicz's conjecture \cite{p06-Faj} is $R(G)\ge\mu(G)$, where
\[
 R(G)=\sum_{uv\in E(G)}\frac1{\sqrt{d_ud_v}},
 \qquad \mu(G)=\binom n2^{-1}\sum_{\{u,v\}\subseteq V(G)}d_G(u,v).
\]
The equality $R=\mu=1$ holds on $K_2$. Earlier results include the tree
case \cite{p06-CPSMean}, unicyclic graphs \cite{p06-YL} and a minimum-degree
condition for graphs of sufficiently large order \cite{p06-Dal}. We first
bound $\mu$ using a maximum-leaf spanning tree and bound $R$ by summing
local inequalities on degree pairs. We then give estimates independent
of a maximum-degree restriction, using the number of cycle edges or the
distribution of vertices around a geodesic. These results concern
the original index $R$.

For use below recall the elementary estimate $\mu\le(n+1)/3$. Pass to
a spanning tree and delete a leaf. Its sum of distances to the other
$n-1$ vertices is at most $1+\cdots+(n-1)$, since the $j$th smallest
distance is at most $j$. Induction bounds the sum over all pairs by
$\binom{n+1}{3}$. A second elementary estimate is
\begin{equation}\label{p06-eq:basic}
 R\ge \frac{n\sqrt{\Delta}}{\Delta+1}\ge\sqrt{n-1}.
\end{equation}
In fact the ratio of the edge contribution to the incidence sum
$1/i+1/j$ at degrees $i\le j$ equals $\sqrt{j/i}/(1+j/i)$. This is
decreasing for $j/i\ge1$ and $j/i\le\Delta$. Summing proves the
first bound. For the second, use $\Delta\le n-1$ and the same
monotonicity. The star lower bound is classical \cite{p06-BE}.

\section{A leaf bound for maximum degree at most fifteen}\label{p06-sec:leaf15}
For a connected graph of order $n\ge3$, write $L$ for the number of
degree-one vertices and $s$ for the largest number of leaves in any of
its spanning trees. Let $J$ denote the subgraph induced by the vertices
of degree at least three and put $e_J=|E(J)|$.
\begin{lemma}\label{p06-lem:leaf15}
Writing $m=|E(G)|$, one has
\[
  s\ge\Delta,\qquad s\ge L,\qquad s\ge L+m-n+1-e_J .
\]
\end{lemma}
\begin{proof}
Extend the incident edges of a maximum-degree vertex to a spanning tree.
Each of its $\Delta$ branches contains a leaf. Every spanning tree
retains all original leaves. In any spanning tree, at most $e_J$ of
the $m-n+1$ omitted edges have both endpoints of degree at least three.
Each other omitted edge has a degree-two endpoint, since pendant edges
cannot be omitted. This endpoint is a new leaf of the tree. Distinct
omitted edges cannot share such a degree-two endpoint.
\end{proof}

\begin{theorem}\label{p06-thm:leaf15}
Every connected simple graph of maximum degree at most fifteen and order
at least two satisfies $R(G)\ge\mu(G)$, with equality only for $K_2$.
\end{theorem}
\begin{proof}
The case $K_2$ is immediate. Put $S=100000$. For $2\le\Delta\le14$
use the first row, and for $\Delta=15$ use the second:
\[
\begin{array}{c|rrrrr}
\Delta&d&s_0&K&B&C\\ \hline
2,\ldots,14&14&2&35235&8348&2886\\
15&15&15&34405&8211&2862
\end{array}
\]
For the endpoint degrees $i\le j$ of an edge set $w=1/i+1/j$,
$p=\mathbf{1}_{i=1}$ and $u=\mathbf{1}_{i\ge3}$.
For either row, the following holds whenever $1\le i\le j\le d$ and
$j\ge2$:
\begin{equation}\label{p06-eq:leaf15dual}
 \frac{S}{\sqrt{ij}}\ \ge\ K w-Bp-C(p+1-w-u).
\end{equation}
An exact integer check is included below. These coefficients give
positive weight to two separate bounds on $s$: the original pendant
edges and the omitted edges that create extra leaves. Summing
\eqref{p06-eq:leaf15dual} and applying Lemma~\ref{p06-lem:leaf15} gives
\[
  S R\ge Kn-BL-C(L+m-n-e_J)\ge Kn-(B+C)s+C.
\]

Let $t=n-s-1$. Distances in $G$ are bounded by distances in a
spanning tree with $s$ leaves. Among such trees the double comet with
balanced end clusters maximizes the Wiener index
\cite[Lemma~2.4 and the observation following it]{p06-CPSMean}. This gives
\begin{equation}\label{p06-eq:leaf15double}
 \mu\ \le\ U(n,s)=\frac{12s(n-1)+3tn^2-t(t^2-1)}{6n(n-1)}.
\end{equation}
Indeed the pendant edges contribute $s(n-1)$ to the Wiener index.
Writing $a=\lfloor s/2\rfloor$, the $t$ internal edges contribute
$\sum_{j=1}^t(a+j)(n-a-j)$. This sum equals
$tn^2/4-t(t^2-1)/12$ for even $s$, with a further subtraction of
$t/4$ for odd $s$. The same sum applies to the path when $s=2$. When $s=n-1$, the spanning tree is a star and
\eqref{p06-eq:leaf15double} holds directly.

Set $x=s/n$ and $F(x)=1/3-x^2/2+x^3/6$. Expansion gives
\[
 U(n,s)-nF(x)
 =\frac{n(x^3+6x+2)-10x-2}{6(n-1)}\le\tfrac32 .
\]
The numerator increases on $[0,1]$ for $n\ge3$ and substitution at
$x=1$ suffices. Write $k=K/S$, $b=(B+C)/S$. For $0\le x\le1/2$,
\[
  k-bx-F(x)\ \ge\ k-\tfrac13-bx+\tfrac5{12}x^2
        \ \ge\ k-\tfrac13-\tfrac35 b^2>\tfrac1{300}.
\]
For $1/2\le x\le1$ use $F(x)\le1/3-x^2/3$. The minimum of
$x^2/3-bx$ on that interval is $1/12-b/2$, again giving a gap
greater than $1/300$ for both rows. This proves strictness for
$n\ge450$.


For small orders we also use \eqref{p06-eq:basic}. The integer checks below
verify the edge inequalities and show that at every admissible triple
$(\Delta,n,s)$ with $n\le450$, at least one of the two certified bounds
for $R$ strictly exceeds $U(n,s)$. Write $Q=\lfloor S\sqrt{\Delta}\rfloor$,
so $nQ/[S(\Delta+1)]$ is a rational lower bound. Both loops run over
integers and use no enumeration of graphs.
{\small
\begin{verbatim}
from math import isqrt
S = 100000
for lo,d,K,B,C in [(2,14,35235,8348,2886),(15,15,34405,8211,2862)]:
    for j in range(2,d+1):
        for i in range(1,j+1):
            p,u = int(i==1),int(i>=3)
            A = (K+C)*(i+j) - i*j*(B*p+C*(p+1-u))
            assert A <= 0 or A*A <= S*S*i*j
    for h in range(lo,d+1):
        Q = isqrt(S*S*h)
        for n in range(h+1,451):
            for s in range(h,n):
                t = n-s-1
                P = 12*s*(n-1)+3*t*n*n-t*(t*t-1)
                V = max((h+1)*(K*n-(B+C)*s+C), n*Q)
                assert 6*n*(n-1)*V > S*(h+1)*P
\end{verbatim}
}
\end{proof}

\begin{corollary}\label{p06-cor:leaf15local}
Suppose every edge whose lower endpoint degree is at most two has upper
endpoint degree at most fifteen. Suppose also that whenever both endpoint
degrees are at least three, the larger is at most five times the smaller.
Then $R(G)>\mu(G)$ whenever $n>2$.
\end{corollary}
\begin{proof}
It suffices to consider $\Delta\ge16$. Use the second row in the preceding
proof, so $s\ge\Delta\ge16$. For an edge with $i\ge3$ its local bound
\eqref{p06-eq:leaf15dual} becomes $1/\sqrt{ij}\ge (K+C)w/S$.
The ratio of the left-hand side to $w$ is at least $\sqrt5/6$ by
$j/i\le5$, and $37267/100000<\sqrt5/6$. Edges with $i\le2$ are covered
by the same finite table. The rest of the argument depends only on
$s\ge15$ and $s<n$.
\end{proof}
More generally, divide the right-hand side of \eqref{p06-eq:leaf15dual}
for the second row by $S$ and denote it by $\phi(i,j)$. Put
\[
 E_* = \sum_{uv\in E}\max\{0,\phi(d_u,d_v)-1/\sqrt{d_ud_v}\},
\]
sorting the degrees before evaluating $\phi$. For every connected graph
with $s\ge15$ the proof gives the sufficient condition
\[
 E_* < \frac{Kn-(B+C)s+C}{S}-U(n,s)
       \quad\Longrightarrow\quad R(G)>\mu(G).
\]
For $n\ge450$ it suffices that $E_*\le n/300-3/2$.

\section{Cost of cycle edges}\label{p06-sec:cycleloss}
Write $b=|E(G)|-n+1$ for the circuit rank. On deleting a nonbridge
$uv$ with degrees $d\le e$, put $x=1/\sqrt d$, $y=1/\sqrt e$.
The net increase in $R$ is $A_u+A_v-xy$, where
\[
 A_u=((d-1)^{-1/2}-x)\sum_{w\sim u,\ w\ne v}d_w^{-1/2}.
\]
In particular $A_u\le x/2$ since
$\sqrt{d(d-1)}\le d-1/2$; likewise $A_v\le y/2$.
Since $uv$ belongs to a cycle, a neighbor of $u$ other than $v$ has
degree at least two. Hence
\[
 A_u+A_v-xy\le f(d)+x\max(0,1/2-x),\quad
 f(d)=(d-2+1/\sqrt2)((d-1)^{-1/2}-d^{-1/2}).
\]
If $d\ge9$ the right side is at most $x-x^2\le 2/9$. For $d=2,\ldots,8$,
successive strict upper bounds in thousandths are
$208,222,210,220,221,219,216$, respectively, each below $2000/9$.
They follow by bounding $j^{-1/2}$ between
$10^{-6}\lfloor\sqrt{10^{12}/j}\rfloor$ and that quantity plus $10^{-6}$.
Every nonbridge deletion therefore increases $R$ by at most $2/9$.

Fix a spanning tree $T$ with the largest possible number $s$ of leaves
and delete the $b$ complementary edges, keeping $T$ at each step.
Since distances increase under edge deletion, $R(G)-\mu(G)\ge
R(T)-\mu(T)-2b/9$. By \cite[Theorem~1.2]{p06-CPSMean} we obtain
\[
 R(G)-\mu(G)\ge\sqrt{s}-2-2b/9.
\]
For $3\le s\le n-2$ the stronger estimate
$R(T)\ge(n-s)/2+\sqrt{s}-0.462$ is
\cite[Corollary~2.3]{p06-CPSMean}. Using $\mu(T)\le nF(s/n)+3/2$ and
$F(x)=(1-x)/2-(1-x)^3/6$ gives
\[
 R(G)-\mu(G)>\frac{(n-s)^3}{6n^2}+\sqrt{s}-2-2b/9.
\]
When $s=n-1$, $G$ has a universal vertex, so $\mu(G)<2$ and
$R(G)\ge\sqrt{n-1}$ by \eqref{p06-eq:basic}. When $s=2$, $\Delta\le2$. Consequently,
for every fixed $b$, with $h=2+2b/9$, the conjecture holds if
$n\ge\max\{48h,2h^2\}$. Indeed if $s\ge n/2$ the square-root term
is at least $\sqrt{n/2}$, and otherwise the cubic term exceeds $n/48$.

\begin{corollary}
Every connected graph of order at least three satisfies $R>\mu$ if its circuit rank is at most
seventeen.
\end{corollary}
\begin{proof}
We give the finite check for $s\ge13$, since smaller values imply
$\Delta\le12$ and are covered by Theorem~\ref{p06-thm:leaf15}.
The preceding estimates cover $n\ge278$. Exclude $s=n-1$ as above.
For $Q_j=\lfloor 10^5\sqrt{j}\rfloor/10^5$, put
\[
 A=Q_{2s}(n-s-1)/(s+2)+Q_s,\quad
 B=Q_{n-1},\quad C=(n-s)/2+Q_s-462/1000-34/9 .
\]
Here $A\le R$ since $\Delta\le s$: on edges of lower degree at least two,
$R_{uv}\ge k(1/d_u+1/d_v)$ for $k=\sqrt{2s}/(s+2)$.
The deficit on each pendant edge is at most $k(1+1/s)-1/\sqrt{s}$,
by monotonicity in the support degree. It is positive for $s\ge13$,
and $L\le s$. Summation gives $R\ge k(n-s-1)+\sqrt{s}$.
Also $B\le R$ by the star lower bound, and $C\le R$ by deleting cycle
edges as above. Thus a sufficient lower bound for $R-\mu$ is
\[
 g(n,s)=\max\{A-U(n,s),B-U(n,s),C-U(n,s),Q_s-2-34/9\}.
\]
Computing its minimum over the $34\,716$ pairs
$14\le n\le277$, $13\le s\le n-2$ in rational arithmetic
gives $327879/3700000>0$ at $(n,s)=(38,19)$. The calculation is:
{\small
\begin{verbatim}
from fractions import Fraction as Q
from math import isqrt
q = lambda j: Q(isqrt(j*10**10),10**5)
best = None
for n in range(14,278):
    for s in range(13,n-1):
        t=n-s-1
        U=Q(12*s*(n-1)+3*t*n*n-t*(t*t-1),6*n*(n-1))
        A=q(2*s)*(n-s-1)/(s+2)+q(s)
        C=Q(n-s,2)+q(s)-Q(462,1000)-Q(34,9)
        g=max(A-U,q(n-1)-U,C-U,q(s)-2-Q(34,9))
        value=(g,n,s)
        if best is None or value<best: best=value
assert best==(Q(327879,3700000),38,19)
\end{verbatim}
}
\end{proof}

\section{Projection onto a shortest path}
The modified index $M=\sum_{uv}1/\max(d_u,d_v)$ satisfies $M\ge D/2$
by Dvo\v{r}\'ak, Lidick\'y and \v{S}krekovski \cite[Theorem~3]{p06-DLS}.
Writing $a\ge b$ for the reciprocal square roots of the two endpoint
degrees, one has $(a-b)^2\le(\sqrt{\Delta}-1)b(a-b)$. The sums of the two
terms $(a-b)^2$ and $b(a-b)$ are $n-2R$ and $R-M$, respectively.
Thus $(\sqrt{\Delta}+1)R\ge n+(\sqrt{\Delta}-1)M$.
For a shortest path $P$ of length $d$, put $t=n-d$ and
$k=\max_{v\in V(G)}\dist(v,P)$. We allow $d$ to be smaller than the
diameter $D$. If $k\ge1$, then
\begin{equation}\label{p06-eq:faj-pathprojection}
 \mu(G)\le \frac{dn}{2(n-1)}+\frac{2k(t-1)}n<\frac d2+2k.
\end{equation}
Indeed, project each vertex to a nearest vertex of $P$. The distance
sum between pairs of projected vertices is at most $dn^2/4$, since
each of the $d$ cuts of the path separates at most $n^2/4$ unordered
pairs. The total cost of going to and from the path, divided by
$\binom n2$, is at most $2k(t-1)/n$. The strict inequality follows
from $d/[2(n-1)]<2k(d+1)/n$.

Every vertex on $P$ has at most two neighbors on $P$, and an outside
vertex meets at most three consecutive positions. Thus $\Delta\le t+1$.
The edge-energy bound $(\sqrt\Delta+1)R\ge n+(\sqrt\Delta-1)M$
and the diameter bound $M\ge D/2$ imply
\[
 R\ge\frac d2+\frac{t}{\sqrt\Delta+1}
 \ge\frac d2+\sqrt{t+1}-1.
\]
Consequently $R>\mu$ whenever $t\ge4k(k+1)$. Also,
$n\ge3t+2$ suffices, since $d/2\ge(n+1)/3\ge\mu$.
For $k=1$ this proves the conclusion for $t\ge8$, and the smaller
values give $\Delta\le8$. In particular, every connected AT-free
graph satisfies the conjecture: Corneil, Olariu and Stewart proved
that such a graph contains a dominating pair, and a shortest path between
its endpoints dominates the graph. See \cite{p06-COS}.

\begin{proposition}
If a connected graph of order at least three admits a shortest path
within distance four of every vertex, then $R>\mu$.
\end{proposition}
\begin{proof}
The case $k=0$ is the path, covered by Theorem~\ref{p06-thm:leaf15}. We may use $k=4$ as an upper bound throughout.
The preceding estimates leave $t<80$, $n<3t+2$ and $\Delta\ge13$,
where the last restriction uses the bounded-degree theorem. There are
268158 integer triples $(t,n,\Delta)$ in this remaining range. The
following reproducible integer calculation proves strictness throughout.
Using $d=n-t$, put
\[
 U=\min\left\{\frac{n+1}3,\ \frac{dn}{2(n-1)}+\frac{8(t-1)}n,
 \ \frac{n+1}3-\frac{\Delta(\Delta-2)(3n-5-\Delta)}{6n(n-1)}\right\}.
\]
The last term follows by expanding $U(n,s)$ in \eqref{p06-eq:leaf15double} and extending the star at a maximum-degree vertex
to a tree, which has at least $\Delta$ leaves, and applying the fixed-leaf
upper bound. Its correction is increasing in the leaf count up to $n-1$
(its derivative has the sign of $3s(2n-2-s)-6n+10>0$ for $2\le s\le n-1$). In each iteration put
$u=\min(U,D)$ and $v=u-D/2$. The integer tests give
$u<\sqrt{n-1}$ or $u<D/2+(n-D)/(\sqrt\Delta+1)$,
using $n>D$. Both are sufficient by the classical $R\ge\sqrt{n-1}$
and the degree-energy bound. Here all rational arithmetic is exact.
{\small
\begin{verbatim}
from fractions import Fraction as F
for t in range(1,80):
    for n in range(t+1,3*t+2):
        d = n-t
        for h in range(13,min(t+1,n-1)+1):
            w = F(n+1,3)
            U = min(w, F(d*n,2*(n-1))+F(8*(t-1),n),
                    w-F(h*(h-2)*(3*n-5-h),6*n*(n-1)))
            for D in range(d,min(n-1,d+8)+1):
                u = min(U,D)
                v = u-F(D,2)
                A = n-D-v
                assert u*u < n-1 or v <= 0 or (A>0 and A*A>h*v*v)
\end{verbatim}
}
\end{proof}

\begin{thebibliography}{9}
\bibitem{p06-BE} B. Bollob\'as and P. Erd\H{o}s. Graphs of extremal weights.
\emph{Ars Combinatoria} \textbf{50} (1998), 225--233.
\bibitem{p06-COS} D. G. Corneil, S. Olariu and L. Stewart.
Asteroidal triple-free graphs. \emph{SIAM J. Discrete Math.}
\textbf{10} (1997), 399--430.
\href{https://doi.org/10.1137/S0895480193250125}{Publisher}.
\bibitem{p06-CPSMean} M. Cygan, M. Pilipczuk and R. \v{S}krekovski.
Relation between Randi\'c index and average distance of trees.
\emph{MATCH Commun. Math. Comput. Chem.} \textbf{66} (2011), 605--612.
\href{https://match.pmf.kg.ac.rs/electronic_versions/Match66/n2/match66n2_605-612.pdf}{Journal PDF}.
\bibitem{p06-Dal} C. Dalf\'o. On the Randi\'c index of graphs.
\emph{Discrete Mathematics} \textbf{342} (2019), 2792--2796.
\href{https://doi.org/10.1016/j.disc.2018.08.020}{Publisher}.
\bibitem{p06-DLS} Z. Dvo\v{r}\'ak, B. Lidick\'y and R. \v{S}krekovski.
Randi\'c index and the diameter of a graph.
\emph{European Journal of Combinatorics} \textbf{32} (2011), 434--442.
\href{https://doi.org/10.1016/j.ejc.2010.12.002}{Publisher}.
\bibitem{p06-Faj} S. Fajtlowicz. On conjectures of Graffiti.
\emph{Discrete Mathematics} \textbf{72} (1988), 113--118.
\href{https://doi.org/10.1016/0012-365X(88)90199-9}{Publisher}.
\bibitem{p06-YL} L. You and B. Liu. On a conjecture between Randi\'c index
and average distance of unicyclic graphs. \emph{Filomat} \textbf{28}
(2014), 767--773. \href{https://doi.org/10.2298/FIL1404767Y}{Publisher}.
\end{thebibliography}

% END SOURCE BODY Fajtlowicz_Average_Distance
\unitstart{Paper 07}{p07}{Threshold forests}{Threshold forests and Fajtlowicz's average-distance conjecture}{}
\makeatletter\@removefromreset{theorem}{section}\@addtoreset{theorem}{section}\makeatother\renewcommand{\thetheorem}{\thesection.\arabic{theorem}}
% BEGIN SOURCE BODY Fajtlowicz_Threshold_Forests
\begin{abstract}
For a finite connected simple graph, write $R$ for the Randi\'c index and
$\mu$ for mean distance over distinct unordered vertex pairs. We prove $R>\mu$
for graphs of order at least three when any one of the following holds:
the maximum degree is at most 100, the circuit rank is at most 155, or some
geodesic lies within distance 18 of every vertex. The first estimate counts
leaves that can be retained when joining components induced by high-degree
vertices. The second combines that estimate with circuit rank, and the
third uses the distances inside and outside a neighborhood of maximum order.
For block graphs and further structural classes we establish the stronger
modified-index inequality $M\ge\mu-1+2/n$. A fractional packing of cycle
shortcuts gives a global lower bound on $R-\mu$ whose sign remains open
in the unrestricted case. Exact integer certificates are included.
\end{abstract}
\section{The conjecture and the estimates}
Let $G$ be a finite connected undirected simple graph of order $n\ge2$.
We write $d_v$ for the degree of $v$, $\Delta$ for maximum degree and $D$
for diameter. Fajtlowicz's conjecture \cite{p07-Faj} is $R(G)\ge\mu(G)$, where
\[
 R(G)=\sum_{uv\in E(G)}\frac1{\sqrt{d_ud_v}},
 \qquad \mu(G)=\binom n2^{-1}\sum_{\{u,v\}\subseteq V(G)}d_G(u,v).
\]
The equality $R=\mu=1$ holds on $K_2$. Earlier results include the tree
case \cite{p07-CPSMean}, unicyclic graphs \cite{p07-YL} and a minimum-degree
condition for graphs of sufficiently large order \cite{p07-Dal}. The main quantitative results are Theorems~\ref{p07-thm:degree100},
\ref{p07-thm:cycl155} and \ref{p07-thm:widthplus}. Each applies independently to
graphs of arbitrary order. We first
bound $\mu$ using a maximum-leaf spanning tree and bound $R$ by summing
local inequalities on degree pairs. We then give estimates independent
of a maximum-degree restriction, using the number of cycle edges or the
distribution of vertices around a geodesic. These results concern
the original index $R$.

For use below recall the elementary estimate $\mu\le(n+1)/3$. Pass to
a spanning tree and delete a leaf. Its sum of distances to the other
$n-1$ vertices is at most $1+\cdots+(n-1)$, since the $j$th smallest
distance is at most $j$. Induction bounds the sum over all pairs by
$\binom{n+1}{3}$. A second elementary estimate is
\begin{equation}\label{p07-eq:basic}
 R\ge \frac{n\sqrt{\Delta}}{\Delta+1}\ge\sqrt{n-1}.
\end{equation}
In fact the ratio of the edge contribution to the incidence sum
$1/i+1/j$ at degrees $i\le j$ equals $\sqrt{j/i}/(1+j/i)$. This is
decreasing for $j/i\ge1$ and $j/i\le\Delta$. Summing proves the
first bound. For the second, use $\Delta\le n-1$ and the same
monotonicity. The star lower bound is classical \cite{p07-BE}.

\section{A leaf bound for maximum degree at most fifteen}\label{p07-sec:leaf15}
For a connected graph of order $n\ge3$, write $L$ for the number of
degree-one vertices and $s$ for the largest number of leaves in any of
its spanning trees. Let $J$ denote the subgraph induced by the vertices
of degree at least three and put $e_J=|E(J)|$.
\begin{lemma}\label{p07-lem:leaf15}
Writing $m=|E(G)|$, one has
\[
  s\ge\Delta,\qquad s\ge L,\qquad s\ge L+m-n+1-e_J .
\]
\end{lemma}
\begin{proof}
Extend the incident edges of a maximum-degree vertex to a spanning tree.
Each of its $\Delta$ branches contains a leaf. Every spanning tree
retains all original leaves. In any spanning tree, at most $e_J$ of
the $m-n+1$ omitted edges have both endpoints of degree at least three.
Each other omitted edge has a degree-two endpoint, since pendant edges
cannot be omitted. This endpoint is a new leaf of the tree. Distinct
omitted edges cannot share such a degree-two endpoint.
\end{proof}

\begin{theorem}\label{p07-thm:leaf15}
Every connected simple graph of maximum degree at most fifteen and order
at least two satisfies $R(G)\ge\mu(G)$, with equality only for $K_2$.
\end{theorem}
\begin{proof}
The case $K_2$ is immediate. Put $S=100000$. For $2\le\Delta\le14$
use the first row, and for $\Delta=15$ use the second:
\[
\begin{array}{c|rrrrr}
\Delta&d&s_0&K&B&C\\ \hline
2,\ldots,14&14&2&35235&8348&2886\\
15&15&15&34405&8211&2862
\end{array}
\]
For the endpoint degrees $i\le j$ of an edge set $w=1/i+1/j$,
$p=\mathbf{1}_{i=1}$ and $u=\mathbf{1}_{i\ge3}$.
For either row, the following holds whenever $1\le i\le j\le d$ and
$j\ge2$:
\begin{equation}\label{p07-eq:leaf15dual}
 \frac{S}{\sqrt{ij}}\ \ge\ K w-Bp-C(p+1-w-u).
\end{equation}
An exact integer check is included below. These coefficients give
positive weight to two separate bounds on $s$: the original pendant
edges and the omitted edges that create extra leaves. Summing
\eqref{p07-eq:leaf15dual} and applying Lemma~\ref{p07-lem:leaf15} gives
\[
  S R\ge Kn-BL-C(L+m-n-e_J)\ge Kn-(B+C)s+C.
\]

Let $t=n-s-1$. Distances in $G$ are bounded by distances in a
spanning tree with $s$ leaves. Among such trees the double comet with
balanced end clusters maximizes the Wiener index
\cite[Lemma~2.4 and the observation following it]{p07-CPSMean}. This gives
\begin{equation}\label{p07-eq:leaf15double}
 \mu\ \le\ U(n,s)=\frac{12s(n-1)+3tn^2-t(t^2-1)}{6n(n-1)}.
\end{equation}
Indeed the pendant edges contribute $s(n-1)$ to the Wiener index.
Writing $a=\lfloor s/2\rfloor$, the $t$ internal edges contribute
$\sum_{j=1}^t(a+j)(n-a-j)$. This sum equals
$tn^2/4-t(t^2-1)/12$ for even $s$, with a further subtraction of
$t/4$ for odd $s$. The same sum applies to the path when $s=2$. When $s=n-1$, the spanning tree is a star and
\eqref{p07-eq:leaf15double} holds directly.

Set $x=s/n$ and $F(x)=1/3-x^2/2+x^3/6$. Expansion gives
\[
 U(n,s)-nF(x)
 =\frac{n(x^3+6x+2)-10x-2}{6(n-1)}\le\tfrac32 .
\]
The numerator increases on $[0,1]$ for $n\ge3$ and substitution at
$x=1$ suffices. Write $k=K/S$, $b=(B+C)/S$. For $0\le x\le1/2$,
\[
  k-bx-F(x)\ \ge\ k-\tfrac13-bx+\tfrac5{12}x^2
        \ \ge\ k-\tfrac13-\tfrac35 b^2>\tfrac1{300}.
\]
For $1/2\le x\le1$ use $F(x)\le1/3-x^2/3$. The minimum of
$x^2/3-bx$ on that interval is $1/12-b/2$, again giving a gap
greater than $1/300$ for both rows. This proves strictness for
$n\ge450$.


For small orders we also use \eqref{p07-eq:basic}. The integer checks below
verify the edge inequalities and show that at every admissible triple
$(\Delta,n,s)$ with $n\le450$, at least one of the two certified bounds
for $R$ strictly exceeds $U(n,s)$. Write $Q=\lfloor S\sqrt{\Delta}\rfloor$,
so $nQ/[S(\Delta+1)]$ is a rational lower bound. Both loops run over
integers and use no enumeration of graphs.
{\small
\begin{verbatim}
from math import isqrt
S = 100000
for lo,d,K,B,C in [(2,14,35235,8348,2886),(15,15,34405,8211,2862)]:
    for j in range(2,d+1):
        for i in range(1,j+1):
            p,u = int(i==1),int(i>=3)
            A = (K+C)*(i+j) - i*j*(B*p+C*(p+1-u))
            assert A <= 0 or A*A <= S*S*i*j
    for h in range(lo,d+1):
        Q = isqrt(S*S*h)
        for n in range(h+1,451):
            for s in range(h,n):
                t = n-s-1
                P = 12*s*(n-1)+3*t*n*n-t*(t*t-1)
                V = max((h+1)*(K*n-(B+C)*s+C), n*Q)
                assert 6*n*(n-1)*V > S*(h+1)*P
\end{verbatim}
}
\end{proof}

\section{Components above a degree threshold}\label{p07-sec:hubs}
Let $s$ be the largest number of leaves in a spanning tree, and $L$ the number of leaves of $G$. We will work with $n\ge3$. A spanning tree extending a maximum-degree star shows that $s\ge\Delta$. Every spanning tree retains the $L$ original leaves.

For $t\ge2$ and $i\le j$ define
\[
 z_t(i,j)=\frac{\mathbf1_{i<t\le j}}{i}
 -2\left(\frac{\mathbf1_{i\ge t}}i+\frac{\mathbf1_{j\ge t}}j\right)
 +2\,\mathbf1_{i\ge t}\left(\frac1{i+1}+\frac1{j+1}\right).
\]
\begin{lemma}\label{p07-lem:hub-components}
For every connected graph of order at least three and every $t\ge2$,
\[
 s\ge2+\sum_{uv\in E(G)} z_t(\min(d_u,d_v),\max(d_u,d_v)).
\]
\end{lemma}
\begin{proof}
Let $H$ be induced by vertices of degree at least $t$, and first suppose $H$ is nonempty. Let $c$ be its number of components, counting isolated vertices. In each component take a spanning tree. Assign each vertex outside $H$ adjacent to $H$ to one of its neighbors in $H$ and retain that edge. The resulting forest has $c$ components, and its $A$ assigned vertices are leaves. If there are missing vertices, extend the forest to a spanning forest by repeatedly adding an edge from the current vertex set to a vertex outside it. Each step creates a leaf and destroys at most one leaf. Join the remaining $c$ components with $c-1$ edges, at a loss of at most $2(c-1)$ leaves. Consequently $s\ge A-2(c-1)$. Each crossing edge of degrees $i<t\le j$ contributes $1/i$ to a sum bounded by $A$, since each low vertex of degree $i$ is incident to at most $i$ such edges.

The rank of a spanning forest of $H$ is $|V(H)|-c$. Give $V(H)$ a uniform random total order. A vertex with an earlier neighbor chooses one such neighbor and retains their edge. The chosen edges form a forest, since the largest vertex of a hypothetical cycle could not choose both incident cycle edges. A vertex of internal degree $h$ has an earlier neighbor with probability $h/(h+1)$. By summing expectations,
\[
 |V(H)|-c\ \ge\ \sum_{v\in V(H)}\frac{d_H(v)}{1+d_H(v)}
 \ \ge\ \sum_{uv\in E(H)}\left(\frac1{1+d_u}+\frac1{1+d_v}\right).
\]
The second inequality follows from $d_u\ge d_H(u)$ at each vertex. Substitution into $s\ge A-2(c-1)$ proves the lemma. If $H$ is empty the edge sum vanishes, and $s\ge2$.
\end{proof}

\begin{theorem}\label{p07-thm:degree100}
Every finite connected simple graph of order $n\ge2$ and maximum degree at most $100$ satisfies $R\ge\mu$. Equality occurs only for $K_2$.
\end{theorem}
\begin{proof}
The range $\Delta\le15$ follows from the earlier maximum-degree theorem. Put $S=10^6$. For each row below use the indicated set of positive weights $(t,C_t)$. The triples in the middle columns give $(s_0,d,K)$, where $s_0\le\Delta\le d$; $B$ is the coefficient of a pendant edge.
\[
\begin{array}{c|r|r|r|r}
 &s_0&d&K&B\\ \hline
 1&16&50&392416&70524\\
 2&51&100&397964&50399
\end{array}
\]
The weights are listed in the certificate below. The edge inequality used in each row is
\[
 \frac{S}{\sqrt{ij}}\ge K\left(\frac1i+\frac1j\right)-B\mathbf1_{i=1}
           -\sum_t C_tz_t(i,j),\qquad 1\le i\le j\le d,\quad j\ge2.
\]
Writing $C=\sum_t C_t$, Lemma~\ref{p07-lem:hub-components} and $L\le s$ give
\[
 SR\ge Kn-(B+C)s+2C.
\]

By \eqref{p07-eq:leaf15double}, the distance bound to compare is $U(n,s)$.
Write $x=s/n$ and $F(x)=1/3-x^2/2+x^3/6$ as above, so that
$U(n,s)\le nF(x)+3/2$.
With $k=K/S$ and $b=(B+C)/S$, calculus gives
\[
 \min_{0\le x\le1}(k-bx-F(x))=k-b-\tfrac13(1-2b)^{3/2}>1/300
\]
for each of the two rows. The last inequality is verified by the exact squaring below. The bound $R>\mu$ follows for $n\ge450$.

The remaining finite checks involve only endpoint degrees and the two integer parameters $n,s$. The minimum certified difference $(Kn-(B+C)s+2C)/S-U(n,s)$ for $n<450$ is $631959/8500000$ in row one and $27854801/133400000$ in row two.
\begin{verbatim}
from fractions import Fraction as Q
S=10**6
T1=[(8,108549),(16,117674)]
T2=[(3,15476),(5,28042),(7,24598),(9,20818),(11,21547),
    (13,19107),(15,10358),(17,26019),(20,27067),(24,35006),
    (48,47502)]
for s0,d,K,B,T in [(16,50,392416,70524,T1),
                   (51,100,397964,50399,T2)]:
    C=sum(c for _,c in T); b=Q(B+C,S); a=Q(K,S)-b-Q(1,300)
    assert 0<b<Q(1,2) and a>0 and 9*a*a>(1-2*b)**3
    for j in range(2,d+1):
        for i in range(1,j+1):
            def z(t):
                low=Q(int(i<t<=j),i)
                high=Q(int(i>=t),i)+Q(int(j>=t),j)
                rank=int(i>=t)*(Q(1,i+1)+Q(1,j+1))
                return low-2*high+2*rank
            v=K*(Q(1,i)+Q(1,j))-B*int(i==1)
            v-=sum(c*z(t) for t,c in T)
            assert v<=0 or i*j*v*v<=S*S
    for n in range(s0+1,450):
        for s in range(s0,n):
            t=n-s-1; P=12*s*(n-1)+3*t*n*n-t*(t*t-1)
            assert 6*n*(n-1)*(K*n-(B+C)*s+2*C)>S*P
\end{verbatim}
\end{proof}
\section{Cost of cycle edges}\label{p07-sec:cycleloss}
Write $b=|E(G)|-n+1$ for the circuit rank. On deleting a nonbridge
$uv$ with degrees $d\le e$, put $x=1/\sqrt d$, $y=1/\sqrt e$.
The net increase in $R$ is $A_u+A_v-xy$, where
\[
 A_u=((d-1)^{-1/2}-x)\sum_{w\sim u,\ w\ne v}d_w^{-1/2}.
\]
In particular $A_u\le x/2$ since
$\sqrt{d(d-1)}\le d-1/2$; likewise $A_v\le y/2$.
Since $uv$ belongs to a cycle, a neighbor of $u$ other than $v$ has
degree at least two. Hence
\[
 A_u+A_v-xy\le f(d)+x\max(0,1/2-x),\quad
 f(d)=(d-2+1/\sqrt2)((d-1)^{-1/2}-d^{-1/2}).
\]
If $d\ge9$ the right side is at most $x-x^2\le 2/9$. For $d=2,\ldots,8$,
successive strict upper bounds in thousandths are
$208,222,210,220,221,219,216$, respectively, each below $2000/9$.
They follow by bounding $j^{-1/2}$ between
$10^{-6}\lfloor\sqrt{10^{12}/j}\rfloor$ and that quantity plus $10^{-6}$.
Every nonbridge deletion therefore increases $R$ by at most $2/9$.

Fix a spanning tree $T$ with the largest possible number $s$ of leaves
and delete the $b$ complementary edges, keeping $T$ at each step.
Since distances increase under edge deletion, $R(G)-\mu(G)\ge
R(T)-\mu(T)-2b/9$. By \cite[Theorem~1.2]{p07-CPSMean} we obtain
\[
 R(G)-\mu(G)\ge\sqrt{s}-2-2b/9.
\]
For $3\le s\le n-2$ the stronger estimate
$R(T)\ge(n-s)/2+\sqrt{s}-0.462$ is
\cite[Corollary~2.3]{p07-CPSMean}. Using $\mu(T)\le nF(s/n)+3/2$ and
$F(x)=(1-x)/2-(1-x)^3/6$ gives
\[
 R(G)-\mu(G)>\frac{(n-s)^3}{6n^2}+\sqrt{s}-2-2b/9.
\]
When $s=n-1$ and $n\ge5$, $G$ has a universal vertex, so $\mu(G)<2$ and
$R(G)\ge\sqrt{n-1}$ by \eqref{p07-eq:basic}. Orders below five and $s=2$ fall under Theorem~\ref{p07-thm:leaf15}. Consequently,
for every fixed $b$, with $h=2+2b/9$, the conjecture holds if
$n\ge\max\{48h,2h^2\}$. Indeed if $s\ge n/2$ the square-root term
is at least $\sqrt{n/2}$, and otherwise the cubic term exceeds $n/48$.

\section{Keeping the rank of the high-degree subgraph}
Let $G$ be connected of order $n\ge3$, size $m$, and circuit rank
$b=m-n+1$. Write $L$ for the number of pendant vertices, $s$ for the
maximum number of leaves over its spanning trees, and $J$ for the
subgraph induced by the vertices of degree at least three. The degrees
$d_v$ below are always taken in $G$.
\begin{lemma}\label{p07-lem:ranklocal}
Put $p_{ij}=\mathbf{1}_{i=1}$ and $z_{ij}=\mathbf{1}_{i\ge3}(1-1/i)$
for ordered endpoint degrees $i\le j$. Then
\[
  \sum_{uv\in E}z_{ij}\ge L+b-s.
\]
\end{lemma}
\begin{proof}
Extend a spanning forest of $J$ to a spanning tree of $G$. Exactly
$|E(J)|-\operatorname{rank}(J)$ omitted edges lie in $J$, where the
rank denotes the number of edges of its spanning forest. All other
omitted edges have an endpoint of degree two, and two omissions cannot
use the same degree-two vertex. These endpoints are new leaves. Thus
$s\ge L+b-|E(J)|+\operatorname{rank}(J)$.
For each nonsingleton component $Q$ of $J$, write $d'_u$ for the
degree in $Q$. Its edges satisfy
\[
 \sum_{uv\in E(Q)}\frac1{\min(d_u,d_v)}
 \le \sum_{uv\in E(Q)}\frac1{\min(d'_u,d'_v)}
 = |V(Q)|-\sum_{uv\in E(Q)}\frac1{\max(d'_u,d'_v)}
 \le |V(Q)|-1.
\]
The last sum of reciprocal maxima is at least one, since the edges
at a maximum-degree vertex alone sum to one. Isolated vertices also
satisfy the bound with both sides zero. Sum over components and use
the preceding spanning-tree inequality.
\end{proof}



\begin{theorem}\label{p07-thm:cycl155}
Every connected simple graph of order at least two and circuit rank
at most $155$ satisfies $R\ge\mu$, with equality only for $K_2$.
\end{theorem}
\begin{proof}
The bounded-degree theorem allows us to assume $s\ge d=\Delta\ge101$.
The earlier deletion and spanning-tree bounds give
$R-\mu\ge\sqrt{s}-2-2b/9$, and strictness for
$n\ge\max(48h,2h^2)$, with $h=2+2b/9$. For $b\le155$ we have
$h\le328/9$, so $s\ge1329$ or $n\ge2657$ suffices. When $s=n-1$,
a universal vertex gives $\mu<2\le\sqrt{n-1}\le R$.
It remains to consider $101\le s\le1328$, $s+2\le n\le2656$.

For a threshold $T$ set
\[
 Z_T(i,j)=\frac{\mathbf1_{i<T\le j}}i
 -2\left(\frac{\mathbf1_{i\ge T}}i+\frac{\mathbf1_{j\ge T}}j\right)
 +2\mathbf1_{i\ge T}\left(\frac1{i+1}+\frac1{j+1}\right).
\]
The high-core forest bound in the bounded-degree proof gives
$\sum_E Z_T\le s$ for every $T$. We keep $T=4,8,16$ and
write $w=1/i+1/j$, $p=p_{ij}$, $z=z_{ij}$ as above.
The exhaustive integer certificate below supplies nine rows with degree
caps $H$. For a row $(H,A,B,P,C,g_4,g_8,g_{16})$ it checks
\[
 \frac{10^5}{\sqrt{ij}}\ge
 Aw+B-Pp-C(p-z)-\sum_T g_T Z_T(i,j)
 \qquad(1\le i\le j\le H,\ j\ge2).
\]
The check multiplies by $ij(i+1)(j+1)$ before squaring.
Summing the inequalities and using Lemma~\ref{p07-lem:ranklocal},
$L\le s$ and $\sum_E Z_T\le s$ gives the bound
\[
 10^5R\ge An+B(n-1)+(B+C)b-(P+C+\textstyle\sum_Tg_T)s .
\]
All rows satisfy $B+C<0$, so the lower bound remains valid with
$b$ replaced by $155$. Every row with cap at least $d$ gives a separate
lower bound.

Put $t=n-s-1$. For the distance upper bound we use the smaller
of the balanced double-comet bound
\[
 U(n,s)=\frac{12s(n-1)+3tn^2-t(t^2-1)}{6n(n-1)}
\]
and the maximum-degree broom bound of Stevanovi\'c~\cite{p07-Stev}.
For $d\ge d_0$ the latter bound is
\[
 \mu\le\binom n2^{-1}\left[
   \binom{u+2}{3}+(d_0-1)\frac{(u+1)(u+2)}2+(d_0-1)(d_0-2)\right],
 \qquad u=n-d_0 .
\]
For completeness, use a spanning tree containing the star of a
degree-$d$ vertex $v$. Replace each branch by a path rooted at an end.
This weakly increases both its internal Wiener index and its sum of
distances to $v$, hence also the total Wiener index. For a spider with
arm lengths $m_i$ the Wiener index is $\sum_i\sum_{k=1}^{m_i}k(n-k)$.
Moving one unit from a shorter arm $b'\ge2$ to a longer arm $a'\ge b'$
adds $(a'+1-b')(n-a'-b'-1)\ge0$. If there are more than $d_0$ arms
we can also move the last unit of a length-one arm to another arm by
the same formula. The resulting broom has one arm of length $n-d_0$
and $d_0-1$ pendant edges, giving the stated bound.

The check additionally uses $R\ge\sqrt{n-1}$,
$R\ge(n-s)/2+\sqrt{s}-0.462-310/9$ from the earlier tree comparison,
and the pendant-edge degree-ratio estimate. For $d\le D$, $D\ge101$,
the latter gives
\[
 R\ge\frac{\sqrt{2D}}{D+2}(n-s-s/D)+s/\sqrt D .
\]
Indeed the ratio of an internal edge weight to its reciprocal-degree
sum is at least $\sqrt{2D}/(D+2)$, while the deficit on each pendant
edge is at most $[\sqrt{2D}/(D+2)](1+1/D)-1/\sqrt D>0$ and $L\le s$.
For $D=s$ this simplifies to $R\ge\sqrt{2s}(n-s-1)/(s+2)+\sqrt{s}$.
In lower bounds with rounded square roots we check the sign of each
coefficient. The integer program first examines $2\,382\,934$ pairs
$(n,s)$. For the surviving pairs it checks $2\,113\,784$ intervals
$d_0\le\Delta\le\min(s,H)$, using $d_0$ for the broom and
$D=\min(s,H)$ in the lower bound. All comparisons are strict.
\begin{verbatim}
from math import isqrt, comb
S=10**5
rows=[
(125,61878,-13163,8420,7558,5321,7443,11938),
(160,62388,-13552,8440,7678,5394,20145,0),
(220,63006,-14023,8469,7819,5487,7315,13698),
(350,63820,-14638,8514,7999,5614,7221,14849),
(600,64630,-15247,8564,8171,5745,7126,15916),
(800,63700,-15247,8564,8171,5745,7126,15916),
(1000,63000,-15247,8564,8171,5745,7126,15916),
(1200,62500,-15247,8564,8171,5745,7126,15916),
(1600,61600,-15247,8564,8171,5745,7126,15916)]
for H,A,B,P,C,*rt in rows:
    assert B+C<0 and min([P,C]+rt)>=0
    for j in range(2,H+1):
        for i in range(1,j+1):
            D=i*j*(i+1)*(j+1)
            v=A*(i+j)*(i+1)*(j+1)+B*D-(P+C)*D*(i==1)
            if i>=3: v+=C*(D-D//i)
            for T,G in zip((4,8,16),rt):
                if i<T<=j: v-=G*(D//i-2*(D//j))
                if i>=T:
                    v+=2*G*j*(j+1)+2*G*i*(i+1)
            assert v<=0 or v*v*i*j<=S*S*D*D,(i,j,H)
Q=10**6
p=[isqrt(j*Q*Q) for j in range(5400)]
pairs=trials=0
for n in range(103,2657):
    den=6*n*(n-1)
    for s in range(101,min(n-1,1329)):
        pairs+=1
        t=n-s-1
        v=12*s*(n-1)+3*t*n*n-t*(t*t-1)
        # tree bound denominator 9000*Q
        tw=4500*Q*(n-s)+9000*p[s]-(4158+310000)*Q
        base=p[2*s]*(n-s-1)+p[s]*(s+2)
        if (tw*den>9000*Q*v or p[n-1]*den>Q*v
            or base*den>Q*(s+2)*v): continue
        D0=101
        for H,A,B,P,C,*rt in rows:
            if D0>s: break
            u=n-D0
            wb=comb(u+2,3)+(D0-1)*(u+1)*(u+2)//2
            bw=12*(wb+(D0-1)*(D0-2))
            upper=min(v,bw)
            # all rows with cap >= d_hi apply
            hi=min(s,H); D0=H+1; trials+=1
            prov=(tw*den>9000*Q*upper or p[n-1]*den>Q*upper
                  or base*den>Q*(s+2)*upper)
            e=hi*(n-s)-s
            if e>=0:
                num=p[2*hi]*e+s*p[hi]*(hi+2)
                prov|=(num*den>Q*hi*(hi+2)*upper)
            for cap,a,b,p0,c,*rt0 in rows:
                if cap<hi: continue
                lb=a*n+b*(n-1)+(b+c)*155-(p0+c+sum(rt0))*s
                if lb*den>S*upper: prov=True; break
            assert prov,(n,s,hi)
assert (pairs,trials)==(2382934,2113784)
\end{verbatim}
\end{proof}
\section{Projection onto a shortest path}
The modified index $M=\sum_{uv}1/\max(d_u,d_v)$ satisfies $M\ge D/2$
by Dvo\v{r}\'ak, Lidick\'y and \v{S}krekovski \cite[Theorem~3]{p07-DLS}.
Writing $a\ge b$ for the reciprocal square roots of the two endpoint
degrees, one has $(a-b)^2\le(\sqrt{\Delta}-1)b(a-b)$. The sums of the two
terms $(a-b)^2$ and $b(a-b)$ are $n-2R$ and $R-M$, respectively.
Thus $(\sqrt{\Delta}+1)R\ge n+(\sqrt{\Delta}-1)M$.
For a shortest path $P$ of length $d$, put $t=n-d$ and
$k=\max_{v\in V(G)}\dist(v,P)$. We allow $d$ to be smaller than the
diameter $D$. If $k\ge1$, then
\begin{equation}\label{p07-eq:faj-pathprojection}
 \mu(G)\le \frac{dn}{2(n-1)}+\frac{2k(t-1)}n<\frac d2+2k.
\end{equation}
Indeed, project each vertex to a nearest vertex of $P$. The distance
sum between pairs of projected vertices is at most $dn^2/4$, since
each of the $d$ cuts of the path separates at most $n^2/4$ unordered
pairs. The total cost of going to and from the path, divided by
$\binom n2$, is at most $2k(t-1)/n$. The strict inequality follows
from $d/[2(n-1)]<2k(d+1)/n$.

Every vertex on $P$ has at most two neighbors on $P$, and an outside
vertex meets at most three consecutive positions. Thus $\Delta\le t+1$.
The edge-energy bound $(\sqrt\Delta+1)R\ge n+(\sqrt\Delta-1)M$
and the diameter bound $M\ge D/2$ imply
\[
 R\ge\frac d2+\frac{t}{\sqrt\Delta+1}
 \ge\frac d2+\sqrt{t+1}-1.
\]
Consequently $R>\mu$ whenever $t\ge4k(k+1)$. Also,
$n\ge3t+2$ suffices, since $d/2\ge(n+1)/3\ge\mu$.
For $k=1$ this proves the conclusion for $t\ge8$, and the smaller
values give $\Delta\le8$. In particular, every connected AT-free
graph satisfies the conjecture: Corneil, Olariu and Stewart proved
that such a graph contains a dominating pair, and a shortest path between
its endpoints dominates the graph. See \cite{p07-COS}.

\subsection*{A geodesic of covering radius at most eighteen}
\begin{theorem}\label{p07-thm:widthplus}
Let $G$ be a connected simple graph of order $n\ge3$. If there is a
shortest path within distance eighteen of every vertex, then $R(G)>\mu(G)$.
\end{theorem}
\begin{proof}
Choose such a geodesic with at least one edge (a one-vertex path can be extended by one edge). Write $d$ for its length, $D$ for the diameter,
$h=\Delta$, $K=18$ and $t=n-d$. Let $f(v)$ be a nearest vertex of $P$.
For each of the $d$ cuts of the path let $N_i$ count the vertices
projecting to its initial segment. The projected distance sum, divided by
$\binom n2$, equals
\[
 \frac{dn}{2(n-1)}-\frac{2}{n(n-1)}\sum_{i=1}^d(N_i-n/2)^2.
\]
Each vertex of the path projects to itself, so the $N_i$ are strictly
increasing integers. The sum of their squared deviations from $n/2$ is
at least $d(d^2-1)/12$. Indeed their pairwise gaps are at least those
of $1,\ldots,d$, and their sum of squared deviations is minimized about
their own mean. The resulting projection bound is
\begin{equation}\label{p07-eq:widthplusproj}
 \mu\le \frac{dn}{2(n-1)}+\frac{2K(t-1)}n-
 \frac{d(d^2-1)}{6n(n-1)}.
\end{equation}
Here the outside paths to and from the geodesic have total average cost
at most $2K(t-1)/n$. One also has $D\le d+2K$ and $h\le t+1$ by the
geodesic neighbor bound. The ordinary projection and degree-energy
bounds already handle $t\ge4K(K+1)$ and $n\ge3t+2$. The bounded-degree
theorem handles $h\le100$. It remains to check
\begin{gather*}
 100\le t<1368,\qquad t+1\le n\le3t+1,\\
 101\le h\le\min(t+1,n-1),\qquad d\le D\le\min(n-1,d+36).
\end{gather*}

For a maximum-degree vertex $z$, put $H=N[z]$, $m=h+1$, $a=n-m$.
Any two vertices of $H$ are within distance two. For fixed $u\in H$ the
sum of the distances from $u$ to the $a$ vertices outside $H$ is at most
\[
 aD-\tfrac12(D-3)_+(D-2)_+ .
\]
To see this, let $e$ be the eccentricity of $u$. A path to a farthest
vertex gives an outside vertex at every distance $3,\ldots,e$. Thus the
sum is at most $ae-(e-3)_+(e-2)_+/2$. This expression is nondecreasing on
the integer interval $0\le e\le D$, since $a\ge D-2$ by the diametral
path neighbor bound. To estimate the distances between outside vertices,
contract the connected subgraph $H$ to one vertex. In the contracted graph
a shortest path lifts to a walk in $G$ at an additional cost of at most
two, because the traversal of $H$ can be replaced by a path of length at
most two. The total distance sum of the contracted graph of order $a+1$
is at most $\binom{a+2}{3}$. Thus, putting $f=(D-3)_+(D-2)_+$, one has
\begin{align*}
 n(n-1)\mu\le {}&2m(m-1)+2maD-mf\\
 &+\min\{D a(a-1),\ a(a+1)(a+2)/3+2a(a-1)\}.
\end{align*}
There is an additional saving on the outside pairs. Write $g=(D-K)_+$,
$r=\min\{g,d+1\}$ and let
\begin{align*}
 L&=[d(d+1)(3D-d-2)-18d(D-1)]_+,\\
 V&=[r(2g-r+1)-6g]_+ .
\end{align*}
In the preceding display, the term $D a(a-1)$ may be replaced by
\[
 D a(a-1)-L/3-(a-d-1)_+ V .
\]
To see this, measure the deficit of each \emph{ordered} pair from the
upper bound $D$. Between distinct vertices of the path the total deficit
is $d(d+1)(3D-d-2)/3$. At most three path vertices belong to $H$, by
shortestness. Removing their incident pairs loses at most $6d(D-1)$,
leaving the lower bound $L/3$ on the contribution from path vertices
outside $H$. Now take any outside vertex not on $P$, projecting to a
path position $j$. For a path vertex at position $i$ its deficit is at
least $[g-|i-j|]_+$. The sum of these bounds over $i=0,\ldots,d$ is
smallest at an endpoint $j$ (the successive difference is
$[g-(j+1)]_+-[g-(d-j)]_+$, nonnegative up to the middle).
That endpoint sum is $r(2g-r+1)/2$. Removing the at most three path
vertices in $H$ costs at most $3g$. Thus for both orderings the
remaining pair contributions total at least $V$ per outside vertex
not on $P$. There are at least $(a-d-1)_+$ such vertices. The two
collections of pairs used for $L$ and $V$ are disjoint.

Finally, $R\ge\sqrt{n-1}$ and $R\ge D/2+(n-D)/(\sqrt h+1)$, the
latter by the degree-energy inequality and $M\ge D/2$. Below all upper
bounds are stored as numerators with common denominator $Q=6n(n-1)$.
The function \texttt{done} checks by squaring whether one of these two
lower bounds is strictly above the given upper bound. Its second argument
may be decreased, or its third argument increased, since the lower bound
is nondecreasing in $D$ and nonincreasing in $h$. The check compares
\eqref{p07-eq:widthplusproj}, $\mu\le(n+1)/3$, the double-comet estimate for
$s\ge h$, the hub estimate above, and $\mu\le D$. It uses only integer
arithmetic.
\begin{verbatim}
K=18
tested=0
for t in range(100,4*K*(K+1)):
    for n in range(t+1,3*t+2):
        d=n-t; Q=6*n*(n-1); den2=Q//2
        W=2*n*(n-1)*(n+1)
        U=min(W, 3*d*n*n+12*K*(t-1)*(n-1)-d*(d*d-1))
        def done(u,D,h):
            V=u-D*den2
            A=(n-D)*Q-V
            return (u*u<(n-1)*Q*Q or V<=0 or A>0 and A*A>h*V*V)
        if done(U,d,t+1):continue
        for h in range(101,min(t+1,n-1)+1):
            U0=min(U,W-h*(h-2)*(3*n-5-h))
            if done(U0,d,h):continue
            m=h+1; a=n-m
            for D in range(d,min(n-1,d+2*K)+1):
                if done(U0,D,h):break
                if done(D*Q,D,h):continue
                f=max(D-3,0)*max(D-2,0)
                g=max(D-K,0); r=min(g,d+1)
                L=max(0,d*(d+1)*(3*D-d-2)-18*d*(D-1))
                V=max(0,r*(2*g-r+1)-6*g)
                w=min(3*D*a*(a-1)-L-3*max(0,a-d-1)*V,
                      a*(a+1)*(a+2)+6*a*(a-1))
                U1=6*(2*m*(m-1)+2*m*a*D-m*f)+2*w
                assert done(min(U0,U1),D,h),(t,n,h,D)
                tested+=1
assert tested == 129342972
\end{verbatim}
\end{proof}
\section{Structural classes for the modified index}
A block graph is a connected graph whose blocks are cliques, with every
bridge counted as a block of order two. In particular the family allows
cliques of unbounded order, joined at arbitrary cut vertices. In this
section $M=\sum_{uv\in E}1/\max(d_u,d_v)$.

\begin{theorem}\label{p07-thm:strongblocks}
Let $G$ be a connected block graph of order $n\ge2$. Then
\[
 M(G)\ \ge\ \mu(G)-1+2/n,\qquad R(G)\ \ge\ \mu(G).
\]
The second inequality is strict for $n>2$. If $G$ has $c\ge1$ cut
vertices, then $2M(G)\ge c+1$.
\end{theorem}
\begin{proof}
A single clique has $M=R=n/2$ and $\mu=1$. Otherwise the graph has a
nonempty set of cut vertices. A shortest path between any two of these
vertices has only cut vertices in its interior: a noncut vertex belongs
to a single block and has all its neighbors in one clique, which would
shorten that path. Every other vertex is adjacent to a cut vertex,
since its unique block meets the rest of $G$ at a cut vertex. Take a
spanning tree of the induced connected graph on all $c$ cut vertices,
and attach each remaining vertex as a leaf to a cut neighbor. Every
cut vertex is internal in every spanning tree, since deleting it leaves
at least two components even in the full graph. This constructs a tree
with exactly $c$ internal vertices.

We now shrink the cliques while preserving their cut vertices. Choose a
noncut vertex $z$ in a block of order $q\ge3$. Its degree is $q-1$;
each neighbor $v$ has degree $d_v\ge q-1$. Deleting $z$ removes the
weights $1/d_v$ on its incident edges. At a remaining edge incident
with $v$, lowering that endpoint degree from $d_v$ to $d_v-1$ can
increase its weight by at most $1/[d_v(d_v-1)]$, even when its other
endpoint also loses a neighbor. In the latter case lower the two
degrees successively and apply the bound to each individual change.
The $d_v-1$ surviving edges at $v$ contribute a total increase at most
$1/d_v$. Summing over the neighbors proves $M(G-z)\le M(G)$.
The clique of order $q$ becomes one of order $q-1\ge2$. Other blocks
and all memberships of cut vertices in those blocks are preserved.
Thus there are still exactly $c$ cut vertices. Repeat until every
noncut vertex of the resulting graph $G_0$ is a leaf. All nonleaves of
$G_0$ are now cut vertices. The reduced graph has at least three vertices,
since it retains a cut vertex.

For any connected graph $X$ of order $N\ge3$ write $m=|E(X)|$,
$L=|\{v:d_v=1\}|$, $\beta=m-N+1$ and
\[
 B=\sum_{uv:\,\min(d_u,d_v)\ge3}\left(1-\frac2{\min(d_u,d_v)}\right).
\]
Summing $1/d_u+1/d_v-2/\max(d_u,d_v)$ gives
$N-2M=\sum_{uv}|1/d_u-1/d_v|$. Inserting $1/2$ between the two
reciprocals gives total $\sum_v|1-d_v/2|=L+m-N$ before subtracting
the triangle-inequality slack. This slack is precisely $B$, since no
two leaves are adjacent. Consequently
\[
 2M(X)=N-L+1+B-\beta .
\]
For $X=G_0$ each clique block of order $q\ge3$ has only cut vertices.
Each belongs to another block and thus has full degree at least $q$.
The clique edges contribute at least
$\binom q2(1-2/q)=(q-1)(q-2)/2$ to $B$. This is exactly the circuit
rank of that block. Distinct blocks are edge-disjoint and their ranks
sum to $\beta$, since gluing connected graphs at a vertex adds
circuit ranks. Bridges contribute no circuit rank. Hence $B\ge\beta$,
and $2M(G)\ge2M(G_0)\ge c+1$.

In the constructed spanning tree the $n-c$ pendant edges each carry
$n-1$ unordered vertex pairs. Each of the $c-1$ internal edges carries
at most $n^2/4$ pairs. Distances can only decrease in $G$, so
\[
 \mu(G)\le \frac{2(n-c)}n+\frac{(c-1)n}{2(n-1)} .
\]
Subtracting this from $(c+1)/2+1-2/n$ gives
$(c-1)(2/n-1/[2(n-1)])\ge0$, proving the claim for $M$.

For completeness, write $y=\sqrt\Delta$. The edge identity
\[
(d_u^{-1/2}-d_v^{-1/2})^2
=(\sqrt{\max(d_u,d_v)/\min(d_u,d_v)}-1)(R_{uv}-M_{uv})
\]
gives $n-2R\le(y-1)(R-M)$. Since $\mu\le(n+1)/3$ and
$M\ge\mu-1+2/n$, we have
\[
 (y+1)(R-\mu)\ge(n-2)
      \left(\frac13-\frac{\sqrt{n-1}-1}{n}\right)>0
 \quad(n>2).
\]
The strict sign follows by squaring $1+n/3>\sqrt{n-1}$.
\end{proof}

For a simultaneous illustration of the scope, take a clique $K_q$
with $q\ge3$, fix one of its vertices, and attach three internally
disjoint pendant paths of length $h$ there. Its maximum degree is
$q+2$ and its circuit rank is $(q-1)(q-2)/2$. Every simple path meets
the interiors of at most two of the three arms, so some leaf is at
distance at least $h$ from the path. Thus the theorem permits
simultaneously unbounded maximum degree, circuit rank and distance
from a best covering geodesic.
\begin{corollary}
Let $H$ be a connected graph in which every nonleaf is a cut vertex and
all blocks of order at least three are regular (their degrees are measured
inside the block). Then $M\ge\mu-1+2/n$ and $R>\mu$ for $n\ge3$.
\end{corollary}
\begin{proof}
In the connected induced subgraph on the nonleaves, take a spanning
tree and attach the leaves to their supports. Every nonleaf remains
internal since it is a cut vertex. Use the identity $2M=N-L+1+B-\beta$
from the preceding proof. If a block has $v$ vertices and internal degree
$r\ge2$, every vertex is a cut vertex, belongs also to another block,
and thus has full degree at least $r+1$. The block contributes at least
$vr/2-vr/(r+1)$ to $B$, which is at least its circuit rank $vr/2-v+1$
since $v\ge r+1$. Sum over blocks and finish as for block graphs.
\end{proof}
\subsection*{A vertex-cover condition for the modified index}
\begin{proposition}\label{p07-prop:fajtcover}
Let $G$ be connected of order $n\ge2$, and let $M=\sum_{uv}1/\max(d_u,d_v)$. If $G$ has a vertex cover $X$ with $2|X|-r_X\le2M$, where $r_X$ is the number of edges in a spanning forest of $G[X]$, then $M\ge\mu-1+2/n$ and $R\ge\mu$, strictly for the latter inequality when $n>2$. This applies in particular to every bipartite graph $G=(X,Y;E)$ such that $d_x\ge d_y$ on each edge $xy$ with $x\in X$. Thus it includes all connected biregular bipartite graphs.
\end{proposition}
\begin{proof}
Extend a spanning forest of $G[X]$ to a spanning tree $T$ of $G$. Exactly $r_X$ of its edges lie in $G[X]$. Since $Y=V(G)\setminus X$ is independent, $\sum_{y\in Y}d_T(y)=n-1-r_X$. Every nonleaf of $T$ in $Y$ contributes at least one to $\sum_{y\in Y}(d_T(y)-1)=|X|-1-r_X$. Thus the number $I$ of nonleaves satisfies $I+1\le 2|X|-r_X\le 2M$. For $n\ge3$, pendant edges of $T$ contribute $n-1$ each to its Wiener index and the other $I-1$ edges contribute at most $n^2/4$ each. Hence
\[
 M-\mu+1-2/n\ge(I-1)\left(2/n-1/[2(n-1)]\right)\ge0.
\]
The modified-index transfer to $R$ gives the claim. On a bipartition with the stated degree ordering, $M=\sum_{x\in X}d_x/d_x=|X|$. For $K_2$ both assertions hold directly.
\end{proof}

More generally, suppose that every edge from $X$ to $V(G)\setminus X$ has its larger degree at the $X$ end, allowing equality. Then
\[
 |X|-M=\sum_{uv\in E(G[X])}\frac1{\min(d_u,d_v)}.
\]
Indeed, each $x\in X$ contributes a total of one when its incident edges are counted with weight $1/d_x$. The boundary terms agree with $M$, and an internal edge is counted at both ends. Thus the proposition applies if $\sum_{uv\in E(G[X])}2/\min(d_u,d_v)\le r_X$. This condition holds in particular when $G[X]$ is a forest and the endpoints of its edges have full degree at least two. The forests may contain adjacent high-degree vertices, and vertices outside $X$ may share arbitrarily many of these neighbors.

The simpler sufficient condition $|X|\le M$ fails for $K_3$, which has minimum vertex-cover size two and $M=3/2$. Its minimum cover has $r_X=1$, so the forest version includes it. Writing $\gamma_c$ for the minimum size of a connected dominating set, a related universal condition $2M\ge\gamma_c+1$ also fails, with an unbounded deficit. Give every vertex of $K_{2,k}$ one private leaf, with $k\ge4$. The $k+2$ original vertices are cut vertices. Each belongs to every connected dominating set, since its private leaf must be dominated and a private leaf cannot connect the rest. Those core vertices themselves form a connected dominating set. Hence $\gamma_c=k+2$, whereas $M=2+k/3$ by summing over the three edge types. The deficit $(\gamma_c+1)-2M=k/3-1$ tends to infinity with $k$.
\subsection*{Graphs obtained by replacing every vertex with a module}
Let $F$ be a connected graph with at least two vertices. Replace each $v\in V(F)$ by a nonempty graph $H_v$. Join every pair of vertices in $H_u,H_v$ when $uv\in E(F)$, with no other intermodule edges. Denote the result by $G$. No regularity or common internal edge pattern is required below.
\begin{proposition}
Put $a=\min_v|H_v|$ and $b=\max_v|H_v|$. If $a^2\ge3b$, then $M(G)\ge\diam(G)$. In particular both the modified-index bound $M(G)\ge\mu(G)-1+2/|V(G)|$ and the original Randi\'c bound hold. This applies when all modules have the same size $t\ge3$.
\end{proposition}
\begin{proof}
Suppose first $|V(F)|\ge3$. On a base edge $uv$, put $j=\max(d_F(u),d_F(v))\ge2$. Each endpoint in $G$ of the corresponding cross edges has degree at most $bj+b-1$. Hence their total weight in $M(G)$ is at least
\[
 \frac{a^2}{bj+b-1}>\frac{3}{j+1}\ge\frac{2}{j}.
\]
Summing the cross edges gives $M(G)>2M(F)\ge\diam(F)$ by the published modified-index diameter theorem. Distances between distinct modules agree with distances in $F$ and distances within a module are at most two. If $\diam(F)\ge2$ this proves the assertion. Otherwise $F$ is complete, so $2M(F)=|V(F)|\ge3$ and $\diam(G)\le2$.

If $F=K_2$, write $p\le q$ for the two module sizes. Both are at least three. If there is no internal edge, $M(K_{p,q})=p\ge3$. If there is at least one, then every full degree is at most $p+q-1$ and the graph has at least $pq+1$ edges, giving $M(G)\ge(pq+1)/(p+q-1)\ge2$. The last inequality is $(p-2)(q-2)\ge1$. Again $\diam(G)\le2$.
\end{proof}
\section{Metric savings from cycle edges}
Fix a spanning tree $T$ of $G$. Write $W(X)=\sum_{\{x,y\}}d_X(x,y)$ and
$\mathcal C=E(G)\setminus E(T)$. For $e=ab\in\mathcal C$ let
$P_e=(v_0=a,v_1,\ldots,v_{L_e}=b)$ be its tree path. Deleting all
edges of $P_e$ from $T$ produces components of orders $w_i(e)$ containing
$v_i$. Define $[z]_+=\max\{0,z\}$ and
\begin{equation}\label{p07-eq:chordsavings}
 C_e=\sum_{0\le i<j\le L_e}w_i(e)w_j(e)\,[2(j-i)-L_e-1]_+ .
\end{equation}
\begin{theorem}\label{p07-thm:chordpack}
One has $W(T)-W(T+e)=C_e$. For every set of real numbers
$\lambda_e\ge0$ satisfying $\sum_{e:f\in P_e}\lambda_e\le1$ on every
edge $f$ of $T$, one has
\[
 W(T)-W(G)\ \ge\ \sum_{e\in\mathcal C}\lambda_e C_e .
\]
In particular the single-edge savings add when the paths $P_e$ are
pairwise edge-disjoint, with no restriction on shared vertices.
\end{theorem}
\begin{proof}
Two vertices attached to $P_e$ at positions $i,j$ can travel between
those positions along the tree arc or the complementary arc of the cycle.
These routes have lengths $|i-j|$ and $L_e+1-|i-j|$. If $i=j$ the
chord gives no improvement. Otherwise the two portions outside the
cycle have the same length on both routes. Summing over the positions
gives~\eqref{p07-eq:chordsavings}.

For the assertion with several edges, fix $x,y$. Write $Q$ for their
tree path. For every chord that shortens this distance, the intersection
$Q\cap P_e$ is an interval of positive length on $Q$. Its endpoints
are the two attachment positions, and replacing the interval by the
complementary cycle arc saves $g_e(x,y)=d_T(x,y)-d_{T+e}(x,y)>0$.
For any collection of such intervals with disjoint edge sets, perform
all of their replacements in their order along $Q$. The result is an
$x,y$-walk in $G$. Even if some replacement arcs meet, its length is
$d_T(x,y)$ minus the sum of the individual savings.

First suppose that all $\lambda_e$ are rational, and take a common
denominator $m$. For this pair $x,y$, replace the interval for $e$ by
$m\lambda_e$ copies. At most $m$ copies cover any tree edge. The copies
can be colored with $m$ colors so that same-colored intervals have
disjoint edge sets. To see this, represent an interval by its first and
last \emph{edge indices} on $Q$, process the intervals by first index,
and reuse any color whose previous interval ends strictly before that
index. If no color were free, more than $m$ copies would cover the first
edge. Thus some color saves at least the average over the $m$ colors,
namely $\sum_e\lambda_e g_e(x,y)$. The result for real coefficients
follows by approximating each coefficient from below by nonnegative
rationals. Sum the inequalities over all pairs $x,y$.
\end{proof}
A sharper bound, computed separately for each $x,y$, follows by taking
the maximum total saving of disjoint intervals on $Q$. The fractional
version has one variable per chord and one capacity constraint per tree
edge. In particular it gives the sufficient condition
\[
 R(T)-\mu(T)+\binom n2^{-1}\sum_e\lambda_e C_e
  \ \ge\ R(T)-R(G)
 \quad\Longrightarrow\quad R(G)\ge\mu(G).
\]

For a maximum-leaf spanning tree, let $S(T)$ be the maximum of
$\sum_e\lambda_e C_e$ over these capacity constraints (zero for a tree).
Combining the preceding theorem with the tree estimate and the bound
on deletion cost from Section~\ref{p07-sec:cycleloss} yields
\begin{equation}\label{p07-eq:combined}
 \boxed{\quad R(G)-\mu(G)\ \ge\ \sqrt{s}-2-\frac{2b}{9}
                           +\frac{S(T)}{\binom n2}.\quad}
\end{equation}
Indeed $R(G)\ge R(T)-2b/9$ and $\mu(G)\le\mu(T)-S(T)/\binom n2$.
The single-edge change in the Wiener index is also studied in
Khalifeh and Esfahanian~\cite{p07-KE}.

\paragraph{The effect of a single edge.}
For a nonbridge $e$ in any connected $G$, put
$A_e=R(G-e)-R(G)$ and $B_e=\mu(G-e)-\mu(G)$. If $A_e\le B_e$, a proof
of the conjecture for $G-e$ also proves it for $G$, since the margin
$R-\mu$ increases by $B_e-A_e$ on restoring $e$. The sign can reverse
even on arbitrarily large unicyclic graphs. Take a triangle $abc$, a
private leaf at each of $a,b$, and a path of length $q\ge1$ attached at
$c$. Put $e=ab$ and $n=q+5$. The three triangle degrees are all three.
The four edges at $a,b$ other than $ab$ yield, upon deleting $e$,
\[
 A_e=2(1+1/\sqrt3)(1/\sqrt2-1/\sqrt3)-1/3
     =\sqrt2+\sqrt{2/3}-2/\sqrt3-1>0 .
\]
Indeed $\sqrt6+\sqrt2>2+\sqrt3$ by squaring. The tree path $a,c,b$
has endpoint fibers of order two; only pairs from opposite endpoint
fibers are shortened. Thus $W(G-e)-W(G)=4$ and $B_e=4/\binom n2$.
For all sufficiently large $q$ one has $A_e>B_e$. The same example
can be arranged at arbitrarily large maximum degree by taking several
new branches at an interior vertex of the path beyond $c$, without
changing $A_e$ or the Wiener saving of four.
\begin{corollary}\label{p07-cor:longnonbridge}
In any graph $G$ with nonbridge $e=ab$, let $L=d_{G-e}(a,b)$. Then
$W(G-e)-W(G)\ge L^3/24$. In particular, if $G$ is an edge-minimal
counterexample to $R\ge\mu$, then every nonbridge satisfies
\[
 d_{G-e}(a,b)^3<\tfrac83 n(n-1).
\]
\end{corollary}
\begin{proof}
A shortest $a,b$-path $v_0,\ldots,v_L$ in $G-e$ satisfies
$d_{G-e}(v_i,v_j)=j-i$ when $i<j$. Restoring $e$ shortens this
distance by at least $[2(j-i)-L-1]_+$. Summing over the path vertices
alone yields $\sum_{h=1}^L(L+1-h)[2h-L-1]_+$. For even $L$ this is
$L(L+1)(L+2)/24$, and for odd $L$ it is $(L-1)(L+1)(L+3)/24$.
Both are at least $L^3/24$ for $L\ge2$. If $L^3\ge (8/3)n(n-1)$,
the savings in mean distance on restoring $e$ are therefore at least
$2/9$, which covers its maximum Randi\'c loss. Edge minimality would now
transfer $R\ge\mu$ from $G-e$ to $G$.
\end{proof}
\subsection*{A necessary bridge in a minimum counterexample}
Call an edge a bridge if its deletion disconnects the graph. Remove all
bridges simultaneously, and let $Q$ be any remaining component containing
an edge. All its edges were nonbridges of $G$. For $v\in V(Q)$ put
$c_v=d_Q(v)$ and $b_v=d_v-c_v$. Let $S_v$ be the sum of the Randi\'c
weights $1/\sqrt{d_vd_u}$ of incident edges in $Q$, and $T_v$ the sum
of the weights of the incident bridges, using full degrees throughout.
Write $a_d=\sqrt{d/(d-1)}-1$ and
$\epsilon_d=1/2-(d-1)a_d=(\sqrt d-\sqrt{d-1})^2/2>0$.
\begin{proposition}\label{p07-prop:bridgemin}
If $G$ is an edge-minimal counterexample, every such $Q$ contains a
vertex $v$ with $b_v>0$ and $T_v/b_v>S_v/c_v$. In particular $G$
has a bridge.
\end{proposition}
\begin{proof}
Deleting an edge $e$ of $Q$ incident to $v$ multiplies the remaining
weights at $v$ by $1+a_{d_v}$, taking one endpoint at a time.
Distinct endpoints have no remaining edge between them. Summing the
changes over $e\in E(Q)$ and assigning half the deleted weight to each
endpoint gives the identity
\[
 \sum_{e\in E(Q)}\big(R(G-e)-R(G)\big)
 =\sum_{v\in V(Q)}\big[-\epsilon_{d_v}S_v+
        a_{d_v}(c_vT_v-b_vS_v)\big].
\]
Indeed the coefficient before rearrangement at $v$ is
$a_{d_v}[(c_v-1)S_v+c_vT_v]-S_v/2$. By edge minimality,
each summand on the left exceeds $\mu(G-e)-\mu(G)>0$, since restoring
the edge shortens the distance between its endpoints. Some $c_vT_v-b_vS_v$
must therefore be positive. Finally, the tree case is already settled,
so an edge-minimal counterexample would have a component $Q$ of this kind.
\end{proof}
This proposition concerns a graph chosen minimal over all counterexamples.
Deleting a nonbridge of a bridgeless graph can create bridges, so the
proposition alone does not establish the conjecture for all bridgeless graphs.
\subsection*{An obstruction to monotone cycle deletion}
For $X=R$ or $X=R-\mu$, call the deletion of a nonbridge $uv$
admissible if $X(G-uv)\le X(G)$. Even among bridgeless graphs there
need not exist an admissible sequence from a graph to a spanning tree.

Take a clique of order $20$, and at each of its twenty vertices attach
ten pairwise disjoint pairs of new vertices. Each pair is joined by an
edge and both its vertices are joined to their common clique vertex.
Every edge lies in a cycle. The order is $420$ and each clique vertex
has degree $39$. Consider a state in which just some pair edges (those
between two new vertices) have been deleted. The order and clique
degrees stay fixed. Call each clique vertex with its twenty new vertices
a cluster. The following exhausts the nonbridges in every such state.

\emph{Clique edge.} Deleting it removes $1/39$ from $R$. At each
endpoint, the reciprocal square-root sum of the other neighbors is at
least $20/\sqrt2+18/\sqrt{39}$. Thus $R$ increases by at least
\[
 2(1/\sqrt{38}-1/\sqrt{39})(20/\sqrt2+18/\sqrt{39})-1/39
 > 1/30 .
\]
Indeed $1/\sqrt{38}-1/\sqrt{39}=1/[\sqrt{38\cdot39}(\sqrt{38}+\sqrt{39})]>1/500$,
and $20/\sqrt2+18/\sqrt{39}>13+2=15$, using $\sqrt{39}<25/4$.
Thus the estimate exceeds $30/500-1/39>1/30$.
Only distances between the two endpoint clusters increase, by one
each. Thus $\mu$ increases by $21^2/\binom{420}2<1/100$, so neither
admissibility condition holds.

\emph{An edge from a clique vertex to a new vertex.} If its pair
edge was already deleted, this edge is a bridge. Otherwise deleting
it increases the weight of the corresponding pair edge from $1/2$ to
$1/\sqrt2$. Other retained edge weights can only increase, so the net
gain in $R$ is at least $1/\sqrt2-1/2-1/\sqrt{78}>7/10-1/2-1/8=3/40$.
Only distances from the affected new vertex can increase, and they
increase by at most one since its sibling provides a detour of one
extra step. Thus $\mu$ increases by at most $2/420<1/100$.

\emph{An intact pair edge.} Its deletion changes $R$ by
$2(1-1/\sqrt2)/\sqrt{39}-1/2<0$ and cannot lower $\mu$.
These are exactly the admissible nonbridges. Deleting any such edge
returns to the type of state considered above. By induction all
admissible sequences delete only pair edges. Even after the last
deletion the original clique is intact, so no admissible sequence
reaches a spanning tree. This obstruction is to the proposed proof
strategy. The graph itself satisfies the conjecture (it is a block graph).

\section*{What remains}
For a shortest path $P$, set $k(P)=\max_v d(v,P)$, and minimize over
all such paths to obtain $k(G)$. If a counterexample exists, then it
satisfies $\Delta\ge101$, $b\ge156$ and $k(G)\ge19$ simultaneously.
It also has $n<\max\{48h,2h^2\}$ with $h=2+2b/9$.
Choose a counterexample with the fewest edges. By Proposition~\ref{p07-prop:bridgemin}, every cyclic component after removal of the bridges is incident to a bridge, with the stated imbalance of incident weights. On each of its nonbridges
one must have $R(G-e)-R(G)>\mu(G-e)-\mu(G)$, and
Corollary~\ref{p07-cor:longnonbridge} bounds the detour length on every such edge.

The sign of the explicit lower bound~\eqref{p07-eq:combined} cannot settle
the question on every graph. On $K_n$ a maximum-leaf tree is a star.
Every chord has $C_e=1$, and summing the capacities on its $n-1$ spokes
gives $S(T)\le(n-1)/2$. Equality is attained for $n\ge3$ by putting
weight $1/(n-2)$ on each chord. Thus the right-hand side is
$\sqrt{n-1}-2-(n-1)(n-2)/9+1/n$, which tends to minus infinity;
the actual margin is $n/2-1$.

This failure persists even when all three numerical cutoffs are exceeded.
In the clique with three length-$q$ arms at a common vertex, take the
clique order also equal to $q$. Then $n=4q$, $\Delta=q+2$,
$b=(q-1)(q-2)/2$ and $k(G)\ge q$. Every choice of tree satisfies
$S(T)/\binom n2\le\mu(T)\le(n+1)/3$. The quadratic term $-2b/9$
therefore dominates the other terms of~\eqref{p07-eq:combined}. The block-graph
theorem proves the conjecture on these examples. An induction using metric
savings would need to bound the actual loss $R(T)-R(G)$, or restrict to a
structurally reduced counterexample for which the available estimates suffice.

\paragraph{Proof status (G1--G6).}
G1 identifies Fajtlowicz's named 1988 conjecture. G2 fixes the index and
the mean over distinct vertex pairs as in Section~1. For G3 we located
published proofs in restricted cases; the literature check did not locate
a resolution in full generality. G4 is established on the stated classes
and remains open for general connected graphs in this work. For G5 the
new cutoffs hold uniformly in $n$ and follow from the stated lemmas and
reproducible arithmetic. For G6 we attribute the underlying tree and
diameter estimates above; priority for the precise extensions has not
been established.
\begin{thebibliography}{99}
\bibitem{p07-BE} B. Bollob\'as and P. Erd\H{o}s. Graphs of extremal weights.
\emph{Ars Combinatoria} \textbf{50} (1998), 225--233.
\bibitem{p07-COS} D. G. Corneil, S. Olariu and L. Stewart.
Asteroidal triple-free graphs. \emph{SIAM J. Discrete Math.}
\textbf{10} (1997), 399--430.
\href{https://doi.org/10.1137/S0895480193250125}{Publisher}.
\bibitem{p07-CPSMean} M. Cygan, M. Pilipczuk and R. \v{S}krekovski.
Relation between Randi\'c index and average distance of trees.
\emph{MATCH Commun. Math. Comput. Chem.} \textbf{66} (2011), 605--612.
\href{https://match.pmf.kg.ac.rs/electronic_versions/Match66/n2/match66n2_605-612.pdf}{Journal PDF}.
\bibitem{p07-Dal} C. Dalf\'o. On the Randi\'c index of graphs.
\emph{Discrete Mathematics} \textbf{342} (2019), 2792--2796.
\href{https://doi.org/10.1016/j.disc.2018.08.020}{Publisher}.
\bibitem{p07-DLS} Z. Dvo\v{r}\'ak, B. Lidick\'y and R. \v{S}krekovski.
Randi\'c index and the diameter of a graph.
\emph{European Journal of Combinatorics} \textbf{32} (2011), 434--442.
\href{https://doi.org/10.1016/j.ejc.2010.12.002}{Publisher}.
\bibitem{p07-Faj} S. Fajtlowicz. On conjectures of Graffiti.
\emph{Discrete Mathematics} \textbf{72} (1988), 113--118.
\href{https://doi.org/10.1016/0012-365X(88)90199-9}{Publisher}.
\bibitem{p07-KE} M. H. Khalifeh and A.-H. Esfahanian. Inset edges effect and
average distance of trees. Preprint, 2020.
\href{https://arxiv.org/abs/2008.05674}{arXiv:2008.05674}.
\bibitem{p07-Stev} D. Stevanovi\'c. Maximizing Wiener index of graphs with fixed
maximum degree. \emph{MATCH Commun. Math. Comput. Chem.} \textbf{60}
(2008), 71--83.
\href{https://match.pmf.kg.ac.rs/electronic_versions/Match60/n1/match60n1_71-83.pdf}{Journal PDF}.
\bibitem{p07-YL} L. You and B. Liu. On a conjecture between Randi\'c index
and average distance of unicyclic graphs. \emph{Filomat} \textbf{28}
(2014), 767--773. \href{https://doi.org/10.2298/FIL1404767Y}{Publisher}.
\end{thebibliography}

% END SOURCE BODY Fajtlowicz_Threshold_Forests
\unitstart{Paper 08}{p08}{Leaf attachment and bridges}{Leaf attachment and bridge gluing in Fajtlowicz's conjecture}{}
\makeatletter\@removefromreset{theorem}{section}\@addtoreset{theorem}{section}\makeatother\renewcommand{\thetheorem}{\thesection.\arabic{theorem}}
% BEGIN SOURCE BODY Fajtlowicz_Leaf_and_Bridge_Extrema
\begin{abstract}
We determine the largest possible increase of mean distance when a single leaf
is attached to a rooted connected graph of fixed order. This sharp local extremum
excludes every leaf with a degree-two neighbor from a minimum counterexample to
Fajtlowicz's conjecture. We also determine the supremum of the loss in total
Randi\'c index when joining two nontrivial graphs by a bridge. The resulting
rooted distance test restricts both sides of every internal bridge in a minimum
counterexample. A complementary spanning-forest estimate establishes a stronger
modified-index inequality on cacti, planar graphs of girth at least six,
triangle-free outerplanar graphs and graphs obtained by subdividing every edge.
These conclusions do not resolve the unrestricted conjecture.
\end{abstract}
\section{The transfer problem}
All graphs are finite, simple, undirected and connected unless a subgraph is
explicitly introduced. Write $n=|V(G)|$, $d_v=d_G(v)$ for degree,
$d_G(u,v)$ for distance and $W(G)=\sum_{\{u,v\}}d_G(u,v)$ for the
sum over unordered pairs of distinct vertices. For $n\ge2$ put
\[
 \mu(G)=W(G)/\binom{n}{2},\quad
 R(G)=\sum_{uv\in E(G)}(d_ud_v)^{-1/2},\quad
 M(G)=\sum_{uv\in E(G)}1/\max(d_u,d_v).
\]
Fajtlowicz asked whether $R\ge\mu$ on all connected graphs~\cite{p08-Faj}.
The tree and unicyclic cases are known~\cite{p08-Tree,p08-Uni}.
The present paper asks two extremal questions independently of the conjecture.
How much can attaching a leaf increase $\mu$? How much can inserting a bridge
reduce the sum of the two original indices $R$? Sharp bounds for these quantities
give reductions applicable to a putative counterexample with the fewest edges.
A third argument uses the rank of a forest to produce an unconditional
graph-class theorem.

Two elementary bounds used below are $R(G)\ge\sqrt{n-1}$ and $\mu(G)\le(n+1)/3$.
For the first, on a degree pair $i\le j$ the ratio of the edge weight to
$1/i+1/j$ is $\sqrt{j/i}/(1+j/i)\ge\sqrt{n-1}/n$.
The incidence terms sum to $n$. For the second, pass to any spanning tree,
then remove a leaf. Its distances in nondecreasing order to the other vertices
are bounded by $1,2,\ldots,n-1$. Induction on the order gives
$W(G)\le\binom{n+1}3$. For $n=2$ one has $R=\mu=1$.
\section{The sharp increase in average distance on attaching a leaf}
For a connected rooted graph $(H,u)$ of order $q\ge2$, form $H^+$ by
attaching a new leaf to $u$.
\begin{theorem}\label{p08-thm:leafresponse}
Writing $p=\lfloor q/2\rfloor$, one has
\[
 \max_{(H,u):\,|H|=q}\bigl(\mu(H^+)-\mu(H)\bigr)
   =\frac{5p+1}{6(2p+1)}
   =\frac5{12}-\frac1{4(2p+1)}.
\]
An extremizer consists of singleton distance layers from its root,
followed by a terminal layer of $p$ vertices that form a clique and
are all adjacent to the preceding singleton.
\end{theorem}
\begin{proof}
Let $W$ be the unordered distance sum in $H$, $\tau$ the sum of distances
from the root, and $s_0=1,s_1,\ldots,s_\ell$ the positive layer sizes.
Put $k_i=\sum_{j=i}^\ell s_j$. We have $\tau=\sum_{i=1}^\ell k_i$. Attaching the leaf adds $q+\tau$ to $W$, giving
\[
 \mu(H^+)-\mu(H)=\frac2{q+1}\bigl(1+\tau/q-\mu(H)\bigr).
\]
Each pair in different layers contributes at least the difference
in its indices, and distinct vertices in the same layer contribute
at least one. Thus $2W\ge 2\sum_i k_i(q-k_i)+\sum_i s_i(s_i-1)$.
Equality is attainable for every set of layer sizes, by filling all
intralayer edges and all edges between consecutive layers. Therefore
it suffices to maximize
\[
 B=\sum_i f(k_i)-\sum_i s_i(s_i-1),\qquad f(k)=k(2k-q-1).
\]
Indeed the mean-distance increase is at most
$2(1+B/[q(q-1)])/(q+1)$.

For fixed $\ell$, if two layers other than layer zero contain more
than one vertex, redistribute their excess vertices between them,
leaving each layer nonempty. For layers $i,j$, $B$ is a quadratic
in the transferred amount with quadratic coefficient $2|i-j|-2\ge0$.
Its maximum occurs at an endpoint. Repetition leaves all excess
$E=q-1-\ell$ in a single layer $j$. Put $t=\ell-j$. The tail counts
then omit the $E$ consecutive integers $t+1,\ldots,t+E$ from
$1,\ldots,q-1$. We must minimize
\[
 C(E,t)=\sum_{k=t+1}^{t+E} f(k)+E(E+1),\qquad
 0\le E\le q-2,\quad 0\le t\le q-2-E.
\]
For $E=0$ this value is zero. Completing the square gives
\[
 C(E,t)=A(E)+2E(t-t_E)^2,\quad
 t_E=(q-1-2E)/4,\quad
 A(E)=E(4E^2+24E-3q^2-6q+17)/24.
\]
Let $D(E)$ be the minimum over allowed integer $t$. For $E\le p-1$
the nearest integer to $t_E$ lies in the allowed range (for $E=0$
the choice is immaterial). Thus $D(E)=A(E)+2E\delta_E^2$,
where $\delta_E=\operatorname{dist}(t_E,\mathbb Z)\le1/2$.
Since $\delta_{E+2}=\delta_E$ and $A(E+2)-A(E)$ is increasing for
$E\ge0$, for $0\le E\le p-3$ we obtain
\[
 D(E+2)-D(E)\le A(p-1)-A(p-3)+1
 =\begin{cases}-(4p+5)/4,&q=2p,\\-2(p+1),&q=2p+1.\end{cases}
\]
For $p\ge2$, $t=1$ minimizes at $E=p-2$ and $t=0$ at $p-1$, giving
$D(p-2)-D(p-1)=-f(1)-2(p-1)=q-2p+1>0$.
For $E\ge p-1$, $t=0$ always minimizes and
$C(E,0)-C(E-1,0)=E(2E-q+1)$ for $E\ge p$. The latter is
nonnegative for both parities. This also treats $p=1$, when $E\ge0$.
The global minimum is at $E=p-1,t=0$. Substituting in
$B=\sum_{k=1}^{q-1}f(k)-C(E,t)$ and summing $k$ and $k^2$ gives the
claimed value. The layer construction attains all the distance bounds.
\end{proof}

\begin{corollary}\label{p08-cor:leafsupport}
In an edge-minimal counterexample of order $n\ge3$, a leaf adjacent
to a vertex $u$ of degree $d\ge2$ must satisfy
\[
 \frac1{\sqrt d}
 -\left(\frac1{\sqrt{d-1}}-\frac1{\sqrt d}\right)
      \sum_{w\sim u,\ w\ne v} \frac1{\sqrt {d_w}}
 < \frac{5p+1}{6(2p+1)},\qquad p=\lfloor(n-1)/2\rfloor .
\]
In particular, no leaf is adjacent to a vertex of degree two.
If a degree-three vertex is adjacent to exactly one leaf, the
unordered degree pair of its other neighbors must be $(2,2)$ or
$(2,3)$.
\end{corollary}
\begin{proof}
The displayed left side is exactly the change in $R$ when the deleted
leaf $v$ is restored. Minimality requires it to be strictly less than
the change in $\mu$. For degree two the left side is at least
$\sqrt2-1$ in the three-vertex path, which satisfies the conjecture
directly; in every other graph its other neighbor has degree at least
two and the increase is at least $1/2$. For $d=3$ with a single leaf,
all pairs except the stated ones satisfy
$\sum_{w\sim u,\,w\ne v}d_w^{-1/2}\le 1/\sqrt2+1/2$.
Their increase in $R$ exceeds $5/12$: use
$1/\sqrt3>5773/10000$ and $1/\sqrt2<7072/10000$, both by squaring,
to bound it below by
$5773/10000-(1299/10000)(12072/10000)>5/12$.
\end{proof}
For finite-order checks the closed form can also be compared directly
with the layer reduction by the following integer certificate.
\begin{verbatim}
for q in range(2,151):
    p=q//2
    f=lambda k:k*(2*k-q-1)
    sums=[0]
    for k in range(1,q): sums.append(sums[-1]+f(k))
    c=min(sums[t+E]-sums[t]+E*(E+1)
          for E in range(q-1) for t in range(q-1-E))
    assert (12*(2*p+1)*(q*(q-1)+sums[-1]-c)
            ==(q+1)*q*(q-1)*(5*p+1))
\end{verbatim}
If a graph of order at least three is enlarged by attaching one vertex
to an existing leaf, its Randi\'c margin increases by more than $1/12$.
Indeed the old neighbor of that leaf has degree at least two, so the
same computation gives an increase of at least $1/2$ in $R$. Repetition
allows arbitrary pendant-path extensions of every graph on which the
conjecture is already established.
Even a base graph $H$ with unknown sign of its margin admits an
unconditional extension statement. If $|H|=q\ge3$, $H$ has a leaf, and
$L\ge1$ vertices are appended by such extensions, then
\[
 R(G)-\mu(G)>R(H)-\mu(H)+L/12.
\]
It therefore suffices to take
$L\ge4(q+1)-12\sqrt{q-1}$, by the elementary bounds
$R(H)\ge\sqrt{q-1}$ and $\mu(H)\le(q+1)/3$.

The surviving degree pair $(2,3)$ cannot be eliminated using only the
single-leaf increment. Take an extremizer $K$ of order $q$ above, rooted
at its singleton endpoint $u$. Attach at $u$ a new neighbor $w$ carrying
two private leaves to obtain $H$. For $q$ large the original path
neighbor of $u$ has degree two and $w$ has degree three. The three new
vertices change $\tau(u)$ by five, $W$ by $O(q^2)$ and the order by
three. In the exact leaf-increment formula these changes are $O(1/q)$:
$\tau(u)=O(q^2)$, $W=O(q^3)$ and their respective coefficients are
$O(q^{-2})$ and $O(q^{-3})$. Hence the increment in $\mu$ when a leaf
is added to $u$ tends to $5/12$, while the increment in $R$ equals
$1/\sqrt3-1/2+1/3=1/\sqrt3-1/6<5/12$, since
$1/\sqrt3<7/12$. For pair $(2,2)$ use one private leaf at $w$ instead.
Its $R$ increment is even smaller. Thus these local configurations
give actual failures of monotone transfer, without producing
counterexamples to the conjecture.
\section{Joining two rooted graphs by a bridge}
Let $A,B$ be disjoint connected graphs of orders $a,b\ge2$, rooted at
$x,y$ respectively, and let $G$ be obtained by adding the bridge $xy$.
Put $n=a+b$ and write $\tau_A(x)=a^{-1}\sum_{u\in A}d_A(x,u)$,
including $u=x$; define $\tau_B(y)$ analogously. Put
\[
 \alpha=\frac{\binom a2}{\binom n2},\quad
 \beta=\frac{\binom b2}{\binom n2},\quad
 \gamma=\frac{ab}{\binom n2},\quad c=1-1/\sqrt2 .
\]
Then $\alpha+\beta+\gamma=1$. Write $R_A=R(A)$, $\mu_A=\mu(A)$,
and use analogous notation for $B$.
\begin{lemma}\label{p08-lem:bridgesharp}
The bridge attachment cost is bounded sharply by
\[
 R_A+R_B-R(G)<c.
\]
Moreover $\mu(G)=\alpha\mu_A+\beta\mu_B+
\gamma(1+\tau_A(x)+\tau_B(y))$.
\end{lemma}
\begin{proof}
Write $p=d_A(x)$, $q=d_B(y)$, $P=p+1$ and $Q=q+1$.
Writing $S_A(x)$ for the sum of the Randi\'c weights on the edges at
$x$ in $A$, its exact loss equals $(1-\sqrt{p/P})S_A(x)$.
Since $S_A(x)\le\sqrt p$, the loss is at most
\[
 u(P)=\sqrt{P-1}-\frac{P-1}{\sqrt P}
     =\frac{1}{\sqrt P}
       \frac{\sqrt{1-1/P}}{1+\sqrt{1-1/P}}<\frac1{2\sqrt P}.
\]
Adding the new bridge therefore gives total loss at most
$u(P)+u(Q)-1/\sqrt{PQ}$. Assume $P\le Q$. If $P=2,3,4$, the last
two terms have negative sum, and direct evaluation gives
$u(P)\le u(2)=c$. If $P\ge5$, put $z=1/\sqrt P$, $w=1/\sqrt Q$.
Both lie in $[0,1/2]$, where $(z+w)/2-zw\le1/4<c$.
To approach $c$, attach an endpoint of a $K_2$ to the center of a star
with $q$ leaves and let $q\to\infty$. The distances between vertices in
opposite components all pass through $xy$, giving the stated formula.
\end{proof}
\begin{theorem}\label{p08-thm:rootsglue}
Suppose $R_A\ge\mu_A$, $R_B\ge\mu_B$ and
\[
 \gamma\big[(\tau_A(x)-R_A)+(\tau_B(y)-R_B)\big]
 \ \le\ \beta R_A+\alpha R_B-\gamma-c .
\]
Then $R(G)>\mu(G)$. In particular, for $n\ge102$, it suffices
that $\tau_A(x)\le R_A$ and $\tau_B(y)\le R_B$.
\end{theorem}
\begin{proof}
Subtract the distance identity from the lower bound on $R(G)$.
The result exceeds
\[
 \beta R_A+\alpha R_B-\gamma-c
   -\gamma[(\tau_A(x)-R_A)+(\tau_B(y)-R_B)] .
\]
For the last assertion use $R(X)\ge\sqrt{|X|-1}$ and
$\gamma\le n/[2(n-1)]$. If $a,b\ge5$, then
$\beta R_A+\alpha R_B-\gamma-c\ge2-3\gamma-c
\ge(n-4)/[2(n-1)]-c>0$ for $n\ge10$.
If $\min(a,b)\le4$, both $R_A,R_B\ge1$ and $\gamma\le8/(n-1)$,
so the same expression is at least $1-16/(n-1)-c>0$ for $n\ge102$.
\end{proof}
It follows that, in an edge-minimal counterexample of order at least
102, at every bridge with two nonsingleton sides the displayed
criterion fails strictly. In particular one endpoint has average
distance to its own side larger than the entire Randi\'c index of that
side. Conversely every graph satisfying $\mu_A\le R_A$ admits a root
with $\tau_A(x)\le R_A$, since averaging $\tau_A$ over $x$ gives
$(1-1/a)\mu_A$. Thus any two graphs satisfying the conjecture can
be joined at suitable roots with preservation of the conjecture
when $a,b\ge2$ and $n\ge102$. Safe roots need not include every vertex.
For example, form a broom by adding $k$ leaves at the far end of a
length-$k$ path, with root at the opposite tip. Then the mean distance
from this root, including the root itself, is
$[k(k+1)/2+k(k+1)]/(2k+1)=3k/4+O(1)$,
while $R=k/2+O(\sqrt{k})$. Joining two such brooms at their roots
shows that even among trees both roots can fail the test.

\begin{corollary}\label{p08-cor:treeofregular}
Take arbitrary disjoint connected graphs $H_1,\ldots,H_k$, each
of order $n_i\ge2$ and total order $n$. Join them following any tree on
the $k$ indices, choosing the two endpoints of each joining edge
arbitrarily. The resulting graph satisfies $R>\mu$ if
\[
 \sum_{i=1}^k (n_i/2-R(H_i)) + c(k-1)<(n-2)/6 .
\]
In particular the conclusion holds whenever all $H_i$ are regular
and $n>2$, with different degrees and different orders allowed.
\end{corollary}
\begin{proof}
Join the pieces one at a time using Lemma~\ref{p08-lem:bridgesharp}. The
result has $R\ge\sum_i R(H_i)-c(k-1)$, which exceeds $(n+1)/3$
under the hypothesis. The universal bound $\mu\le(n+1)/3$ finishes.
For regular components the sum is zero, $k\le n/2$, and $c<1/3$.
\end{proof}
For a family with all three numerical parameters unbounded, follow a
three-armed tree of arm length $h$ with one $K_{q,q}$ at every tree
vertex. Use any fixed port in each copy for all incident bridges.
The maximum degree is at least $q$, the circuit rank is
$(3h+1)(q-1)^2$, and no shortest path enters the interiors of all
three arms. Its distance from some terminal copy is at least $h$.
\section{A forest question on the vertices of degree at least three}
Let $G$ be finite, connected, and simple, of order $n\ge3$. Let $L$
count its leaves, $b=|E(G)|-n+1$, and $J=G[\{v:d_G(v)\ge3\}]$.
Write $r=|V(J)|-c(J)$, taking $r=c(J)=0$ for the empty graph. All
weights below use degrees in $G$. Set
\[
 P=\sum_{uv\in E(J)}\frac2{\min(d_u,d_v)},\qquad
 B=\sum_{uv\in E(J)}\left(1-\frac2{\min(d_u,d_v)}\right).
\]
\begin{theorem}
If $P\le r$, then $M\ge\mu-1+2/n$ and $R>\mu$.
\end{theorem}
\begin{proof}
Extend any spanning forest of $J$ to a spanning tree $T$ of $G$.
Exactly $|E(J)|-r$ of the omitted edges are internal to $J$.
Each other omitted edge has an endpoint of degree two: an edge
incident to a degree-one vertex is a bridge. This endpoint becomes a
leaf of $T$, and distinct omitted edges cannot use the same degree-two
endpoint. Hence $T$ has at least $s=L+b-|E(J)|+r$ leaves.
The identity
\[
 2M=n-L+1+B-b
\]
follows by summing $n-2M=\sum_{uv}|1/d_u-1/d_v|$.
Indeed inserting $1/2$ between $1/d_u$ and $1/d_v$ gives a total of
$\sum_v|1-d_v/2|=L+|E(G)|-n$, before subtracting the triangle
inequality slack, which is precisely $B$.
Thus if $T$ has $I$ nonleaves, $2M-(I+1)\ge r-P\ge0$.
Counting unordered pairs crossing edges of $T$ gives
\[
 \mu(G)\le\frac{2(n-I)}n+\frac{(I-1)n}{2(n-1)}.
\]
Subtracting from $(I+1)/2+1-2/n$ leaves
$(I-1)(2/n-1/[2(n-1)])\ge0$. This proves the assertion about $M$.
For completeness put $y=\sqrt\Delta$. The edge identity
\[
 (d_u^{-1/2}-d_v^{-1/2})^2
 =(\sqrt{\max(d_u,d_v)/\min(d_u,d_v)}-1)(R_{uv}-M_{uv}),
\]
where $R_{uv},M_{uv}$ denote the edge weights,
yields $(y+1)R\ge n+(y-1)M$. Using $\mu\le(n+1)/3$ proves
$(y+1)(R-\mu)\ge (n-2)(1/3-(\sqrt{n-1}-1)/n)>0$.
\end{proof}
The arithmetic condition has an alternative form. If
$t_v=d_G(v)-d_J(v)$, it reads
\[
 \sum_{v\in J}\frac{t_v}{d_v}
 \ \ge\ c(J)+\sum_{uv\in E(J)}
  \big|d_u^{-1}-d_v^{-1}\big|.
\]
Indeed $2/\min(i,j)=i^{-1}+j^{-1}+|i^{-1}-j^{-1}|$.
The reverse strict inequality is necessary for any counterexample.

\begin{corollary}
The strong modified-index inequality holds whenever
$2|E(J)|\le3r$. In particular it holds for all cactus graphs, all
planar graphs of girth at least six, and all triangle-free outerplanar
graphs. In these classes $R>\mu$ for $n\ge3$.
\end{corollary}
\begin{proof}
Every term of $P$ is at most $2/3$. Treat components of $J$ by blocks,
counting bridges as two-vertex blocks. Rank and size are additive over
blocks. A bridge uses one edge and one unit of rank. For a cactus
every other block is a cycle of length $l\ge3$, so
$2l\le3(l-1)$. For planar girth at least six, a nonbridge block with
$v$ vertices and $e$ edges has $f=e-v+2$ faces, each bounded by a
cycle of length at least six. Thus $2e\ge6f$ and $2e\le3v-6$.
For triangle-free outerplanar graphs, embed a nonbridge block with all
vertices on the outer face. The boundaries are cycles since the block
is 2-connected. The outer face has length $v$, each bounded face has
length at least four, and $2e\ge v+4(f-1)$ gives $2e\le3v-4$.
All three classes are inherited by subgraphs, so the counts apply to
$J$. Singleton blocks have zero rank and size.
\end{proof}
For examples simultaneously beyond all three previous numerical cutoffs,
glue $q\ge3$ otherwise disjoint hexagons at one common vertex $z$
and attach three internally disjoint pendant paths of length $h$ there.
Then $b=q$, $\Delta=2q+3$ and every path misses the interior of some
arm, giving geodesic covering radius at least $h$. In particular
$q=156,h=19$ gives $b=156$, $\Delta=315$, $k\ge19$ and a nonclique
block.

In a graph obtained by subdividing every edge of any connected simple graph
at least once, $J$ is independent, hence $P=r=0$. This permits graphs of
arbitrarily high degree and unrestricted topology before subdivision.
\begin{corollary}
The strong modified-index bound also holds if every block is either a clique,
a cycle, a planar graph of girth at least six, or a triangle-free outerplanar
graph. The block types may be mixed freely.
\end{corollary}
\begin{proof}
A single clique is direct; every nonsingleton clique block in a graph
with more than one block contains a cut vertex. Repeatedly remove a noncut
vertex $z$ from a clique block of current order $q\ge3$ until that block
has order two or has only cut vertices. Its neighbors have degree at least
$q-1=d_z$. A neighbor $v$ of degree $d_v$ loses the weight $1/d_v$ from
the removed edge. Decreasing its degree on any one of its $d_v-1$ surviving
incidences can increase that edge's weight by at most $1/[d_v(d_v-1)]$.
Lower endpoint degrees successively on edges affected twice. Hence
$M(G-z)\le M(G)$. All earlier cut vertices remain cut vertices, since
the block retains at least two vertices and all incidences with other
blocks.
In the reduced graph, each remaining clique block of order $q\ge3$
contains only cut vertices. Their full degrees are at least $q$, so
the clique is contained in $J$ and its total contribution to $P$ is
at most $\binom q2(2/q)=q-1$, exactly its graphic rank.
In the remaining block types the induced high-degree subgraph has
$2e/3\le r$ by the preceding argument, and bridges also pass this test.
Ranks of the intersections of $J$ with the blocks add, since every cycle
lies in one block. The theorem therefore supplies a spanning tree $T'$
of the reduced graph with $I'+1\le2M(G')$. Restore the deleted vertices
in reverse order, each as a leaf at a retained cut vertex of its clique.
Such a vertex was internal in $T'$ and stays internal. Hence the
completed spanning tree still has $I'$ internal vertices and
$I'+1\le2M(G')\le2M(G)$. The tree distance estimate proves the result.
\end{proof}

\section{The remaining bridge condition}
Suppose a counterexample of order at least 102 with the fewest possible edges
exists. A leaf then satisfies Corollary~\ref{p08-cor:leafsupport}, and every bridge
with two nonsingleton sides violates the quantitative inequality of
Theorem~\ref{p08-thm:rootsglue}. The forest test additionally requires $P>r$.
The preceding research note~\cite{p08-Notes} gives the further necessary conditions
$\Delta\ge101$, $|E|-|V|+1\ge156$ and distance at least 19 from every
geodesic to some vertex, as well as the existence of a bridge.
A graph may have only pendant bridges, in which case the nonsingleton-side test
is vacuous. Neither the rooted excesses at the other bridges nor the remaining
leaf supports have been ruled out in general. This paper claims no priority for
its exact extremal statements.
\begin{thebibliography}{9}
\bibitem{p08-Faj} S. Fajtlowicz. On conjectures of Graffiti.
\emph{Discrete Math.} \textbf{72} (1988), 113--118.
\href{https://doi.org/10.1016/0012-365X(88)90199-9}{Publisher}.
\bibitem{p08-Tree} M. Cygan, M. Pilipczuk and R. \v{S}krekovski.
Relation between Randi\'c index and average distance of trees.
\emph{MATCH Commun. Math. Comput. Chem.} \textbf{66} (2011), 605--612.
\href{https://match.pmf.kg.ac.rs/electronic_versions/Match66/n2/match66n2_605-612.pdf}{Journal}.
\bibitem{p08-Uni} Z. You and B. Liu. On a conjecture between Randi\'c index
and average distance of unicyclic graphs. \emph{Filomat} \textbf{28} (2014),
767--773. \href{https://doi.org/10.2298/FIL1404767Y}{Publisher}.
\bibitem{p08-Notes} Keio and Shaitetsu. \emph{Threshold forests and Fajtlowicz's
average-distance conjecture}. Research note, 13 November 2026.
\end{thebibliography}

% END SOURCE BODY Fajtlowicz_Leaf_and_Bridge_Extrema
\unitstart{Paper 09}{p09}{Subdividing a bridge}{Subdividing a bridge and joining weighted graph cores}{}
\makeatletter\@removefromreset{theorem}{section}\@addtoreset{theorem}{section}\makeatother\renewcommand{\thetheorem}{\thesection.\arabic{theorem}}
% BEGIN SOURCE BODY Fajtlowicz_Bridge_Subdivision_and_Weighted_Cores
\begin{abstract}
We determine, at each order, the maximum increase in mean distance when a
bridge is subdivided. From order six onward it is attained by a bridge
between two cliques of orders differing by at most one. The maximum is
strictly smaller than one half, which excludes adjacent degree-two vertices
in a minimum counterexample to Fajtlowicz's conjecture. A second reduction
eliminates two exceptional single-leaf configurations when all bridges are
pendant. Finally, a spanning-tree condition on two component graphs
guarantees the original conjecture after gluing at arbitrary vertices and
attaching any number of additional leaves. This includes pieces in which
the chosen gluing points are arbitrarily far from their own typical vertices.
The unrestricted conjecture remains open in this work.
\end{abstract}
\section{Two extremal questions}
For a finite connected simple graph $G$ of order $n\ge2$, write $d_v$ for
degree, $d_G(u,v)$ for distance, $W(G)=\sum_{\{u,v\}}d_G(u,v)$ and
$\mu(G)=W(G)/\binom n2$. Set
\[
 R(G)=\sum_{uv\in E(G)}(d_ud_v)^{-1/2},\qquad
 M(G)=\sum_{uv\in E(G)}1/\max(d_u,d_v).
\]
Fajtlowicz's conjecture asserts $R(G)\ge\mu(G)$~\cite{p09-Faj}.
Earlier results cover trees and unicyclic graphs~\cite{p09-Tree,p09-Uni}.
Our first question concerns distance only. Over all $n$-vertex graphs
and all bridges, how much can the mean distance increase when the bridge
is replaced by a two-edge path? The second question allows an arbitrary
probability measure on a component. How large can its expected pair distance
be compared with the Randi\'c index, and what happens to the bound at a bridge?

\section{The exact increase under subdivision}
Write $G_e$ for the graph obtained by subdividing a bridge $e$ once, and let
$B_n=\max_{G,e}(\mu(G_e)-\mu(G))$ over graphs of order $n$ with a bridge.
\begin{theorem}\label{p09-th:sub}
For $n=2,3$ one has $B_n=1/3$, and for $n=4,5$ one has $B_n=11/30$.
For all $n\ge6$, with $\eta=n\bmod2$,
\[
 B_n=\frac{n^3-n^2-4n+8-(n-5)\eta}{2n(n-1)(n+1)} .
\]
A bridge between $K_{\lfloor n/2\rfloor}$ and $K_{\lceil n/2\rceil}$
attains this value. In particular $B_n<1/2$ at every finite order and
$B_n\longrightarrow1/2$.
\end{theorem}
\begin{proof}
Let $A,B$ be the sides rooted at the endpoints, of sizes $a\ge b$, with
$n=a+b$ and $d=a-b$. For a rooted side $X$ let $T_X$ be the sum of
distances to its root and $W_X$ its unordered distance sum. Before
subdivision $W(G)=W_A+W_B+bT_A+aT_B+ab$. The old pairs crossing the
bridge gain $ab$ in total, and distances to the new vertex sum to
$n+T_A+T_B$. Thus, on putting
$J=\tfrac12n(n-1)(n+1)(\mu(G_e)-\mu(G))$, one obtains
\[
 J=2n(n-1)+(n-5)ab+F_A+F_B,\quad
 F_X=(2|X|-n-1)T_X-4W_X .
\]
Let $J_0(d)$ denote the same quantity for two complete sides. Direct
substitution gives
\[
 4J_0(d)=n^3-n^2-4n+8-(n-5)d^2 .
\]
On $B$ the coefficient of $T_B$ is $-d-1$, and the clique minimizes
both $T_B$ and $W_B$. Only the larger side could improve upon this benchmark.

Write $k_i$ for the number of vertices in $A$ at distance at least $i$
from the root. Then $T_A-(a-1)=\sum_{i\ge2} k_i$, and
\[
 W_A-\binom a2\ \ge\ \sum_{i\ge2} k_i(a-k_{i-1}).
\]
Indeed a pair in layers $j<l$ is counted $[l-j-1]_+$ times on the
right. The excess of its distance over one is at least that large.
For the indices with $k_i>0$ the numbers $a-k_{i-1}$ are distinct
positive integers. Writing $[z]_+=\max(z,0)$ yields
\[
 J-J_0(d)\le n\sum_{u\ge1}[d-1-4u]_+
            \le n(d-1)^2/8\quad(d\ge1).
\]
The final bound is the integral of the decreasing positive part on
$[0,\infty)$. For $d=0$ the excess is nonpositive. At even $n\ge10$,
$n-5\ge n/2$ so the larger-side excess cannot outweigh the $d^2$
penalty in $J_0$. At odd $n\ge11$ and $d\ge1$ use
$n(d-1)^2/8\le(n-5)(d^2-1)/4$ for the same conclusion. Equality
is attained at the smallest possible $d$, with both sides complete.

For clarity the finite end of the calculation is included below.
At each fixed list of layer sizes the corresponding clique layers,
fully joined across successive layers, minimize $W_X$ and preserve
$T_X$. The program enumerates every positive composition of $|X|-1$,
hence maximizes $F_X$ exactly. The singleton is included separately.
\begin{verbatim}
from fractions import Fraction as Q
from itertools import combinations
def best_side(m,n):
    if m==1: return 0
    best=None
    for k in range(m-1):
        for cuts in combinations(range(1,m-1),k):
            x=(0,)+cuts+(m-1,)
            s=[1]+[v-u for u,v in zip(x,x[1:])]
            T=sum(i*c for i,c in enumerate(s))
            W=sum(c*(c-1)//2 for c in s)
            W+=sum(c*s[j]*(j-i) for i,c in enumerate(s)
                   for j in range(i+1,len(s)))
            f=(2*m-n-1)*T-4*W
            best=f if best is None else max(best,f)
    return best
for n in range(2,13):
    J=max(2*n*(n-1)+(n-5)*a*(n-a)+best_side(a,n)
          +best_side(n-a,n) for a in range(1,n))
    actual=Q(2*J,n*(n*n-1))
    expected=(Q(1,3) if n<=3 else Q(11,30) if n<=5 else
          Q(n**3-n*n-4*n+8-(n-5)*(n%2),2*n*(n*n-1)))
    assert actual==expected
\end{verbatim}
For $n=6,7,8,9$ the maxima are respectively $41/105,17/42,53/126,77/180$,
agreeing with the balanced-clique value. For $n\ge6$, the numerator of
$1/2-B_n$ with the displayed common denominator equals
$n^2+3n-8+(n-5)\eta>0$.
\end{proof}
If $i,j$ are the old endpoint degrees, the exact change in $R$ on
subdividing $e$ is
\[
 \Delta R=\frac1{\sqrt{2i}}+\frac1{\sqrt{2j}}-\frac1{\sqrt{ij}}
 =\frac12-(1/\sqrt2-1/\sqrt i)(1/\sqrt2-1/\sqrt j).
\]
In particular $\Delta R=1/2$ if one endpoint has degree two.
Every minimum counterexample must survive suppressing a degree-two
vertex on a bridge path. Suppression restores a simple connected graph
since the new edge cannot already exist. Therefore no such counterexample
has adjacent degree-two vertices on a bridge. More quantitatively,
if the suppressed graph has $n$ vertices then
$(1/\sqrt2-1/\sqrt i)(1/\sqrt2-1/\sqrt j)>1/2-B_n$ is necessary.
\section{Eliminating the two exceptional leaf supports on a cycle}
Let $G$ have the fewest edges among hypothetical counterexamples.
An edge $ab$ which lies in a cycle can be deleted while preserving
connectedness. If $R(G-ab)\le R(G)$ then
$R(G)\ge R(G-ab)\ge\mu(G-ab)>\mu(G)$, a contradiction.
\begin{lemma}\label{p09-lem:cyc23support}
No cycle contains adjacent degree-two vertices. Suppose a degree-three
vertex $u$ is incident with exactly one leaf and its other two neighbors
have degrees $(2,2)$ or $(2,3)$. Then every edge from $u$ to one of its
degree-two neighbors is a bridge.
\end{lemma}
\begin{proof}
For adjacent degree-two vertices on a cycle, let $d_1,d_2$ be the
degrees of their other neighbors, necessarily at least two. On deleting
the edge the index change is
\[
 (1-1/\sqrt2)(1/\sqrt{d_1}+1/\sqrt{d_2})-1/2
 \le\sqrt2-3/2<0 .
\]
For the second statement let $w$ be a degree-two neighbor and $y$ its
other neighbor. Suppose that $uw$ is on a cycle. Then $wy$ is also
on this cycle. The first assertion gives $d_y\ge3$.
Write $z$ for the other nonleaf neighbor of $u$.
Deleting $uw$ changes the index by
\[
 (1/\sqrt2-1/\sqrt3)(1+1/\sqrt{d_z})
 +(1-1/\sqrt2)/\sqrt{d_y}-1/\sqrt6 .
\]
Since $d_z\ge2$ and $d_y\ge3$, this is at most
$1/\sqrt2+1/2-3/\sqrt6<0$.
For the last inequality multiply by $\sqrt6$ to obtain
$\sqrt3+\sqrt{3/2}<3$, which follows on squaring from
$\sqrt2<3/2$. Thus $uw$ must be a bridge.
\end{proof}
The argument permits $y=z$, since the two retained edges at that
vertex are distinct. In fact the second conclusion requires only
$d_z\ge2$. Together with the single-leaf response bound of \cite[Corollary 2.2]{p09-Prev}, the specified two neighbor pairs exhaust all possibilities.
Consequently, in a minimum counterexample in which all bridges are
pendant, every support of a unique leaf has degree at least four.
\begin{corollary}
In an edge-minimal counterexample both neighbors of any degree-two
vertex lying on a cycle have degree at least four. A degree-four or
degree-five vertex with exactly one leaf cannot be adjacent on a
cycle to a degree-two vertex. If all bridges are pendant, every leaf
support has degree at least four.
\end{corollary}
\begin{proof}
A degree-three vertex on a cycle has at most one leaf. The preceding
deletion estimate applies whether it has zero leaves or one, since
of its two other neighbors at least one lies on the cycle and has
degree at least two. This proves the first assertion. For the second,
write $d\in\{4,5\}$ for the support degree. On deleting its edge to
a degree-two vertex the change is at most
\[
 (1/\sqrt{d-1}-1/\sqrt d)(1+(d-2)/\sqrt2)
       +(1-1/\sqrt2)/2-1/\sqrt{2d}<0,
\]
where the other neighbor of the degree-two vertex has degree at least
four. For an exact check, lower bound all individual reciprocal square
roots by $10^{-5}\lfloor\sqrt{10^{10}/j}\rfloor$, and upper bound by
this number plus $10^{-5}$. Substitution with outward bounds gives
less than $-1/100$ for $d=4$, and less than $-1/250$ for $d=5$.
Finally a degree-three support with two leaves is joined to the rest
through a nonpendant bridge, except for the four-vertex star. A
degree-two support is already excluded by the sharp leaf response.
\end{proof}
\begin{theorem}
Let $u,w,v$ occur consecutively on a cycle of an edge-minimal
counterexample, where $d_w=2$. Put $m=d_u\le d_v=M$, and let $l$
count the leaves adjacent to $u$. Then
\[
 \frac{l}{m-1}>1+\sqrt2-2\sqrt{m/M}\ \ge\sqrt2-1,
 \qquad M<2m .
\]
At the larger endpoint the analogous lower bound is
$1+\sqrt2-2\sqrt{M/m}$.
\end{theorem}
\begin{proof}
Set $a=1/\sqrt2$, $b=1-a$, $x=l/(m-1)$. On deleting $uw$
the Randi\'c change is at most
\[
 (1/\sqrt{m-1}-1/\sqrt m)(m-1)(a+bx)
        +b/\sqrt M-a/\sqrt m .
\]
This change is strictly positive by minimality. The product of the
first difference with $m-1$ equals
$\sqrt{1-1/m}/[\sqrt m(1+\sqrt{1-1/m})]<1/(2\sqrt m)$.
It follows that $a+bx>2a-2b\sqrt{m/M}$ and
$a/b=1+\sqrt2$ gives the claim. Since $u$ has another neighbor
on its cycle, $x<1$, forcing $\sqrt{m/M}>1/\sqrt2$.
Deleting $wv$ gives the symmetric bound.
\end{proof}
\begin{proposition}
Let $uv$ be a nonbridge in a minimum counterexample. With full
degrees $i=d_u,j=d_v$, let
$A_u=(i-1)^{-1}\sum_{z\sim u,\ z\ne v}d_z^{-1/2}$ and define
$A_v$ symmetrically. Then $\sqrt j A_u+\sqrt i A_v>2$.
In particular, suppose all bridges are pendant and let $S$ be the
nonempty set of vertices incident to leaves. If every vertex of $S$
has degree at least the largest degree of a nonleaf outside $S$,
then the nonleaves outside $S$ are independent.
\end{proposition}
\begin{proof}
The exact change of $R$ on deleting $uv$ is
\[
 (1/\sqrt{i-1}-1/\sqrt i)(i-1)A_u+
 (1/\sqrt{j-1}-1/\sqrt j)(j-1)A_v-1/\sqrt{ij}.
\]
It must be positive. Use the strict coefficient bounds from the
previous proof and multiply by $2\sqrt{ij}$ to obtain the assertion.
For the corollary let $X$ be the nonleaves outside $S$. If $G[X]$
has an edge, choose $uv\in E(G[X])$ minimizing $d_ud_v$. This edge
is a nonbridge since all bridges are pendant. Each neighbor of $u$
other than $v$ lies in $X$ or $S$ and has degree at least $d_v$,
by minimality of the product or by the assumption on $S$. Thus
$A_u\le1/\sqrt j$. The symmetric estimate at $v$ yields a
contradiction.
\end{proof}
As one consequence, an all-pendant-bridge minimum counterexample
cannot have all its leaves at one vertex of maximum degree. Otherwise the
other nonleaves are independent by the proposition, and the core
contains no cycle. This would make $G$ a tree, for which the
conjecture holds.

Combining Lemma~\ref{p09-lem:cyc23support} with Theorem~\ref{p09-th:sub}, every
degree-two vertex in any minimum counterexample is isolated within
the subgraph induced by vertices of degree two. If all bridges are
pendant, there is no leaf support of degree three, and every nonleaf
degree-two vertex must lie on a cycle.

\begin{corollary}\label{p09-cor:bridgebalance}
There is no degree-three leaf support with exactly one leaf in a
minimum counterexample having only pendant bridges. More generally,
if such a support $u$ occurs in a minimum counterexample of order
$N\ge102$, a degree-two neighbor $w$ lies on a bridge path $u,w,y$.
Suppress $w$. The two sides of the resulting bridge in the
$(N-1)$-vertex graph have orders $a,b$ satisfying
\[
                |a-b|<\frac{N-1}{\sqrt2}.
\]
Equivalently each side contains more than $(2-\sqrt2)/4$ of those
vertices.
\end{corollary}
\begin{proof}
There is at least one degree-two neighbor by the preceding reduction.
Its two incident edges are bridges, and $d_y\ge3$ by the isolation
of degree-two vertices and the absence of degree-two leaf supports.
For $n=N-1$ and $d=|a-b|$, the proof of Theorem~\ref{p09-th:sub} gives
$J\le J_0(0)-(n-10)d^2/8$, even when $n$ is odd. The case $d=0$
follows directly. The index gain on restoring $w$ is strictly greater
than $1/\sqrt6$, since $d_u=3$ and the coefficient of $1/\sqrt{d_y}$
is $1/\sqrt2-1/\sqrt3>0$. Minimality would require the distance
gain to exceed this. Writing $U_0=2J_0(0)/[n(n^2-1)]$, we obtain
\[
 \frac{d^2}{n^2}<
 \frac{4(n^2-1)(U_0-1/\sqrt6)}{n(n-10)}
 <\frac{4(1/2-1/\sqrt6)}{1-10/n}<\frac12 .
\]
Here $U_0<1/2$ and $n\ge101$; the last numerical inequality follows
already from $1/\sqrt6>2/5$ and $n\ge100$. Suppression leaves a
simple connected graph with fewer edges, as in Theorem~\ref{p09-th:sub}.
\end{proof}

\section{Pair distances with freely chosen attachment masses}
For any connected graph $H$, let $Q(H)=\max_p\sum_{x,y}p_xp_y d_H(x,y)$
where the maximum is over probability measures on $V(H)$. For every
spanning tree, each edge contributes $2t(1-t)\le1/2$ for the mass $t$
on either side. Consequently $Q(H)\le(|H|-1)/2$.
A spanning tree with $s\ge2$ leaves and $I=h-s$ internal vertices gives
the sharper bound
\[
 Q(H)\le (I-1)/2+2-2/s\qquad (h=|H|\ge3).
\]
Indeed the leaf masses have sum $b\le1$ and their total contribution is at
most $2b(1-b/s)\le2-2/s$ by Cauchy--Schwarz. All other edges are internal.

Denote by $s(H)$ the maximum number of leaves in a spanning tree, and
$I(H)=h-s(H)$. The condition used for a component is
\begin{equation}\label{p09-eq:core}
           2M(H)\ge I(H)+1 .
\end{equation}
Then with $s=s(H)$, $I=I(H)$,
\begin{equation}\label{p09-eq:score}
 R(H)\ge (I-1)/2+\sqrt s,\qquad
 R(H)-Q(H)\ge\sqrt s-2+2/s>0 .
\end{equation}
To verify this, put $y=\sqrt{\Delta(H)}$. On an edge with degrees
$i\le j$, $(i^{-1/2}-j^{-1/2})^2=(\sqrt{j/i}-1)(1/\sqrt{ij}-1/j)$.
Summation gives $(y+1)R\ge h+(y-1)M$. Extending the full star at a
maximum-degree vertex to a spanning tree gives $s\ge\Delta(H)$.
Substitution of~\eqref{p09-eq:core} now gives
$R\ge(I+1)/2+(s-1)/(y+1)\ge(I-1)/2+\sqrt s$.

Examples satisfying~\eqref{p09-eq:core} include all connected regular graphs
of order at least three and the sparse or mixed block classes of
\cite[Section 4]{p09-Prev}. For an explicit rank criterion, let $J$ be the
subgraph induced by the vertices of full degree at least three and let
$r=|V(J)|-c(J)$, where $c(J)$ counts components. The inequality
$\sum_{uv\in E(J)}2/\min(d_u,d_v)\le r$ implies~\eqref{p09-eq:core}:
extend a spanning forest of $J$ to a spanning tree. Each of the
$b-|E(J)|+r$ omitted edges outside $J$ supplies a distinct new leaf
of old degree two, where $b=|E(H)|-h+1$. With $L$ the original leaf
count, the identity $2M=h-L+1-b+
\sum_{uv\in E(J)}(1-2/\min(d_u,d_v))$ proves the claim. One obtains
the identity by summing $h-2M=\sum_{uv}|1/d_u-1/d_v|$ and inserting
$1/2$ between the two reciprocals. In particular all cacti pass the test,
since $2|E(J)|/3\le r$ by adding the ranks of their edge and cycle blocks.

\paragraph{Attaching leaves.}
Add $t_v$ leaves at each $v\in V(H)$ to form $G$, put $L=\sum t_v$,
$N=h+L$ and $p_v=(1+t_v)/N$. Then
\begin{align}
 \mu(G)&=\frac{N}{N-1}\sum_{x,y}p_xp_y d_H(x,y)+2L/N,\label{p09-eq:proj}\\
 R(G)-R(H)&\ge\sum_v(\sqrt{d_v+t_v}-\sqrt{d_v})
 \ge\sqrt{N-1}-\sqrt{h-1}.\label{p09-eq:gain}
\end{align}
The first identity projects each leaf to its support. Among distinct
vertex pairs each new pendant edge contributes exactly $2/N$. To prove
the second, put $z_v=\sqrt{d_v/(d_v+t_v)}$. The old edge weight loses a
fraction $1-z_uz_v\le(1-z_u)+(1-z_v)$. The old incident weights at $v$
sum to at most $\sqrt{d_v}$, and the new edges add $t_v/\sqrt{d_v+t_v}$.
This gives the first lower bound. For the last one replace every $d_v$
by $h-1$ and use the decreasing increments of the square root to combine
the $t_v$. If $N\ge17$, this gain is at least $2L/N$ since
$N\ge4\sqrt{N-1}$. The remaining normalization cost in~\eqref{p09-eq:proj}
is at most $Q(H)/(N-1)\le1/2$.

\begin{proposition}\label{p09-prop:single}
Every core satisfying~\eqref{p09-eq:core}, decorated with arbitrary numbers
of leaves at arbitrary vertices, satisfies $R\ge\mu$.
\end{proposition}
\begin{proof}
If $s\ge4$, the second bound of~\eqref{p09-eq:score} is at least $1/2$.
Together with~\eqref{p09-eq:proj}--\eqref{p09-eq:gain} this proves the result
for $N\ge17$. For $s\le3$ the core has maximum degree at most three.
If the final maximum degree is at least $101$, the gain at that vertex
alone is at least $\sqrt{101}-\sqrt3>5/2$, covering both the normalization
and the pendant cost. In the remaining small-degree cases, including
$N<17$, we use the bounded-degree theorem of~\cite[Theorem 3.2]{p09-Threshold}.
\end{proof}

\begin{theorem}\label{p09-th:two}
Take any two graphs $H_1,H_2$ of order at least three satisfying
\eqref{p09-eq:core}. Join any chosen vertex of the first to any chosen vertex
of the second by a bridge, obtaining $H$. Attach arbitrary numbers of
additional leaves at arbitrary vertices of $H$ to form $G$. Then
$R(G)\ge\mu(G)$. There is no condition on either chosen vertex.
\end{theorem}
\begin{proof}
Write $s_i=s(H_i)$, $I_i=I(H_i)$,
$C=(I_1-1+I_2-1)/2$ and $c=1-1/\sqrt2$. In each side take a
maximum-leaf spanning tree and connect the two trees by the bridge.
For any probability measure on $H$, the internal tree edges on the two
sides contribute at most $C$ to the expected distance. At most two leaf
edges have a gluing vertex as their leaf, each costing at most $1/2$.
All remaining leaf edges still isolate distinct singletons, so their
total contribution is at most two. The bridge costs at most $1/2$. Hence
\[
 Q(H)\le C+7/2,\qquad
 R(H)>C+\sqrt{s_1}+\sqrt{s_2}-c .
\]
For the latter inequality apply~\eqref{p09-eq:score} and the sharp
bridge-cost lemma~\cite[Lemma 3.1]{p09-Prev}. Here is its short derivation.
If the new endpoint degrees are $P\le Q$, each endpoint loses at most
$u(P)=\sqrt{P-1}-(P-1)/\sqrt P<1/(2\sqrt P)$ on its previous edges.
The new edge adds $1/\sqrt{PQ}$. For $P=2,3,4$ the net loss is
less than $u(P)\le u(2)=c$. For $P\ge5$, it is at most
$(x+y)/2-xy\le1/4<c$ where $x=1/\sqrt P$, $y=1/\sqrt Q$.

If $\sqrt{s_1}+\sqrt{s_2}\ge4+c$, these estimates give
$R(H)-Q(H)>1/2$. Equations~\eqref{p09-eq:proj}--\eqref{p09-eq:gain} pay
for all extra leaves whenever $N\ge17$. Otherwise $s_i\le8$
because $s_j\ge2$. Hence $\Delta(H)\le9$ after the join.
If $\Delta(G)\ge101$, the gain at the largest support is at least
$\sqrt{101}-3>7$. Even the bounds $R(H)-Q(H)>2\sqrt2-c-7/2>-1$,
$Q(H)/(N-1)\le1/2$ and $2L/N<2$ suffice. All remaining cases follow
from the same bounded-degree theorem.
\end{proof}

Both attachment points can fail the earlier root test. On each side take
a fixed four-cycle, attach at the same cycle vertex a path of length $k$ and
$k$ private leaves, and use the far tip of the path as the gluing vertex.
The core is a cactus. The mean distance from the tip, including the root
itself, is $3k/4+O(1)$ as $k\to\infty$, by summing the path distances
and the $k$ distances $k+1$ to the leaves. The index is
$k/2+O(\sqrt k)$ from the path, cycle and pendant edge weights. Thus both
tips eventually exceed their side's index. Theorem~\ref{p09-th:two} uses
the common probability budget on the remaining leaves to avoid both
root-transmission bounds.

\begin{corollary}
Let $k\ge2$ graphs satisfying~\eqref{p09-eq:core}, each with maximum
spanning-tree leaf count $s_i\ge4$, be joined along an arbitrary
tree on these $k$ pieces, at any vertices. Arbitrary extra leaves may
again be attached at any vertices. The resulting graph has $R\ge\mu$.
\end{corollary}
\begin{proof}
Put $C=\sum_i(I_i-1)/2$. For the joined core $H$ write $P$ for the
number of distinct chosen gluing vertices which are leaves of their
respective spanning trees. Then $P\le2(k-1)$. Counting edges in the
joined spanning tree exactly as above gives
$Q(H)\le C+(k-1+P)/2+2$. Repeated bridge insertion gives
$R(H)>C+\sum_i\sqrt{s_i}-c(k-1)$. Subtracting, we obtain
\[
 R(H)-Q(H)>\sum_i(\sqrt{s_i}-2)+(1/2-c)(k-1).
\]
For $k\ge4$ this is greater than $1/2$. If $k=3$ and some $s_i\ge5$,
it also exceeds $1/2$. These cases follow by the leaf extension argument
above (orders below 17 use the bounded-degree result).
The case $k=2$ is Theorem~\ref{p09-th:two}. Finally, if $k=3$ and
$s_i=4$ for all three, $\Delta(H)\le6$. When the final degree reaches
101 at a support, the additional gain there exceeds
$\sqrt{101}-\sqrt6>5/2$; otherwise the bounded-degree theorem applies.
\end{proof}

\section*{Remaining case}
The absence of consecutive degree-two vertices is necessary for any
minimum counterexample. A single degree-two bridge subdivision with
both ends of degree at least three may still have an index increment
below $B_n$. The leaf-density bound on a cycle leaves higher-degree
supports with sufficiently many leaves. For a join with general
components, inequality~\eqref{p09-eq:core} can fail. For example, give
every vertex of $K_{2,k}$ one private leaf, with $k\ge4$. Its $k+2$
original vertices are cut vertices and are precisely the internal
vertices of every maximum-leaf spanning tree. Its modified index
is $2+k/3$, so $2M<I+1$. This example restricts the method and leaves
the original global question open.
\begin{thebibliography}{9}
\bibitem{p09-Faj} S. Fajtlowicz. On conjectures of Graffiti.
\emph{Discrete Math.} \textbf{72} (1988), 113--118.
\href{https://doi.org/10.1016/0012-365X(88)90199-9}{Publisher}.
\bibitem{p09-Tree} M. Cygan, M. Pilipczuk and R. \v{S}krekovski.
Relation between Randi\'c index and average distance of trees.
\emph{MATCH Commun. Math. Comput. Chem.} \textbf{66} (2011), 605--612.
\href{https://match.pmf.kg.ac.rs/electronic_versions/Match66/n2/match66n2_605-612.pdf}{Journal}.
\bibitem{p09-Uni} Z. You and B. Liu. On a conjecture between Randi\'c index
and average distance of unicyclic graphs. \emph{Filomat} \textbf{28} (2014),
767--773. \href{https://doi.org/10.2298/FIL1404767Y}{Publisher}.
\bibitem{p09-Threshold} Keio and Shaitetsu. \emph{Threshold forests and
Fajtlowicz's average-distance conjecture}. Research note, 2026.
\bibitem{p09-Prev} Keio and Shaitetsu. \emph{Leaf attachment and bridge gluing
in Fajtlowicz's conjecture}. Research note, 2026.
\end{thebibliography}

% END SOURCE BODY Fajtlowicz_Bridge_Subdivision_and_Weighted_Cores
\unitstart{Paper 10}{p10}{Vertex deletion and leaf covers}{Vertex deletion and leaf covers for the Randi\'c index and average distance}{}
\makeatletter\@removefromreset{theorem}{section}\@addtoreset{theorem}{section}\makeatother\renewcommand{\thetheorem}{\thesection.\arabic{theorem}}
% BEGIN SOURCE BODY Fajtlowicz_Vertex_Deletion_and_Leaf_Covers
\begin{abstract}
We study Fajtlowicz's inequality $R(G)\ge\mu(G)$ for the average distance
of a connected graph. An identity summing all vertex deletions gives
$a>\sqrt{n-1}/5$ cut vertices in any hypothetical vertex-minimal
counterexample of order $n$. A sharper estimate also accounts for all
pendant-edge weights. Independently, the conjectured inequality holds
after attaching at least one private leaf at every vertex of any connected
core of average degree at most eight, including every connected planar
core. For cores of arbitrary average degree $\dd$, it suffices to attach
at least $k\ge2$ leaves per vertex with $4k^2-k\ge\dd$.
We also give a gluing inequality at four-way cut vertices.
These results leave the unrestricted conjecture unresolved.
\end{abstract}

\section{A deletion identity}
Graphs are finite, simple and undirected. Write $n=|V(G)|$, $d_u=\deg_G u$,
$w_{uv}=1/\sqrt{d_ud_v}$, and $R=\sum_{uv\in E(G)}w_{uv}$. On a connected
graph of order at least two, $\mu=W/\binom n2$ with $W$ the unordered
sum of all pair distances. For $K_1$ put $\mu=R=0$.
Fajtlowicz's conjecture is $R\ge\mu$~\cite{p10-Fajt}. Published proofs
include trees~\cite{p10-Trees}, unicyclic graphs~\cite{p10-Uni} and a degree
condition for sufficiently large graphs~\cite{p10-Dalfo}.
A vertex-minimal counterexample below is a connected counterexample all of
whose connected proper induced subgraphs satisfy the conjecture.
Every hypothetical counterexample contains one of these.

For $n\ge3$ let $P$ denote the sum of the weights of pendant edges.
For $e=uv$ and $z\notin\{u,v\}$ put
$a_{u,z}=\sqrt{d_u/(d_u-1)}$ if $uz$ is an edge and $a_{u,z}=1$ otherwise.
The first case implies $d_u\ge2$. Write $R(G-z)$ also for disconnected
deletions, summing over the surviving edges.
\begin{lemma}\label{p10-del}
On every connected graph of order $n\ge3$ one has
\[
 \sum_z\bigl(R(G)-R(G-z)\bigr)
 =R(G)+\frac P2+
 \frac12\sum_{uv\in E(G)}w_{uv}
       \sum_{z\notin\{u,v\}}(a_{u,z}-a_{v,z})^2 .
\]
\end{lemma}
\begin{proof}
For a fixed $uv$, its retained weight upon deletion of a third vertex is
$w_{uv}a_{u,z}a_{v,z}$. The sum of $a_{u,z}^2$ over all such third
vertices equals $n-1-\mathbf1_{d_u=1}$, since exactly $d_u-1$ of them
are other neighbors. Expanding the square difference gives the identity
edge by edge. For $n\ge3$ each pendant edge has exactly one leaf endpoint.
\end{proof}

Let $K$ be the set of cut vertices, $a=|K|$, $S=V(G)\setminus K$,
$b=|S|=n-a$ and $D=\diam G$. Let $R_S,R_K$ sum the original weights
over edges internal to $S,K$, respectively.
\begin{theorem}\label{p10-cut}
Every vertex-minimal counterexample has $a>D/5>\sqrt{n-1}/5$.
If $a\le b$ it satisfies the stronger bound
\begin{equation}\label{p10-eq:cut}
 D\left(1-\frac{2ab}{n(n-1)}\right)+P
 < 2a+\sum_{z\in K}\frac{d_{G[K]}(z)}{d_z}\le3a .
\end{equation}
\end{theorem}
\begin{proof}
A deleted vertex $z$ costs at most the sum of its incident weights,
since surviving edge weights can only rise. Subtract these bounds for
$z\in K$ from Lemma~\ref{p10-del} to get
\[
 \sum_{z\in S}(R(G)-R(G-z))\ge R_S-R_K+P/2 .
\]
Put $T_z=\sum_y d(z,y)$ and $\tau_z=T_z/(n-1)$.
Deletion can increase distances among retained vertices. Consequently
for $z\in S$ the decrease in $\mu$ is at most
$2(\tau_z-\mu)/(n-2)$. Minimality and $R(G)<\mu(G)$ imply that the
decrease in $R$ is strictly smaller than this bound for every $z\in S$.

For $a\le b$, write $W_S,W_K$ for unordered internal distance sums in
the original metric and $Z$ for the cross sum. Then
\[
 n(n-1)\sum_{z\in S}(\tau_z-\mu)=2aW_S+(a-b)Z-2bW_K .
\]
By the triangle inequality $(b-1)Z\ge aW_S$ (sum over each $z\in K$
and unordered pair in $S$). Since $b\ge2$ and $W_S\le b(b-1)D/2$,
the right hand side is at most $ab(n-2)D/2$. Hence
\[
 R_S-R_K+P/2<\frac{abD}{n(n-1)}. \tag{1}
\]

For any vertex set $X$ we claim that
\begin{equation}\label{p10-res}
 \sum_{uv\in E(G-X)}\frac1{\max(d_u,d_v)}\ge D/2-|X|.
\end{equation}
Add $t$ fresh leaves at each $x\in X$. The old distances are preserved.
The modified-index diameter theorem~\cite[Theorem 3]{p10-DLS} gives at least
$D/2$ for the new sum of reciprocal maximum endpoint degrees.
As $t\to\infty$, original edges incident to $X$ contribute zero in the
limit, and the new edges at $x$ contribute $t/(d_x+t)\to1$.
This proves~\eqref{p10-res}, also for $X=\varnothing$ by the same cited theorem.
Apply it with $X=K$ to get $R_S\ge D/2-a$. On an edge,
$w_{uv}\le(1/d_u+1/d_v)/2$, so
$2R_K\le\sum_{z\in K}d_{G[K]}(z)/d_z\le a$.
Substitution in (1) proves~\eqref{p10-eq:cut}. Also $D\le n-1$ and (1) imply
$D/2-a\le R_S<R_K+a\le3a/2$, hence $D<5a$. If $a>b$, already $D<n<2a$.
Finally, on an edge the ratio of its weight to $1/d_u+1/d_v$ is
$\sqrt{x}/(1+x)$ with $x$ the larger degree divided by the smaller.
It is at least $\sqrt{n-1}/n$. Summing gives $R\ge\sqrt{n-1}$, and
the counterexample hypothesis gives $R<\mu\le D$.
\end{proof}
When the cut vertices are independent the last sum in~\eqref{p10-eq:cut}
vanishes. The conclusion concerns a vertex-minimal induced counterexample.
Deletion can create new cut vertices in a previously 2-connected graph.

\section{Leaf covers of a core}
Let $H$ be a connected graph of order $h\ge3$, maximum degree $\Delta$,
and average degree $\dd$. At each $v\in V(H)$ attach $t_v\ge0$ new
private leaves, obtaining $G$. Let $N=h+L$, $L=\sum_v t_v$ and
$p_v=(1+t_v)/N$. For a graph metric define
$Q(H)=\max_p\sum_{x,y}p_xp_y d_H(x,y)$ over probability measures.
Let $s$ be the largest possible number of leaves in a spanning tree of
$H$, and put $I=h-s$. For any tree edge separating mass $t$ from its
complement the distance contribution is $2t(1-t)\le1/2$. Leaf-edge
contributions together are at most two, as their singleton masses sum
to at most one. The $I-1$ internal edges therefore give
\[
 Q(H)\le (I-1)/2+2,\qquad Q(H)\le(h-1)/2.
\]
The $\Delta$ edges at a maximum-degree vertex extend to a spanning tree
whose branches each contain a leaf. In particular $I\le h-\Delta$.
The exact projection identity is
\[
 \mu(G)=\frac N{N-1}\sum_{x,y}p_xp_y d_H(x,y)+\frac{2L}{N}.
\]
Each pendant edge accounts for $N-1$ unordered pairs. Since
$Q(H)/(N-1)\le(h-1)/(2(N-1))\le2h/N$, it follows that
\begin{equation}\label{p10-distbd}
 \mu(G)\le Q(H)+2\le\min\{I/2+7/2,(h+3)/2\}.
\end{equation}
Adding a leaf at a vertex of current degree $d\ge1$ increases $R$ by
at least $\sqrt{d+1}-\sqrt d$: the old incident edges have total
weight at most $\sqrt d$ and retain a fraction $\sqrt{d/(d+1)}$.
Summing over additions also gives $R(G')-R(G)\ge\sqrt{N'-1}-\sqrt{N-1}$
if the order rises from $N$ to $N'$, by $d\le N-1$ at each step.
We may thus compute a lower bound on $R$ with only the minimum required
loads.

\begin{theorem}\label{p10-covers}
Suppose every vertex of $H$ receives at least one new leaf. Then
$R(G)\ge\mu(G)$ holds in either of the following cases:
\begin{enumerate}
\item $\dd\le8$;
\item each $t_v\ge k$ for an integer $k\ge2$, and $\dd\le4k^2-k$.
\end{enumerate}
The loads can be different and arbitrarily large. All connected simple
planar cores of order at least three satisfy the first condition.
\end{theorem}
\begin{proof}
We first handle $\Delta\le6$. Give every core vertex one leaf.
The contribution of its leaf and half of each incident core edge is at
least $[1+d/(2\sqrt7)]/\sqrt{d+1}\ge3/4$ for $1\le d\le6$, where $d$
is its original degree. Hence $R(G)\ge3h/4$. This is at least
$(h+3)/2$ when $h\ge6$. For $h=3,4,5$ we may assume $2h\le N\le16$,
since $\sqrt{N-1}\ge4$ at $N\ge17$. The finite check at the end compares
$3h/4+\sqrt{N-1}-\sqrt{2h-1}$ with
$N(h-1)/(2(N-1))+2(N-h)/N$.

Suppose now $\Delta\ge7$. For the second assertion retain just $k$ leaves
at each vertex. These pendant edges alone, by convexity of
$x\mapsto(x+k)^{-1/2}$, supply
\[
 R(G)\ge\sum_{v\in H}\frac{k}{\sqrt{d_v+k}}
 \ge\frac{hk}{\sqrt{\dd+k}}\ge h/2 .
\]
Since $I\le h-\Delta\le h-7$, (\ref{p10-distbd}) suffices.

For the first assertion retain one leaf at every vertex. Set
$U=\{v:d_v\ge38\}$, let $A$ be the vertices outside $U$ adjacent to it,
and $B=V(H)\setminus(U\cup A)$. If $U\ne\varnothing$ then
\[
 I\le |B|+3|U|-2 .
\]
Indeed, take a spanning tree in each component of $H[U]$, assigning each
vertex of $A$ as a leaf to one of its neighbors in $U$. Extend the forest
over all remaining vertices, adding at each step a leaf and losing at most
one. Joining its $c\le |U|$ components loses at most $2(c-1)$ leaves.
This gives $s\ge |A|-2(c-1)$ and proves the claim.

Set $c_0=23/1000$, $\rho=13/20$ and
$f(d)=[1+d/(2\sqrt{38})]/\sqrt{d+1}$. Assign a leaf weight to each vertex
and half of the incident core-edge weights to each vertex of $B$. This
does not exceed the available total. A vertex in $B$ receives at least
$f(d)$, since all its neighbors have degree at most 37. Adding
$c_0(d-8)$ to each assignment gives the following lower bounds:
\[
 \begin{array}{c|c|c}
   \text{set}&\text{expression bounded below}&\text{lower bound}\\ \hline
   B&f(d)+c_0(d-8)&1/2\\
   A&(d+1)^{-1/2}+c_0(d-8)&13/40\\
   U&(d+1)^{-1/2}+c_0(d-8)&17/20 .
 \end{array}
\]
The finite checks below prove the first two rows for $1\le d\le37$ and
the third at $d=38$. The third expression increases for $d\ge38$ by
differentiation. The sum of the added terms is nonpositive because
$\dd\le8$. If $U$ is empty this gives $R\ge h/2$, enough for $\Delta\ge7$.
Otherwise the weighted average of
$I\le |B|+3|U|-2$ and $I\le h-\Delta$, with weights $1-\rho,\rho$,
gives
\[
 R\ge |B|/2+13|A|/40+17|U|/20
 \ge I/2+1-\rho+\rho\Delta/2 .
\]
The last two terms exceed $7/2$ for $\Delta\ge38$.
This proves the first assertion by~\eqref{p10-distbd}; the planar corollary
uses $|E(H)|\le3h-6$.
\end{proof}
The proof also permits holes in the cover. For $\Delta\ge38$ let
$Z=\{v:t_v=0\}$. Relative to the one-leaf baseline, removing each
missing leaf costs at most $1/\sqrt{d_v+1}$, since surviving weights rise.
Keeping the exact adjustment $c_0(8-\dd)h$ in the preceding proof gives
the sufficient condition
\[
 \sum_{v\in Z}\frac1{\sqrt{d_v+1}}
 \le \frac{13\Delta}{40}-\frac{63}{20}
        +\frac{23}{1000}(8-\dd)h.
\]
For planar cores this allows positive densities of uncovered vertices
when the right hand side is large. This condition can also be checked
when $\dd>8$.

\paragraph{Integer verification of the finite inequalities.}
The function $x(j)$ below rounds $1/\sqrt j$ down; $y(j)$ rounds
$\sqrt j$ down. The omitted tail of either rounding is less than $1/S$.
\begin{verbatim}
from math import isqrt
from fractions import Fraction as F
S = 10**6
def x(j): return F(isqrt(S*S//j), S)
def y(j): return F(isqrt(S*S*j), S)
c = F(23,1000)
for d in range(1,38):
    assert (1+d*x(38)/2)*x(d+1)+c*(d-8) >= F(1,2)
    assert x(d+1)+c*(d-8) >= F(13,40)
assert x(39)+30*c >= F(17,20)
for d in range(1,7):
    assert (1+d*x(7)/2)*x(d+1) >= F(3,4)
for h in range(3,6):
    for N in range(2*h,17):
        lo = F(3*h,4)+y(N-1)-y(2*h-1)-F(1,S)
        hi = F(N*(h-1),2*(N-1))+F(2*(N-h),N)
        assert lo >= hi
\end{verbatim}

\section{Compensation at a cut vertex}
For a connected graph $J$ put $b(J)=R(J)-Q(J)$.
Suppose $v$ is a cut vertex of $J$ with $m_v$ components in $J-v$.
Choose one neighbor of $v$ in each component and write $\nu_v$ for their
uniform measure. For every $z$, the expected distance from $\nu_v$ to $z$
is at least $d(v,z)$: at most one component permits a decrease of one,
and all other representatives add one. Moving any original atom at $v$
to $\nu_v$ cannot reduce a cross term. Moving all cut-vertex atoms
independently thus gives, for every probability measure $p$,
\begin{equation}\label{p10-atom}
 Q(J)\ge \sum_{x,y}p_xp_y d_J(x,y)
       +\sum_{v\ {\rm cut}}2(1-1/m_v)p_v^2 .
\end{equation}
For a same-vertex term at $v$ the two independent representatives belong
to different components with probability $1-1/m_v$, at distance two.
For two distinct buckets, apply the preceding pointwise inequality first
to one bucket and then to the other.

\begin{theorem}
Connect $r\ge2$ disjoint connected graphs $H_i$ along a tree of $r-1$
bridges, producing $J$. Suppose each chosen endpoint splits its
\emph{original} $H_i$ into at least four components when removed.
Endpoints can be reused. Then
\[
 b(J)>\sum_{i=1}^r b(H_i)
       +\left(\frac9{20}-\frac1{\sqrt5}\right)(r-1).
\]
In particular, $\sum_i b(H_i)\ge1/2$ implies $R(J)\ge\mu(J)$.
This last conclusion survives arbitrary additional private leaves
whenever the final order is at least 17.
\end{theorem}
\begin{proof}
Given a probability distribution on $J$, project outside mass onto its
entry vertex in each piece. The pair distance is the sum of internal
projected distances and bridge crossing indicators, since outer
excursions cannot shorten any internal distance. For a bridge with side
masses $t,1-t$, the projected masses at its endpoints in the respective
pieces are at least $1-t,t$. At a reused endpoint the incoming masses
are disjoint and $(\sum_j t_j)^2\ge\sum_j t_j^2$.
By~\eqref{p10-atom}, the internal expected distances therefore save at least
$\tfrac32(t^2+(1-t)^2)$ per bridge compared with the sum of the $Q(H_i)$.
Subtract the bridge distance contribution $2t(1-t)$: the net saving
is at least $1/4$ per bridge.

On adding a bridge whose new endpoint degrees are $P,Q$, the old
weight loss at an endpoint is at most
$\sqrt{P-1}-(P-1)/\sqrt P<1/(2\sqrt P)$.
The new edge contributes $1/\sqrt{PQ}$. Thus the net index loss is
less than $(x+y)/2-xy$ where $x=P^{-1/2},y=Q^{-1/2}$.
Here $x,y\le1/\sqrt5<1/2$, and this expression is increasing in each
argument. Its maximum is $1/\sqrt5-1/5$. Combining the two changes
proves the result.

For a core of order $h$, $Q/(h-1)\le1/2$. Thus $b\ge1/2$ gives the
ordinary unweighted conclusion. If leaves are attached and $N\ge17$,
their index gain is at least $\sqrt{N-1}-\sqrt{h-1}\ge2(N-h)/N$
by $N\ge4\sqrt{N-1}$. This pays for the extra pendant contribution
in the projection formula, while $Q/(N-1)\le1/2$.
\end{proof}

\section*{The unresolved implication}
Vertex deletion yields the cut-vertex condition~\eqref{p10-eq:cut} in any
minimal obstruction, and leaf cover and cut-vertex gluing give direct
proofs for the stated classes. No argument here forces an arbitrary
graph with many cut vertices and relatively small pendant weight into
one of those classes. Paths already show that the cut-count and
pendant-weight restrictions can coexist in graphs satisfying the
conjecture. Their conjunction alone gives no contradiction.
\begin{thebibliography}{9}
\bibitem{p10-Fajt} S. Fajtlowicz, On conjectures of Graffiti, Discrete Math.
72 (1988), 113--118.
\bibitem{p10-Trees} M. Cygan, M. Pilipczuk and R. \v{S}krekovski,
Relation between Randi\'c index and average distance of trees,
MATCH Commun. Math. Comput. Chem. 66 (2011), 605--612.
\bibitem{p10-Uni} Z. You and B. Liu, On a conjecture between Randi\'c index
and average distance of unicyclic graphs, Filomat 28 (2014), 767--773.
\bibitem{p10-Dalfo} C. Dalf\'o, On the Randi\'c index of graphs, Discrete
Math. 342 (2019), 2792--2796.
\bibitem{p10-DLS} Z. Dvo\v{r}\'ak, B. Lidick\'y and R. \v{S}krekovski,
Randi\'c index and the diameter of a graph, preprint KAM--DIMATIA
2010--496, Theorem 3;
\href{https://iti.mff.cuni.cz/series/2010/496.pdf}{authors' text}.
\end{thebibliography}

% END SOURCE BODY Fajtlowicz_Vertex_Deletion_and_Leaf_Covers
\unitstart{Paper 11}{p11}{The weighted cut-vertex reduction}{A weighted cut-vertex reduction for the Randi\'c index and average distance}{}
\makeatletter\@removefromreset{theorem}{section}\@addtoreset{theorem}{section}\makeatother\renewcommand{\thetheorem}{\thesection.\arabic{theorem}}
% BEGIN SOURCE BODY Fajtlowicz_Weighted_Cut_Vertex_Reduction
\begin{abstract}
Fajtlowicz conjectured that the Randi\'c index of a connected graph is at least its average distance. We introduce a stronger inequality for arbitrary vertex probability distributions and prove that its pointwise deficit splits exactly at every cut vertex after padding both pieces with leaves. Its universal validity reduces to a single two-connected block with at most one private leaf at each vertex. We prove the resulting inequality for every block of diameter at most two, every regular block, certain irregular blocks of larger diameter, and every block of order at most eight. The bounded-order result uses an exact reproducible computation. These classes may be mixed within one graph of arbitrary order. The case of an unrestricted two-connected block with private leaves remains open in this argument.
\end{abstract}
\section{The inequality and the graph classes}
All graphs are finite, simple, undirected, and connected, of order $n\ge2$.
Put $d_v=\deg_G(v)$, let $d_G(x,y)$ be graph distance, and write
\[
 R(G)=\sum_{uv\in E(G)}\frac1{\sqrt{d_ud_v}},\qquad
 M(G)=\sum_{uv\in E(G)}\frac1{\max(d_u,d_v)},\qquad
 \mu(G)=\binom n2^{-1}\sum_{\{x,y\}}d_G(x,y).
\]
Fajtlowicz's average-distance conjecture is $R(G)\ge\mu(G)$~\cite{p11-Fajt,p11-Trees,p11-Uni}. For a probability vector $p$ on $V(G)$ and independent $X,Y$ with law $p$ let
\[
 L_G(p)=\mathbb E_p d_G(X,Y)+2\sum_x p_x^2,\qquad
 \Psi(G)=\max_p L_G(p)-1,\qquad Q(G)=\max_p\mathbb E_p d_G(X,Y).
\]
The diagonal coefficient $2$ is chosen because a pendant edge carrying probability $t$ contributes $2t(1-t)$ to expected distance. Its sum with the diagonal term $2t^2$ is linear in $t$.
A cyclic block is a maximal two-connected subgraph of order at least three.
A private leaf is a new degree-one vertex joined to a specified old vertex.
\begin{theorem}\label{p11-thm:main}
Suppose each cyclic block $H$ of $G$ satisfies at least one of the following conditions, with $h,\delta,\Delta$ its own order, minimum degree and maximum degree:
\begin{enumerate}
\item $\diam(H)\le2$, or $H$ is regular.
\item $h\le8$.
\item $\diam(H)=3$ and at most two vertices of $H$ are incident to edges of $G$ outside $H$.
\item $H$ is Hamiltonian and $h\delta/(\delta+\Delta)\ge h/4+1$.
\item $\delta\ge24$ and $\Delta\le2\delta$.
\item $M(H)\ge Q(H)+1$.
\end{enumerate}
Then $M(G)\ge\Psi(G)$ and $R(G)\ge\mu(G)$. Distances, degrees, and both indices in the conclusion are evaluated in the full graph. There is no restriction on the bridges, the number of cyclic blocks, or how many blocks contain the same vertex.
\end{theorem}
The proof of the bounded-order clause uses the exact certificate in the appendix. All other clauses have analytic proofs. For instance, every chain of diameter-two blocks satisfies the conclusion even when its diameter, the degrees at its junctions, and the sizes of its blocks are all unbounded.

\section{Why the weighted statement implies the original conjecture}
\begin{proposition}\label{p11-prop:transfer}
For any individual connected graph $G$, the inequality $M(G)\ge\Psi(G)$ implies $R(G)\ge\mu(G)$.
\end{proposition}
\begin{proof}
Using the uniform probability vector in the hypothesis gives
\[
 M\ge(1-1/n)\mu-1+2/n.
\]
In particular the coefficient here is $(n-1)/n$: the independent draws in $L$ include the diagonal. Write $k=\sqrt\Delta$. For degrees $a\le b$ on an edge,
\[
 (a^{-1/2}-b^{-1/2})^2
 =(\sqrt{b/a}-1)\bigl((ab)^{-1/2}-b^{-1}\bigr).
\]
Summing yields $n-2R\le(k-1)(R-M)$, whence
\[
 (k+1)(R-\mu)\ge n-2\mu-(k-1)\bigl(1+(\mu-2)/n\bigr).
\]
We recall $\mu\le(n+1)/3$. Pass to a spanning tree, whose distances dominate those of $G$. In a tree of order $j$, the $i$th smallest distance from a fixed vertex to another vertex is at most $i$. Its transmission is therefore at most $j(j-1)/2$. Remove a leaf, preserving all other distances, and sum this bound over $j=2,\ldots,n$ to obtain $\binom{n+1}3$ for the unordered distance sum.

The right side above decreases with $\mu$ and with $k$. Substitution of $(n+1)/3$ and $\sqrt{n-1}$ gives the lower bound
\[
 \frac{n(n-2)-(4n-5)(\sqrt{n-1}-1)}{3n}.
\]
For $n\ge5$ its numerator is nonnegative: square the two positive sides of $n(n-2)+4n-5\ge(4n-5)\sqrt{n-1}$. The difference is $(n-5)(n-2)(n^2-5n+5)\ge0$.
For $2\le n\le4$, the ratio of $(ab)^{-1/2}$ to $a^{-1}+b^{-1}$ is $\sqrt{x}/(1+x)$ for $x=b/a\in[1,n-1]$ and is at least $\sqrt{n-1}/n$. Summation gives $R\ge\sqrt{n-1}\ge(n+1)/3$ for these orders.
\end{proof}

\section{An exact law at each cut vertex}
Let $v$ be a cut vertex, and let $A,B$ partition $V(G)\setminus\{v\}$ into two nonempty unions of components of $G-v$. Write $a=|N(v)\cap A|$, $b=|N(v)\cap B|$. Let $G_A$ be the subgraph on $A\cup\{v\}$ with $b$ new leaves at $v$, and form $G_B$ using $a$ new leaves. Put $D_G(p)=L_G(p)-1-M(G)$.
\begin{theorem}[Exact splitting]\label{p11-thm:split}
For each distribution $p$ on $G$ there exist distributions $p_A,p_B$ on these padded sides with
\[
 D_{G_A}(p_A)+D_{G_B}(p_B)=D_G(p).
\]
In particular, if both pieces satisfy $M\ge\Psi$, so does $G$.
\end{theorem}
\begin{proof}
Every original edge keeps its endpoint degrees on its side. The new leaf edges across the two graphs have total $M$-weight $(a+b)/(a+b)=1$. Hence $M(G_A)+M(G_B)=M(G)+1$, the classical degree-padding identity~\cite[Lemma~4]{p11-DLS} also used in corpus manuscript 087~\cite{p11-VMRC}.

Put $\alpha=p(A)$, $\beta=p(B)$ and $\gamma=p(v)$. Keep all masses of vertices of $A$ on side $A$, retain $\gamma$ at its copy of $v$, and put $\beta$ on one new leaf there. On side $B$, keep the masses on $B$, give its copy of $v$ zero mass, and put $\alpha+\gamma$ on one new leaf. Both laws have total mass one.

First contract the new leaves to $v$ on both sides. The sum of the two projected expected distances equals the original expected distance: pointwise, project $x,y$ onto each side by mapping the opposite side to $v$. The two projected distances add to $d_G(x,y)$, including when one vertex is $v$. Undoing the contractions adds $2t(1-t)$ for each occupied new leaf of mass $t$. The squared masses on ordinary vertices sum to the original sum of squares, because just one copy of $v$ retains $\gamma$. Adding the new diagonal terms gives $2\sum t=2(\beta+\alpha+\gamma)=2$. Thus $L_{G_A}(p_A)+L_{G_B}(p_B)=L_G(p)+2$. Subtracting the constants and the $M$-weights proves the identity.
\end{proof}

\begin{lemma}[Support and deletion]\label{p11-lem:support}
There is a maximizer of $L_G$ whose distinct positive atoms have distance at least three. Deleting an unoccupied pendant leaf from a graph of order at least three keeps $L_G(p)$ fixed and can only decrease $M$.
\end{lemma}
\begin{proof}
Choose a maximizer of minimum support. Transfer mass between two positive coordinates at distance $d$, fixing their sum. The coefficient of the squared transfer parameter is $4-2d$. If $d\le2$, at least one endpoint of the interval of transfers attains the maximum and eliminates an atom. For deletion, consider the reverse attachment at a vertex of original degree $j\ge1$. Each old incident edge loses at most $1/[j(j+1)]$ and the new edge contributes $1/(j+1)$. Distances among surviving vertices stay fixed.
\end{proof}
\begin{theorem}[Reduction to one block]\label{p11-thm:reduction}
The assertion $M\ge\Psi$ on all connected graphs of order at least two is equivalent to its validity on every two-connected core with at most one private leaf per core vertex. The same reduction holds within any prescribed class of allowed cyclic blocks.
\end{theorem}
\begin{proof}
Split whenever each of $A,B$ contains an edge of $G-v$. Both padded sides have strictly fewer edges: for $G_A$ the decrease is precisely the number of edges inside $B$. Their cyclic blocks come from the original graph. Induction on the number of edges and Theorem~\ref{p11-thm:split} reduce the assertion to pieces admitting no such split.

If there are two cyclic blocks, take a separating vertex on the path between them in the block--cut tree. On each side a cyclic block contains an edge avoiding that vertex. With just one cyclic block, every attached component is a tree meeting the block in one vertex. If any branch extends two edges away from the block, split at that attachment point. Thus an unsplittable piece with a cycle is a single block with private leaves. In a minimum-support maximizing distribution at most one leaf at each vertex is occupied. Remove all surplus leaves using Lemma~\ref{p11-lem:support}.

For a tree of diameter at least four, a middle vertex of a four-edge path gives the required split. A tree of diameter at most two is a star (including $K_2$), with $M=\Psi=1$. A remaining tree of diameter three consists of adjacent centers with $a\ge b\ge1$ leaves. A separated support has at most two sites and two such sites are distance three apart. Hence $\Psi=3/2$, while $M=1+b/(b+1)\ge3/2$. This handles all terminal trees. The converse follows by taking the stated one-block graphs themselves.
\end{proof}

\section{Blocks that satisfy the weighted inequality}
\subsection{Blocks of small diameter}

\begin{theorem}\label{p11-th:diameter-two}
Let $H$ be a finite connected simple graph of minimum degree at least two and diameter at most two. Add any number of private leaves at any vertices, forming $G$. Then $F(G)\le M(G)+1$.
\end{theorem}
\begin{proof}
Write $F(G)=\max_p L_G(p)=1+\Psi(G)$. By Lemma~\ref{p11-lem:support}, choose a maximizer whose distinct atoms are at distance at least three, and remove all but one leaf at each loaded site, keeping the occupied leaf if there is one. Write $t$ for the number of loaded sites after this simplification.

A separated maximizing support contains at most one original vertex.
Its distance to a supported leaf is at most three, and distances
between supported leaves are at most four. In a distribution with mass $m$
on the new leaves and the remainder at one old vertex, the excess above
$2$ comes only from distances exceeding two, since the collision reward
is two on the diagonal. For $t\ge1$ the excess is at most
\[
 2m(1-m)+2\left(m^2-\sum_{\ell\ {\rm leaf}}p_\ell^2\right)
 \le 2m-2m^2/t .
\]
The same estimate applies to an all-leaf measure by setting $m=1$.
It follows that $F\le 2$ for $t=0$, $F\le5/2$ for $t=1$, and
$F\le4-2/t$ for $t\ge2$.

We first treat $t\le2$. In a minimum-degree-two graph $M(H)\ge2$
unless $H$ consists of an adjacent pair joined to a nonempty independent
set of common neighbors. To see this, take the two vertices of highest
degrees $d_1\ge d_2$. Edges incident to the first total one; counting
only previously uncounted edges at the second gives at least $1-1/d_2$.
If the two are nonadjacent, one gets two immediately. Any edge outside
the pair of stars adds at least $1/d_2$. If there is none, two-connectivity
is unnecessary: minimum degree two forces all remaining vertices to be
common neighbors. In the exceptional graph with $b\ge1$ common neighbors
one has $M(H)=2-1/(b+1)$. For $b=1$ this is a triangle: when $t\le2$,
$F\le5/2$, as at most two tips at mutual distance three can be selected
and no original vertex can be supported together with a tip.
For $b\ge2$, if any common neighbor receives a leaf, the first such
attachment, performed before any other, increases $M$ by $1/3$ since
its old neighbors both have degree $b+1\ge3$. Thus $M(G)\ge2$.
If leaves occur only on the adjacent pair, each loaded site is universal
in $H$ and the same $F\le5/2$ argument applies. These facts give
$M(G)+1\ge F(G)$ in all the cases $t\le2$. We will also use
$M(G)\ge2$ when $t=3$: the same argument applies to the exceptional
graph with $b\ge2$, and the fully decorated triangle has $M=2$.

For $t\ge3$, order the core vertices by full degree in $G$ and
take the three largest, forming a set $T$. Write $d$ for the third
largest degree, $e$ for the number of edges inside $T$, and
$j$ for the number of its loaded vertices. Since every loaded core
vertex has full degree at least three, $d\ge3$.
Count edges incident to $T$ in descending degree order, assigning each
edge to the first endpoint visited. At a vertex of degree $d_i$,
each edge to a still unvisited vertex has weight $1/d_i$. At most
$e/d$ in total is lost from the three unit contributions due to edges
already counted at earlier endpoints. Each of the $t-j$ pendant
edges at sites outside $T$ has weight at least $1/d$. Thus
\[
 M(G)\ge3+(t-j-e)/d .
\]
For $t\ge6$ this gives $M(G)\ge3$. For $t=5$ it gives
$M(G)\ge3-1/d\ge8/3$, sufficient since $F(G)\le 18/5$.

If $t=3$ and $e+j\le5$, the same bound gives $M(G)\ge7/3$, while
$F(G)\le10/3$. If $e+j=6$, all three loaded sites form a triangle.
All supported distances are then at most three, so $F(G)\le3$;
we already established $M(G)\ge2$ for $t=3$.

Finally, if $t=4$ and $e+j\le5$, then $M(G)\ge8/3$, which suffices
for $F(G)\le7/2$. Otherwise three of the loaded sites form a triangle.
In any separated support all pairs which could be at distance four
involve the tip of the fourth site. Let $x$ be the mass at that tip,
allowing $x=0$. Relative to the constant three on every ordered pair,
this gives
\[
 F(G)\le3+2x(1-x)-\sum_v p_v^2
 \le3+2x-3x^2\le10/3.
\]
The degree bound gives $M(G)\ge3-2/d\ge7/3$, completing the proof.
\end{proof}

\begin{proposition}\label{p11-prop:diam-three}
Let $H$ have minimum degree at least two and diameter three. Attach
any number of private leaves at a single vertex, or attach none.
The resulting graph $G$ satisfies $F(G)\le M(G)+1$.
\end{proposition}
\begin{proof}
We first show $M(H)\ge7/3$. Let $d_1\ge d_2\ge d_3$ be the three
largest degrees. If $d_3\ge3$, take their three vertices $T$.
The ordered-star count used above gives $M(H)\ge3-e/d_3$,
where $e$ counts the edges within $T$. Each additional edge entirely
outside $T$ increases this bound by at least $1/d_3$.
This proves the claim if $e\le2$ or if any edge lies outside $T$.
In the remaining case $T$ is a triangle and a vertex cover. Every
outside vertex meets at least two vertices of $T$, and any two such
neighbor sets intersect. All pairs then have distance at most two.

If $d_3=2$, choose the two highest-degree vertices $x,y$. Edges
entirely outside them weigh $1/2$, while their two stars contribute
at least $2-1/d_2$. Two outside edges suffice. With at most one such
edge and $xy$ present, every vertex outside $\{x,y\}$ is a common
neighbor of $x,y$ except possibly the ends of that one edge. Each
exception meets a hub and the two exceptions are adjacent, so the
diameter is at most two. If $xy$ is absent then the stars contribute
two; one outside edge suffices. Without that edge the graph is
complete bipartite with one part of size two and hence has diameter two.

By the support reduction it is enough to consider one added tip.
All pair distances are at most three except for possible distances of
four involving this tip. If its mass is $x$ (zero without a tip),
the collision correction yields
\[
 F(G)\le3+2x(1-x)-x^2\le10/3.
\]
Leaf monotonicity gives $M(G)\ge7/3$, as required.
\end{proof}

\begin{proposition}\label{p11-prop:diam-three-two}
Proposition~\ref{p11-prop:diam-three} remains valid with arbitrary
private leaves at two distinct sites of the diameter-three core.
\end{proposition}
\begin{proof}
Retain one leaf at each site and choose a separated maximizing measure.
Let the tip masses be $x,y$ and $m=x+y$. Relative to the constant
three, a tip--old pair has excess at most one and the tip--tip pair
has excess at most two. Each diagonal entry is two, so
\[
 F(G)\le3+2m(1-m)+4xy-x^2-y^2
      \le3+2m-3m^2/2\le11/3.
\]
Let $T$ be the three core vertices of largest full degree, with smallest
degree $d$ among the three. If $d\ge3$, let $e$ be the number of core
edges inside $T$, $c$ the number entirely outside $T$, and $r$ the
number of loaded sites outside $T$. The ordered-star estimate gives
$M(G)\ge3+(r+c-e)/d$, enough when $e-r-c\le1$.
If $e-r-c\ge2$ and $c=0$, then $e\ge2$, $T$ is connected and
all outside vertices have at least two neighbors in $T$. Any two
outside vertices share a neighbor, and each member of $T$ lies within
two steps of every outside vertex, contradicting diameter three.
If $e-r-c\ge2$ and $c\ge1$, then $e=3$, $c=1$ and $r=0$.
All outside vertices have at least one neighbor in the triangle $T$.
Thus each vertex of $T$ has core eccentricity at most two.
Since both loaded sites lie in this triangle, $G$ has diameter at most
three. Now $F(G)\le3$ and $M(G)\ge M(H)\ge7/3$.

If $d=2$, then the two loaded sites are the two highest-degree core
vertices, since their full degrees are at least three. Their stars
contribute at least $2-1/3$, and every core edge entirely outside them
weighs $1/2$. If there are two such edges we have $M(G)\ge8/3$.
Otherwise the diameter-three hypothesis forces exactly one such edge
and the two sites to be nonadjacent, by the two-hub argument of
Proposition~\ref{p11-prop:diam-three}. Thus $M(G)\ge2+1/2$.
If only one tip receives positive mass, the earlier bound $F(G)\le10/3$
applies. If both tips receive mass, every old vertex lies at distance
at most two from one of them: each site has its own tip and every
other core vertex has a neighbor among the two sites, since it has
core degree two and there is only one edge outside them.
Hence no old vertex can belong to the separated support. On two
points at distance at most five the energy is at most
$2+2(5-2)\max_{0\le x\le1}x(1-x)=7/2\le1+M(G)$.
\end{proof}

The diameter-three result applies in Theorem~\ref{p11-thm:main} because each cut operation can add leaves to a surviving block only at a vertex already incident to an external edge. Thus the set of its external attachment sites can only shrink throughout the induction.

\subsection{Regular blocks of unrestricted order}
Let $H$ be connected of order $h\ge2$.
If $s\ge2$ is the number of leaves of a spanning tree of order $h\ge3$,
then
\begin{equation}\label{p11-eq:regbag-q}
 Q(H)\le (h-s-1)/2+2-2/s .
\end{equation}
Indeed, distances are bounded by those in the tree, and an edge separating
masses $t$ and $1-t$ contributes $2t(1-t)\le1/2$. The $s$ leaf edges
together contribute at most $2b(1-b/s)\le2-2/s$, where $b\le1$ is the
total leaf mass. There are $h-s-1$ other edges.

For any collection of private leaves attached to the vertices of $H$, the
resulting $H'$ satisfies $M(H')\ge M(H)$ and $\Psi(H')\le Q(H)+1$.
For the former claim, a newly attached leaf at an old degree-$j$ vertex
contributes $1/(j+1)$, and its $j$ old edges lose at most
$1/j-1/(j+1)$ each. For the latter claim project each new leaf to its
support. In the expectation for two independent draws every new leaf
of mass $t$ contributes the additional distance $2t(1-t)$. Its diagonal
correction $2t^2$ cancels the quadratic term. The remaining correction
from core atoms $p_v$ is at most $2\sum_{v\in V(H)}p_v$; including the
constant $-1$ leaves a total addition of at most one. Consequently
$M(H)\ge Q(H)+1$ suffices for the entire decorated block.

Suppose $H$ is $d$-regular. Then $M(H)=h/2$. If a spanning tree has
$s\ge4$ leaves, the gap $M(H)-Q(H)$ is at least
$s/2-3/2+2/s\ge1$ by~\eqref{p11-eq:regbag-q}. For $d\ge4$ the star of all
edges incident to a vertex extends to a spanning tree, and its $d$
branches each contain a leaf. The same four-leaf conclusion holds in
every connected cubic graph other than $K_4$, by the following local
argument. If some edge $uv$ lies on no triangle, its endpoints have
four distinct other neighbors; their four incident edges together with
$uv$ form a tree with four leaves. Extending this tree one vertex at a
time preserves its leaf count or increases it. Otherwise every edge lies
on a triangle. At any vertex $u$ the subgraph on its three neighbors
has no isolated vertex. If it is a triangle, the cubic connected graph
is $K_4$. If it is a path $a,b,c$, then $u$ and $b$ already have their
three neighbors. The third neighbor $x$ of $a$ is distinct from $c$,
but the edge $ax$ can lie on a triangle only through $u$ or $b$.
Neither can be adjacent to $x$, a contradiction. In the exceptional
graph $K_4$, $Q=3/4$ and $M=2$.

For $d=2$ the graph is $C_h$. With $k=\lfloor h/2\rfloor$, use the $h$
cyclic intervals of $k$ consecutive vertices. Exactly $2d(x,y)$ intervals
separate a given pair $x,y$. If their masses are $t_i$, then
$\sum_i t_i=k$ and the expected distance is
$\sum_i t_i(1-t_i)\le k(1-k/h)\le h/4$. Thus $M-Q\ge1$ for $h\ge4$.
The graph $K_3$ is covered by Theorem~\ref{p11-th:diameter-two}, and $K_2$ with leaves is a star or a double star as above.

The same test includes irregular blocks. Write $\delta$ and $\Delta$
for the minimum and maximum degrees of $H$. On an edge with endpoint
degrees $i\le j$, the ratio of $1/j$ to $1/i+1/j$ is
$i/(i+j)\ge\delta/(\delta+\Delta)$. Thus
$M(H)\ge h\delta/(\delta+\Delta)$. A Hamiltonian cycle gives
$Q(H)\le h/4$ by the interval proof above, so
$h\delta/(\delta+\Delta)\ge h/4+1$ suffices. For example, any
Hamiltonian block with $h\ge12$ and $\Delta\le2\delta$ passes.

Another example requires no Hamiltonian cycle. If $\delta\ge24$ and
$\Delta\le2\delta$, then $M(H)\ge h/3$ again. Choose each vertex
independently with probability $p=\log(\delta+1)/(\delta+1)$, and add
every undominated vertex. The expected resulting size is at most
$h[p+(1-p)^{\delta+1}]\le h(1+\log(\delta+1))/(\delta+1)$, so a
dominating set $U$ with this size exists. Assign every vertex outside
$U$ to an adjacent center in $U$. Adjacency of the resulting star
clusters connects their centers by paths of length at most three.
Joining along a spanning tree of the cluster graph gives a connected
dominating set of size at most $3|U|-2$. It extends to a spanning tree
with at most $3|U|-2$ internal vertices, by attaching outside vertices
as leaves. Consequently
\[
 Q(H)+1\le \tfrac32|U|+\tfrac32
 \le h\cdot\frac{3(1+\log25)}{50}+\tfrac32 < h/3 .
\]
The second inequality uses that $(1+\log x)/x$ decreases for $x>1$.
For the last one, use $h\ge25$ and $\log25<13/4$, which follows, for
example, by summing the first eight terms of the exponential series
at $13/4$. This proves the decorated assertion also for these blocks.
The preceding calculation also gives the general sufficient condition
\[
 \frac{\delta}{\delta+\Delta}\ \ge \frac{3(1+\log(\delta+1))}{2(\delta+1)}+\frac{3}{2h}.
\]

\section{The exact calculation through eight core vertices}
For each two-connected $H$ of order at most eight, and $K_2$, we verified $M\ge\Psi$ for every one-leaf-per-site decoration. We describe how the finite computation proves an inequality over all real probability vectors.

Fix a set of marked sites (vertices with a leaf). At each marked site merge its center mass and all of its leaf masses onto one leaf. Within the site the kernel $d_G(x,y)+2\mathbf1_{x=y}$ is at most two, and becomes two after the merger. Distances to every other site can only increase. Thus a maximum can be sought on at most one representative per core site, using the leaf at a marked site. Choose a minimum-support maximum there. By the mass-transfer argument, all positive representatives are at pairwise distance at least three. Write $X$ for these sites and $A$ for the marks among them. Deleting any marks off $X$ leaves the candidate energy unchanged and can only lower $M$. It therefore suffices to test all $A\subseteq X\subseteq V(H)$. On that face the symmetric kernel is
\[
 K_{ii}=2,\qquad K_{ij}=d_H(i,j)+\mathbf1_{i\in A}+\mathbf1_{j\in A}\quad(i\ne j).
\]
At a relative interior maximum $Kp=\lambda\mathbf1$, $\mathbf1^\top p=1$, and $p^\top Kp=\lambda$. At a minimum-support maximizer the bordered system is invertible. Indeed, a null vector $(z,c)$ satisfies $Kz=c\mathbf1$, $\mathbf1^\top z=0$; multiplication by $p^\top$ gives $c=0$. A nonzero such $z$ would preserve the energy along $p+tz$ up to a boundary face. Thus the code checks Cramer's rule using integer determinants, discards singular and nonpositive solutions, and tests $\lambda-1\le M$. On two sites at distance $d\ge3$ the candidate value $\lambda-1$ is $d/2$. The singleton value is one.

The connected representatives are generated from $K_1$ by adding a vertex adjacent to each possible nonempty subset. Completeness follows by removing a leaf of a spanning tree. The canonical code uses degree refinement and minimizes the adjacency encoding over permutations inside each stable color class. A test for connectivity after each vertex deletion selects the blocks. One pruning rule comes from $Q(H)\le\diam(H)(1-1/h)$ by Cauchy--Schwarz. If a smaller mark set already pays $1+\diam(H)(1-1/h)$ in $M$, every extension of it passes. All comparisons are exact.

\begin{center}
\begin{tabular}{rrrrr}
\toprule
$h$ & Connected graphs & Blocks incl. $K_2$ & Marked instances & Tested faces\\
\midrule
2 & 1 & 1 & 4 & 4\\
3 & 2 & 1 & 8 & 6\\
4 & 6 & 3 & 48 & 43\\
5 & 21 & 10 & 320 & 277\\
6 & 112 & 56 & 3,584 & 1,655\\
7 & 853 & 468 & 59,904 & 21,951\\
8 & 11,117 & 7,123 & 1,823,488 & 404,991\\
\midrule
Total & 12,112 & 7,662 & 1,887,356 & 428,927\\
\bottomrule
\end{tabular}
\end{center}
Marked instances count $2^h$ subsets on a fixed representative of each block. Symmetries can identify some of these decorated instances. The final column counts the candidate faces tested after pruning. There were no failed comparisons. The preceding results, Theorem~\ref{p11-thm:reduction} and Proposition~\ref{p11-prop:transfer} prove Theorem~\ref{p11-thm:main}.

\section{Structure of a least weighted obstruction}
\begin{proposition}\label{p11-prop:weighted-obstruction}
Suppose $M\ge\Psi$ fails and choose a connected counterexample $G$
with the fewest edges (on at least two vertices). By Theorem~\ref{p11-thm:reduction} and its strict edge decrease, $G$ consists of a two-connected graph $H$ with new private leaves attached.
Fix a distribution $p$ maximizing $L_G$ with minimum support.
Every leaf has positive mass; each core vertex bears at most one leaf.
Every vertex of $H$ of intrinsic degree two bears a leaf, and these
intrinsic degree-two vertices form an independent set. Every core
vertex without a leaf has a core neighbor of strictly smaller full degree.
The support of $p$ has at least three points at pairwise distance at
least three.
\end{proposition}
\begin{proof}
Moving mass between two positive support points at distance $d\le2$
produces a quadratic with second derivative $4(2-d)\ge0$. Its value
cannot be smaller at both endpoints of the transfer interval; an
endpoint reduces support. No zero-mass leaf can remain, because removing
it preserves all relevant distances and cannot increase $M$. Indeed,
at its neighbor of degree $j\ge2$, the deleted edge costs $1/j$ and
at most $j-1$ surviving edges can gain $1/[j(j-1)]$ each.
In particular at most one of the leaves at a vertex can remain.
A singleton support gives $L_G=2\le M+1$, since the edges at a
maximum-degree vertex contribute one to $M$. For support on a pair at
distance $d\ge2$ the maximum is $1+d/2$; the modified-index diameter
bound $M\ge\diam(G)/2$~\cite[Theorem~3]{p11-DLS} handles this case.

For a nonbridge $vu$ with $i=d_v\le j=d_u$, suppose that all neighbors
of $v$ have degree at least $i$. On deleting $vu$ the surviving edges
incident to $v$ keep their weights. The change at each of the at most
$j-1$ other incidences of $u$ is at most $1/[j(j-1)]$. The removed
edge had weight $1/j$, so $M(G-vu)\le M(G)$. All distances weakly
increase, contradicting edge minimality. Each edge of $H$ is a
nonbridge. Thus a core vertex without a leaf has a neighbor of lower
full degree, necessarily in the core. In particular a vertex of
intrinsic degree two must be loaded: otherwise its full degree is two
and its two core neighbors both have degree at least two.

All core vertices now have full degree at least three. If two adjacent
vertices of intrinsic degree two existed, their full degrees would be
three and each would have precisely one leaf. On deleting the edge
between them, the two leaf edges gain $1/2-1/3$ apiece. Each endpoint's
other remaining core edge keeps its weight since its other endpoint
has full degree at least three, also when these two core edges have
a common other endpoint. The total change is $2(1/2-1/3)-1/3=0$,
again a contradiction. This proves the claims.
\end{proof}
A descending-degree walk from an unloaded core vertex of full degree
$d$ must reach a loaded vertex in at most $d-3$ steps, since the least
full degree in $H$ is three. Thus at full degree three every core
vertex is loaded; at degree four an unloaded vertex needs a loaded
degree-three core neighbor. No open ear in $H$ can have two consecutive
internal vertices of intrinsic degree two.

\begin{corollary}\label{p11-cor:weighted-order}
In Proposition~\ref{p11-prop:weighted-obstruction}, if the core has diameter
$D$ and the full graph has order $N$, then $N\le D(D+3)$.
In particular a least weighted obstruction of core diameter three
would have at most 18 vertices.
\end{corollary}
\begin{proof}
Put $s=D+1$ and take the $s$ core vertices with largest full degrees
as $T$, with minimum degree $d$ among them. Let $e$ be the number of
core edges inside $T$, $c$ the number entirely outside $T$, and $r$
the number of loaded sites outside $T$. Counting the stars on $T$ in
descending degree order gives $M(G)\ge s-e/d+(c+r)/d$.
Since every distinct pair in $G$ is at distance at most $D+2$, its
weighted pair objective is at most $D+2$, including the diagonal.
Failure thus forces $M(G)<s$, so $c+r\le e-1\le\binom{s}{2}-1$.
By Proposition~\ref{p11-prop:weighted-obstruction}, each unloaded core
vertex outside $T$ has a core neighbor of lower full degree, also
outside $T$. There are therefore at most $2c$ such vertices. The core
contains at most $s+2c+r$ vertices and there are at most $s+r$ leaves.
Hence $N\le 2s+2(c+r)\le s(s+1)-2=D(D+3)$.
\end{proof}


In particular, the remaining minimum weighted obstruction has a single irregular core with at least nine vertices and diameter at least three. When the core has diameter three, it must have at least three loaded sites. To settle the original conjecture through this reduction it suffices to rule out all such obstructions; a failure of the stronger weighted inequality alone would not refute Fajtlowicz's inequality.

\section*{Methodological sources}
Dvo\v{r}\'ak, Lidick\'y and \v{S}krekovski used cut-vertex degree padding and locally minimal vertices in their diameter proof~\cite{p11-DLS}. Manuscript 087 of the Verifiable Mathematical Reasoning Corpus uses that padding for radius. Manuscript 094 tracks edge contributions by endpoint degree under local deletions. Manuscript 106 uses exact rational certificates after structural reduction. Manuscript 161 uses the edgewise squared-degree comparison in Proposition~\ref{p11-prop:transfer}. Their graph arguments address radius. The weighted distance functional, asymmetric assignment of cut mass, terminal estimates, and stationary-face verifier here are derived for average distance.
\begin{thebibliography}{9}
\bibitem{p11-Fajt} S. Fajtlowicz, On conjectures of Graffiti, Discrete Math. 72 (1988), 113--118.
\bibitem{p11-Trees} M. Cygan, M. Pilipczuk and R. \v{S}krekovski, Relation between Randi\'c index and average distance of trees, MATCH Commun. Math. Comput. Chem. 66 (2011), 605--612.
\bibitem{p11-Uni} Z. You and B. Liu, On a conjecture between Randi\'c index and average distance of unicyclic graphs, Filomat 28 (2014), 767--773.
\bibitem{p11-DLS} Z. Dvo\v{r}\'ak, B. Lidick\'y and R. \v{S}krekovski, Randi\'c index and the diameter of a graph, KAM--DIMATIA Series 2010--496, Theorem 3 and Lemmas 4--5. \href{https://iti.mff.cuni.cz/series/2010/496.pdf}{Authors' text}.
\bibitem{p11-VMRC} QEEDDD, \emph{Verifiable Mathematical Reasoning Corpus},
\href{https://github.com/QEEDDD/Verifiable-Mathematical-Reasoning-Corpus/tree/38348353416f1f7fc50cecb93c3b554a00260c13}{snapshot 3834835}, manuscripts 087, 094, 106, and 161.
\end{thebibliography}
\clearpage\volappendix
\section{Reproducible verifier}
Save the lines between the two \texttt{Verbatim} markers in the accompanying TeX source as \texttt{bagcheck.py}. Run \texttt{python3 bagcheck.py} with Python 3.10 or later and assertions enabled. It uses only the standard library.
\smallskip
\begin{Verbatim}[fontsize=\scriptsize]

"""Verify all decorated two-connected blocks of order at most eight.
Python 3.10+, standard library only.
"""
from collections import deque
from fractions import Fraction
from itertools import combinations, permutations, product


def canonical_key(graph):
    n = len(graph)
    colors = tuple(neighbors.bit_count() for neighbors in graph)
    while True:
        signatures = [
            (colors[i], tuple(sorted(colors[j] for j in range(n)
                                     if graph[i] >> j & 1)))
            for i in range(n)
        ]
        palette = {sig: j for j, sig in enumerate(sorted(set(signatures)))}
        refined = tuple(palette[sig] for sig in signatures)
        if len(set(refined)) == len(set(colors)):
            break
        colors = refined
    classes = [[i for i in range(n) if refined[i] == color]
               for color in sorted(set(refined))]
    best = None
    for groups in product(*(permutations(group) for group in classes)):
        order = sum(groups, ())
        code = sum(((graph[order[i]] >> order[j]) & 1)
                   << (i * (i - 1) // 2 + j)
                   for i in range(n) for j in range(i))
        if best is None or code < best:
            best = code
    return best


def extend(graphs, n):
    representatives = {}
    for graph in graphs:
        for mask in range(1, 1 << (n - 1)):
            new = tuple(row | (((mask >> i) & 1) << (n - 1))
                        for i, row in enumerate(graph)) + (mask,)
            key = canonical_key(new)
            if key not in representatives:
                representatives[key] = new
    return list(representatives.values())


def determinant(matrix):
    if not matrix:
        return 1
    if len(matrix) == 1:
        return matrix[0][0]
    total = 0
    for i, val in enumerate(matrix[0]):
        if val:
            minor = [row[:i] + row[i+1:] for row in matrix[1:]]
            total += (-1)**i * val * determinant(minor)
    return total


def distances(graph):
    n = len(graph)
    result = []
    for u in range(n):
        d = [-1] * n
        d[u] = 0
        queue = deque([u])
        while queue:
            v = queue.popleft()
            for w in range(n):
                if d[w] == -1 and graph[v] >> w & 1:
                    d[w] = d[v] + 1
                    queue.append(w)
        result.append(d)
    return result



from math import lcm
from functools import lru_cache

@lru_cache(None)
def face_value(rows):
    # Maxima can be taken on a face with invertible bordered equations.
    t = len(rows)
    if t == 1:
        return Fraction(1)
    if t == 2:
        return Fraction(rows[0][1], 2)
    matrix = [list(row) + [-1] for row in rows] + [[1] * t + [0]]
    den = determinant(matrix)
    if not den:
        return Fraction(0)
    cofactors = [(-1)**(t+j) *
                 determinant([row[:j] + row[j+1:] for row in matrix[:-1]])
                 for j in range(t + 1)]
    if any(c * den <= 0 for c in cofactors[:-1]):
        return Fraction(0)
    return Fraction(cofactors[-1], den) - 1


def is_block(graph):
    n = len(graph)
    for v in range(n):
        active = ((1 << n) - 1) ^ (1 << v)
        seen = active & -active
        frontier = seen
        while frontier:
            bit = frontier & -frontier
            frontier ^= bit
            neighbors = graph[bit.bit_length()-1] & active & ~seen
            seen |= neighbors
            frontier |= neighbors
        if seen != active:
            return False
    return True


def check_decorations(graph):
    n = len(graph)
    q = 1 << n
    full = q - 1
    scale = lcm(*range(1, n+1))
    d = distances(graph)
    diameter = max(max(row) for row in d)
    deg = [x.bit_count() for x in graph]
    edges = [(i, j) for i in range(n) for j in range(i) if graph[i] >> j & 1]
    base = sum(scale // max(deg[i], deg[j]) for i, j in edges)
    # For variable marks A, project back to sites. Theta <= Q(H)+1
    # and Q(H) <= diam(H)*(1-1/n) by Cauchy--Schwarz.
    ub_num = diameter*(n-1)+n
    if n*base >= scale*ub_num:
        return 0, 0
    # Compute M(A) in integer units by first and second differences.
    linear = [scale // (v+1) for v in deg]
    paired = [[0]*n for _ in range(n)]
    for i,j in edges:
        a,b = deg[i],deg[j]
        def w(ii,jj): return scale // max(ii,jj)
        u = w(a+1,b)-w(a,b)
        v = w(a,b+1)-w(a,b)
        linear[i] += u
        linear[j] += v
        paired[i][j] = paired[j][i] = w(a+1,b+1)-w(a,b)-u-v
    weights = [base]*q
    for mask in range(1,q):
        bit = mask & -mask
        i = bit.bit_length()-1
        rest = mask ^ bit
        weights[mask] = weights[rest] + linear[i] + sum(
            paired[i][j] for j in range(n) if rest >> j & 1)
    # Certificate can compare ALL admissible supports X at once with
    # M(A intersect X). Other marks only increase M, and leave the
    # face's metric unchanged.
    positions = [tuple(i for i in range(n) if mask >> i & 1)
                 for mask in range(q)]
    fac = 0
    def visit(x, cand, marks):
        nonlocal fac
        if x:
            bound = weights[marks]
            if n * bound >= scale*ub_num:
                return
            X = positions[x]
            rows = tuple(tuple(2 if i == j else d[i][j] +
                               ((marks >> i) & 1) + ((marks >> j) & 1)
                               for j in X) for i in X)
            val = face_value(rows)
            assert val*scale <= bound, (graph,X,marks,val,Fraction(bound,scale))
            fac += 1
        # Every extension respects smallest support vertex ordering.
        while cand:
            bit = cand & -cand
            v = bit.bit_length()-1
            cand ^= bit
            for a in (0,1):
                if any(d[u][v]+((marks >> u) & 1)+a < 3 for u in positions[x]):
                    continue
                visit(x|bit,cand,marks|(a*bit))
    visit(0,full,0)
    return 1, fac


if __name__ == '__main__':
    graphs=[(0,)]
    print('order connected blocks decorations difficult_blocks tested_faces', flush=True)
    for n in range(2,9):
        graphs=extend(graphs,n)
        blocks=[g for g in graphs if is_block(g)]
        checks=[check_decorations(g) for g in blocks]
        print(n,len(graphs),len(blocks),(1<<n)*len(blocks),
              sum(c[0] for c in checks),sum(c[1] for c in checks),flush=True)
\end{Verbatim}

% END SOURCE BODY Fajtlowicz_Weighted_Cut_Vertex_Reduction
\unitstart{Paper 12}{p12}{Clipped distances}{Clipped distances and the Randi\'c average-distance conjecture}{}
\makeatletter\@removefromreset{theorem}{section}\@addtoreset{theorem}{section}\makeatother\renewcommand{\thetheorem}{\thesection.\arabic{theorem}}
% BEGIN SOURCE BODY Fajtlowicz_Clipped_Distances
\begin{abstract}
We study Fajtlowicz's inequality between the Randi\'c index and average distance on a finite connected graph. The remaining diameter-three block case with unrestricted attachment sites follows from a stronger statement: a capped distance inequality on every graph, including graphs of arbitrarily large diameter. We prove a universal cap of $12$ and deduce the ordinary inequality on arbitrary assemblies of blocks of diameter at most $10$. An independent degree argument covers every block of maximum degree at most thirteen, of unrestricted diameter. We include the mathematical reductions and exact verifiers. The conjecture for all connected graphs remains open in this work.
\end{abstract}

\section{Statement and conventions}
Let $G$ be finite, simple, undirected, connected, and of order $n\ge2$. With vertex degree $d_v$, write
\[
 M(G)=\sum_{uv\in E(G)}\frac1{\max(d_u,d_v)},\qquad
 R(G)=\sum_{uv\in E(G)}\frac1{\sqrt{d_ud_v}},\qquad
 \mu(G)=\binom n2^{-1}\sum_{\{x,y\}}d_G(x,y).
\]
Fajtlowicz conjectured $R(G)\ge\mu(G)$~\cite{p12-Fajt,p12-Trees}. For $C\ge2$ (including $C=\infty$) put
\[
 K_{G,C}(x,y)=
 \begin{cases}
 2,&x=y,\\
 \min\{C,d_G(x,y)\},&x\ne y,
 \end{cases}\qquad
 F_C(G)=\max_{p}\sum_{x,y}p_xp_y K_{G,C}(x,y).
\]
The maximum runs over all probability vectors on $V(G)$. Write $\Psi(G)=F_\infty(G)-1$ and $Q(H)=\max_p\sum_{x,y}p_xp_y d_H(x,y)$. A cyclic block is a maximal two-connected subgraph of order at least three. A decoration attaches arbitrary numbers of private leaves at core vertices. A loaded site bears a leaf. The degrees and diameter of a block in a hypothesis refer to that block in isolation.

\begin{theorem}\label{p12-th:main}
Every connected graph of order at least two satisfies $F_{12}(G)\le M(G)+1$.
Consequently $R(G)\ge\mu(G)$ if $\diam(G)\le12$, or if every cyclic block independently has diameter at most $10$ or maximum degree at most thirteen. Individual blocks may also be replaced by any block whose every decoration satisfies $M\ge\Psi$, including the degree classes below and regular blocks of arbitrary order~\cite{p12-previous}.
\end{theorem}
A chain of blocks covered by the theorem may have arbitrarily large total diameter. The proof uses the capped inequality on each decorated piece and the ordinary (uncapped) gluing inequality for the chain.

\section{Why clipping commutes with the reduction}
\begin{lemma}\label{p12-lem:split}
For a cut vertex $v$, partition $G-v$ into unions of components $A,B$, both nonempty. Pad the induced piece on $A\cup\{v\}$ with one new leaf at $v$ for each original neighbor of $v$ in $B$; define $G_B$ symmetrically. Then $M(G_A)+M(G_B)=M(G)+1$ and
\[
 F_C(G)+2\le F_C(G_A)+F_C(G_B)
\]
for every $C\ge2$, including infinity.
\end{lemma}
\begin{proof}
Every old edge retains its degrees, and the weights of all the new edges sum to one, as in Dvo\v{r}\'ak, Lidick\'y and \v{S}krekovski~\cite[Lemma~4]{p12-DLS}. For the law on the $A$ side keep $A$ pointwise and map all other vertices, including $v$, to a single new leaf. Do the analogous operation on the $B$ side. For two points in the same open side, the sum of the projected kernel entries is the old entry plus two. Two copies of the same original point have sum four. For $v$ and an open-side point at distance $a$ from $v$, the sum is $2+\min(C,a+1)\ge2+\min(C,a)$. Across open sides, the distances $a,b$ from $v$ are at least one. The required inequality is
\[
 \min(C,a+1)+\min(C,b+1)\ge2+\min(C,a+b).
\]
With no truncation on the left this follows by addition, while a truncated summand contributes $C$ and the other contributes at least two. Take expectation for a maximizing law on $G$.
\end{proof}

\begin{lemma}\label{p12-lem:minimal}
If $F_C\le M+1$ fails, any counterexample of least edge count is a two-connected core $H$ with at most one leaf per site. In a minimum-support maximizing law, distinct occupied vertices have distance at least three and every leaf is occupied. Each unloaded site has a core neighbor of strictly smaller full degree. In particular all full degrees on the core are at least three. Sites of full degree three have one leaf and form an independent set. The same reduction holds within an a priori upper bound on the intrinsic maximum degree of each cyclic block.
\end{lemma}
\begin{proof}
Choose the split in Lemma~\ref{p12-lem:split} whenever each open side contains an edge. Each padded side then has fewer edges. Two distinct cyclic blocks have such a separator on their path in the block--cut tree. For a single cyclic block with any branch of length at least two, split at its attachment. The terminal trees have diameter at most three, since a four-edge path supplies a split. A star has $M=1$ and $F_C\le2$. A double star with $a\ge b\ge1$ leaves at the two centers has $M=1+b/(b+1)\ge3/2$. Its separated supports have at most two points at distance at most three, giving $F_C\le5/2$.

For separation, shift mass between two positive coordinates at distance at most two with their total fixed. The quadratic coefficient is $4-2\min(d,C)\ge0$, so an endpoint of the interval eliminates a coordinate without loss. Adding a leaf at a vertex of degree $j$ pays $1/(j+1)$ and reduces each of at most $j$ old edges by at most $1/[j(j+1)]$. Thus removing a zero-mass leaf keeps the value of the law and can only decrease $M$. Sibling leaves cannot both be occupied.

If an unloaded core vertex $v$ has no lower-degree neighbor, take a core edge $vu$ with $d_v\le d_u$. After deletion, every other edge at $v$ retains its weight. The gains at $u$ sum to at most $1/d_u$, exactly the weight removed. Every core edge is a nonbridge, so the deletion stays connected and all capped distances can only increase. This contradicts minimality. An unloaded degree-two vertex is impossible, since every core neighbor has degree at least two. It follows that all full core degrees are at least three, and an unloaded full-degree-three vertex is now impossible as well. For two adjacent full-degree-three sites, deletion of their common edge removes $1/3$ and the two leaf edges each gain $1/6$. All other edge weights stay fixed. This again contradicts minimality. Every operation preserves an imposed bound on the intrinsic maximum degree of cyclic blocks, since any new cyclic block is contained in an old one.
\end{proof}

\section{Two bounds from spanning trees}
Project a law on a decorated block onto its core. A leaf of mass $t$ contributes $2t(1-t)$ additional expected distance and $2t^2$ diagonal correction. An old core atom $p_v$ contributes $2p_v^2\le2p_v$. Summing gives $F_C(G)-1\le Q(H)+1$. For a core of order $h$, a spanning tree with $s$ leaves gives
\[
 Q(H)\le (h-s-1)/2+2-2/s.
\]
Indeed each nonpendant edge separates masses $x$ and $1-x$ and contributes at most $1/2$ in the tree. For the $s$ leaf edges, with total leaf mass $b\le1$, their sum is at most $2b-2b^2/s\le2-2/s$. A star at a vertex of intrinsic degree $\Delta(H)$ extends to a spanning tree with at least $\Delta(H)$ leaves.

We also record a bound for every two-connected $H$. If its diameter is $D$, two internally vertex-disjoint paths between diametral endpoints form a cycle of length $c\ge2D$. Extend the cycle to a spanning connected unicyclic subgraph and traverse the cycle once and every attached tree edge twice. The resulting closed walk has length $m=2h-c$. Choose one occurrence of every site. Distances in $H$ are bounded by circular distances along the walk. On a discrete circle of length $m$, take all cyclic intervals of length $k=\lfloor m/2\rfloor$. A pair at circular distance $r$ is separated by $2r$ intervals. If the interval masses are $t_i$, the expected circular distance is $\sum_i t_i(1-t_i)\le k(1-k/m)\le m/4$. It follows that
\[
 Q(H)\le\min\left(D(1-1/h),\,(h-D)/2\right)
 \le\frac{h(h-1)}{3h-2}.
\]

\paragraph{Small cores without graph enumeration.}
In any least obstruction of Lemma~\ref{p12-lem:minimal}, write the core's full degrees in decreasing order, here indexed from one, and let $s$ be the largest index with $d_s\ge s$. Assign edges touching the top $s$ sites to the earlier endpoint. Their sum is at least $s-\binom{s}{2}/s=(s+1)/2$, because each already-assigned incidence costs at most $1/s$ at the later site. Each remaining vertex supplies its pendant edge if loaded, or an edge to a strictly lower-degree vertex otherwise. These edges are distinct, all outside the top set, and have weight at least $1/s$. Indeed the remaining full degrees are at most $s$. Thus
\[
 M(G)\ge B(h):=\min_{1\le s\le h}\bigl(h/s+(s-1)/2\bigr).
\]
For $h=3,\dots,10$ these lower bounds are respectively
$2,7/3,8/3,3,13/4,7/2,15/4,4$.
On the same orders the maxima of $1+\min(D(1-1/h),(h-D)/2)$ over the integers $1\le D\le\lfloor h/2\rfloor$ are
$5/3,2,5/2,8/3,3,7/2,11/3,4$.
Hence no least obstruction has $h\le10$. This calculation also applies when the reductions are restricted to a bound on intrinsic maximum degree.

\section{A degree certificate for blocks of arbitrary diameter}
\begin{theorem}\label{p12-thm:thirteen}
If every cyclic block of $G$ has intrinsic maximum degree at most
thirteen, then $M(G)\ge\Psi(G)$. Each such block with arbitrary
private leaves is therefore certified and may be freely mixed with
any other individually certified blocks.
\end{theorem}
\begin{proof}
Choose a counterexample of minimum edge count within this class, which
is preserved by cut-vertex padding, deleting a leaf, and deleting a
nonbridge: every new cyclic block is contained in an old one.
The cut and descent reductions give a terminal graph with a
two-connected core $H$ of order $h$, diameter $D$ and maximum intrinsic
degree $K\le13$. The small-order argument above permits $h\ge9$. Each core vertex has type $(k,a)$
with $k=\deg_H(v)$ and load $a\in\{0,1\}$, and full degree $d=k+a$.
Every full degree is at least three and each unloaded core vertex
has a core neighbor of smaller full degree. In particular $(2,0)$
and $(3,0)$ are excluded. Expanding $2/\max(d,e)$ gives
\[
 B=h-2M(G)=\sum_{ij\in E(H)}
 \left( |1/d_i-1/d_j|-\frac{a_i}{k_id_i}-\frac{a_j}{k_jd_j}\right).
\]

\emph{Metric bounds.} For the diametral cycle of length $l\ge2D$ above we have $Q(H)\le(2h-l)/4$ and $Q(H)\le D(1-1/h)$. When $l<4D$ another bound uses the total mass $b$ on the chosen cycle.
The conditional expected distance between two draws on the cycle is
at most $l/4$. The off-cycle diagonal alone has probability at least
$(1-b)^2/(h-l)$ if $h>l$. Consequently
\[
 D-Q(H)\ge
 \left((D-l/4)^{-1}+(h-l)/D\right)^{-1}\quad(l<4D).
\]
Indeed the right side is the minimum of
$(D-l/4)b^2+D(1-b)^2/(h-l)$; at $h=l$ take $b=1$.
Finally breadth-first layers from a diametral endpoint have at least
two vertices in each interior layer, since an interior singleton
would separate the diametral vertices. Their baseline sizes are
$1,2,\ldots,2,1$, summing to $2D$. The closed neighborhood of any
vertex lies in at most three successive layers and hence has at most
$6+(h-2D)$ vertices. Thus $2D\le h-K+5$.

\emph{Rational certificate.} For actual $K$ in the four ranges
$2$--$9$, $10$--$11$, $12$, $13$, respectively, the verifier below
certifies
\[
 B\le ch,\qquad c\in\{198/1000,253/1000,84/300,2851/10000\}
\]
in the same order. Specifically the integers $X_i,Y_i,m,C$ in its
four rows satisfy $Y_i\ge0$, $Y_i=0$ on loaded sites, $c=C/m$, and
\begin{align*}
 m\left(|1/d_i-1/d_j|-a_i/(k_i d_i)-a_j/(k_jd_j)\right)
 &\le X_i+X_j-Y_i\mathbf 1_{d_i>d_j}-Y_j\mathbf 1_{d_j>d_i},\\
 k_i X_i-Y_i &\le C.
\end{align*}
Sum the first inequality over the core edges. Every nonzero $Y_i$
is paid at least once by the smaller neighbor of its unloaded vertex.
The second inequality then proves $mB\le Ch$.

If $h\ge2/(1/3-c)$, the bound $Q(H)\le h/3$ gives
$h(1-c)\ge2(Q(H)+1)$. Below this threshold the parameter check
in the verifier uses precisely the integer ranges proved above and
the minimum of the applicable metric upper bounds. In all cases it
gives the same conclusion. Thus
$2M(G)\ge h(1-c)\ge2(Q(H)+1)\ge2\Psi(G)$, a contradiction.
The order of the entire graph and its articulation degrees do not
enter this proof.
\end{proof}

\section{A capped inequality on all graphs}
\begin{theorem}\label{p12-crit:cap12}
For every connected graph $G$ with at least two vertices the cutoff
$K=12$ satisfies $F_K(G)\le M(G)+1$. In particular the
original weighted inequality holds if $\diam(G)\le12$. Every
two-connected core of diameter at most ten satisfies the untruncated
inequality with arbitrary private leaves. Such cyclic blocks can be
mixed without restriction using the ordinary cut reduction.
\end{theorem}
\begin{proof}
Consider a counterexample with the fewest edges to the cutoff assertion.
The cut, leaf, and nonbridge reductions give one two-connected core $H$
with at most one leaf at each site. Let $h=|V(H)|$ and
$d_0\ge\dots\ge d_{h-1}$ be the full degrees of its vertices in
this order. It suffices to test $h\ge9$, $\Delta(H)\ge8$, since Theorem~\ref{p12-thm:thirteen} already handles the larger range $\Delta(H)\le13$ (in particular all $h\le14$). Keeping the weaker floors in the finite verifier is harmless. Every unladen vertex has a strictly lower-degree core neighbor.
Every full degree is at least three, and every vertex of full degree
three has exactly one leaf and two core neighbors of larger degree.

For $s\le h$ take the first $s$ core vertices as $T$, and let $e,c$
count core edges inside and entirely outside $T$, respectively. Let
$r$ count laden vertices outside $T$. Each unladen vertex outside $T$
chooses an edge to a strictly lower-degree core neighbor outside $T$.
These edges have distinct higher endpoints; thus $h-s\le c+r$.
Summing the first $s$ stars, then unused core edges and leaves, gives
$M\ge s+(c+r-e)/d_{s-1}$. Since $M<11$ in any obstruction, taking
$s=11$ when $h\ge11$ yields $h\le11+\binom{11}{2}-1=65$.

\paragraph{An upper bound.}
Let $p$ be any distribution on the decorated graph, and project it onto
the core by moving mass on a leaf to its site. For sites of the two
independent draws let their core distance be $t$. The full kernel,
including its diagonal value two, is at most $2+\min(t,10)$. For
$j=1,\ldots,10$, let $P_j$ be the maximum cardinality of a subset of
mutually $j$-separated core vertices. The probability that $t\ge j$ is
at most $1-1/P_j$. Indeed, in the graph joining pairs at distance at
least $j$, maximize the ordered edge probability over distributions.
For two distinct positive nonadjacent sites the objective is affine
under transfer of mass between them, so some maximizing distribution
is supported on a clique. A distribution on at most $P_j$ points has
collision probability at least $1/P_j$.

Set $w=2\lfloor(j-1)/2\rfloor+1$. When there are at least two
$j$-separated sites, the disjoint balls of radius $(w-1)/2$ around
them have at least $w$ vertices each. Indeed two internally disjoint
paths between a center and another site give two disjoint initial
segments of length $(w-1)/2$. If $j$ is even there is also no edge
between the balls. At least two vertices lie outside their union:
otherwise the graph is disconnected or a single vertex separates balls.
Consequently $P_j\le B_j=\max\{1,\lfloor(h-2\mathbf1_{2\mid j})/w\rfloor\}$.
This yields the first of the following independent upper estimates:
\[
 F_{12}(G)-1\ \le\ U=\min\left\{
 11-\sum_{j=1}^{10}1/B_j,\quad
 1+\frac{h(h-1)}{3h-2},\quad
 \frac{h-\ell+5}2-\frac2\ell\right\},
\]
Here $\ell=\max(8,d_0-1)$. For the second estimate use the two-connected distance bound for $Q(H)$.
For the third, extend a maximum-degree star to a spanning tree with at
least $\ell$ leaves and use the spanning-tree estimate for $Q(H)+1$.

\paragraph{Degree comparisons.}
Let $k_i$ count the core neighbors of vertex $i$ with smaller index.
Assign each edge to its first core endpoint, including any private
leaf at that endpoint. Then
\[
 M=h-\sum_{i=0}^{h-1} k_i/d_i,\qquad
 0\le k_i\le b_i=\min(i,d_i-1).
\]
Each vertex either has a private leaf or a lower-degree neighbor,
which explains the $d_i-1$ bound. Writing $n_z=\#\{i:d_i=z\}$,
we may further decrease $b_i$ to $\min(b_i,h-n_3)$ at degree three.
For every full degree $z$, at least
$g_z=\max(0,n_z-\sum_{y<z} y n_y)$ vertices must be laden:
an unladen one requires a distinct incidence at a lower-degree core
vertex. In this formula the contribution of degree three to the sum
may be replaced by $2n_3$. Also take $g_h=n_h$, since full degree $h$
forces a leaf. Thus
\[
 \sum_i k_i \le B=\left\lfloor
 \tfrac12(\textstyle\sum_i d_i-\sum_z g_z)\right\rfloor.
\]
Maximizing $\sum k_i/d_i$ under this budget and the individual caps
is greedy, starting at the smallest degree. Another lower bound on
$M$ uses any initial segment and the inequality $h-s\le c+r$:
\[
 M\ge s+\frac{h-s}{d_{s-1}}-\sum_{i<s}\frac{\min(i,d_i-1)}{d_i}.
\]
The verifier compares these lower bounds with $U$ using exact integer
arithmetic for every nonincreasing degree sequence between $3$ and $h$
with $d_0\ge8$, for $9\le h\le65$.

For efficiency it also uses this bound to prune degree prefixes. Fix an
integer $q$ between $3$ and $10$. Put $f_i=\mathbf1_{d_i\in\{3,h\}}$.
The looser budget $\sum k_i\le\sum_i(d_i-f_i)/2$ and the individual
caps imply
\[
 M\ge h-\sum_i\left[
 \frac{d_i-f_i}{2q}
 +\max(0,1/d_i-1/q)\min(i,d_i-1)\right].
\]
Indeed subtract $\sum k_i/q$ before applying the individual caps.
After a prefix ending in $x$, each unassigned summand is bounded above
by the maximum over $3\le z\le x$ of the same expression with
$d_i=z$ and $\min(i,z-1)$ replaced by $z-1$. This safely certifies all
extensions whenever the resulting lower bound is at least $U$.
The program also drops an entire order when the lower bound $B(h)$ from Section~3 exceeds or equals the degree-independent upper bound. The program is given in the appendix. The code scales by $\mathrm{lcm}(1,\ldots,h)$ and rounds the second
entry in $U$ upwards to this grid before comparison. No degree sequence
fails. This proves the cap theorem, and decorations of cores with
diameter at most ten have total diameter at most twelve.
\end{proof}

\section{From the weighted bound to average distance}
Write $\tau_C=\binom{n}{2}^{-1}\sum_{\{x,y\}}(d_G(x,y)-C)_+$ and $k=\sqrt{\Delta(G)}$. The capped theorem implies, for every graph,
\[
 \begin{split}
 (k+1)(R-\mu)\ \ge\ n-2\mu
 -(k-1)\left[1+\frac{\mu-2}{n}
            +\left(1-\frac1n\right)\tau_C\right],
 \qquad C=12.
 \end{split}
\]
Indeed the uniform law gives $M\ge(1-1/n)(\mu-\tau_C)-1+2/n$.
On an edge with endpoint degrees $a\le b$, the identity
\[
 (a^{-1/2}-b^{-1/2})^2=(\sqrt{b/a}-1)\bigl((ab)^{-1/2}-b^{-1}\bigr)
\]
sums to $n-2R\le(k-1)(R-M)$. Substitution proves the displayed inequality. Thus any individual graph, even beyond the block classes, is certified whenever this right side is nonnegative.

If its diameter is at most $C$, then $\tau_C=0$. The same estimate with $\tau_C=0$ holds whenever the untruncated inequality $M\ge\Psi$ is proved by gluing. To check nonnegativity, recall that $\mu\le(n+1)/3$: pass to a spanning tree, delete a leaf repeatedly, and note that at order $j$ the sum of distances from the leaf is at most $1+\cdots+(j-1)$. The $i$th nearest other vertex is at distance at most $i$. This gives the bound after summing and dividing by $\binom n2$. Substituting $\mu\le(n+1)/3$ and $k\le\sqrt{n-1}$ suffices if
\[
 n(n-2)+4n-5\ \ge\ (4n-5)\sqrt{n-1}.
\]
The squared difference is $(n-5)(n-2)(n^2-5n+5)\ge0$ for $n\ge5$. For $2\le n\le4$, $(ab)^{-1/2}/(1/a+1/b)=\sqrt{b/a}/(1+b/a)\ge\sqrt{n-1}/n$, so $R\ge\sqrt{n-1}\ge(n+1)/3$.

\paragraph{Scope and remaining cases.}
The ordinary kernel agrees with the clipped kernel on every decorated block of intrinsic diameter at most $C-2$. Lemma~\ref{p12-lem:split} with $C=\infty$ and induction on the number of edges then give the unrestricted mixed-block statement, even if the assembled graph has diameter greater than $C$. In a globally least weighted obstruction the unique core therefore has $D\ge C-1$ and $\Delta(H)\ge14$. The breadth-first layer bound $h\ge2D+\Delta(H)-5$ gives $h\ge 2C+7$. In fact every graph whose cyclic blocks have at most $2C+6$ vertices is certified: each such block separately satisfies the diameter or the degree condition. The two stated bounds can also be mixed block by block. For such a least obstruction we also have the lower bound $M\ge\sqrt{2h}-1/2$ directly from $B(h)$ by the arithmetic--geometric mean inequality. Its ordinary kernel is at most $D+2$, so failure forces $M<D+1$ and
\[
 2D+9\ \le\ h\ <\ \frac{(D+3/2)^2}{2}.
\]
At the first remaining diameter $D=11$, the core must have between 31 and 78 vertices, inclusive, before any further degree or packing tests. The graph classes give sufficient conditions for the original conjecture. We have not ruled out a weighted obstruction with unrestricted diameter and degree.

\begin{thebibliography}{9}
\bibitem{p12-Fajt} S. Fajtlowicz, On conjectures of Graffiti, Discrete Math. 72 (1988), 113--118.
\bibitem{p12-Trees} M. Cygan, M. Pilipczuk and R. \v{S}krekovski, Relation between Randi\'c index and average distance of trees, MATCH Commun. Math. Comput. Chem. 66 (2011), 605--612.
\bibitem{p12-DLS} Z. Dvo\v{r}\'ak, B. Lidick\'y and R. \v{S}krekovski, Randi\'c index and the diameter of a graph, KAM--DIMATIA Series 2010--496, \href{https://iti.mff.cuni.cz/series/2010/496.pdf}{authors' text}.
\bibitem{p12-previous} Keio and Shaitetsu, \emph{A weighted cut-vertex reduction for the Randi\'c index and average distance}, companion research manuscript. For its methodological sources see \href{https://github.com/QEEDDD/Verifiable-Mathematical-Reasoning-Corpus/tree/38348353416f1f7fc50cecb93c3b554a00260c13}{Verifiable Mathematical Reasoning Corpus, snapshot 3834835}, entries 087, 094, 106 and 161.
\end{thebibliography}

\clearpage\volappendix
\section{Reproducing the degree certificate}
Copy the following block from the single TeX source to \texttt{degree13.py}. Run \texttt{python3 degree13.py} with assertions enabled. The four rows report respectively the smallest and largest intrinsic degree in a group, the large-order cutoff, the number of ordered type comparisons, and the number of integer metric cases. There are 1,404 type comparisons and 12,649 metric cases in total.
\begin{Verbatim}[fontsize=\scriptsize]
from fractions import Fraction as F
from math import ceil

# cap, m, C, edge potentials X, descent credits Y (indices start at 0)
rows = [
    (9, 1000, 198,
     [98,42,125,18,68,1,34,-13,11,-10,8,8,22,22],
     {2:305,4:142,6:9}),
    (11, 1000, 253,
     [124,68,151,37,87,21,54,6,30,-6,12,-11,3,3,14,14,23,23],
     {2:357,4:187,6:74}),
    (12, 300, 84,
     [40,24,49,14,29,9,19,5,12,1,6,-2,2,-1,2,2,5,5,7,7],
     {2:113,4:61,6:30}),
    (13, 10000, 2851,
     [1424,868,1701,512,1012,346,679,203,441,84,262,-15,
      123,-99,12,-11,80,80,155,155,219,219],
     {2:3957,4:2212,6:1224,8:236}),
]
start = 2
for cap,m,C,X,Y in rows:
    c = F(C,m)
    types = [(k,a) for k in range(2,cap+1) for a in (0,1)
             if k+a >= 3 and (a or k >= 4)]
    assert len(X) == len(types)
    for i,(k,a) in enumerate(types):
        d = k+a
        assert k*X[i]-Y.get(i,0) <= C
        assert Y.get(i,0) >= 0 and (not a or Y.get(i,0) == 0)
        for j,(l,b) in enumerate(types):
            e = l+b
            debt = abs(F(1,d)-F(1,e))-F(a,k*d)-F(b,l*e)
            rhs = X[i]+X[j]-Y.get(i,0)*(d>e)-Y.get(j,0)*(e>d)
            assert m*debt <= rhs
    stop = ceil(2/(F(1,3)-c))
    cases = 0
    for K in range(start,cap+1):
        for h in range(max(9,K+1),stop):
            for D in range(1,min(h//2,(h-K+5)//2)+1):
                for l in range(2*D,h+1):
                    q = min(F(D*(h-1),h), F(2*h-l,4))
                    if l < 4*D:
                        q = min(q,D-1/(1/(D-F(l,4))+F(h-l,D)))
                    assert h*(1-c) >= 2*(q+1), (K,h,D,l)
                    cases += 1
    print(start,cap,stop,len(types)**2,cases)
    start = cap+1
\end{Verbatim}
\clearpage
\section{Reproducing the distance certificate}
Save the following block as \texttt{cap.cpp}, compile with \texttt{g++ -std=c++17 -O3 cap.cpp -o cap} on a compiler supporting 128-bit integers, and run \texttt{./cap}. Each output row gives the core order, the number of completed sequences reaching the final comparison after prefix pruning, and the number of failures. The run reached 12,121,898 complete sequences and returned zero failures. A failure also exits with a nonzero status. The counts refer to degree sequences examined after pruning, not distinct graphs.
\begin{Verbatim}[fontsize=\scriptsize]
#include <algorithm>
#include <array>
#include <cstdio>
#include <numeric>
#include <vector>

using Wide = __int128_t;
using std::max;
using std::min;

constexpr int cutoff = 12;
constexpr int max_order = 65;
constexpr int first_degree = 8;
constexpr int max_q = 10;

struct Certificate {
    int h;
    Wide scale = 1;
    Wide metric_upper;
    long long checked = 0;
    long long failed = 0;
    std::vector<int> degree;
    std::array<std::array<Wide,max_order+1>,max_q+1> tail_term{};
    std::array<std::array<Wide,max_q+1>,max_order+1> prefix_sum{};

    explicit Certificate(int order) : h(order), degree(order) {
        for (int z=2; z<=h; ++z)
            scale = scale / std::gcd(z,static_cast<int>(scale%z)) * z;
        metric_upper = (cutoff-1)*scale;
        for (int j=1; j<=cutoff-2; ++j) {
            const int width = 2*((j-1)/2)+1;
            const int packing = max(1,(h-(j%2==0 ? 2 : 0))/width);
            metric_upper -= scale/packing;
        }
        const Wide cycle_upper = scale+(scale*h*(h-1)+3*h-3)/(3*h-2);
        metric_upper = min(metric_upper,cycle_upper);
        for (int q=3; q<=max_q; ++q) {
            for (int z=3; z<=h; ++z) {
                Wide term = dual_term(q,z,z-1);
                tail_term[q][z] = max(tail_term[q][z-1],term);
            }
        }
    }

    Wide dual_term(int q,int z,int earlier_cap) const {
        const int forced_leaf = (z==3 || z==h);
        return scale*(z-forced_leaf)
             + 2*max(0,q-z)*(scale/z)*earlier_cap;
    }

    Wide upper() const {
        const int leaves = max(first_degree,degree.front()-1);
        return min(metric_upper,
                   (h-leaves+5)*(scale/2)-2*(scale/leaves));
    }

    void check_budget(Wide target) {
        ++checked;
        std::array<int,max_order+1> counts{};
        int sum = 0;
        for (int z : degree) {
            ++counts[z];
            sum += z;
        }
        int forced = 0;
        int lower_incidences = 0;
        for (int z=3; z<=h; ++z) {
            forced += (z==h ? counts[z] : max(0,counts[z]-lower_incidences));
            lower_incidences += counts[z]*(z==3 ? 2 : z);
        }
        int remaining = (sum-forced)/2;
        Wide lower = h*scale;
        for (int i=h-1; i>=0; --i) {
            int z = degree[i];
            int spend = min({i,z-1,remaining});
            if (z==3) spend = min(spend,h-counts[3]);
            lower -= spend*(scale/z);
            remaining -= spend;
        }
        if (lower < target) ++failed;
    }

    void visit(int i,int highest,Wide debit) {
        for (int z=(i==0 ? first_degree : 3); z<=highest; ++z) {
            degree[i] = z;
            const int s = i+1;
            const Wide target = upper();
            const Wide next = debit+min(i,z-1)*(scale/z);
            if (s*scale+(h-s)*(scale/z)-next >= target) continue;
            bool passed = false;
            for (int q=3; q<=min(max_q,degree.front()); ++q) {
                prefix_sum[s][q] = prefix_sum[i][q]+dual_term(q,z,min(i,z-1));
                Wide lower = 2*q*h*scale-prefix_sum[s][q]
                             -(h-s)*tail_term[q][z];
                if (lower >= 2*q*target) {
                    passed = true;
                    break;
                }
            }
            if (passed) continue;
            if (s==h) check_budget(target);
            else visit(s,z,next);
        }
    }

    void run() {
        Wide lower = h*scale;
        for (int s=1; s<=h; ++s)
            lower = min(lower,h*(scale/s)+(s-1)*(scale/2));
        if (lower < metric_upper) visit(0,h,0);
    }
};

int main() {
    for (int h=9; h<=max_order; ++h) {
        Certificate certificate(h);
        certificate.run();
        std::printf("%d %lld %lld\n",h,certificate.checked,certificate.failed);
        std::fflush(stdout);
        if (certificate.failed!=0) return 1;
    }
}
\end{Verbatim}

% END SOURCE BODY Fajtlowicz_Clipped_Distances
\unitstart{Paper 13}{p13}{Iterated cutoffs}{An iterated distance cutoff for the Randi\'c index}{\subtitlevolume{Block diameter, maximum degree, and the average-distance conjecture}}
\makeatletter\@removefromreset{theorem}{section}\@addtoreset{theorem}{section}\makeatother\renewcommand{\thetheorem}{\thesection.\arabic{theorem}}
% BEGIN SOURCE BODY Fajtlowicz_Iterated_Cutoffs
\begin{abstract}
We prove a weighted distance inequality through cutoff $34$. This proves Fajtlowicz's average-distance inequality whenever each cyclic block has diameter at most $32$. Independently, a degree certificate handles every block of maximum intrinsic degree at most 25, with unrestricted diameter. These conditions can be mixed block by block. Two exact verifiers are included in the source. The original conjecture for arbitrary connected graphs remains open in this argument.
\end{abstract}

\section{Statement and starting results}
Let $G$ be finite, simple, undirected and connected, with $n\ge2$. For degrees $d_v$ and graph distance $d_G$ put
\[
 R(G)=\sum_{uv\in E(G)}(d_ud_v)^{-1/2},\quad
 M(G)=\sum_{uv\in E(G)}1/\max(d_u,d_v),\quad
 \mu(G)=\binom{n}{2}^{-1}\sum_{\{x,y\}}d_G(x,y).
\]
For independent draws $X,Y$ with vertex probability vector $p$ define
\[
 L_C(G,p)=\mathbb E_p\min\{d_G(X,Y),C\}+2\sum_v p_v^2,
 \qquad F_C(G)=\max_p L_C(G,p).
\]
The cutoff $C$ can be $\infty$. Equivalently, the kernel has diagonal entries two and off-diagonal entries $\min\{d_G(x,y),C\}$. Write $Q(H)=\max_p\mathbb E_p d_H(X,Y)$. A cyclic block is a maximal two-connected subgraph with at least three vertices. Degrees and diameters in a block hypothesis refer to the block alone. A private leaf at a site is a new vertex with that site as its only neighbor.
\begin{theorem}\label{p13-th:main}
Every connected graph of order at least two satisfies $F_{34}(G)\le M(G)+1$. In particular $R(G)\ge\mu(G)$ for graphs of diameter at most $34$. The same conclusion holds if each cyclic block independently has diameter at most $32$ or maximum degree at most 25. Neither the number of blocks nor the degrees at their common vertices are restricted.
\end{theorem}

The proof starts at cutoff two and establishes its degree estimates independently. The cutoff-twelve approach in the preceding manuscript~\cite{p13-cut12} motivates the induction. Fajtlowicz's proposed comparison is $R\ge\mu$~\cite{p13-Fajt,p13-Trees}.

\paragraph{The minimal-counterexample reduction.}
For completeness, here is the reduction used for both the capped and uncapped assertions. At a cut vertex $v$ split the components of $G-v$ into nonempty sides $A,B$. Pad each side at $v$ with one new leaf per opposite neighbor, as in~\cite[Lemma~4]{p13-DLS}. Every old edge keeps its degrees, and the new weights sum to one. On the $A$ side map every point outside $A$, including $v$, to one new leaf. Do the same on the $B$ side. For any original pair the two projected kernel entries sum to at least the original entry plus two. For opposite sides at distances $a,b\ge1$ from $v$ this is
\[
 \min(C,a+1)+\min(C,b+1)\ge 2+\min(C,a+b).
\]
For a pair within a side the extra entry is two. For a pair involving $v$ its projected distance on the other point's side increases by one; coincident pairs give four. Taking expectation proves that the two padded assertions imply the original.

Split whenever each open side contains an edge. Both padded pieces then have fewer edges. Two cyclic blocks have such a separator, and a tree branch of length two off a single cyclic block allows such a split at its attachment. Terminal trees are stars or double stars. A star has $M=1$, $F_C\le2$. For a double star with $a\ge b\ge1$ leaves at the centers, $M=1+b/(b+1)\ge3/2$, and $F_C\le5/2$, since a separated support has at most two points at distance at most three.

Here a separated support is obtained by shifting mass between two sites at distance $d\le2$. At fixed total mass the squared-transfer coefficient is $4-2\min(C,d)\ge0$; some endpoint of the interval eliminates an atom without loss. Choose a maximizing law of least support. Adding a new leaf at old degree $j$ contributes $1/(j+1)$ and each of its $j$ old edges loses at most $1/[j(j+1)]$. Thus removing an unoccupied leaf preserves the law's value and cannot increase $M$. In a least-edge exception the sole remaining cyclic core $H$ therefore has at most one leaf at each site, and every leaf is occupied. Every core edge is a nonbridge. Deleting one cannot decrease distances, so it must strictly increase $M$ in an exception. If an unloaded vertex $v$ had no neighbor of lower full degree, deleting any core edge $vu$ would leave all other edges at $v$ unchanged and permit gains at $u$ totaling at most the removed $1/d_u$. Thus every unloaded vertex has a lower-degree core neighbor. Every core vertex has full degree at least three, because an unloaded degree-two vertex would violate this condition. Full-degree-three sites must then be loaded. Deleting an edge between two such sites removes $1/3$, while the two leaf edges gain $1/6$ each. Hence these sites are independent. These reductions preserve any imposed upper bound on the intrinsic degrees of each cyclic block.

\paragraph{Projecting a law onto the core.}
For any decoration, project its leaf masses onto their sites. Every leaf of mass $t$ contributes $2t(1-t)$ above the projected expected distance. Its diagonal adds $2t^2$. The remaining diagonal terms at old atoms $p_v$ are at most $2\sum p_v$ over those atoms. It follows that $F_C(G)-1\le Q(H)+1$. A spanning tree of $H$ with $s\ge2$ leaves gives $Q(H)\le(h-s-1)/2+2-2/s$ for $h=|V(H)|$. A nonpendant bridge contributes at most $1/2$ to expected distance. The $s$ pendant contributions total at most $2b-2b^2/s\le2-2/s$, where $b$ is their total mass. The bound also applies when $s$ is a lower bound on the available number of leaves.

\section{Iterating the cutoff}
Let $C$ be an integer and use the diagonal value $2$ in the capped quadratic kernel $F_C$.
We first prove the steps through thirty. We use only the degree-at-most-thirteen consequence of the unrestricted degree theorem in Section~\ref{p13-sec:degree}. Its proof uses the elementary metric estimates below, independently of all cutoff induction steps.
\begin{theorem}
For $2\le C\le30$ one has $F_C(G)\le M(G)+1$ on every
connected graph $G$ with at least two vertices.
\end{theorem}
\begin{proof}
It is enough to establish the integer steps $C-1\longrightarrow C$ for
$3\le C\le30$. For $C=2$ one has $F_2\le2$, while the edges at a maximum-degree vertex alone contribute one to $M$.
Suppose the new assertion fails and choose an exception with fewest edges.
By padded splitting and deletion it has a two-connected core $H$ and at
most one private leaf at each site. Write $h,D,K$ for the core order,
diameter and maximum degree. Every full site degree is at least three;
each unloaded site has a neighbor of strictly smaller full degree. The
unrestricted degree theorem gives $K\ge14$. By the inductive hypothesis
the exception must contain two vertices at distance at least $C$.
Consequently $D\ge C-2$. The breadth-first layers from a diametral endpoint
have baseline sizes $1,2,\ldots,2,1$ by two-connectivity, summing to
$2D$. A maximum-degree vertex and its neighbors occupy at most three
successive layers. Comparing with their total baseline of at most six
gives $K\le h-2D+5$. In particular $h\ge2C+5$.

Let $s$ be the largest index with $d_s\ge s$ for full core degrees in
nonincreasing order, indexed from one. The top $s$ stars give at least
$s-\binom s2/s$ in $M$. Each vertex outside them supplies a distinct
additional edge, its own leaf edge if loaded or an edge to a strictly
lower-degree neighbor outside the top set otherwise. These edges have
weight at least $1/s$. Thus
\[
 M\ge B(h):=\min_{1\le s\le h}\left(\frac hs+\frac{s-1}{2}\right).
\]
If $h\ge C(C-1)/2$, the right side is at least $C-1$: for integer $s$
the required inequality is $(s-C)(s-C+1)\ge0$. Since $F_C\le C$ this
excludes such orders. Therefore the certificate needs just
\[
 2C+5\le h<C(C-1)/2,\quad
 14\le K\le h-2C+9,\quad C-2\le D\le\left\lfloor (h-K+5)/2\right\rfloor .
\]

\paragraph{Packing and contraction bounds.}
If $j$ is an integer at least one, let $P_j$ be the largest size of a
$j$-separated set in $H$, and $r=\lfloor(j-1)/2\rfloor$. Two independent draws
from any probability distribution on $H$ have probability at most
$1-1/P_j$ of being at distance at least $j$: shift mass between nonadjacent
sites of the distance-threshold graph to reduce the support of a maximizing
law to a clique. The quadratic form is affine on each such transfer.
The disjoint radius-$r$ balls of two or more separated sites have at
least $2r+1$ vertices each, by two internally disjoint paths to a second site.
For even $j$ there is no edge between the balls, and at least two vertices
must be outside them, by two-connectivity. Moreover, when $r\ge1$, the
closed neighborhood of a vertex of maximum degree meets at most one
of the radius-$(r-1)$ balls; otherwise two centers have distance at most
$2r<j$. Each of the remaining balls contains at least $2r-1$ vertices.
Thus an upper bound $b_j$ on $P_j$ is the smaller of
\[
 \max\left(1,\left\lfloor\frac{h-2\mathbf1_{2\mid j}}{2r+1}\right\rfloor\right)
 \quad\text{and, for $r\ge1$,}\quad
 1+\left\lfloor\frac{h-K-1}{2r-1}\right\rfloor .
\]
For two vertices projecting to sites at distance $t$ in $H$, the capped
kernel entry is at most $2+\min(t,C-2)$, even if they coincide. Summing the
tail probabilities gives $F_C-1\le C-1-\sum_{j=1}^{C-2}1/b_j$.

We also have $F_C-1\le1+Q(H)$, where $Q$ is the largest expected ordinary
distance of two independent draws. In particular
$F_C-1\le(h-K+5)/2-2/K$ by extending a degree-$K$ star to a spanning tree
with at least $K$ leaves. A second upper bound contracts the closed neighborhood of a maximum-degree vertex to a single site $o$. The resulting connected simple graph $J$ has $m=h-K$ vertices. Every vertex other than $o$ can be deleted while preserving connectivity, because $H$ is two-connected. For $m\ge2$ add one new vertex adjacent to $o$ and to every neighbor of $o$. The graph $J^+$ on $N=m+1$ vertices is two-connected: deleting $o$ or its new twin leaves a copy of $J$, and deleting any other vertex preserves connectivity. Folding the twin onto $o$ proves that $J$ embeds isometrically in $J^+$. In any two-connected graph of order $N$, two internally disjoint paths between diametral vertices give a cycle of length at least twice the diameter $d$. Traverse that cycle once and a spanning forest attached to it twice. This closed spanning walk has length at most $2(N-d)$. On a circular walk of length $s$, choose one occurrence for each vertex and take all cyclic intervals of $\lfloor s/2\rfloor$ consecutive positions. Twice the circular distance counts intervals separating a pair. If the interval masses are $b_i$, the expected circular distance is $\sum_i b_i(1-b_i)\le s/4$. Hence $Q(J)\le Q(J^+)\le\min(d,(N-d)/2)\le N/3$. If $m=1$, $Q(J)=0$. Contracting the hub neighborhood saves at most two along any individual shortest route, since its boundary sites can be joined through the hub in at most two steps. The statement includes endpoints inside its neighborhood and endpoints mapping to the same site. Thus $Q(H)\le Q(J)+2$ and $F_C-1\le3+(h-K+1)/3$.

\paragraph{Lower certificates and exhaustion.}
Write $a_v\in\{0,1\}$ for the load, $k_v=\deg_H(v)$ and $d_v=k_v+a_v$.
A row $(K_*,m,t,X,Y)$ of the checker applies whenever $K\le K_*$;
the types $(k,a)$ occur in lexicographic order up to $K_*$ with
$k+a\ge3$ and either $a=1$ or $k\ge4$. The checker verifies
$Y_i\ge0$, $Y_i=0$ at loaded types, $k_iX_i-Y_i\le t$, and
\[
 m\left(|1/d_i-1/d_j|-a_i/(k_i d_i)-a_j/(k_j d_j)\right)
 \le X_i+X_j-Y_i\mathbf1_{d_j<d_i}-Y_j\mathbf1_{d_i<d_j}
\]
for every unordered type pair. The sum of the expression in
parentheses over core edges equals $h-2M$. Every type with nonzero
$Y$ has a lower-degree neighbor, so summing yields $M\ge h(1-t/m)/2$.
Together with $B(h)$ this gives the lower bound used in the checker.
The displayed finite ranges exhaust all possible parameters. For each
tuple the verifier compares the better lower bound with the smallest
applicable upper estimate. It checks rational inequalities without
floating-point arithmetic. No tuple fails, completing each induction
step.
\end{proof}

\section{Unrestricted distances when the intrinsic degree is at most 25}\label{p13-sec:degree}
\begin{theorem}
Every connected graph whose cyclic blocks have intrinsic maximum degree at most 25 satisfies the ordinary weighted inequality. Each such block may be freely mixed with every other class for which all private-leaf decorations are certified.
\end{theorem}
\begin{proof}
Apply the previously established cut and deletion reduction in the class with this degree bound. Deleting an edge cannot increase the maximum degree in any surviving cyclic block. Write $H$ for the terminal core, $h$ for its order, $K$ for its intrinsic maximum degree, and $D$ for its diameter. Each edge of $H$ must strictly increase $M$ upon deletion, since deleting a nonbridge cannot decrease any kernel distance.

Give a vertex the type $(k,a,t)$, where $k$ is its intrinsic degree, $a\in\{0,1\}$ its leaf indicator, $d=k+a$, and $t$ the number of neighbors (including its leaf) of strictly smaller full degree. Then $t\ge1$, $a\le t\le k+a$, and a vertex of full degree three has $t=1$ and $a=1$. A vertex has exactly $t-a$ lower-degree core neighbors. For a core edge with endpoint types $(k,a,t),(p,b,u)$ and full degrees $d<e$,
\[
  M(G-xy)-M(G)=\frac{t}{d(d-1)}-\frac{e-u}{e(e-1)}>0.
\]
The change counts the surviving neighbors below the original degree at each endpoint, remembering that $x$ itself counted as a lower neighbor at $y$. For $d=e$ the change is $(t+u-d+1)/[d(d-1)]$ and strict increase requires $t+u\ge d$. These conditions and the exact lower-neighbor count rule out the pairs excluded in the verifier.

Put $B=h-2M$. Directly from the degree identity above, $B$ is the sum on core edges of
\[
 w_{ij}=|1/d_i-1/d_j|-a_i/(k_i d_i)-a_j/(k_j d_j).
\]
The two rational certificates give $B\le ch$ with $c=24472/100000$ for $2\le K\le14$, and $c=38932/100000$ for $15\le K\le25$. For scale $m=100000$ and integer entries $X_i,Y_i$ the supplied verifier tests
\[
 mw_{ij}\le X_i+X_j-Y_i\mathbf1_{d_j<d_i}-Y_j\mathbf1_{d_i<d_j},
 \qquad k_iX_i-(t_i-a_i)Y_i\le mc .
\]
Each $Y_i$ is charged exactly $t_i-a_i$ times when summing the first inequalities. This requires no sign restriction on $Y_i$. The ordered type lists and compressed integer arrays are included in the verifier. The packing bound below and the lower bound $2M\ge(1-c)h$ finish the argument.

\emph{Metric bound on the hub quotient.}
Use the preceding isometric embedding into the two-connected $J^+$ on $N=h-K+1$ vertices when the quotient has at least two vertices. We first sharpen the expected-distance bound for any two-connected graph on $N$ vertices, with diameter $d\le N/2$. Put $w_j=2\lfloor(j-1)/2\rfloor+1$. The disjoint-ball estimate proved above gives $P_j\le\lfloor N/w_j\rfloor$ for $j\le d$. The closed-walk argument gives $Q\le(N-d)/2$.

The distance-threshold clique argument gives $Q\le\sum_{j=1}^d (1-1/P_j)$, where $P_j$ denotes the maximum size of a $j$-separated set. Set $f(0)=0$, $f(x)=1/\lfloor1/x\rfloor$ on $0<x\le1/2$. Since $w_j\ge j-1$, we have $1/P_j\ge f((j-1)/N)$. The error of the left endpoint sum for this monotone function on $[0,d/N]$, multiplied by $N$, is at most $f(d/N)\le1/2$. With $x_0=197/500$ one has the rational bound
\[
 x_0-\int_0^{x_0} f(t)\,dt
 \le x_0-\sum_{m=1}^{32}
 \frac{(\min(x_0,1/m)-1/(m+1))_+}{m}<303/1000.
\]
If $d/N\le x_0$, monotonicity of $x-\int_0^x f(t)dt$ proves $Q\le303N/1000+1/2$. If $d/N\ge x_0$, the tour bound $Q\le(N-d)/2$ gives the same result. Apply this estimate to $J^+$ and use $Q(J)\le Q(J^+)$. A one-vertex quotient satisfies it directly. Lifting back to $H$ gives
\[
 Q(H)+1\le 7/2+303(h-K+1)/1000\le 303h/1000
 \le h(1-c)/2\le M,
\]
in the second degree range, where $K\ge15$ and $c=38932/100000$. This suffices for every decoration by the projection inequality. For the first degree range, apply the same mass-packing and tour proof directly to $H$, using $h$ in place of $N$. Indeed all disjoint radius-$r$ balls around packing sites have at least $2r+1$ vertices when there are two or more sites, by two internally disjoint paths to another site. For even $j$ the balls have no joining edges, so at least two vertices lie outside their union by two-connectivity. Thus in $H$ we may put 
$P_j\le \max(1,\lfloor(h-2\mathbf1_{2\mid j})/w_j\rfloor)$.
The tour bound is $(h-D)/2$, by one diametral cycle. Hence $Q(H)\le303h/1000+1/2$. At $h\ge21$ this gives $Q(H)+1\le h(1-c)/2$ for $c=24472/100000$. For smaller $h$ sort full core degrees $d_1\ge\cdots\ge d_h$ and let $s$ be the largest index with $d_s\ge s$. The top $s$ stars contribute at least $(s+1)/2$, losing at most $1/s$ at each repeated incidence. Each outside vertex supplies its private leaf edge or an edge to a strictly lower-degree outside vertex. The $h-s$ chosen edges are distinct and each contributes at least $1/s$. Therefore
\[
 M\ge \min_{1\le s\le h}\{h/s+(s-1)/2\}.
\]
The code checks $3\le h\le20$, $1\le D\le\lfloor h/2\rfloor$ directly with the better of this lower bound and $h(1-c)/2$, against $1+\min((h-D)/2,D(1-1/h),\sum_{j\le D}(1-1/P_j))$, using the stated upper bounds on $P_j$. This proves the decorated assertion throughout the first degree range as well.
The verifier checks 13,741 and 138,793 admissible unordered pairs from 214 and 665 local types. All arithmetic comparisons are rational.
\end{proof}

\section{Combining the estimates at cutoffs 31 through 34}
\begin{theorem}
For every connected graph of order at least two, $F_{34}(G)\le M(G)+1$.
\end{theorem}
\begin{proof}
Starting at cutoff 30, suppose the assertion first fails at $C\in\{31,32,33,34\}$. A least-edge exception must now have $K\ge26$, $D\ge C-2$, and
\[
2(C-2)+26-5\le h<C(C-1)/2,\qquad 26\le K\le h-2(C-2)+5 .
\]
Use the order and local degree lower bounds as before. If $K$ is beyond the last available degree row the trivial lower bound zero is valid. For $F_C-1$, in addition to the packing estimate in the original core use the three bounds
\[
 \tfrac72+303(h-K+1)/1000,\qquad (h-K+5)/2-2/K,\qquad 1+(h-C+2)/2.
\]
The first follows from the hub quotient, the second from a spanning tree with at least $K$ leaves, and the third from the diametral cycle in $H$. At cutoff 31 the first program's joint check eliminates all pairs, including 50,354 for which it uses the original packing comparison after the three displayed bounds.

At the next three cutoffs use a further inequality. Put $N=h-K+1$ and define
\[
 b_j^+=\max\left(1,\left\lfloor
 \frac{N-2\mathbf1_{2\mid j}}{2\lfloor(j-1)/2\rfloor+1}
 \right\rfloor\right).
\]
Any kernel entry between decorated vertices is bounded by $4+\min(d_{J^+}(x,y),C-4)$ after projection to the augmented quotient: passing from the decoration to $H$ costs at most two, lifting a path in $J$ costs at most another two, and $J$ embeds isometrically in $J^+$. Diagonal entries two obey this bound as well. The earlier two-connected packing estimate and $\diam(J^+)\le\lfloor N/2\rfloor$ give
\[
 F_C-1\le 3+\sum_{j=1}^{\min(C-4,\lfloor N/2\rfloor)}(1-1/b_j^+).
\]
A singleton quotient obeys the same inequality directly. The verifier uses an additional degree row through $K=320$ whose integer potentials pass all 202,566 unordered local comparisons, together with the preceding bounds. For each of the three steps $C=32,33,34$ every pair in the displayed finite ranges passes. Thus no first failure occurs through cutoff 34.
\end{proof}

\section{Original distances and the remaining range}
For $C=34$ the capped quadratic form equals the ordinary one on graphs of diameter at most 34. It also agrees on every decoration of any core of diameter at most 32, since two terminal leaves can add only two to a distance. Apply the cut rule with $C=\infty$ to assemble the certified pieces. This proves the mixed-block weighted assertion. It includes arbitrarily large graphs with arbitrarily large overall diameter.

For clarity we recall the conversion to Fajtlowicz's original inequality. If $F_\infty(G)\le M+1$, the uniform law gives $M\ge(1-1/n)\mu-1+2/n$. Put $k=\sqrt{\Delta(G)}$. On each edge with endpoint degrees $a\le b$,
\[
 (a^{-1/2}-b^{-1/2})^2=(\sqrt{b/a}-1)((ab)^{-1/2}-1/b).
\]
Summing gives $n-2R\le(k-1)(R-M)$ and hence
\[
 (k+1)(R-\mu)\ge n-2\mu-(k-1)(1+(\mu-2)/n).
\]
The bound $\mu\le(n+1)/3$ follows by passing to a spanning tree and repeatedly removing a leaf. At order $j$, that leaf's distances to the other vertices, in increasing order, are at most $1,2,\ldots,j-1$. Substitution of this mean-distance bound and $k\le\sqrt{n-1}$ shows that the last display is nonnegative if
$n(n-2)+4n-5\ge(4n-5)\sqrt{n-1}$. For $n\ge5$ the squared difference is $(n-5)(n-2)(n^2-5n+5)\ge0$. For $2\le n\le4$ each edge has ratio $(ab)^{-1/2}/(1/a+1/b)\ge\sqrt{n-1}/n$, hence $R\ge\sqrt{n-1}\ge(n+1)/3$. This establishes Theorem~\ref{p13-th:main}.

\paragraph{Obstruction parameters.}
A connected graph failing the original conjecture would contain a cyclic block outside both proved classes. Such a block must have $D\ge33$, $K\ge26$ and $h\ge2D+K-5\ge87$. In particular, the original inequality holds whenever every cyclic block has at most 86 vertices. If a least-edge counterexample to the stronger unrestricted weighted statement exists, its terminal core must further satisfy $M<D+1$, hence $h<(D+2)(D+1)/2$ by the order bound.

When $D=33$ the decorated diameter is at most 35. The cutoff-35 diagnostic leaves precisely 893 integer pairs $(h,K)$ under these estimates, with $159\le h\le500$ and $86\le K\le336$. In fact a least-edge counterexample to the unrestricted weighted inequality with this core diameter would also be a least-edge counterexample to cutoff 35: its own distances are at most 35 and every smaller graph satisfies the unrestricted inequality. Thus its parameters must be among the reported pairs. Their survival does not establish existence of any such graph. At this first diameter and for any fixed probability vector the only omitted term at cutoff 34 has the form
\[
 L_\infty(G,p)-L_{34}(G,p)
 =2\!\sum_{\substack{\{u,v\}\text{ loaded}\\d_H(u,v)=33}}p_{u'}p_{v'},
\]
where $u',v'$ denote the private leaves at their respective sites. The general case with unrestricted $D$ and $K$ remains open in this argument.

\paragraph{A condition without diameter or maximum-degree bounds.}
The same contraction argument gives another sufficient condition. Suppose $t\ge1$ sites $v_i$ have pairwise disjoint closed neighborhoods and degrees $k_i$. Contract these $t$ sets. The resulting graph has $m=h-\sum_i k_i$ vertices and at most $t$ articulation vertices, all contracted sites. A shortest route in its block--cut tree crosses at most $t+1$ blocks. If its diameter is $d$, the same closed-walk construction along a diametral route saves at least $2d-2(t+1)$ over a doubled spanning-tree tour. The quotient thus has expected distance at most $\min(d,(m+t-d)/2)\le(m+t)/3$. Lifting adds at most $2t$ to each distance; the singleton quotient also satisfies the estimate. In particular, every decoration of $H$ satisfies the uncapped bound whenever
\[
 3M(H)\ge 3+h-\sum_{i=1}^t(k_i-7).
\]
Indeed $F_\infty(G)-1\le Q(H)+1$, and adding leaves cannot decrease $M$. For example, take a cyclic chain of $t\ge3$ disjoint cliques of order $b\ge9$. Join each neighboring pair by two independent edges, using four distinct attachment vertices in every clique. The graph stays connected after any one vertex deletion. Choose one unattached site per clique; their closed neighborhoods partition the core. Internal clique edges alone give $M(H)\ge t(b-1)/2\ge4t$, while the sufficient threshold above is $1+8t/3$. Thus this criterion applies even as both $t$ and $b$ increase without bound.

\begin{thebibliography}{9}
\bibitem{p13-Fajt} S. Fajtlowicz, On conjectures of Graffiti, Discrete Math. 72 (1988), 113--118.
\bibitem{p13-Trees} M. Cygan, M. Pilipczuk and R. \v{S}krekovski, Relation between Randi\'c index and average distance of trees, MATCH Commun. Math. Comput. Chem. 66 (2011), 605--612.
\bibitem{p13-DLS} Z. Dvo\v{r}\'ak, B. Lidick\'y and R. \v{S}krekovski, Randi\'c index and the diameter of a graph, KAM--DIMATIA Series 2010--496, \href{https://iti.mff.cuni.cz/series/2010/496.pdf}{authors' manuscript}.
\bibitem{p13-cut12} Keio and Shaitetsu, \emph{Clipped distances and the Randi\'c average-distance conjecture}, preceding research manuscript, Theorem 1.1 and Section 4. Earlier degree-padding and local-certificate methods are collected in \href{https://github.com/QEEDDD/Verifiable-Mathematical-Reasoning-Corpus/tree/38348353416f1f7fc50cecb93c3b554a00260c13}{VMRC, snapshot 3834835}, manuscripts 087, 094, 106, and 161.
\end{thebibliography}

\clearpage\volappendix
\section{Reproducing the exact calculations}
The single source contains two complete Python programs, including the compressed arrays of integer coefficients. Extract the lines between each pair of \texttt{Verbatim} markers into a separate file and run each with Python 3.9 or later, with assertions enabled. Both programs use only the standard library.

The first program checks the local rational potentials and iterates the distance bounds, starting at cutoff two. Its first 23 rows entail 378,220 unordered comparisons. For cutoffs 13 through 30 the loop visits 3,345 pairs $(C,h)$ and checks 330,279 triples $(C,h,K)$ after whole-order pruning. The final checks combine the metric and degree bounds and report the unresolved parameters at the next core diameter.

\begin{Verbatim}[fontsize=\scriptsize]
"""Exact verifier for the iterated capped distance theorem. Run Python 3.9+."""
from fractions import Fraction as F
from functools import lru_cache

def verify_potential(row):
    K,scale,target,X,Y=row
    types=[(k,a) for k in range(2,K+1) for a in (0,1)
           if k+a>=3 and (a or k>=4)]
    assert len(X)==len(Y)==len(types)
    for i,(k,a) in enumerate(types):
        d=k+a
        assert Y[i]>=0 and (not a or not Y[i])
        assert k*X[i]-Y[i]<=target
        for j in range(i+1):
            l,b=types[j]
            e=l+b
            right=X[i]+X[j]-Y[i]*(d>e)-Y[j]*(e>d)
            # Positive common denominator k*l*d*e for the edge debt.
            left=abs(d-e)*k*l-a*l*e-b*k*d
            assert scale*left <= right*k*l*d*e
    return F(target,scale)

@lru_cache(None)
def order_lb(h):
    return min(F(h,s)+F(s-1,2) for s in range(1,h+1))

def pack(h,k,cap):
    bound=F(cap-1)
    for j in range(1,cap-1):
        r=(j-1)//2
        size=max(1,(h-(2 if j%2==0 else 0))//(2*r+1))
        if k and r: size=min(size,1+(h-k-1)//(2*r-1))
        bound-=F(1,size)
    return bound

def verify_caps(rows,last):
    table=[(row[0],verify_potential(row)) for row in rows]
    for cap in range(3,last+1):
        orders=pairs=0
        for h in range(2*cap+5,cap*(cap-1)//2):
            orders+=1
            b=order_lb(h)
            if b>=pack(h,0,cap): continue
            for k in range(14,h-2*cap+10):
                pairs+=1
                coefficient=next((c for m,c in table if k<=m),F(1))
                lower=max(b,h*(1-coefficient)/2)
                upper=min(pack(h,k,cap),F(h-k+5,2)-F(2,k),3+F(h-k+1,3))
                assert lower>=upper,(cap,h,k)
        print(cap,orders,pairs,flush=True)

# Packed JSON rows contain [K, scale, target, X, Y] in that order.
import base64, json, zlib
encoded = """
eNrtnAuSJDluRC8UZUaC/7OM7f2vIX8ORta0ptd2Je2ORqPpLmdF1iczgwQBhwOsX36p+6mF
f0+bpZfnlzrnXE+NUuOJ0vp41tZQg3GsNZ9Ta3l66/uZK87T6lhPP709Ued42hj9qa3pt7u+
NkO/O5q+dB5++WuUqW+P56vqi70OY59q1Fb2HaZ+2EOUEh7+9vxSHr2ynq3rc+y6jz7XPYOv
613Nps9tlPGUf/b/3/SsEe8c9NLH/sxB6+vopnZtz2kaamecZ+jFiubjGUO/uk5pmo5Wn1nr
elpr5+kr9hPcapv6Wl1HT7S4rXr0PDzDnL5TXqA8axZNiL7Fs37Fbry4MXTzYOsJQa1R79B2
v8PkDjxsvTjDO1dzaxk1J/oNz1mNyePa99r6vPsZ/nKt//yU/Wb6xmf6gjt7p0+zpYWLifWU
yjgY16xa2eI5qFt3NaYm7Jwlc9KN96LlG+10vekVekp9LTqG1ZZudQ/NmJ+h2aiWjGfxYjE1
6/3x4ujbozVNZGFesLXpWTL6GMaS6YMzplF5uRzaKncYo96BH2V4Z3ZNW2GLZSuMYA31G6v0
ZmtsekflGbPs/97U/maaPztVO2B+T/OoemehjRqeVY2DcWuW9C50rXmtmhh9D4Ptz8Lq9Bwy
rBPs+7OfvrHEoy3Yjl4h2taEawXZXRVL1tP2YtPNzaKH2puTnyhakqbnnSeetvV8mvDQE2sm
v5beg74MeivG7NPYmj+grTTu4B/10Pq8Q5eJ5MA2Z3hXQA4Jq5U3KNh09LCty/zL4nPDDvQ2
TtjGV+v/84X49YK08lmQuWR6nwUZhxk8mpIqu8TyGPdkb8pi5TKKXYpsfGw8CDbSN3MTcpON
vTCqvtZCz9QrzmNqn+tOtaEqO2AXmxr7C+PUyumh9sLpTJ582Ay9AsvZtK9k/fJkTYvxpXf3
8NRfteDhqiFHZwxcPFtUpgu2PBKoRd/MAZ+bQ7MPYujsbQ+6dQ+fJTqRLpmlLPaI1Ut0drcb
KofHh1u0dax/3RL9sFTfHn6XXj9LNYscX3QCim6KcdoscSs12CcrNHNnar9pkfRw6zuj4Kz6
lscJhZzRlsOeHuKn7AMbzj5wSq3wQmcSDwgFgx/hPSjE2qjlwJjOUBSdBbPApLVpG85ea7Pw
X+y/Ngma+vHKB+8L9FmMqdcCaxeDmQfanuMOde07BIvqoZd6h3GHd+1OsyNTANvLa6d4y5qd
4ce1n1zL0hr7a6ziaFx2+9ev4Wcd+2cdj5bwex3bZroX1hjBuBhPEBrkCDUlA1e/ZWezEUrl
n4kREw4jHuIwO/W1tuQmu5y43N9mL8roYzr8KJrLrenhIlrIjvRwMxHsC+1uwvL2vMiv4cgU
14bsnPjUHO0I6ytXVFuUjaz9qK0P8Xh4w4BbAUM2Atj5YMubAfkNQ2va7+DQ5gF7zKGNdYdu
N9rnZ2GX+ZL8fmcBFaCPF1iMzws7Vk8+pW2rz1PW7/Vf59+zsD8s8HwXWCsS30HOo9wHNhqD
CV+Mx5exsfnVMHZ53GcOWe6W81S8k3fUfRKvtf7yttow7PnuyCe6J/eqefMztMEKdDnkwElr
PrTDIEh6dR4unB4Bb+kJNPEr2RfcToEX+9ECQYg1bQMyq6WFs3zZHpfWmoip/+YYAgwBdAiQ
IG9kYJEAMgeOFhD82g3bwBgCHx7pyMc7nBw+K74rW1SGl2548XZZ8ZkRU1TGKy6yZlZYbBAt
A+fv8N8r/6E3IjR9/3zlF2x3MR5fNhj3Otx9gW3O1SCUcnsyYvlC5xiHbxR9rQfUEgtpQ4vW
2R0KOYPIOGGVeohX0S7U3p6HQKUfEbVr5DB6g4fkJcyxCAsK4DK+afpeeQXMcUGdBuZATG4O
u4dflQGQJMFL8eaJkIGArmcGA6sVYJdg6ZnAllEB/NGxGbR3+M4IAlaQgwMYw8ukPiF5ezu3
teyocWa2CdtV8le7eU2MQ3M5ROrBRNsnnN/HJF6z6B+Spd05f+4QWlEEkkNuTkQY5di1dWr3
PtVsT+WsImDwqkMiSoY5ilZuio/JEBaTjt+VF4Eaw+u73bCDrB62blas19xiU0qMIZz2Q/X1
DUR7uSXiALkdQfkQsHgJItMZvIqsRdv2ZVsiq3LEMkgCKy4C3yQL6SYsRrRmNEUdIE9mDAVy
MPXaAAoKiETg21J4Xx4qs5XDJ9cJZthDK93DJ8PZzmiaJjbz7TKd6Ygk2EJWb/YebQaPz7a7
0E5zVIm6fz9L+cFiPlxP67h/4kh4ywrpmnuPxaM9wm6OIaRLqyj928pNdUOaZHkTLViYl7EG
eIPB+rejSVCafYgmCwFhXwbYoe04hIZnCNhETPzHXgR6IviQq6i4J70h3mF1JHHwViiBW5Db
Tf3UsQwg24GY6A5xceYnvBm9or79JVNCRhjYNx8YFMAxAiUSxohiTAUFwHMamiywdauAXML5
BNnKHXa9g/N5D2Gf6IEsW8NHfDjNcUTz7AxM/sumUY8/1+uCtE2bDap0+xrdr+0uyvj9LegH
S+rflnR6/7nv6fDp7bF4bBtfo/yNmTqkwXKdSqrRwsLyR+f5FOoHwtG0jnQQPpazwJkSmlZK
caQcJCRNs8JRIzodfMHEbTcYyspUXXFGNlUglZ2VX5h0mhhJuzc/q7OdCKHcOZw4HCg9lB3p
OUfm91juYx83TVNhMHrOASWzmgIIieHgltBtGEPxEkyIa6/2e4buA3xsSy7FONuwiJADJn8H
bsVD+AbiE8Sci2FZfZm67jNsYeJBKQnWMm1Zqb7o6Sq/t5LvdCgA9j//9+zrY2PrY2Nyq3/H
xryc22PxqLxZ3mp1T5dix+okBgF/RGpYqH6WIebS17TIECI0GBKiQYhsGFC3zBasOWmvpjGw
BFLV5kTWGckk90NrI+aRptZJ4Djs9g4rWuZmCDidzIPYqIiyIGLe1NrRrLqso+Ft0Ji6qRKJ
Si/8XkmLmhisPBlk/PmSfSx/VIjM8iYw2kggMANCvOEAv6ytgiWrBqglxqwGHMBAMFyWoNY7
rHeojpjxK2HwtT6lxtvWt9KvHQgDc7OTQu1habA6XmB9x7RbxIFvj2aCFUxM+aP8xxTHN9VS
xtx+Gjg7sUxcltGMVxZArLIx1GbCehAxICXOZc5BG2BBNiwLAUhLxCzKyQ0ku4Zh9bGsQ8lG
lfihmrDEOL6V4QvHh0C+0CMRr5Bt8Lwo6EQpi8LbRAxHNj8SiyW4EmgsgR3yCJFLz+u3tHBR
SvplhEoQyfvRA1gg1F0ZCLK5rBH/9HzBF3mgrYJQSSxEjcEnohL0hFKNxKnGGNtApDcI1sLC
RwI44ln2yQAR1pCtGNqjQOty3uE7FpNPaHjtU8HFAfQwk/pcShZMWksFZ59m7zi2/aBs3Ong
GWb8yjrsW1ukKjD+OFb6g7V+B+eFAPxTa4UqyW85WfYoGqeJDIyhWhZycQfuN1HAd00Pgpva
JGPYLbl7RwIYXlHIisKOzFO/okUm48et2qOGdQVHaKw/jgN2xTjQbhENgvgq7kn83tS4YAq7
Y6jk7rM5qXf9hFs0RyfxCxMKOTQyT/3Q1K+QmgzZ7kRFEu3AwbaS5sfKQftIKQLYiT1faBn6
JMeGTW5LUgAjNyi4kJbI+4O3RofeDyjmGfsY3JhxEh+T3svYoiYGuQkOV89ljJ5gJ7GViisc
HhxGZnsNWnHCtRatpBX/Mrf1i0bpjcypx9UxMuxrhv24DFe5FgI4CU1PctD+gCb9UxM/HxPf
MMmfmjhiXrWkR4hiVxesZVjSRsHZhUknl1tE5k0K4LQ3s+AOS2jIIWYJuM1ObPXTa4IrnhrX
LJbWCJaNtLrhzILiRUPZkIFOMuOGWAJ9IMXByYZrudvslHraxNjQMSqFVjHJ+erQy9WStwTs
e6uDbPZY1ZNJLbvrqRQNAZZM2m82sthHxbmTKneueW4ydLa0CatilJw37FX8Qs+CG+alEy7K
NtI1482YoEKgR028qTdZFBAXA2xhQ4THOM1AkjRWMbQUCb1vw9VLsruduNIuXsNIdvxuA/l4
K/JnmeU2sy/ucriKWw8/zyTvvTOvyq87CSCbdViYKeW07orwavWPvxl+tjnmp+o8Sx8/3xwe
wz8EFUDFM1FmMyhesgLNobRacVUOsVN3oyiHUQ2kALn+nvW4gXbWIcMDuturxVsLiMiLVLHk
+pWxkdi24qoavh/OYi0iWIkgEQpzWV5DfpqHBPedOjPGy5u1AkDdQmkf2jPMn99Irg3FqsUS
pDbfMUcVZdmDOkel4PHwdufgOjMzhwl4uDYqaSObjrJGc8FS75FQxL6RhzzmUIpgEKAvvtn0
CW2cD/y6Q1RCNwSiX5xu4K4NJFFeL3bizRbLxa07Qr8ASqxB1QS8bAlhBvQLhA5h0yDhJomT
YCUBpTIB2dZoSHJvAbmd4X3y+SwTcVgZCKcOoC55KANduZ+SV8lYuvkS3pDKkoOOVnZnmaRk
Vez/3rb6e1ttfaiWdsKY/3irddJF1wuRHlyWoiJjoy5UOtbAM5NGzXnoMEL5PrKouSFgbdNM
Qx8N+vsg9HSHKKiwLqi0Uj+hot7ZPnL+FLRQcBErtQkpkK6ZeYRceli+lGFjvoMMl0pqwG8q
zR9RiUmLqo13kmv1/EblxStOXfuP8m0hZHTvP5LElgUpbrSnPO66fzHpwdqSDCHrTIspi0xc
t42iAoly75B3Afu7N+uF1I8pRXl3ihw+fkUiFaRqumjvWq94HGxNnyBwfDjTIh1IjJE4iWgX
FLiYQt2O0XvCRFZ8lxYvQITtv4p8rsyVT5mCqrfx7lzdh6F7Mux1WNdiiIIkUMIAVRowV+Il
jWUm4iKbRd6dK4f2w2dZnIsZs6yUqJu1aK1gzUxnpuLYemY+NXoWN5yYz+1Gks4bZdrdIhVz
/Tl28D+zw/d3MBXZGz9VoTyyw7F4d4VMy3JlLrO2gZ9tKRLTKEHWSh4z0ea0lrQnUB6jXdA0
06zThaWxHJhcH6E3KVzmaLQTuU5X0UQRhbz/HVkmlfLSvPXdD0PAnZaqeEj7HVXQmPXVsJS4
4dKtvZO4I8/zG5UXTxV+ogCtadl0XLLaXMmtrvRRHznuf+Eh3yjH+g5XxQ9pjzR/XVm61fZc
CK/9MkK2J8wVNwcJbSVb/Nj7nXY/IiXK1rImy6tTYYTRMpkUqXg5F2zQDk/KEjU1W2RP0kF3
iuAhvmAKOGD5Tj42hXHngcYoiVfFeFNF9Dhw9bYgtqNEw/cBywBeBZgqA8LIq3pcXY52GINW
GfLpshJxoe1n6GcMmBXtDiUSVCkAZgRGS6ySkN0CNEWjzgRe1553JKiJfLqSlLukfvd+tu7J
3e7shozsYREds75HXcYxvYdVZvGtuDF+DzdR2n2IaE23p+7bVeHms1r/f3iR/4q3scFed7N+
6m7Kj+6mudczfybLGg2zOK7JV3oPteFDKRNSHjW8zV6a0LrVrehs9wYcWCZtVhCKg/eBhyh9
smnSIAChICltxykfhGKiqdDr1BwcXS2B9DbEhnC7FvXeWB/JnIAoCtlxQTxE2eaZw7/Bi4d7
CA53W1P+ITxDKDBbBGScEcK3azjuIcRp0T9Zh5s9rGY1tyNsl6HY0ec2/5AIpJcqrkAX97DZ
YnHSDwIsPUqQYhoptpv+3N720HON77Wbsj6Ky0K8J8RTMlrWf2j467gqmmhwV2x0kolg0nRn
dB8m3WKamILH75VkgztC3jg3CZ9Z63YtwxXHB3pPDqTleb74Lbm0/fDfyf5DFIEn8FFd5CIx
6gl3PKGtXZCrULv3/JNtjQRFD0DeZV0se+o+/o6GMGCN73zUXdpvDRjmyvwJUMswbqW1W5MI
+gIT4yQum4InGzBaG+HNh26FlpU2eiToLBpWaRI0I4J9cRKzXNwKL6ZktAuiUX/ltXZzlv/8
mcaH1I3Fn6wvTD+uq7jhFlmhXFaVOsPu2REYt94x+LFVTKtGz+787RYj99QU29dfTvHf5Wa/
e2eW3EL5h6zuBzdrjaSHW2i8PevpZNTZTFKdtC0ObcjHloIUTcIub8mLlWy1m31hxOQN9qd0
wQ77DWjZKJYDKH3TX9vhTB3v3RBqe7HMWIdblWi/DApCeNsCNwqL7RudDm+L9Hj6S/j4jeDF
A7OPZue76ZauV4AP/EWYwJJP+gSEOKGPm1BWouEWT1nRG/SQJBDni/SWBwKIC26Cr/j72p0s
bk+jtfzyacDCx5lAuv/WkqcJWMkG/azI22ffUnjNYwzFRSBUcvq/jr31rsslTJJb82zRIu52
scotE0uoEB4b+mJ3AMmkfoGLwZ3gOkq1RyePhoASMGkdQB9jhezOaIijOQlPzsxT4IKLQw0p
FUOrmWS8O10oyxvbsmu6WTfGImldNzyTtPqAEYW6bHHmGh6LykunI20s5KDuZ8FGmDtnubTG
KA7o7T6sCh8ZDfSsncDuX3XNOBWqg4j6nQ2/8YG42Wa+gZ6duMYqicuH36MLga2AehEXtzUL
zdnQQhmfeFITu1zsBDVFOHOZidoTL7+GVBg70WditARsxEh9+hOHkB3BjUcfXk6jCXiz+roS
booEJ9FmosfFTtCPAm6v+YfXv3GM1QLaYomdYG8b27DiYVBfePm4ZiQbkX/zefgITVj7QV1a
qa3tbDHTimWW3yP1uTjZulhXfr2k7n3CIsGmPcZaotvV6cZ2rdQhr87sbXfBfrma9Nf/P2VU
jPodFQcVkV84PWDJjwAhc7ZbrybkbqamUeP+TJ5B6pnr92wpoJTZ3YxroYrs/Qw/O0o0JZqJ
ErjIRiaVXKWZ/bZEUJAfbVgHrHkIbqBT9Lm8W9iOtO0sE0qT7lh5REuOGkewV/Zpu+TZOIPR
7EkGR0dQBRqOJI77AAmStDS2UqF3MGCasWK6KkH+0t2liLfN2cBRkVoFtxUc1ws7MPqPoiLz
Fwoczujo9gu2pusYldMFqahuqyxOZ7KbHi8bdrvM7My6DEkAOQ1BwtVL17+7jxrtdpOd5u5e
fiT8Deb9aqNu94jsPLqHIZwFuVXFBWVULrt/r2ghKSjVZyjoXsUsvJbUT5Qa0bWHYkIUIfMg
26A70vWFnfUG1pUaHormKA7KVjXH1SOcQhGcqQuObNW1KkrNMMzFOSmALxrJkjiip+twTdPX
bhdsmR+4DI+QQ0dC5LXpVKT/xmBq+vOGz61vPLhxwmX9ntcEdK4jA3rmOdhGXgeWhaZPgzq5
k+PZgDfmddt5TcCHX9UM+LQdcJ3GMaxec+0+Ph9WhVSR/rW8ZqlHHnb1tWtbPtSS18XL+ZgZ
QTi5oq7D2SwuphNFZpRkigv0ahODTvbNBSckaeRyPslnMYmegx/md/ip4d+o/mDCIB/W5fwx
/TquHafOfrmGs9BNe6NsxR/uG/KxLWMn6KE1dgLbNHaCt2rsxFtf+/CUdrES90jrJ89d9SL7
5eouF/1iJy6veTuEwzlAzbUzTuLlO6yt17dc9IudoBxhvBWGk6DuYry1wZOgU8aYiVku2sW8
uPrj7TaNNS9uno5nMsbFy6/qRb94+Va5aBeXf1F9AXaa9ZuP+UQIOImXn2Ejxkp8+Fq7GBcv
fysX7eLtFMtzl5wOS8RFv1gXt+cxed/Lm1pZN5dO/TKGz08G/26fdR7tmBYoteL15u/3oFfM
bJfZ9/jPzBN9da7M47f5l5KS40OaTud381mi1f3ac5qqjW1loB8XTXptqSPk2+rjqqjuxdl/
CaJ//f8z88v9zS/3Kha3m3s4GPlLBm5DTP13WP7c4363ppTQTp7tchM/u+74Tx8URI6NBOJf
XBuhm8NEi8Nfy81RdCWt6jwYfjHotnW9mQPAw+ffaD/tbnyNlCTdobpTeumovu4uF92i/kbl
Bs7UZv5tgA219KE2H4riof8GCGqpf4MXag4eFMKzJgdXCtr9Y/oEFRLNcNesm0a2H+57vgp9
OeLs2+/rs3gcTcqKvcNPcQ9wd8uTH8I+3Zh87p+I0EPCIaTBZJQjrVSD6awZPrrElDf3StGq
EFfP8V8q8bEc/8mBYZHAZPSU97C0D9TNevuMmw/zmW66Puhw3h3PoaXVD88rzld3uVlfg3XA
cRzYk5uiCtHOcLmpWyHHemuLcW6zegkz0+NW1vyTF+4VjvxbGPTsUUHgT3xw4MOdMNW9Tm4Y
5XCAO9JOdgktV00ir2tL8Z+uDHqfEP+1SEQGX1NXmVde4horozhA133NHgB34XNcal82O93S
kcp1+I97UPb1tVUB9BzsDKtDguJ61rxGCeRaO8LXmCesDCZsU80/ItN8gKZftot9Rl5jkutG
Va6Jvly3G4mxUa55jrjdUGQ0V535qDWvevOqObCU5VpPXnOaYVzmw7X//s1lSlyjO3JN4sW1
1aSSXVMlz3K5i2pHXpu9YZX3et5rrINrV27nZc8lm3h7v8dKsb6Z11Hy2vojnunktSsmJ1l1
NrNv62U0UXEhO9IFf6uCihAXcF7XXFASyVFQ3AYFDi5YEytxpAVUSkyloZFcwB2LC9FKkPgs
z/V8yaE8X5ob/fex6OKP7o9k4+5oSU6efWep9lEYn/fYPC//6n+XpXd/5Hut/uj+yDsp/sij
JXmfyd5dGU8Ov4/zDHf403g/mB1PVb3oF+vmJ+Xi6ozoBsbL88vFy/fnxcv768XNg1590qVO
hNJ6cXOkt771yQvqxZsfzIudmOXizRfGxRVjsSxb2ps/9It5sROffCIu3rziiruf/KJcvHnG
zQFvn2BlLvHw2TPm44uJfnHrdPwRDePNR94ehTcvGRfrYic+eUpcvPnKrfvFujiJVi7iol+8
+cy6ePOa2x/xyW9eDXhcvPnOToxyUS/e/KdfjIt1sROfvKhevPlRvxgX62InVrm4+vPt9noP
mzuPN958al+8edWrV+Pt3nxD0TWyDjh86k4rG3meZGef5M5zUAptO8937jwgvLI7S5Z4//hU
O5l31NRz457/dCk7z4E64dhOX051V/MePo28jrXfFX7KOXueOS55ssoHlHv2NHdO0vpvqdy/
Fpa5UErFwxnPyT8cNPtf1PWv/3/9//OmZX/7DybHm1E=
"""
rows = json.loads(zlib.decompress(base64.b64decode(encoded)))
verify_caps(rows,30)

# One additional induction step after the unrestricted degree-25 theorem.
# Also record the uncovered numerical pairs at the next ordinary diameter.
def check_joint_step(cap, diagnostic=False):
    table = [(row[0], F(row[2], row[1])) for row in rows]
    dmin = cap - 2
    unresolved = []
    comparisons = 0
    for h in range(2*dmin + 26 - 5, cap*(cap-1)//2):
        base = order_lb(h)
        if base >= pack(h, 0, cap):
            continue
        for degree in range(26, h - 2*dmin + 6):
            c = next((value for top,value in table if degree <= top), F(1))
            lower = max(base, h*(1-c)/2)
            upper = min(F(7,2) + F(303,1000)*(h-degree+1),
                        F(h-degree+5,2) - F(2,degree),
                        1 + F(h-dmin,2))
            if lower >= upper:
                continue
            comparisons += 1
            if lower < pack(h, degree, cap):
                unresolved.append((h, degree))
    if diagnostic:
        assert unresolved
        print("frontier:", len(unresolved),
              min(h for h,k in unresolved), max(h for h,k in unresolved),
              min(k for h,k in unresolved), max(k for h,k in unresolved))
    else:
        assert not unresolved, unresolved[:1]
        print("joint cutoff:", cap, comparisons, len(unresolved))
check_joint_step(31)
# Additional row for the next three caps. This JSON uses the same row format.
extra_encoded = """
eNrtl9GR5DAIRFPZAPaqJECAYrna/NM4dSOxKdzHTJk3NGNhWwbZ81dlfH/Ngc/3V9bX3+kq
64TrS8zXPup+7fqSvaV3mXO5ni/dZkinASUTe27NiVgEjrAsTogpIg2RMPiiZ2Awp5+ciCCH
n8yI+OHyMETWSbdkH3/FPBlsOyIqJ4+5g3v7oQ34Mc6vNmMhovn9pdtPfhNcqoYHInmy6Ur8
GnnGqgkiPs4olRBEOHbqRIT7c7J02vlVMpBn4HwkMBeSeSLiCwxkkEXf3TinjHBPjYqcsxXZ
gxH6czKCezA1GDnnJsPBaZjWvbHnwNzOrdhzDNyXRM4TSfgYNXmGMzBXJ4JbGCMZwZ6+6OOI
03EV0x1c9etEziUVQTZjtiWIG87tRDYrgD6Pq8ygm3HM54kgv2RFUDhnUlhM2FOEVbOQf7K6
hKMmR52LAVEbJ0IO/qpa5cUIso3Fmg6cw5CKzKpyRk7+Hc7ACe+l0ChLHU9iZjZuNe/iRiEl
rjtRImnwUDgpAEo6MYIycHWUERM1fZLGQnWjCAJzVRKzWXI86UhP6diZkhXkbAYDcKNc8smJ
nVH6jotxtMvCOVMuVAHLemHGFppgod9K4uAlMUslpz6JOlpIb3teaTmuZLWWxAmVNHhobd7/
kjjTkrMl+xIlaLgjlLrHlVwuFMfgKqKJ5kMVlkR1lMS0l9T9JLuTErVUcq4nRz7JA0HK1isF
fU8pqJySuIPCZsQKJNXG+SSqqqS11Jao1JKouJJzP4l5Ftzz07xPYoop5/YrJztjs4HtSXZa
orC5zlI6+5fNqk+ypymXPGnsW8RMnqz+hWQvsXkln6yuhJz55GTn4bi4eSWr9xgbV/K82WRS
AqUIP7V8nDp8FCZ8dgxW/CgfJw8fRQUft3uxYMtHicHHY4flNsvHFRlroXx0Np4wmHrlCou7
iTkuv9ZVXHH58/rXLa+cP5gnbsYtuHHAS/CSwHYZFkraupZlVaF1cjS/tstMrq1rUYYljqbX
1rUsw9JL02vrWpbhgUHTa+talOW4JtfsWlzbZehTml5b17Ls1kItvA1reCMfOI0FaVjDG78l
ORrSsIY3fgt7NKRhjdXo9tD9wOduQRvW8EY+8HlUmA1trIY3ftt3NGZDG9bwRjT2A3rpQhrW
WI1o5APK4GI2tGGN1YhGPtQ6RsyGNKyxGt7Ixr7g29uFNLRhDW9EIxv7ge9aBWlowxreiEY2
di/yozEb2rDGangjGtnYDzob0tCGNVbDG9HIxn7A+/jFbEhDG9ZYDW9EIxv7AW1w8HP+Zgwu
qadQMHX0NPx6A5fAZ52jQOjpuiNijDvi/DO4I5xvmOyxnXfEecm+nmre/c4s3bGnX7W84VGZ
N/84wLFRv2ViFrgQ4vk0+FZX+8RcdcTzb6YO40MrzzL0H9+V8BiEo/dg5/9FJVTWyeC7RCU8
hfguQSvhvudlo9J9ts/22T7bZ/ts//v28/MP7OaN9w==
"""
row320 = json.loads(zlib.decompress(base64.b64decode(extra_encoded)))
table = [(row[0], F(row[2],row[1])) for row in rows]
table.append((row320[0], verify_potential(row320)))

@lru_cache(None)
def quotient_pack(n,cap):
    bound = F(3)
    for j in range(1,min(cap-4,n//2)+1):
        width = 2*((j-1)//2)+1
        packing = max(1,(n-(2 if j%2==0 else 0))//width)
        bound += 1-F(1,packing)
    return bound

def check_quotient_step(cap,diagnostic=False):
    dmin = cap-2
    unresolved = []
    checked = 0
    for h in range(2*dmin+26-5,cap*(cap-1)//2):
        base = order_lb(h)
        if base >= pack(h,0,cap):
            continue
        for degree in range(26,h-2*dmin+6):
            c = next((val for top,val in table if degree<=top),F(1))
            lower = max(base,h*(1-c)/2)
            n = h-degree+1
            upper = min(F(7,2)+F(303,1000)*n,
                        3+F(n,3),
                        F(h-degree+5,2)-F(2,degree),
                        1+F(h-dmin,2),quotient_pack(n,cap))
            if lower>=upper: continue
            checked += 1
            if lower<pack(h,degree,cap): unresolved.append((h,degree))
    if diagnostic:
        assert unresolved
        print("frontier:",len(unresolved),
              min(h for h,k in unresolved),max(h for h,k in unresolved),
              min(k for h,k in unresolved),max(k for h,k in unresolved))
    else:
        assert not unresolved, unresolved[:1]
        print("quotient cutoff:",cap,checked,len(unresolved),flush=True)

for cap in (32,33,34):
    check_quotient_step(cap)
check_quotient_step(35,diagnostic=True)
\end{Verbatim}

\subsection*{Unrestricted distances through degree 25}
The second program verifies the inequalities for 214 and 665 vertex types in the two degree ranges, including all 13,741 and 138,793 admissible unordered type pairs, respectively. It checks the rational integral estimate and the finitely many small orders directly.

\begin{Verbatim}[fontsize=\scriptsize]
# Python 3; extract and run with assertions enabled.
from base64 import b85decode
from fractions import Fraction as F
from math import ceil
import json
import zlib

# Degree groups: lower K, upper K, integer cost, compressed integer potentials.
data=[
(2,14,24472,
    'c-pO0363Kn4EzvC_=dMc^Z#$G%g{=+J2`eV-AN^YDT>18b|Yq%Hc3~bAv7Dgdz=m@;zrYr8rvA'
    'us^bndLL4!7o9F?h1_OiCbk5aO)gHjtv^`Kkf5B1Hbc8h-<ljP#nyh#ZqDwZ=WPc7j&Qtmu80H'
    '(M=+kdH24g>fj0}qW`R6@megq6h3$E!sMtlpo;E<2LbSb<TEuI5CeUoZFG`s5yc=678!JW%8Ph'
    'Z1J%FsCQNUGs)fWF$TXJnRy!P2aOKucCWD3$<HCisrZPxZ}wDh@@`$nt5V4zC?Oz3KrI%K2Vr+'
    'arBpry5{9Rb$j5)4vr*=*qaN226Ib(Yg?tOay6G3v&zvkmS>{(jBF!t(uHxv0s`AS3x)AbC(jv'
    '3uG+lp{2E2F>F;jD{Tq;gu3PxxIb0aeHM=`hLgzGE|#^F%%1yYQwT?OozdfiPQKhIg3!Iqq^j9'
    'bo6J9Z*1pzyU1PoOrsM3ksA~NzbEZn|sV=`uSM_l992IMTyLyhIH;cNA<S@4ER~uM#)edcN&D%'
    'RXN4?K_Ol<1zfCM<sVU~xiIL1BgOnxb%$V_3d`UH-ZG~$}0!bQ2se&`-Q9932OCr4p|t3z!e?F'
    'T`S1wjzO_k~2<WWOrw6J_iMWwg(0K^av+8ERk#l)1onTt~mDulVRcQu*Xa^@zezsTiINt%fx@H'
    '-Q4w`1<+-XTpiL'
),
(15,25,38932,
    'c-qZZ3vS~&47`vHh@?K$S&F^?H6CXuyM4P&5j(|h<Ay10i+_$rZ*MYdZJX%qwn_6*H|f2PO}2c'
    'ZRm<44l@fQhIn=gX!TYlAOU-+flPgL-8DsbdA*Hn*=2|US79ah;H1~Npvi0h^b!+8!F-75`5Bp'
    'q5IObb}trXJcV-bdCIR3Set$Nq3N+joNAs-EYDXOWYFthrMh}V(hXGJJRKD9@>o)xjCS-&RM8Z'
    '8mZB@sVpi|ZmRYDMGCen)I7Htvd8J!kZ4F~2R)<Nlh+M8rH?{$>W()V(MI8>WRuzxBoFwcFn?U'
    'XA$k;o1XQzWSnb#!Fw-NI%7ImQDN9LY}%*`kyy<jqZC^$g`G89iLc=ZQ+uT=dVQ7KTXDXToMly'
    'J-oc>r*@^^5OVDr-9LS>%N#d^T)p~~2loB8xN&uT^k@(IaaY`Og?!|JE6rXK@@sqBp$DX-Y%O3'
    'HbtBh`eHkF)Ga6z&Ai3iE09HO3#f}Gt56M~hJ2|Xax9#LrX(s+^*I+d>T09>>^&qYCZ7RqPJ|V'
    'TwZ70FP!ZB;W652|8LvPK^Bsf^Uq4=zoqXxSR_CkhH!Hw!+UDFc_#<5S}958xEg7*r1Am_;l9K'
    'oa(RF1JR3o+1uX;3Yw<dQ=%26w4viqZ&O=Yn*QS}iiS6w-BX+;KA#hrbSPw38_9EXp45q6BF!b'
    'EBgi=E&5jo|;oB^wDl)cN?)4Mn<+Skx5c7nj){-y9qqi5_e4+a@~|LQjl^lNMsvWQ^WfY>mUD-'
    'cI`PfK|OZ|NnBYY8Qn=6R|X@wC{yuzW|ON@%d@#=RlE2qZ7bZ->gs_q`Cl}s;8hrGYIFz6tWIu'
    'V+@7dJv87@i`wmH1%F=Fn74>|NB8!wBHa}ZV)Y{>#l&DDhoVH-JvBA;#@wQjM22T?a0-H_^S;q'
    'iHeFmDs6>j3N5^kbyDyZyI`AW*AVFx)JcRLdw%XxI#^K2E7@Mkp0&j0g4RcB(F92SlDbtvre{c'
    't*U>PlK>CuZ1Bm)ZWA-xYN_UeWluj8y$_e@Hp<PU=UOYhu?h&TVvY=<)7s0qAc7?+5jU<A(B@H'
    '@3sDa~d4or5%l<m1(1X9S5!~U@52816vBftH0|jKF&CD#=-=uXWTfw9np|KGJ*tyJoGM%UWlCL'
    '6DpU*Sq*U{Xg1`pp?xiu!FpBpP9r56hZV~zt+rXE$nMEoOZh+!S0`mzR&`C@g91vg-dZK{6NfL'
    'waztb&6xPD}r0FW}(~uCkW)StfmNa5IF?2y&!cI>6LWC6^%g&^wA=2RMWlD>)u+1@VwE%*xW89'
    '5s5PgBB<1VHHtrgZ!i+?#la`3xhI><mJpxB%Uus77m2E#N~kh3rtZ1e2@dD#6o!E%saNp&><`5'
    'OOj44Huocb6%NiIY=9U}F+=`fGN-A?gIwmtwDkpvu^lT@COt0f7L!gg=Bn(3bEA{U81Slw|1P!'
    '@o}g4o2Tmj!C5L37b4?7u*ly#4e_#vH!`jG-jS&!RLY<DLwfq3Bh^oThb!jqzKxEea8v|cW$<{'
    '67&eZ8NJoc4r$z38qcTg^yg<Kr>lJ9Ad|URFe#*{qfeZhLQg)zt=8B}+UOw<4+VT6QUqq?JV15'
    '*p^2ic?@Jf%TatmoZt+%564!_}M)IJ%PXeZehZd43@f<o|zWw_39~}O}?E'
),
]
x0=F(197,500)
I=sum(max(F(0),min(x0,F(1,m))-F(1,m+1))/m for m in range(1,33))
assert x0-I < F(303,1000)

for k0,K,C,compressed in data:
    scale=100000
    c=F(C,scale)
    X,Y=json.loads(zlib.decompress(b85decode(compressed)))
    types=[(k,a,t) for k in range(2,K+1) for a in (0,1)
           for t in range(1,k+a+1) if k+a>=3 and (a or k>=4)
           and a<=t<=a+k and (k+a!=3 or t==1)]
    assert len(X)==len(Y)==len(types)
    edges=0
    for i,(k,a,t) in enumerate(types):
        d=k+a
        assert k*X[i]-(t-a)*Y[i]<=C
        for j in range(i,len(types)):
            p,b,u=types[j]
            e=p+b
            if ((d>e and t==a) or (e>d and u==b) or
                (d<=e and t-a==k) or (e<=d and u-b==p)):
                continue
            if d<e and t*e*(e-1)<=(e-u)*d*(d-1): continue
            if e<d and u*d*(d-1)<=(d-t)*e*(e-1): continue
            if e==d and t+u<d: continue
            fee=abs(F(1,d)-F(1,e))-F(a,k*d)-F(b,p*e)
            assert scale*fee<=X[i]+X[j]-Y[i]*(d>e)-Y[j]*(e>d)
            edges+=1
    alpha=F(303,1000)
    assert (1-c)/2>=alpha
    if k0>=14:
        assert F(7,2)+alpha*(1-k0)<=0
    else:
        stop=ceil(F(3,2)/((1-c)/2-alpha))
        for h in range(3,stop):
            lower=max(h*(1-c)/2,
                      min(F(h,s)+F(s-1,2) for s in range(1,h+1)))
            mass=F(0)
            for D in range(1,h//2+1):
                w=2*((D-1)//2)+1
                P=max(1,(h-2*(D%2==0))//w)
                mass+=1-F(1,P)
                upper=min(F(h-D,2),F(D*(h-1),h),mass)
                assert upper+1<=lower, (h,D)
        print("small orders checked through",stop-1)
    print(k0,K,len(types),edges,c)
\end{Verbatim}

% END SOURCE BODY Fajtlowicz_Iterated_Cutoffs
\unitstart{Paper 14}{p14}{Large degrees and mean distance}{The Randi\'c index, large degrees, and mean distance}{\subtitlevolume{A weighted cutoff of 40 and an ordinary diameter bound of 53}}
\makeatletter\@removefromreset{theorem}{section}\@addtoreset{theorem}{section}\makeatother\renewcommand{\thetheorem}{\thesection.\arabic{theorem}}
% BEGIN SOURCE BODY Fajtlowicz_Diameter_and_High_Degree_Sites
\begin{abstract}
We prove the Randi\'c average-distance inequality for connected graphs of diameter at most 53. A weighted distance inequality with cutoff 40 covers all cyclic blocks of intrinsic diameter at most 38. At unrestricted diameters we prove the weighted inequality when the vertices of intrinsic degree at least 26, joined at distance at most two, form an eleven-colorable graph in each block. In a least weighted obstruction the high-degree sites must also spread over distances growing at least as the square root of their number. The starting point is the preceding cutoff-34 manuscript. The three finite calculations for the extensions are reproduced in the appendix.
\end{abstract}
\section{Statements and inputs}
All graphs are finite, undirected, simple and connected, with at least two vertices. The ordinary average $\mu$ is taken over unordered pairs of distinct vertices. At full vertex degrees $d_u,d_v$, put
\[
 R(G)=\sum_{uv\in E(G)}\frac1{\sqrt{d_ud_v}},\qquad
 M(G)=\sum_{uv\in E(G)}\frac1{\max(d_u,d_v)} .
\]
For two independent vertices $X,Y$ with the same probability vector $p$, set
\[
 L_C(G,p)=\mathbb E_p\min\{d_G(X,Y),C\}+2\sum_v p_v^2,
 \qquad F_C(G)=\max_p L_C(G,p).
\]
Allow $C=\infty$ with its usual meaning. Write $Q(H)=\max_p\mathbb E_p d_H(X,Y)$. A cyclic block means a maximal two-connected subgraph with at least three vertices. Degrees and distances denoted by $k_v$, $K$ and $D$ inside a block are intrinsic to that block. A decoration attaches private leaves to core vertices.
\begin{theorem}\label{p14-main}
Every graph satisfies $F_{40}\le M+1$. Fajtlowicz's $R\ge\mu$ holds whenever the whole graph has diameter at most 53. It also holds whenever each cyclic block individually satisfies one of these two conditions:
\begin{enumerate}
\item the block has diameter at most 38,
\item the sites of degree at least 26 can be partitioned into at most eleven sets, within each of which distinct sites have distance at least three.
\end{enumerate}
There is no upper bound on the number of blocks or on their articulation degrees.
\end{theorem}
\paragraph{Earlier results used.}
The general conjecture is due to Fajtlowicz~\cite{p14-fajt}. The tree case was proved in~\cite{p14-tree}. Our preceding manuscript~\cite{p14-prior} proves $F_{34}\le M+1$ on all graphs, and $F_\infty\le M+1$ for blocks of intrinsic degree at most 25 with arbitrary private leaves. It proves that the unrestricted weighted inequality implies $R\ge\mu$ by evaluating the uniform law and comparing degree weights. Its padded cut-vertex reduction glues all independently certified decorated blocks. These results supply the starting point here.
\paragraph{Minimality and estimates.}
For an edge-minimal weighted obstruction, the same reduction gives a two-connected core $H$ with at most one private leaf at each site. Each core site has full degree at least three, and each unloaded site has a core neighbor of strictly smaller full degree. Degree-three sites are loaded. Deleting any core edge must strictly increase $M$. These statements also apply to classes preserved by blockwise edge deletion and padding. Write $h=|H|$. The previous manuscript gives
\[
 M\ge B(h)=\min_{1\le s\le h}\left(h/s+(s-1)/2\right),
 \qquad h\ge 2D+K-5,\qquad F_C-1\le Q(H)+1.
\]
A least obstruction to passing from $C-1$ to $C$ must have $D\ge C-2$, since the two private leaves can add at most two. It also has $h<C(C-1)/2$, since otherwise $B(h)\ge C-1$. Another input is $Q(J)\le0.303|J|+1/2$ for every two-connected graph $J$. All decimals in this paper denote exact rationals.
\paragraph{Local degree comparisons.}
Let $a_i\in\{0,1\}$ record the load at site $i$ of intrinsic degree $k_i$, and $d_i=k_i+a_i$. Then $h-2M$ equals the sum over core edges of
\[
 w_{ij}=|1/d_i-1/d_j|-a_i/(k_i d_i)-a_j/(k_j d_j).
\]
A verifiable row for the coarse types $(k_i,a_i)$ consists of rational potentials $x_i,y_i$ such that $y_i\ge0$, $y_i=0$ on loaded sites, $k_i x_i-y_i\le c$, and
\[
 w_{ij}\le x_i+x_j-y_i\mathbf1_{d_j<d_i}-y_j\mathbf1_{d_i<d_j}.
\]
Every site with nonzero $y_i$ has a smaller neighbor, so summation yields $M\ge h(1-c)/2$. For finer types $(k_i,a_i,l_i)$ the value $l_i$ counts all neighbors of strictly smaller full degree. The number of such core neighbors is exactly $l_i-a_i$. The vertex test becomes $k_i x_i-(l_i-a_i)y_i\le c$, allowing signed $y_i$. The admissible pair tests in the appendix follow from the exact deletion changes of~\cite{p14-prior}. At full degrees $d<e$ and lower-neighbor counts $l,u$, the change in $M$ on deleting the edge is $l/[d(d-1)]-(e-u)/[e(e-1)]>0$. At $d=e$ it is $(l+u-d+1)/[d(d-1)]>0$.
\section{From cutoff 34 to cutoff 40}

\paragraph{A diameter retained under contraction.}
Assume the universal cutoff theorem at $C-1$ for some $C\ge35$.
A least-edge counterexample to the cutoff at $C$ has a single decorated two-connected core of order $h$, intrinsic maximum degree $K\ge26$ and diameter $D\ge C-2$. Write $N=h-K+1$. Contract the closed neighborhood of a vertex of degree $K$ to $o$, obtaining $J$, and augment $J$ by one true twin of $o$, obtaining $J^+$ as in the cutoff-34 paper. If $x,y$ are diametral in the core then
\[
 d_{J^+}(\pi x,\pi y)=d_J(\pi x,\pi y)\ge D-2.
\]
Indeed a shortest quotient route lifts with at most two additional edges. The quotient has at least two vertices since $D\ge33$. Hence the two-connected graph $J^+$ on $N$ vertices contains a cycle of length at least $2(D-2)$, obtained from two internally disjoint paths between those projected vertices. Traversing this cycle once and the edges of an attached spanning forest twice gives a closed spanning walk of length at most $2N-2(D-2)$. Circular averaging bounds the expectation between independent draws on this walk by one quarter of its length. Finally lifting the quotient costs at most two in core distance and projecting decoration to core costs an additional one in the $F_C-1$ bound. We conclude
\[
 F_C(G)-1\le 3+\frac{h-K+3-D}{2}.
\]
This diameter-dependent estimate holds for every decorated two-connected core with $D\ge3$ and at every cutoff, including $\infty$. Combined with the rank lower bound $B(h)$ it gives an unrestricted exclusion whenever $B(h)\ge3+(h-K+3-D)/2$ in a least obstruction.

\paragraph{Exact verification.}
For $C=35,\ldots,40$ the preceding argument and the prior reduction restrict the check to
\[
 2(C-2)+21\le h<C(C-1)/2,\qquad 26\le K\le h-2(C-2)+5.
\]
For each order first compare $B(h)$ with $1+\sum_{j\le C-2}(1-1/P_j)$, using the two-connected packing bound in the core. At remaining pairs use the maximum of $B(h)$ and the applicable rational degree-row lower bound, and the minimum of the diameter-dependent bound, the capped quotient and core packing bounds and the prior untruncated quotient metric bound $7/2+303N/1000$. For clarity, the ordinary packing cap on $P_j$ is $\max(1,\lfloor(n-2\mathbf1_{2\mid j})/(2r+1)\rfloor)$, $r=\lfloor(j-1)/2\rfloor$, where $n$ is the order of the two-connected graph. Beyond $j>\lfloor n/2\rfloor$ the tail probability is zero. For core packing use also $P_j\le1+\lfloor(h-K-1)/(2r-1)\rfloor$ when $r\ge1$. For quotient packing truncate at $C-4$ and add $3$. The integer edge and vertex comparisons above verify each local row, and the included data are exact integer potentials. All six cutoff steps pass.

There are 38 rows and 1,502,496 unordered local type comparisons. The six cutoff steps examine 947,003 pairs $(h,K)$ after whole-order pruning. At cutoff 40, $F_{40}$ agrees with the unrestricted functional on every decoration of a core of diameter at most 38. Padded splitting then yields the first block condition of Theorem~\ref{p14-main}.

\section{Blocks with eleven separated classes of high-degree vertices}
Let $H$ be a cyclic block. Put $T=\{v\in V(H):k_v\ge26\}$, $t=|T|$ and $S=\sum_{v\in T}k_v$. Form a graph on $T$ by joining two sites if their distance in $H$ is at most two.
\begin{theorem}
If this graph admits a proper coloring with eleven colors in each cyclic block, the original connected graph satisfies the unrestricted weighted inequality $F_\infty\le M+1$. In particular this applies when each block contains at most eleven vertices of intrinsic degree greater than 25.
\end{theorem}
\begin{proof}
The condition is closed under the cut and deletion operations: new blocks are subgraphs of old blocks, intrinsic degrees can only decrease and distances inside a surviving block can only increase. In a least exception the original reductions therefore apply to its terminal decorated block $H$. Let $h=|V(H)|$. The previous degree theorem handles $t=0$.

Use the supplied scale-$m=10^8$ degree-25 row and write $x_i=X_i/m$, $y_i=Y_i/m$ on each ordinary type $i=(k,a,l)$ with $k\le25$, full degree $d=k+a$ and $l-a$ lower-degree core neighbors. This row satisfies the ordinary edge and vertex comparisons with $c=39300050/10^8$, together with the following comparison for every ordinary type capable of having a nonlower core neighbor:
\[
 l-a<k\quad\Longrightarrow\quad
 1/d-a/(kd)-x_i\le g:=16419997/300000000 .
\]
This is checked by an exact rational comparison in the appendix. Give a high site of full degree $e$ the two potentials $x=g-1/e$, $y=0$. An edge from an ordinary site to a high site has $d\le e$; the preceding bound gives the required edge inequality for its debt $w=|1/d-1/e|-a/(kd)-a'/(k'e)$. On a high--high edge $d\le e$, one has $1/d-1/e\le2g-1/d-1/e$ because $1/d\le1/26<g$. All ordinary--ordinary inequalities are verified by the supplied row with the original local degree tests. Summing the potentials and counting lower-degree neighbors as in Section 1 yields
\[
 h-2M\le c(h-t)+gS-26t/27,
 \qquad M\ge Ah+B t-gS/2 ,
\]

where $A=(1-c)/2=0.30349975$ and $B=1/2-A+13/27$. Fix one color class $I$. Its closed neighborhoods are disjoint. Contract these into single vertices and suppress loops and multiple edges. Create one true twin of each contracted vertex, cloning its present closed neighborhood successively. The augmented graph has $N=h-\sum_{i\in I}k_i+|I|$ vertices. Folding the twins to their originals fixes all pre-existing distances. Deleting any uncontracted vertex preserves connectivity of the quotient because its preimage is a single vertex of $H$; deleting either copy of a contracted vertex leaves its other copy. Thus the augmentation is two-connected whenever it has at least three vertices. The expectation bound of the previous paper gives $Q\le303N/1000+1/2$ there. For orders one and two this is immediate. Lifting a simple quotient path increases its length by at most $2|I|$. The projection of the private leaves adds at most one to the upper bound for $F_\infty-1$ relative to $Q(H)$. Consequently
\[
 F_\infty-1\le Q(H)+1\le1.5+0.303h-
       \sum_{i\in I}(0.303k_i-2.303).
\]
For an empty class use the same metric bound directly on $H$. The best of the eleven classes saves at least $(0.303S-2.303t)/11$. Thus
\[
 M-(Q(H)+1)\ge
 (A-0.303)h+(B-2.303/11)t+(0.303/11-g/2)S-1.5 .
\]
The coefficient of $S$ is positive. Substituting $S\ge26t$ gives a coefficient $B-2.303/11+26(0.303/11-g/2)>3/8$ on $t$. This suffices for $t\ge4$. For $1\le t\le3$ use the single maximum-degree site $K$ as $I$. Then $S\le3K$, so the difference is at least
\[
 (A-0.303)h + B t +(0.303-3g/2)K-3.803>0,
\]
since $K\ge26$. This contradicts minimality.
\end{proof}
Exactly the same argument applies when the conflict graph has fractional chromatic number at most eleven: an optimal fractional coloring supplies a distribution on independent sets covering each site with probability at least $1/11$.

\paragraph{High sites confined to a region of small diameter.}
Consider a least-edge counterexample to the unrestricted weighted assertion; keep the notation $T,t$ from the coloring argument. Let $U=\bigcup_{v\in T}N_H[v]$, $r=|U|-t$, $z=h-|U|$, and $V_T=\sum_{ij\in E(H[T])}1/\min(d_i,d_j)$, where $d_i$ are full degrees. For an ordinary type put
\[
 \delta_i=(1/d_i-a_i/(k_i d_i)-x_i)_+ .
\]
The same finite row verifies $(k_i-l_i+a_i)\delta_i\le579/1000$. Use the ordinary potentials at normal--normal edges. At an ordinary--high incidence assign $-1/e$ to the high end of full degree $e$, and keep its load correction $-a/(ke)$; the remaining excess at the ordinary end is at most $\delta_i$. Each affected ordinary site meets at most $k_i-l_i+a_i$ nonlower neighbors. On a high--high edge of degrees $d,e$, retain the contributions $-1/d,-1/e$ and both load corrections, leaving precisely $2/\min(d,e)$ extra. At each high vertex the sum of the negative charges and its load corrections is $-k_v/d_v-a_v/d_v=-1$. Summing gives
\[
 M\ge A z+(A-0.579/2)r+t-V_T,\qquad A=0.30349975 .
\]
Order $T$ by full degree and orient each induced edge from smaller to larger rank. Thus $V_T=\sum_{v\in T}d^+(v)/d_v$. The $j$th site from the end has outdegree at most $j-1$, with $d_v\ge26$. Putting $m=\min(t,26)$ gives
\[
 t-V_T\ge m-m(m-1)/52 .
\]

If $T$ has diameter $\delta$ measured in $H$, any pair in $U$ has distance in $H$ at most $\delta+2$. Identify all of $U$ to one vertex, even if its induced subgraph is disconnected, and duplicate that vertex as a true twin. The augmented quotient on $z+2$ sites is two-connected for orders at least three: deleting any other site preserves connectivity by taking the image of $H$ with its unique preimage removed. Deleting either twin preserves a copy of the unaugmented quotient. A shortest quotient route lifts with at most $\delta+2$ extra edges, since the entrance and exit may be joined by an ambient shortest path in $H$. For quotient orders one or two the same expectation estimate is immediate. Consequently $Q(H)+1\le0.303z+\delta+4.106$. Since $A>0.303$ and $A-0.579/2>0$, there is no such least obstruction if $m-m(m-1)/52\ge\delta+4.106$. Thus a least exception with $t\ge12$ has $\delta\ge6$, and one with $t\ge22$ has $\delta\ge10$. These are conditions on a global least weighted obstruction. An arbitrary input block with small high-site diameter may fail to preserve that hypothesis under edge deletion, so this conclusion alone does not assert an entire hereditary class.

The retained term $t-V_T$ has a second lower bound that grows indefinitely with $t$. Indeed it equals the total $M$-weight of the edges incident to $T$, including private leaf edges: summing $1/d_v$ over incidences at each high vertex gives one; on high--high edges the overcount is precisely $1/\min(d,e)$, and an ordinary neighbor never has larger full degree. Order the full degrees of these $t$ sites decreasingly, and let $s$ be the largest index with $d_s\ge s$. For each of the top $s$ sites in order, retain edges to all neighbors except earlier sites in this order. Each retained edge has weight $1/d_i$, so the $s$ stars contribute at least $s-\binom{s}{2}/s=(s+1)/2$. Each remaining high site supplies its private leaf edge if loaded, or its edge to a strictly lower full degree otherwise. These $t-s$ edges are distinct from each other and from the retained stars. Their weights are at least $1/s$, since $d_{s+1}\le s$ whenever $s<t$. Hence
\[
 t-V_T\ge B(t):=\min_{1\le s\le t}\left(\frac ts+\frac{s-1}2\right)
 \ge\sqrt{2t}-\tfrac12 .
\]
A global least weighted obstruction therefore needs
\[
 \delta+4.106>B(t)\ge\sqrt{2t}-\tfrac12,
 \qquad t<(\delta+4.606)^2/2 .
\]

\section{Extra whole-graph diameters from the uniform distribution}
The original mean-distance problem admits a second gain beyond the weighted cutoff. Write $d=\diam G$, $n=|G|$, $k=\Delta(G)$ and $q=\binom n2$. Assume the weighted theorem is available at $2\le C<d$. Put $b=\binom{k+1}{2}$ and $m=n-k$. Define
\begin{align*}
U={}&\min\{(n+1)/3,\, U_2,\, U_3\},\\
U_2={}&
2+\frac{(m^3-m)/6+km(m-1)/2}{q},\qquad U_3=d-\frac{(d^3-d)/3+(d-2)b+k-3(d-1)}{q},\\
g={}&\sum_{j=1}^{d}(d+1-j)\bigl(d\min(j,C)-Cj\bigr)+(d-C)(b-6),\\
A(t)={}&\frac{n-1}{n}\left(\frac Cdt+\frac{g}{dq}\right)-1+\frac2n.
\end{align*}
The following tests suffice at any diameter in this range:
\[
\frac{n\sqrt k}{k+1}\ge U\qquad\hbox{or}\qquad
\frac{n+(\sqrt k-1)A(U)}{\sqrt k+1}\ge U.
\tag{*}
\]
Indeed $\mu\le U$. The first upper bound is the path bound. Contract a maximum-degree closed neighborhood to a site $o$. The resulting connected graph on $m$ vertices has Wiener sum at most $(m^3-m)/6$ and rooted transmission at $o$ at most $m(m-1)/2$: a spanning-tree leaf induction proves the first estimate and increasing distance layers prove the second. The site $o$ acquires $k$ extra copies in the distance sum. Every pair of vertices costs at most two extra steps before projection, proving the second bound. For the third bound, sum the deficits $d-d_G(x,y)$ over the pairs on a diametral path and the pairs in a maximum-degree closed neighborhood. The path gives $(d^3-d)/3$. The neighborhood gives at least $(d-2)b+k$, since every pair is at distance at most two and $k$ pairs use edges to the center. At most three vertices of the path lie in the closed neighborhood (it has diameter at most two and the path is isometric). The duplicated contribution from at most three pairs is at most $3(d-1)$.

Similarly $e(j)=d\min(j,C)-Cj$ is nonnegative. Summing it over the path gives the explicit sum in $g$. Each of the $b$ neighborhood pairs contributes at least $d-C$, and the at most three duplicates contribute at most $6(d-C)$. Thus the mean clipped distance is at least $(C/d)\mu+g/(dq)$. The weighted assertion evaluated at the uniform law gives $M\ge A(\mu)$. The edge degree inequalities from the preceding manuscript imply $R\ge n\sqrt{k}/(k+1)$ and $(\sqrt{k}+1)R\ge n+(\sqrt{k}-1)M$. Indeed, for endpoint degrees $a\le b'$, the ratio $1/\sqrt{ab'}$ to $1/a+1/b'$ is at least $\sqrt{k}/(k+1)$; the other bound follows by summing
$(1/\sqrt a-1/\sqrt{b'})^2=(\sqrt{b'/a}-1)(1/\sqrt{ab'}-1/b')$.
The difference $[n+(\sqrt k-1)A(t)]/(\sqrt k+1)-t$ decreases in $t$, proving (*).

For the parameter check we may assume $n\le d^2$, $k\ge26$, $k\le n-d+2$. The degree-25 theorem covers smaller maximum degrees. For $n\ge d^2+1$ the elementary bound $R\ge n\sqrt{k}/(k+1)\ge\sqrt{n-1}\ge d$ suffices. The third restriction follows by counting a diametral path and the closed neighborhood, whose overlap has at most three vertices. Evaluating (*) by integer squaring at $C=40$ certifies every whole-graph diameter through 53. Exactly 1,127,740 instances require the second test after pruning. Together with the weighted and block arguments this proves Theorem~\ref{p14-main}.


\section{First remaining cases}
An ordinary counterexample would have overall diameter at least 54. Some cyclic block must have diameter at least 39, and the distance-two graph of its sites of degree at least 26 must have fractional chromatic number exceeding eleven. In particular such a block has at least twelve high sites. The more specific distance restrictions on high sites above concern a least-edge counterexample to the stronger weighted inequality.

At core diameter $D=39$ the decorated diameter is at most 41. Every proper smaller graph than a globally least unrestricted weighted counterexample obeys the cutoff-41 inequality, and on this obstruction the cutoff would agree with unrestricted distance. Hence the parameter bounds used in the cutoff-41 diagnostic apply. Under the bounds implemented in the first program, 166 order--degree pairs remain, with $119\le h\le331$ and $26\le K\le221$. These are necessary parameters. Feasibility of any corresponding graph is undetermined. The only term missed by cutoff 40 in this diameter is $2\sum p_{u'}p_{v'}$ over unordered pairs of leaves whose attachment sites are at distance 39. Higher core diameters remain open as well.
\begin{thebibliography}{9}
\bibitem{p14-fajt} S. Fajtlowicz, On conjectures of Graffiti, \emph{Discrete Mathematics} 72 (1988), 113--118.
\bibitem{p14-tree} M. Cygan, M. Pilipczuk and R. \v{S}krekovski, Relation between Randi\'c index and average distance of trees, \emph{MATCH} 66 (2011), 605--612.
\bibitem{p14-prior} Keio and Shaitetsu, \emph{An iterated distance cutoff for the Randi\'c index}, preceding research manuscript. Sections 1--5 establish the reductions, two-connected metric bound, unrestricted degree-25 result and cutoff 34.
\end{thebibliography}
\clearpage\volappendix\section{Exact certificates}
Extract the three \texttt{Verbatim} listings from this single TeX source into separate files. Each runs with standard Python 3.9 or later, using assertions. The third takes the arguments \texttt{40 53}. No coefficient fitting or external library is needed. The first reports the unresolved cutoff-41 parameters as a diagnostic.
\subsection*{Cutoff induction}

\begin{Verbatim}[fontsize=\scriptsize]
"""Standard Python 3.9+. Run with assertions enabled."""
ENC="""
eNrtnQt2Y8eObCfEWiv/n7Hc1fOfxsMOAIcyXdX2674/35alSFFliaTOyUQCgQDyb3/r9VUL/73G
Hr28/lbXWvtV5ynr1UYd5VXrYlyM59T2qq2O+dp1lted9gTzzvva2/5tXPvd1ed6jdrHazb7tz7a
fo3e7elOq68+93m1du3b2+2FTrdv+zmvOuxH6uX5y+bba6+xeHl+f4/XKfXV6nrNcl4/qv2WPfmP
Zq+/Xz/s2WoRhv1fMM8Vtr0ZcO3/gVpmi6GeEUO3X/BhtB7DtKvB8F+vv5WXPe+t/dpXe//27uy1
5pj8e73tDr6O2fn3u0+3L7PYu7af2u1V/t4f/2Vvqc/nvt1d3vdtzdXsyhy7vrV1xs14R7X/2ba9
5b3a1e/e1xr2lk9rdqfGWK917FnHqv01t/1bvzYFxql2CzuvU+a2v2dzR+0u2IW3bw+veeuwby+X
a/Kt3St7bbuYy57AZggvbC9m9/OezQu8eHp7k/bsdjv7q61rd7TZjWqvH9tmZHXYqwnDbjKY9sxg
2/OCY38JuGMIdltrDFU3mKGNHQNz0IfRiobn3h57c3Zvd7t839axV7evpW/uYl12vfhaT+Ne7zn4
5952ef3DP3S3b97tWZfdgeduM7bJ/6jNJqjdbcarh73ZPT2FG1FGtRtw7AqfZQtsHlucu9kst1+t
r2VL7DWYxtP+4ldfthrHYOnWZd+eYjdrF/u2zWk32n7Eboq9ozv41uZbXY37fllJo/K2eM3De7A3
e8blbjYzDPb6dvuG/cyatoDtV4et8GsP7fl+2LPzYz+49aPzWc8UmGuA3wTTFjRY9v7AXkM4ZojA
tbkFaukthphTDDbpfGhMVw32p2nICWEroWpCYI9Yxn7jW21a3dUmsBZ9x9bZXzAWE8Km/9V0uf+E
efE5R0Z75ohZsvPTOdLbsj/7aCwa+7x2yXrvWN3K0tpYea4tVnjzh8mMu1XHxpmtsLtxOzfPDGwf
xW5MsavdrpmXPuwJGwukF5tgrdqPNAxBPbaybc5dNpcZlr4zcVZlel+mi73GOSw0s+/LfuoWm61a
7YONxN4Kk8JelDdjr2j/+4fdEZtEP9iDOp9mgIRuLwCY+8Amt7DMogGeU7CLBbBl4NrrXF1NVpEP
p8bAzudD42/yYQ0NOXkWG5tdCbvOWJV22CSZPFdf6/YtotoeyY5hq1M7iP29RValzFf5V31oJo33
TMKs/m4mVdY5e/TRWDR2bOaZ2mKPLYDdbMHeYtZitcZWbrPLno+lb/9mU8e+nXbTB/PMFtyymzoH
d8rWe7E9qy/si11M7JBNIrkQm+2isPV0JtG2K2rXi42bO7+Xb3s2xUrXFPe7c2QGbMIs+3Uu87D7
vexJeU79KC862BFf3NIfk6lvexEWkk/bV4W2u9AD9mcIbHxgYdoGG1zA/g7wzC0zvILZPGCb1Y2B
KR8Df4qGpj+grWdm2ZXVzBral5ptW5phthnL8JxalmaW/aD2q1r5vb1kxYZNXOZZXf+6+fXMsfXM
scEi/OkcW3gSR2PROErRdeT6FJssu9sOf1lGtrcUfFBb5tNsim0KTCpWve3ptrnZKjZHqbK5XbxR
1pnNnb4Pl7EdfBx8EmzUscnR8B0a5sLewNLaxZ+xu3P11ng3vIY5x1Xv64Xra7evy5Zcm5zLbBV7
J8/Jj/KiA2/zYjRtbtgfak7Yj86fZGaMac8nf051r1jAOIPRmoA9FI5j2TsA2ya5YM8Pjm2M4NrW
JZwmVHcJNNwaQx3PcK6GnHO2m7g147IXvPaxNOeW/GCzAUVb4TQnQ3PO//30OuVazab/e8vr3+JD
k+88k882s/LzyXdwR47GonFUNonDDmR23yYfi/A2XJBhN2/j2sxlf+7ai9lmltM2usFsMwMhi3Yq
3s/Eo7YNTA52w0G2rY1Nkx9hv2lyfDF/NlcqrvbB1GH5MAjssJXXMG/Evp34LuU2j2Dk+rHsNz6Q
PS3PqjdkL7tsvplbaIZkae/sZoo02at5XeY22TS8OEL6DLRyHLalAvsjBPx4MDGhXf6Cwy4D2PZX
CzkdzTwCOajA3sOVI1KfYY0Yvnhr8grn45mNqc3TIsYpj8yiH4VjduHkop919NVurCZea5qeh5cq
WAP9tEVzr/Lv9MG0ZJ+JaWkBR/+pBzdwqswGaRZqHHKarq5atzm32VTNA7HNbBV2oMufvfCn7N/M
cBCwb11JHH6cpItlmljHYRGYvYT9W2F3nOzA2/0oduDL7rcIw+1Gt9Exz2ZyW8VdYl+pvIa9SYwl
xqR5TM1tJeq0iW0mie9EILSX3tJmr7TI01biPLyWme5J9DUIsV7MO8ICNnSL+O3bwTc24ZnzOGUG
XC3Qh4MIT7AbDeY8gvkiDrxGjDWGGNjTApwDcEsgIolL9MraKIrFNLydwtY05BztV3PR1ndjDlr0
rGjBPJjqJvN2bdPzaEM2F+Xyc5crYfvzGNrkuzMI/0pv8L+dre94Y7Ht/HS2VsWJGqvGMVn6hVDR
XDabBNvmqF1g7IYZp1Mu1M02cwJXMHH8FIiMaSZsdu4oO5eFJUzPg1PDrmk+tW3t+IanaWvXW2Av
x3OENbKQpOHmQT5Nj+Ds2wPVhZnBc6idGTwVLXbxBXbPt3gE20lbEX3BC1u8ZG6LTabLJFuaWIo0
7BeaAmXxTK8mE+pBLS7yNiu7Mkoh7IRRwavF+Y1opZ8usALAxG0GdrHAstcWzhW2/VnCdRxYOjF1
WyDmEOw1gVnakcPsOYQf8EzgdoYmcNVebiH/1Fw0J2uK87Lr5bzI0p7fqwgRgiWM865yR0d4o+W2
f78p/Jup/A54tlm/n0/lrvmjsWocMDG3YdcgiMwBYGoTT9sWZJOEGNQsLYEmnlrDBDNPbfM1J07G
CTdvwKV0dq3Bc3X2/q4w5zKVFfVgyNtVEFSxc/byTXQcMUvFJWm8ht6GWWRsbiVgYSqzFqpMdWEu
t4npVJBmU9FemkDeHBibPITHU8yMRQq22HBI3ZJqUtfmpBzgIbSrvQ/7whJg2b582bi9Fk4XYFMA
70qwBQ4mJDEIspY/TLiOxzrbRAbmVDqggw3XnkuwOyiwKbArlHZywE+yIae2+Ws+tbvcVdvnjhge
M9YKyu+QzYbz1uS1K6zvy9RSsAiTXx97eMDV57/v1P7NFH/HWwdf66dTfDJejVXjOIpXCXFqGBms
SJHpq3Cvtl3ZFbQVz0bODy5IUbPsZq2maEVFPfgWm3i541v0rRhW3m4L53fZltBxThrsvcXYmE+8
gKF1hwXXBrIVreNbkAeo24k7In6oJMXpXcGlpivGGkK5wEDIKprltE1422S0gIV3a1OBjabJuR1Y
bTjQwgY1FaB1+TGvH7ZGbC3+gCPHi39h5ASY70tG4Aj6E9sXSx6k5TD3GkyCVLiGUhzsJNpNhrD7
dZwqHHsuwS6OgDN9WbjTsYtg93fkAMltwzPpCWmY9IQOTPqrr20Q1vLuIVGV43HeoLPYZLnFYJmb
wqS3nVER4Cz133/S/2by7/fkx4z9fPJv0k8aq8ZhbgpEoJJENi82tvySkrK7tHAChlPnzqQzdVYn
wCPambBBo21yYlWRH2kMi0HM7NhEhyDrWPA+nFHCgVDGDdoBs9xskpMj4U0twm18/U0aA27BAqJG
DIjT4s6EJinfXjdWONsEWgqrmEzsVRAjrOLp7gvOPqSq3gePLwyIPT4E7LhMtlAwhTb5zTVzJ0b5
D/07y2ywLoeDxN6Gr7iO62uLPwlMu0DCdCzWH7BVLjzzfzvgnw1cOsF+R7hVsNDbQQoBvI2/uL/H
5JsTqVl/8Z9wFquIVotfukz79YxeXduZjH7ledf4udNFfEwCacgi+fDjr7YE3pTHgYz66RJYeBNX
Y9U4xBnBblTRYAc7eiF6zT0/XJkNlW7TgWSChdwQcUz8o2l1odqY+Ljd/SyCTgzqmEogdBkiHAOI
XWiFCqdP4rdh8Jq78ewO0MhNCQvclKadfRKOKR9NPrLCL5t3w7fKZkCS8Bu8ttJPR7GkzV6bg2Sd
p80hUh78ebxV25FebPZsavDUECTkqeFkW+Ed40zlBvBj2NvH0GNoHaxB+EO2LqAcN7T2EPjTBaW9
7PqQaAVE5OycSaBAIYFnPUwHpAqOD38LsIsl2B4DLnw7GNdximBbYdXwrIfhjN5l1ymK7TWzxxBH
YgZGZJ2th6t1MYrvDtpnCY+VLJ3cT/320pO1v86CeBbFO616+Jt/uih2k4lgrBonIZN5m/IAbD2c
gifK1TDfhOCLXCvO/4aPs53eTFtnKXCTJrSIxYkwMUUUAc4B7rGtNczlYBXgLeNNsAq6+xbSTyxx
g7gZzLOi/Cte0FEaBB9gSVNBTDeY9wu/r8Fdo6Bwgle/UdzJs9nEHsIMYW+BjiG4RDBhAYtSGQpw
RUET4HbM5qt7dp+1+oIAMS/BjMO0gACG5PXDngWahZd28A4gj+wFwRPsriOMVh3bMSHuAVbFQFwj
PEujC3CiAsuaxCXpAWDvWyD9zK4NMQr6dayAp2GehVHl4NhdUPxqN6I7W+gzvt7rC2KRIRC94f+u
bJPEMUX5larlNVwr0Otfa2Hk4lhvNvLiAvx8cSDcuRqrxgkXdolXbXEcGSp2Zbsnnlm2S0lihMUx
YOQ6/q/dkUGGYi54ZdaK+eZyV1kcl9VAoFyQokD22fI5vho64YWtBqhp8jfkwdoiQi2HxYEnczbB
aJVB2yLftLGRppceBPepaodTXlMJTW6qmV48QGmF8Icv6xry5sakFAm+pEjpuFJ4cofH7BfKziGV
UvA8lIi0SBky5/WDuIZI2vwYPuGXwNuTagIMoHCKkDoY8zAFwnVhF4E0n2DPJWQY3YtjLQElldCn
I7KTT2TRAmZCBPNCDblMynGC8nqU0Kv7VeY6NlfNmPPlMojt2UiPG8gM8f/N5LF6Vtdu03fV/yzr
L7hOnrVS32tl/EKfM5REuhqrxklCySIsfHVIzWOGBZLZpliFQGKSWoAKYwgXOuFCSbVbECmJFu6U
UksQSJfZyHKSwG6QHbUg1L5lR2joKDoyqwZL367ZcHNRkOBoKiLH4jUszmZa4vptfKqrxcGWggMS
oQZsGc9c9Ru8uJnNHnobfHN8ctQY+CgdPx//sMj3R1wmcgnj33gMiQXrlRsKO5Kn2zHbtlSI/qX3
i0x8HQ4ikM0a7w7yWAZoHeE2gayCcFbIzaqD3P6GDRiOZ81UYdu6EwiYFDRFJr8dh71b4TqwDsLY
DgWLT9qqTNFMdpuUtjKv2NNW5ikrNVA8TWX/7IzqGMMz+Ed01W1aWmuF5+aSQ9Jw5a/6oXXzziFc
BH0/S8QOsZVXY9U4xWHu6iEvjn1lj7nQUDgxykYSie+BcXSfWsLRwn2H3xT5yD49irJN8kPsRwbP
18XiwO7aQoGCJdm1cF1YZFJVtqW0qbSWMFBXWTDWw2Z5oDSQVayQXpUZb+GK2Fflw6bSIegiq4xr
xUkiMXaZZaSgmJkDRwwPkpmpndCMACu6KpkropiQSJuhmF5xVxJcIpElaWALiK2HTwzJeihaXBdA
MgR0sV79S7TSHMg0AVKd20ULC1xB8jShpnrWT+Yl2nAgpZlOL4OHzZoBGC6DyBbQtwPlwpu2LT1c
sxuu2QnX7K7mSk2nb7dTVVW3VDNGDJf0r0islBQ2l9ZlnfcvvH6eNdTfa4gN96d7z0WMeTVWjZNc
13URAZfsNGeIIduRRB/ZRjj7rVh1DUzkRA/FlgOFNDyA1VyZuDgkCfDOR1UeA0nVdl51ozQYkLwV
zveIe8HFgeRFL8hrVFZCQ42MihUXStpA5XhlEfnWAwqRT80DmIhnlBAn3MdmFGICuWzc8+KM1+Fx
aI4P4fxRuC/Krmg/JNaKbQg3xRVfECAW7ZNiNv9tvsRscakcLGg2UpQSAGUFNF8oVvvSRdGmJ0hF
CRF8hNmvYx8BIyRkpEPKkexHCsrwDoAUJ+a61QD3EOwmPPnrJMkQSs0nQ23X1gnh69sSrI9Tuy5u
L12UmQW8zZfTOnLlmif6rvPEcuklEaoeHvW/9nJ6ltQjrl+FHOzPlpTGVlx6La3zhIYSD2DuEcFn
1zypMoTwllpDpJc28fSEe1pLEYHqFRYzA1UkcslBuKAMo4UW+ElXYsVFdKH0BDQtCQrkGNL7tilh
NpuTHL0pvdDkWxFl7FUqrED2UKWynVVREKzxUuZbkkM5elLkSse/IgdeMArEAlXZE94uu5J2KJSP
M5SQkoxPVWz0IpHTS0o7BWBFy0xaOrZTtB0/+J/QalgKPsnzKWXpuGLAMREOaHFyRennoZvj9dpx
pCKzBMZ0nCasuh3JOqcQBH9BPkOASG+j4GuOqDV4WOhn31oOc73vUy3SLWw6X7/aFNE+NZvHSiWk
4rs79zxcMlKby5bu3V5ooCSkbK7iXP1M+Q9YZs9S2++ldn+meVKir2g74HFDeWAxiO0N96q8gwD+
qLyl4P1sUm8H53Ah8t0wuxNyeaHWmIQMk7BmMHlnu9KHd6ZrQ+4v54AsDYwXcW4n2OjjtJCJwnbr
Rjgbp/Q8r9HkfilVehv7wVB8oozyDsm75Mlip6fYaTxCCm9cRNqdR4Z8xy2Uqg96t8JbMUEJPJCV
umOFE7Yk3EJNUF2OjCYAHhFCpEmCH4uuIu5j0bH8LJ6CWER7eoKICOcw+brtaPwlYA8B/ZKA0cEN
fZzESHESerKTmWUSiNmm0m7Cs7uN7diOXQJJdWfQhSMJR1KHo19HSFvPWUI4j8/iq719/dqqS1cn
nCiLb3twtacTEyKtWHxOYFAwItm0+5pzVqfJ5ULuf7V6+h+yCh8ifKmC5w83vMaEnfMq4QwZTCoZ
xkK5du4oNvmQrlikKjcJZtvcbMpuyQ5xeQh2xlZZAP5jU00HWwb89SDZ0MmXcxvxpU7oEY9CNJxI
pBwN76YhwG68hoKwiMnI5UiAeCRjJQhbIgSbUlush6WiE8iGInnikOpc0kpMC7z+cSGqbD3Rmez+
dafskEU6oqKZ7mQzqWKcLaI0vEaySTiFmAeuF4lVMlF4mmzsXDAFa5FUcvUtjDpR4nKgmWTHTP2N
jNB8a3IlE/d1Dx6dDWEjSPIjSmC4JQLuxhlvEmR42ElmWEDMhH4X+QSYvvlvzMWGSCwOpO6AvDhY
x8HafaTk/UQiNr/aNFIi1nwGJw5Zx0V20jW8QRw2Ut7iCOV9nq4fm1fFld0FPZHE+g/8YIHu+l6g
+xcE428WaBelvMTQVSlYoNyOdAIFxfAeiBqJhxd8+d5UFsChL1bLZBdZCutwgiaihUHIaLNo+ooc
XXVYiNmvuOpOtAMJf0lnKU3FwtDNI4VCvMeKlCN3KLSZKMlWUvTKX22tV3HiLEt+o/Li/C0vVd9V
ElK+Y2o94QqbaVaMD1k/JX+Ejdyaje6abmY71kjSfOgLVYXACdkfwa4/Q6nGKrsvV9az+zWP80S6
yPf6QQxpS9We+AUH8pK8ZDqoPQXptWYdB5spNG1uoMj/VQIQLGXbjhDVP2xLG47QTLwXbDD9ozjw
xkEy/uhMwSiOKJDciP7YNO0yC704QvZ8tvAs2CZ92/O19ele7NEu2CR/5q4d1w3NMUMPup3U97ra
4yVHq8uNHcVX9z71P3PB/mbhPuzMUlj2hwtXcrfJGrd9oUkcTqZmSdM2KO/BHTlk1PQcG/HKRP+y
SOVOJrAX9HJfpzKqyk5diu9JenppL7NnoKPopIf7UhiJygFZU0PeaEEmJVAKMlnMOJPmQpVQ/3cJ
qJs8X3QZbFFUBDtls1XCj7VZCiiHMiDqHYBMVER+GV4RcCnYlELjbhewnYJI3IMytLTsDzv2LaqI
qbyh/nd4hbrkRoMGAKpO748cSQJrXlBhBFlGLqJ5v5mFqypR4hN9P2/0Onirp/muDYieAZ6+h+IO
CRtxhIejb0dEoQ/N0wJRuPx2iCOxLWd+QvA7WONnfHGMI4OHLwUiOt3oQ+744iAHVOeY7Km5BKJm
nq9tSL9E4NK9rHlog93XHeM5ff12L4Ou4qeLzzreoLZji1z0y83rvs5/9mJ+FvR6L2j46z9c0FPt
NK74QkjQKtrz7BDfX5pcKO2KIG/gVFKmQgp6XSlFJTNUGUuRrrC4YERRsequBzngDrc6YJK6F6eS
uShnu+Cw4fN2ZaOn5FVT+5ImdUmnGf6oIL5VVfMZqjebTpPXI6JTlROsW9z9Ol3iAVkjFbwq/vnd
IhUlGj1dLvhkFR/6hEU9sRHM31jZihSXr2y8VDZA0jK4zooMmlM+etNLHm/1TIwyHfDVqkDGhEmn
iOKyiVmqSiSyooocbSGppWeRFwdkGYB9hqlu04FcDeDwnPsUlD+pkVSz24URenE8i70KEORCO46k
oB5vOwRfSUVlNCwd5viy6KPUoxYHVVVtP4udopkvX9tUot3u9PDF7sXDdgW9aMPcIl/sq/hi9yj5
dm3a27fyWdxFd+XLHus/f7X/buW/g+SOjf/Fyq/PyhfRuryycchfxd0+UmeVpZUvjYkYSBWLU4pR
1YUBRTH05aKoTembCfGjpT7LFkGMrJa9aCu/QOWcOjOwD3T81Q4B1thHuzJ9FLw3XqORgZDiXoF8
GypF1ko4qu5BwOJFpsTRXR0+Rm7qYm/VSkVZTlgr76FxXArFXy5vWN+iIFmuSCYivGJ1hhKcCDFh
d2CxsAdz+UZIVmaaO4FREQHARorWsmpWK/LtCumbBPcq1ZYDTpwwVIiCFfK+EKoNUIANkUWw/hLd
cRwhYn5sAYIEQGwCsoJF2qN3fWGX6OdLmhQmTmxc5HdOEyAqhZQd8P5ApE+hD0Qh0JiilS8OvRd1
r+MZrCcSp3gB9OXIfFAo4J5KmQJ1/TgAK2jo/Nqi6mUNp6Pr8kKA0z2VM6Maprv8oIq6JO9zxYGt
5sSYK0CjQma0/n/LKHwaiPP29Qfe5h+6BihlbKFKhrrkDRMrH980SexstsiLTGjRLuVgFhZTfWMl
JLBWX5ZZ1Bmpqn0TygZySTjBA9K7085lqG0FOrp+vZdPY1OTO4tUiNwrtQ3etYX+LW1LQqc+CFgO
NoKm9UfvqKayF565KffKi1eJSYlPK9IZz9FObwuU4UCTBBWTIDFqc5ecAIAqW++G0V20Sbx6vSEM
xCyKNWWcVG4mYtHJLJpKIaglT4ZPtLCES/64YhuVNXTXcl9xES9JnBRtDMXrEqJQykAHEDUX6+LI
SSsQavCJO1Ml4xOW57MqJL6YxuEIievbcFRHVkVEQhiFeb9fSuIIYQAEBSCndebbgFDK3+5TPcQt
EbJ0js1CG8Z2wLUClLgAaa/acC2BzLxgJlFAbRKKE8GdjTQg88gNeL42UhEYkL2it830grjpX1cJ
7Xh1Sr1VdzJKlSE5VY1MVlOfrdFkZVpTM6WmzNj3h9uUd7gx8BH/0KYob7S0LRbVynXc71skUYSb
3ZIHsO8vNGqHOtDVvIZvI2xVCSapMgk94MnYgGZTHxXIZfabjhJLXf3kNopY6Nogqd/00iYtPGjh
zms04tx2VOGhItau8Hp5H4FGcWpD2CWNl5ezepcxuR/Yj31LdLtYklFJVSVdIpu/errJDYhKVc06
0evN5VZXLU9wltUizEUUqK7kbKv8dGtzVWQ+9Terw8pYQzpbr97Y3n2qNs9vq1DqUMriJa+yC0oA
4hcRLqGMVxNDXu1FvGC2xYwdTGL4JGhtSIZjkjIH9xiYKjxRiq7gEdchvgNpZ/kiFpZhJ5m5hKHe
feVpHzNUvcjWMB2ZG8i2Mi2QiXMyN8reeBSjhCobT+bysiBd7ZPoOjEdj+cyBUV8IIVh2R/EKY3H
0HSJhp+v7XrntXU9V5CG5Hgizqa6eyjDc3a2oUTRuXslZ26nLMRVRNlu87Ld276NzE+Nzn2k/2v2
X/TH+ZqU13K3mEDZv6MVScsFdQmghpRQH9+fIHrT7OqwKhabtJQxVN0zd9jl4ALkr2yYdaSXQ3JE
Slw6TOFQASVbmW3wsHBs8ZCDFhJEt9GGCKwrYSadpqcI1bykqIaeLVol6tO7cpAsxO+Hw3SXppao
LSZosWXYQwI9SE2I4YQ9qUMLlkequ8RTq9L9K2dQlZ9QDq9MqRrN/EDmeoM5xfBQuM0dHTVxwhpw
4TE5JGfolYb6m54mfTgncpxolPXFsvJv3RuioilABtFcLedxUXO9jbc8hOxU2wtxoBK82RfivJeq
bvikouzdAUjdXYAnaCQ4lej0cXlcQpoqgYckzULoMRwIKwAy7XOeVq2K89oXPmVtx2OZglchXwPY
mEAmRagYAlwn9HKlOVLqQ1itLgehOsjYKjoVMUcArUiE2h3qSJplD324xXm+9hLuimvkmrrYqrtL
c8s01AGLlifTu7y41Ke4CPUW9dIIlmZ6LaoS3hIefvtAf85UvaV6iyr9P/KPuvblRXxDiSgLF7r0
SkQqyduBj7k0QtpEZod+Boult8kszCPiToUETEkyHgO3anYksYO+pchGBg7MmEX7MAl52Sa1zipq
EdCj7vVCKjav+lZDhKH0Gl7wCU4We9CURYR1lGvUVOiqBqwiauB8qoiWpaZIsDVTLWckoRdRc5Ky
UbtGPLmKViCNVo+MjBst5VhpmayOutCqVJOSXK1KG7+2Sv+2LxylZYrqJrJHDdGYKlF2c98J9kpy
xqOSQtJMMEYiWiWiqorW5EOpuH14/xyctu2ZVzIzy1M4VJDZCrankCjXW0BTMP+a+qT0jJjuOkRX
Kdh2cC3Ou2n0U0SSDHCUsr8tWXGMwAqQ0t1fSt2fYK47oonP26J1B+Wg+FHZ3OcJ7oYjulVNLt5u
T6u/t+8V6V/pwbjg25GM8hP0dQftCIEHf2nR1MT169emWiyabJ5o8Xej3Y9bJW6yyv5K6DKm1z82
5YXukNXaW0HdvO5yterFYcvZ5PNtt/4nds5zF27odv1FA5ev7LPyuPEzVQuok4u4XgdMlHKamifQ
CZeyxYMBW/gZW53hqJNdS/1KUAZJC3KV0SDp6qUATYmQ5U3D+yE4VMv3KVlDPZ5N7ov1jkdkC12V
WmrOq+auKrdHvAyfy8RRzaYX/buAYknGjBCfhUw00Iq6GBfvx6RyXCygbFyNsk5MbYUX8Sq1IRuH
34K7WJmT1Xu89BXVbIqQ4aSu3N9FvwxaaXtXaVpCZFNNOhLx6mjikQp7O64IG9UNUT6smsxSPKdo
C2WplJIrlJPDfaMVRaPREV08HNcC9hzzyDsdnu+SEgyzu9R9aUon3bzxEsJpDCFSshdfZ7hyIezM
cgVcQXUNDNpb3u1RT2V1an4MYkDcGWWDx6FX3GoMJTwuXhjG1LWgCAKPqxf58UydpYF8XD7vicAO
IZBmg9lS/+f2FpS27ciyCf5IkEFqdIh/G0z14Hr3VqCvdD3vFFwPZA+eKaSBVHPO33+1rcHz5700
N5DDBWrriB8nR+J9dHr0jY/+OqWINTtdhbNrq7/IrNVdP++pdrw34P12AP++BrS+JXRqzP2HQe1X
A9olFGWxq1WS2QmU2V0ZbuhpKg/OlMXEcfRsr9I3xFQYgYWAVaW6c6lMQRWIVEtRdDQILoeOc1BP
BPUlKSpx0GoaiqZYZ2oT7G1cWYrEwqRm21KDCPVn2GqH2aIant9oVQIeGDe4djVlV+8KdZFz2ylT
qtz+8SoovpVRrSU6yyGYq3itOM18u6OiEc/YT9JoujoKHI8q6Vdwb/LG0Nm0qu7I2ZlYKrTlZ3Tw
PpQU1CqEK2C1wvpz1baqh6WGQ8WkSpmFFsK5cpiCsUKNc0KdE9QYJdD0HFOLAJR41VvajuIy91pc
tXdEICmhqMCerZHjCeAF8RBfw4sOOXfk5RZA/4rS9iUJXlavlEALPPbWhbxPutEjk3faESq0fwmp
o5ysERGc+u4pk1mFlChkdmF2R4baj4MaOkN6EQBKX8ETekeVdPRqerIQpDj3eLIRSCcln0yysAaS
NMzq0CQPM9352Od0aJNMLIFIgz52WlWlaYfr8UaqH1/bcftrV6xE5xq3u7v597O5oLivEVkID83L
qt6JQPLEPdUyeBU1OxvT2/AXcYQ8BTb624j+U411e4f1d/9MZfXfGesl0TAtoyzAVxdA2ZjqbW9U
0s86xSgjD96YnoVMUA0rpmQ49H6dZCfmFgtFRAvFr8rxQQ+RgbjQpVe0GBwwYl1+MOVgXasOLX2v
qicqUuuru7HkwOQjr1Q5eFVFnddY6+rMQPco3B81r27wpCHPUJi/JM9QFIvd3lJf1qQquQS0e/b4
X73yWbGellWnLLV21EEtnjPZ+lZHWuisHrnFJCFK9neQWgrqsvI2vGs3HeXVxvu4LVchKiQj5Tve
tU3E5taGKzJNpyPouI6mMgNEHhCcUkaSspbi+epQGM7ESEU0FffssLJoKndAao4X6mUPWEblX6Y6
h1Hx6GG8rl440lJnhyPNZFFbMJV/KK0rzo+O0vjMGFceqEvF6wdMCQcUvexPCufaz6dxXZm72Oq6
/daVsI2onCQEpdO1KW+H2zeARyjO2weSyLd3EzG2a6A0U2xd+61LobZTldbMufKl7jizQVkZMgIz
EJ2XoGMEnPqz30xGCXVtMhrQN0BdoL5UfKGeBVoM+51VohoN0JZEcIf+aYKZnK6agYDjyJY29hqO
cPSfDSUqV5whyQ2itHV+8pWSPe8OH1mk0r2eea/QuXg3ebtDLnhtqwaj4Rq4K//9TDG2KoxXj/Ll
8tcoC3M1u0cI93un+PffXcazu+y6+p8gjb/uLteTzUrMN7l3RZToOVn6Rn7nXBe4qWkMSoyiggwI
5CXlJ1OaeU92bBYVKh8lZfl26cQ8wmElbxGzuKjXsybFj8pRU9GuIuip2hKlw5aULlgBuhegA2uX
aht1+jmStmFnVCynVlWkVCS9aP1EL/+mdvp0kVWKrF41Npdklm1DpyNJwHyUZoKT3RLiqWmCKnS1
0ahRA1dbB6l0yQe90frw7jVqR025CBd2Si3NbqPcmQzpzs1HyTP67Jzov08C/15vhwsRTt6O5apu
CBKwLN+EEFrSSpR6Ainrjo59omxsqb18k+pQDbSKGrm6i8rdwcKhvaRYQqq+UCXOFXJ/54DF5qBK
PGJRnQIhSCPiEtsVrvyJLJ0KAqcXtOMZtOCjS1AyStCpQwkpC38szqqrcLDSfpZHNO3i7fFAhYf0
q+japrRvHQZ9qypfyDgcej5j71IzeuKxmfqD5MxdTyFZtZBZwGhWxmQWuN06qCdwHcpClv7ujPHs
bNUxSiAopmif/IQ8oW9473Q1EJw8xSpActfxpqK8+ci7RCOzj9GD/FFzUwwNooQjayfJrgpqA6o2
me6AkIkVoo3nszPGSRMEoyB3SFrWAVwFIXfKEbgOlK5AfbPLFwqsBCL0enZQ2MQ8bKydoTMkfvl1
TD/BrpwbR0R1731Tmgu0ujcV6KXHEYfRyt8bnd+h2Og0F376CVLz6mfH9mMQp/cO6X7wSg0K7Huv
+T+zn+5HL7I7RNH/eD9VoNJw5y2ckGocw17Vx6ZK+anAgkQXEkaYjgVLtogwpuo7YbcnDRQmqf7h
xawIs0VdsDUP/OxBk7BOdnOooemu6gikYlApu2GNsOBTuYsuu05CUL19nHrTzioNpJpD6IwRqDd2
W2Vqp0g47PtQS4gr8Qlm8UaVTaParalyVC3upTvRaXXwZLHHsvWtnXusDiFb0XAmtCgI+Ks60irA
Y83W7hq5G40l1JlCWYXuBxeW2Hz9kNSiFi8rWioVNVdSKKOzGSVOOx4Cqn8qyeDSVHJ71HSp+7kM
K07RoU6RQJCebPSugOBTc47pZ0Ox/3LiLMni6sI7ckarOtHHEUb43HpcvQWaaCdRpc1pKPKqRBja
j+ujJPF0w/XoRudT8ljqRFX6cgN4rKhJWafjtJoOs7i+J68W6Y9ob96dkVc2Bj61hBJGrdiXt34e
caiLugt0z9Z0b9ciBrQqcc2JG9quYYK5Gi/le19Khmj31pkx3pACgY4XKeoYmdzF9QlbV6IpIvwo
LQn4xAPViXL0xfCixmxQLV1ke+/s7TgyO64JNN5FzKixwVoOHbHYnzOXxPCqAWzu+DUQJ4ZlO1M5
mFVtD4TMtj8eQLT28SP3aK/iwEVdfjyw8HgEJbAcj2cQ5GipgfQQotvwExunp7Ac6kBUv3gM0Voo
az1mCYzAddDqSAg6IpNeobzMtq2PRxHJsCRj355FKDQzSaYDVfe7Z1h6GtTEgVRypseRJK4n1dIz
WFMHVv/y69jasBFUTj85zZWVK+rD7IsXc3tvbzpDeS1I657McrXA9TMkzCx7jKwK8VUVX8/q+u5y
/awVuRqtzGhjxJfvNNf3x9/BB/EUvvsgKzSrQ1EKI3N9Sfi01K56SSRe4v9Wr3AdOh+pTLWZarKe
CoXI6Z6h/nlL9RE60lqyQdLrZI+bqrN0fp2KuZeOzlWwTw5bx4Xh+qvZDRTqUKsyTskZVWfxLR3D
t6LRGy3fzDKwYUF7boyQ9jGdq6wTscWF8sxdG5Oa4qghPEU8KvprS40nVDCHumC59P5Ea7iuNowE
oSpwlwABMrRp41NvDgXiV7052OSOdJ/qoa3zvouOICneft5lF9uDwhlt5RBf+zGpSOkreZY6VCPU
ZjgnbOchS8Pj8YO4dfatiFoRLP0RqF0xq1FYSHNPzJGUERJWzjiLT7WHKiHw6stJzlFa3O2dEHUa
StOZUXD5KD0grg9qj+4JSR1nrGN4lo5bUD3WiaNJqGKsrmFfcYyEuITrZ2ktuAQV4VxPZ0ml4QU6
XrRzvC+WOjY0lzAgb+g6plaPxZqSe1zxmAm0QlvGY3aC/pZVqDJZvuv0uHXFTqRzi/Ar/Ihjeupo
l6tq6OH9wmrspjd6faleo/tuy47M41H9sfyi6l1SNFWUp3D/qLiMRDdmx+PtBLx0dzsSrSv4BHxP
YoHRvViCx9Uzr3I/1dZN91KHPULvcJgtfDsVYJTQvlTJ8lLbC3gNNmiWGsw9zsWSR4UHoT64dkNe
OmT+xTHPL/t7Xna/iVflreiTbPFLRx6gL1QMLy/NfS0Vg/GiNVK+ZPxb+l3SHXoS2DMCHDFw3rXn
mRhOX0zU11ajeEeUlyXrIjU9i6UEwkdTQZ3Xvjr8OMFH4RgHVotBEj8Xvhsrcu8vPlwyN+nLtcAK
3LizwebQ+gOkUCh9vExgZ5VKasMfny9FRZHgJiXf9xcfsAey1WOyQukTRserhyWqgRHYDjrzCT0Q
PmMqOjORTpB2xtuHPNchRVp72kxmh6036xSSAlaaED4mEYcOC0tfM1gpWhEI14H/LvRA+KDqxOKH
RDtaYAQi7ZWJ/sdH7YGsTdaZUukDtuPcjNhf3ZXuhzc3T6Cr34PkB3EQWI/OeHV4Ar1HI5A4sbTF
uXd1e2OQ4gkSW5hil25T1d9xFfteyoas48fbF73EbPIFx7hemHzuu62ApM2KS5SIP98c0vfH/3G/
tr/92qPxb5zq6Wd7Kn9AfO7j1dhU1OA/U6dka5xETgLFha+0WNfJxTowB4dSv7h1aCzLe2OfN0+z
qJrZTackIMihrG+qYqZ681QVnZOr1zngytwvKduXn/Q1KD8YqI91UkW/Tr8VL4ro+GQdX6jjZOnA
HJ0w23X4iX5DTV0kkNVxSUpOoTPWcd9tp5y4IRczX1cbBd+yJeBU01uQjUbf4qioHbFam4pluTql
d8cZwQiLm+rSr/a8/eS+2Lo5cbSqiPT4Udc7hRf1Ruc6dcFSdaqkc1LN0c+pTv2P2iIVNtUJVqyd
xCqS021JBdT6WKSM7K4SQCP5OwmamxJlDLf4MX+RN9MJAzqvUFrmq/KtJkVMFnfoaDRVfZTs0VGc
2ZvuJ/vJUNR/eBsCjmfWcQvbWxzzWCcpVT9emN5FeMscB6fOj9v70XAQA7IbdNN0gSwwU9WPAVAv
m+aeNI9b8cfFH3st1vCeOMelgCsORJBAF+HuiKM7Gh52lWDMa79qCIDxtJsfm6xy+SYtth5LX7D9
8DX1U2mqwtA54Wi0T5wbLpkhirDjj9Vr6UZWJ9pN9+RkTmSFbnA5V144j1luKvcnhUubFmWdumcJ
eTy8pNj39R1cVHrq1fuOeBm1Py5q7eTNDUp67sFkqo9DdSaTJO5u4S/N6GkSGUgeZ80gWUoeQ2Hz
GG+5hQ5dxdjbHxd1aQy9SouDT1dk/Up4/DP8zOuZTNaMgkLvk6/H8leP69pn1hfGsRRDJVle9iNu
MgTrPEbBpT7/cfhquf5YUzl78k91Oi7S8qh7BA90RN5LR3fr9EoaqfBgeRkRVTJiUNEkiU2Fe0bt
iRYS+eNLWhdSq0PtynHF0GwrvnAdKI7xS6WWL6VNXwo91YaGr3xh/GF/oH3Ykyoxy0+UaPqqTq9R
iKlsaYiNvEjcaV9P4aruPo4i4OTOFhpUT+1GoKKjQDxc8eOZPeXrxebeYDYo46JPpYL9bM8IZtRN
K0KaoU8FNn5qlffDiCBHRHOGOlFr8IQ8I7ACIZGqJZChUJzsk01va97ypKtbIKRVGSo9NHYJtMAI
RAiVJ5+kJEsKuPpuB/aEVqHVfUKsEkhaPBLg6hscCYC5nqMgnsL9JxTLZHlof5NOV4psfwnRaqAF
IrEeJx09oZvq+o4H1UKGcj2QId0M7ECGeCVQAxnyZfJ+BjIEPIEMBUOzLHPQvoSGI7ACESpmOiAl
btnEOA8JT4HAE0quQISUEFXCDdOUIWamFaJH4xNyzsAKRAiazczbdaQQQeavvdMST4gaEry+AhGy
Zk9XlOTgCWFroAUypB2BGcgQN6R92dDhCXml/Y4QjkYHXlbsDRHsjveQsDktv7vT6dtDOdKqXqwX
UrfpR3zQjKc5nR9BlndjpZQhDmATO3+vYrc74khDCaeP92TZO1rCKRGwpt7Ach5/Dp2pM66qmkcf
3nBSPV57q5F5WN7meUR355efXF1efiLp98f3x/fH98c/OcYe5R1jd6QHf7OQ9aoKnDByTm+PCDde
VD66th+c485E91M8rkrnxeCKtZbQsbJh3yZtKCEWz+8dH3UwiWRV66r8CVJ5qIi+qWZRpwbjvKFF
RHQxmxqBq5iU8k/24rHknSIJ0xHDqyneJkahQSE6jYHv0qld7VeNn7tyBDrmnRjkiGCEQVQtkWIN
BdgqM2Jfu6o5Y6c+OgbFQw1dn6bGPDCvBN0q3d959tBQU1nFG0XyMLZMCQDl2+ochHlDFuODt27u
uto59ONDVV2uD2PEcNT8QIO0/S7wz0H9Kn3gxvhQZgxL3ak19JuDzhwrO1JXPsgH1KCmt8pa6egW
DYrkNciXdKdyxaDp4oOqLLzU4uQgol9DyUFCKB/WMxAc+aDmUBp0/KFkITe7Nsg79EEEgFc0zxzU
o1eDvGxFTvfmcGYOOiNAg7qIaRCF5DxSy0EhmAYlNfxIHgfqVED53NaBwI48mjzPYaYNFqiBOKeZ
tSb4+RE6l9oRLepXdeQhgnn6WR5o24KF4FASoTq8K+ZzmCd6akd1xBHreVDuhhgTossmHYsAeUEh
zp/2TjHOZIDs4AvrITQhm+/RlVHIvgzZnC870hCXCd0xszQ5DlnMcuU8ryPLmLMXV5Y3q19kd42V
kGd9RNvx7EgMgwJzEuXSeHbCOYERKI7sXprlfDr466igx+FH2sPFObojmxfSMEVojqzyyPLBrP7I
vhhZFUJDNpDng2f5YW0OZPVC9jPrAYnslVt1jEB16FzA83Qcevp0pAZXbe/2uxxSxz03Z4XWfff3
SC1vHuOQGt/sCvtof+NcltQEE7eP/e7NNpojm0JmW3mqhPuX/iLZdTYPP3sUQVGt8yiFQqOcZZ+P
oqgHQmmU2uZUIGW56KN9rkLXbqJThwMrMAKpcAoNdSqfssooDm3r7EBCdWS5alYlpZIquu16EUR/
a7ZXaPuyminLXrPJ5qPQqoHo25JVUKkFR4AnpMIreuilZnyEjjCrqKi0EDLrFhpzZKNCZOcyW5dV
WE8WbwYyu9cDqUSLLGBq2bOa68kWRtlvG4FUtGV2MfrVZNYxlW9PNjK6GT9Zyh6IsmL4KpBZTZhH
IbOdUW2WbSDy8NgnO5oKvOIMZWRPpS8QMqu6AjOQWdceyGxsKvyKI7O1cehztqV4Z3NnILO8UUPg
x7h9yQKHkjCOZ3myxHsHViCzyOMDmWVugajOyyx0KhezfUZmqR9ErQPsrJDZ7FQ+tg9k1rs4siow
s+IPQkGZ/Yye7Pn4QA9klr1+IMrQk5p4qIpEUhhJacwPJPXRP5AUSVIm5bdIauWhWhJJwawPJFUz
PpCUTvtAUj/lt0hq6EFSRvsDSS3ND4wPJDXVPlA/EJRWUlwPzgf2B9YHZmB8oH+gfaB+oPwWSd05
lfdNTX1/fH98f3x/fH/8e1N78wu1t4hc/8a5ohKcyYVy1fOKPprQQ7jHdQ9JLgmt6Pi7dA62Osh5
i2N1FSHQvKv5ebFqfKEWx7QSd2UNfADd9TjvZKEC2dBpq+rsAmlntto4rigF1wlnhbC7u+aXbxUp
qh+yTnLzYqcdx7/g6/ar1lFqrblUZILCACcdDUWni0+nprl7UEEURnZO/aM6r9bR9vgpjRwHxqXB
mVaXkhYCcgiztuTR6ngI6YL9IEf1qFAxz4yhjxKDfGkf5OHIxVkjB/c2KBhaz+BHterc752DnzQn
kWrLYTwDlU0adLyyD73GwIElMUgNq6H3HEoOU0XrGkj9+lBPDF/pvp2Dp4x76thVaXdXDjyVD9I4
a+jPUGsMOgPaB+mzNMyRQ1851JtDeer11AjQuwGWHPg1H3SStAZJgjQo3S7JgerCNKgzl4Z1cpgz
BxXfewV+ycFP1BOnvXz4NfEXQIEPRsBPS/duYKAGioNyBkHCDc4DGQ5E+oDmh4A2bEJx9OZgQQIq
5QEyJcHPEclz8/K4HDVxA3nE1u8IwDh3LwnAPFfzIQB3oDpY98Be09GFpdOCz3NOJwe0OqoDPbtQ
HXmuZ54LSJZAiEavFFGO6KwD+gpUR54G0poj+x8mEcjhisKviMDueIjAGYj+XRwZIXRHnkoipeB0
2dVq74a22bI7+zZKsL/fjW+HN12QPEsYgTg+Kfs/5rmo2RI828dkn8j/NSEYxfa/IwSjnU0cdP4U
62e7myziz/Mfsx1O9reks9EKSZkwAz0QDYOzLyatP4QkBqPtzi+JwWw+8P9LDIaMjSN+hRFogeLI
lu3ZHugpWYxjsKKU8U0MZkuhGeiBkMT9jiBcgSyZ/F8ShHT92+VPEIQtkARhlHg+BGH0NX0Iwuiz
9zSViBLRhyjcgSQKkyBMuX4ShdE39XdEYZSi/t2Jwiw7+BVRGOcTP4Thr4jCLHP4s4RhEoVZNvFX
Iwz/XsRh/8CfJQ4/CcQ/Ig4/CcQP4vCXBOIncfgrArF/4JNALL/FP5xA/CASHwLxk0jcH/gkEP9V
ROInofgrIvGPCMXyW/xpQvFXROI/mlD8DbH4BfsD6wPzA+MD/QPtA/UD5TdIrePvcT6wP7A+MH+B
8YH+gfaB+guU3yL6eP8e5wP7A+sXmL/A+ED/BdoH6i9Qfo59P3B+gf0LrA/MX2D8Av1n+CbAvz++
P74/vj++P74/vj/+bELiv/4f3+z+JQ==
"""
from base64 import b64decode
from fractions import Fraction as F
from functools import lru_cache
import json
import zlib

def check_row(row):
    cap,scale,target,X,Y = row
    types = [(k,a) for k in range(2,cap+1) for a in (0,1)
             if k+a>=3 and (a or k>=4)]
    assert len(types)==len(X)==len(Y)
    for i,(k,a) in enumerate(types):
        d=k+a
        assert Y[i]>=0 and (a==0 or Y[i]==0)
        assert k*X[i]-Y[i]<=target
        for j in range(i+1):
            u,b=types[j]
            e=u+b
            debt=abs(d-e)*k*u-a*u*e-b*k*d
            rhs=X[i]+X[j]-Y[i]*(d>e)-Y[j]*(e>d)
            assert scale*debt<=rhs*k*u*d*e
    return cap,F(target,scale)

@lru_cache(None)
def rank_bound(h):
    return min(F(h,s)+F(s-1,2) for s in range(1,h+1))

@lru_cache(None)
def pack_terms(n,last):
    total=F(0)
    for j in range(1,min(last,n//2)+1):
        r=(j-1)//2
        P=max(1,(n-2*(j%2==0))//(2*r+1))
        total+=1-F(1,P)
    return total

def hub_pack(h,k,C):
    total=F(1)
    for j in range(1,C-1):
        r=(j-1)//2
        P=max(1,(h-2*(j%2==0))//(2*r+1))
        if r:
            P=min(P,1+(h-k-1)//(2*r-1))
        total+=1-F(1,P)
    return total

def check_step(C,rows,diagnostic=False):
    D=C-2
    bad=[]
    orders=pairs=0
    for h in range(2*D+21,C*(C-1)//2):
        orders+=1
        B=rank_bound(h)
        if B>=1+pack_terms(h,C-2):
            continue
        for k in range(26,h-2*D+6):
            pairs+=1
            N=h-k+1
            c=next((val for top,val in rows if k<=top),F(1))
            lower=max(B,h*(1-c)/2)
            upper=min(3+F(N-(D-2),2),3+pack_terms(N,C-4),
                      F(7,2)+F(303,1000)*N)
            if lower<upper and lower<hub_pack(h,k,C):
                bad.append((h,k))
    print(C,orders,pairs,len(bad),flush=True)
    if diagnostic:
        assert bad
        print("range",min(h for h,k in bad),max(h for h,k in bad),
              min(k for h,k in bad),max(k for h,k in bad))
    else:
        assert not bad,bad[:1]

rows=[check_row(row) for row in json.loads(zlib.decompress(b64decode(ENC)))]
assert rows==sorted(rows)
for C in range(35,41):
    check_step(C,rows)
check_step(41,rows,diagnostic=True)
\end{Verbatim}

\subsection*{Eleven classes and local degree bounds}
\begin{Verbatim}[fontsize=\scriptsize]
from fractions import Fraction as F
import json,zlib,base64
s,c,g=100000000,39300050,[16419997, 300000000]
enc=(
    'c-p<13v%7K4n#MkTg5-XwM@<axA9H`)OBLVjy-XnQx(aQNP<8EXn_9bpSss9rPQnTTJq-kq}'
    'JW}8CShp9eLYqdp>Kuo$s#gy5*}jnwj}EY@Ly%uQTn{eT}-7UCl@1r(bi_99YXN_fqC72h3^'
    '3wdDQz7(M6O{+zUVdnI;nmN#R1-t1-5bMiIsz1jYGSX*bJ(T?nUwzVWLGmPfd=2h94@gX}k?'
    'tGz-zFS#?VGPuh-Kz2Vy5_1sUEfSn`8?iyFp2chR^1^o7O9P4^=!h=zGM`B_Jw{^q3*+v>|;'
    'yGznZ7F_G+HlOWcfNy}X)-|It69VZ~-h-{&@&rFP$QPGsF}#!VY~!>hi#$BmCs(uO^CdSCBZ'
    'QRl|a0!$kFWmaor*1kWQwO-T9ZmV2gb1U8rt#9V@9;(rAi7+cG6w*e|Z@bwHVZ4-agLS6@X7'
    'M__0h|_Kd^>Mjo)~`4%hpoPU(Cx+FzmeM-13!6yDnJ8&C1?OPB;7P+@Q;?H*y*CA@8!WeC-R'
    'xw|Hf9*u-{{cUrTC=OOGz^KZxO8u2o{9E~R*PCr(uv#x7+djH&MX={5*ExXiRhF{wfK(&GOz'
    '|EW^V?$k@{@HiJ*NUxM+u^e;Q_3o<zPpUMl?D`^%Tzq*1dGosv#NUm-e3&uU<@2__e-wdd|r'
    'mHlqN7*-NrJ2-zPS%16=B6;K3Tv${=uni(Y*Wy;^83Ua<@^7IgpKa&s`o@F=cc{be~YukXLQ'
    'T*Iico146cxHr7ac7dWY5ChF@;>op`9av*lS2!0Wmc9qdYw7_`MjWX{j@<AgFtM?rUM-SCk5'
    '@$85|7?nmYv1>Z(a7esO)6uVV+=JPOP`{dxZL-y+@2&#8H5MSx<cOF0oMB$hbz@%NwwDGTg`'
    '~6WE;Cm&u{NzdT-BwfA?HU)zXJ>pw3a@f}rb69KFN!d&5gk1-Kr)vok0(~XE#O=);;B70Hju'
    'E-70K%u_h6Gzv~s5@iV8Y?3=@h(>{`_Ar!Ujo}Nq-Z3PN%%bZ^0gF%>te($7TYyK{U1qzIsh'
    'ElXlx28I*egqibk?(NHI}zSKub?PZA7K5dQ%g7;O|xVhx7_Qo$Y+QqWFOhO2t6z&m=y`&cnn'
    'ARyv>nqnfcUbRm0T=V66%Oq^-J!|92@5@AV(s$n}lkAfkAt@8D8dB{g@U@LI8B!dSbhO(*MO'
    'Q}KNV)e6vPRC3T|r8`pJe2rsK~A~4%h4Yt`=sEn?934w4yktDGD2nR|Mb;-HrdMPy@4bF{a;'
    '1rQ|v5-;#<?@Jys4#B?a=AYHnoBAVi$jcmS1pHPu*Yk@A_kO|lW7Evl6#fd+o_AV|mR*P{H%'
    'CyNGtL!04$E;P)H@sKhR|WP0RK)dSM4zOAc!i)#)K)5UAoY3S&<LwOg?qqQEAV=K{WIZ!G@w'
    '*M?Xm+Ythh8dl1?^gZIF5E5!U6GK1@E4c2%aD%pLzgyksURjPUR^G7#GCNfXMbX(~iMp|;=Q'
    'FFCsOMbfMF^CH;-dFzcxeIzaGM%oz|lW?i=kJgFO^UK#;zRD`V89AY<ois)wC&eRY)rzt&25'
    '<#bN`MyOV{U~z2jWwqKe#2Tgj!49Vfr0ANA~E@$>5%^|M}shpIyKA8SB8S<N9j-2MUHwefw3'
    'Cvci<63Y10_!_uAjO$X~pOvzOwaVxM0dB)m1%B=<8xukG1y@Cp1(BEp}{wF|_iaxy_WTeRwz'
    'gZ8D96sccbxr0YjIZ!Suy3_T5R>Hy|8pte8;pqO*lp{3)k|6z!emu1a9{+KWXL?+!Wk)KX6T'
    'BNMte}<D0W0)u*3x!)WR!zlm+}wKDkt%K&;rq<kzlZyhm}_$6(W!7)N}+6o6fnN<}7n5RG6D'
    'a25g#>4B1m=VO@ikSYl>Hx$Zye2=_6UPQ@*zSi$WC#eNj-3&&pzZ3mYoy@oY+BJ$QJOjx|LL'
    '><lnMPjZy0S=gC6O>nh9xuQYCR``@x27j(m81{4@Y0`1Co_`aIxR6Y{9p_zVcpK!$VQS>8t3'
    'YMRLhtxr2o@A}>e+23g=XY!tze+2hK~D7|n2{l0IedWo7(mFSdcv*`+yNLM&_wlhf0;)%TP3'
    'Nuueg1N4wl=%K=py$ry6hg^bQ7$n)D{9m-EeikSerSU1E2qv};cXJfxOy-Fb{-P8l9_{@CX2'
    'L8AbsD(H}Gm%Exxl8N@ClhTRZps#hpCJUXC;W?Q8w_`sYEuiXbpHKPnYl-aVe3ys@Hbr?8|f'
    'O9p_DwS4rrdv*`?IEI*$3jTGzYuS(c;k!u8qIQ(Als)`Vx7K_@C=*B+__(`w5B2*f%9-zFz$'
    '%ur*4jsi`!b+XPVU%1-z{&%81e>;<}S-~mv<*m3z6)99{Dl_)L`VurF!(ok%azeC;|Nv-6*v'
    'mhFU_R$f*&kI;dZahw0}QW6igNRokStKN@^D12F?`*hEs4wg+)I`;I|d<cee$a*%$maqk~s*'
    'EGxB)srZ1rcxRGYN~tkkv*yPi>X>6#hJ2Ohpn0FnR|G<HQbgR4VQxk5_?VADAf-DL#lf=MA4'
    'R$Wv-uRgOSNTAI^4n3JG}PMJ;GeNCoC8L|2U%8j51kc9w7t)f3gyDiH0w@nC<l#ognRA5|5L'
    'DyIk@0ou2dSA0IjeopQYT>YcT8D3Lbetgzwy&boKY+*Og7N1#$(txX|ZVy|C#{@h?AT$q@eW'
    '&5>o}c6|(Jc6Q74pcAt=+4TeJyeFiB$-SSj*K;zn#$%o)wba0_QO501k!zIsxaHT9hDn27O2'
    'bEl+hC#3N08w*eo~uVhvGXt5g2nEJzF)J^@%i<ynsx{KQRAmQc4)zkuR;I`gTkp*`knJkhZf'
    '$;74*zPDGian#q$Y6IZ&)~?4r!V+;e^D>$+*rWBDmOltq`Xat{HG=5ie0@ISe>*ex8vXCJLj'
    '-CxUtKjm~u#9aMT_Jo3=drNQ<r&xPr9p2y+y~ISE-b*MHpYI(GWr$|9p-kndfY1#7}$&5TqC'
    'lGhWtFlIRKQqjZ0YS7_oms*g*Zt$PTQ8#qQNhVYg>kjoNy;#kg$K4_PK!6F`P`meC0n|6e(f'
    '&+;IXwJiiXW==2#VLJ&((<fID3{Pp^lcJ$BtW>&Rul2vRddEHyAUMeAMcBAd}H{dSn&mJYn~'
    '7_X;0x1yN~AE>KN$XvEKg%%FKdb*#w2-y_J_>ATRga5(a>x|bq_V3d^4_K^p!v~(UKsX`4pa'
    'l-wP%403cM2jkjuNComSu)<vT$)H%Ly;u4lwlyi4a&!hBwby=$nKi6zAMQY{=Jgioy#%QMx;'
    'PpnPvzQDI8U>FxZb5LN~)0gX)tyRfVnGh&<hs;*=d&p25k-1wDv)oQRZOTz96LJ~P5WmxAZr'
    'Q$@QWTZWJ9q?!tHxP6b}^kccY73GP)D_7qF=1`xB8cfHeOCGB}Q?jG2LtFK0X5BjXQ`5n9LQ'
    '=NM<36L6C05Ednv|btEh9W*Mj<ogf;EzWx0B1EP%^Phmeew5624FG@DaUD^}{|b9^bUXTthq'
    'vU>><X@_-$*j-urIJrXK!v94Mj@hOEV)8Vu&$uZO<W7?FZcy3wbwBb5L5hNP_exPel-t#;*Y'
    '~j^c9uO09_BvTDDEYN_I)A-=u*DxSglnGOoj4SmY*A5vQcx7>YEr!t^h|#Bo|N*YU5^gJU`X'
    '7EshAKZ?_26TYSX4xnOK%MLd*V3xNIC<J)#B+x(p$yaae9eM~3PR$Zev(Qq)*1<IwixsnILQ'
    '#$XayH3X@gbpjs0TiUFXa})V@{(h#lXJEf0$cxv$K=G-|(~$gv$duGc-R$OzKZ{Hm;L#|1sH'
    'MFnNe6PMNP4~TP*qTbUZFib^g2izx|yctlQseyQ7=fuXiWiC=Z@B2#KwpLr^>8JYl+b&CN0r'
    'jFd{|;BUaKBL)np``iR`3CYO#p)Xnk`@HJpxsNlD0`gA;{^3RHo4Zz87#fO#<CS(%#Ko!NTD'
    'vZ{bGPzE6q75!oEA6N!lBH%Qx)uGHd8s=c`6=gFFGBAoRl@bv7gubKoJ7fPxFB-2h;QtwDQ9'
    '^Pr}hYAu%+z5WA13()XaS-Ve8bHP1US|cIg&Ps*GH?WQ|U558>L1kcQJ{9K-t}k*wnyXC9&C'
    'G(FMu$8r;Ws{W|J^~Vt(OAV-J{Cx$WKC!r~QOx9ky2NP#P|4FQKwYCE2Glx26?le1jiAz0I7'
    'Q3do|uV0tN}Cm#~jZM0RE}-c@G+9$%||zD;HvxGL@7tRsd$~VCXaJkGUcWE?D#0{~u^SHGg`'
    'kBHJHq0*-^azp#lFh28~j<Igg7L?lldyZ5O&dp0wAJ7mi(sd<(&AT&cM9+Qf^n<7VMU<5yd4'
    'OhtxM5^>`%3!x_YK1ARfY%=Gh{FY|%i29MnFV?)kkCO_w%1x#RxKrHNm_ccZ5_XPxlW^RB<D'
    'GGt>2aP=|A1Qyt`fcct>h%I9561z92R6FuG3yIXmc6pm80&2sH50Du`g|rB_f_a)lrb4=Y9A'
    'kgq&o6z2u<h(bD3LajPv_UQB*#d{HG@ERS8p56TW@BaW{wl5<'
)
T=[(k,a,t) for k in range(2,26) for a in (0,1) for t in range(1,k+a+1) if k+a>=3 and (a or k>=4) and a<=t<=a+k and (k+a!=3 or t==1)]
X,Y=json.loads(zlib.decompress(base64.b85decode(enc)));g=F(*g)
count=0
cr=0
for i,(k,a,t) in enumerate(T):
 d=k+a
 assert k*X[i]-(t-a)*Y[i]<=c
 if t-a<k:
  cr+=1
  gap=max(F(0), F(1,d)-F(a,k*d)-F(X[i],s))
  assert (k-t+a)*gap <= F(579,1000)
  assert F(1,d)-F(a,k*d)-F(X[i],s)<=g
 for j in range(i,len(T)):
  p,b,u=T[j];e=p+b
  if ((d>e and t==a) or (e>d and u==b) or (d<=e and t-a==k) or (e<=d and u-b==p)):continue
  if d<e and t*e*(e-1)<=(e-u)*d*(d-1):continue
  if e<d and u*d*(d-1)<=(d-t)*e*(e-1):continue
  if e==d and t+u<d:continue
  w=abs(F(1,d)-F(1,e))-F(a,k*d)-F(b,p*e)
  assert s*w<=X[i]+X[j]-Y[i]*(d>e)-Y[j]*(e>d)
  count+=1
a=(1-F(c,s))/2;b=F(1,2)-a+F(13,27);u=F(303,1000)
coeff=b-(2+u)/11+(u/11-g/2)*26
assert a>=u and u/11>=g/2 and coeff>F(3,8) and (u-3*g/2)*26>4
print('scale,c,g:',s,c,g,'type-pair checks',count,'normal-high',cr,'a,b,per-site',float(a),float(b),float(coeff))
\end{Verbatim}

\subsection*{Ordinary mean distance}
\begin{Verbatim}[fontsize=\scriptsize]
def ordinary(cap,start,stop):
 count=0
 for d in range(start,stop+1):
  cases=0
  path=d*(d-1)*(d+1)//3
  gain=sum((d+1-j)*(d*min(j,cap)-cap*j) for j in range(1,d+1))
  for n in range(d+1,d*d+1):
   q=n*(n-1)//2
   lim=n-d+2
   if lim<26: continue
   if n*n*lim>=(d*(lim+1))**2:
    continue
   for k in range(26,lim+1):
    if n*n*k >= (d*(k+1))**2: continue
    m=n-k
    hub=k*(k+1)//2
    loss=path+(d-2)*hub+k-3*(d-1)
    ud=q*d-loss
    # We store the common-denominator bound U=H/J.
    H=min(q*(n+1),3*ud,6*q+(m*(m-1)*(m+1)//2)+3*k*m*(m-1)//2)
    J=3*q
    if n*n*k*J*J>=(H*(k+1))**2:continue
    # At mean U the clipping lower bound for M equals A/B.
    B=n*d*J
    A=(n-1)*(cap*H+3*(gain+(d-cap)*(hub-6)))-(n-2)*d*J
    # The transfer lower bound exceeds U=H/J
    # exactly when X >= sqrt(k)*Y.
    X=n*B-A-H*n*d
    Y=H*n*d-A
    if Y>=0:
     ok=X>=0 and X*X>=k*Y*Y
    else:
     ok=X>=0 or X*X<=k*Y*Y
    if not ok: return (d,n,k)
    cases+=1
  count+=cases
  print(cap,d,cases,flush=True)
 print('all pass',cap,stop,count)
if __name__=='__main__':
 import sys
 C=int(sys.argv[1])
 assert ordinary(C,C+1,int(sys.argv[2])) is None
\end{Verbatim}

% END SOURCE BODY Fajtlowicz_Diameter_and_High_Degree_Sites
\unitstart{Appendix A}{notes}{The ten research problems}{Ten problems from the Verifiable Mathematical Reasoning Corpus}{}
% Body of the original NOTES.md, typeset for this collection.
Research notes by Keio and Shaitetsu, companion to \emph{Degree cuts and the Randi\'c--radius inequality}. The separately issued four-page manuscript gives the graph proof for Problem~1. These notes report original proof arguments and finite checks. They have not been independently refereed, and they make no assertion of historical priority.

\section{Reproducing the calculations}
Use Python 3.10 or later with assertions enabled. The Chowla scripts read the expressions and Lean interval data from the upstream checkout. In the directory containing the scripts, run
\begin{Verbatim}[fontsize=\small]
git clone https://github.com/clambro/lambert-nonquadraticity upstream
git -C upstream checkout 008e3136a95e1e1ea605837ca9742b8613b2034a
python keio_ch_six_envelope.py upstream
python gieseking.py
python baker.py
python beukers.py
python denominators.py
python heat_window.py
\end{Verbatim}
The first Python call checks both the improved Chowla interval and the obstruction. The Baker calculation requires SymPy, and the heat-window calculation requires mpmath. The numbered corpus references in these notes use the \href{https://github.com/QEEDDD/Verifiable-Mathematical-Reasoning-Corpus/tree/8a3f6895d0b9b9f125e43157e4a003a2a2b51704}{VMRC snapshot 8a3f689}.

\section{Problem 2. Chowla's question for rational Lambert bases}
Set $F(t)=\sum_{n\ge1}(t^n-1)^{-1}$. For coprime integers $a>b>0$, put $\theta=\log b/\log a$. Use the positive determinant, product-formula criterion and six-presentation clearing bound in the upstream work. The fixed parameters are $(500,27,516,41)$, with decay scale $S=371952500/3$. On the additional cell $86/19\le x\le527/116$, the sixth expression saves $500-110x$ over the baseline, as checked at every change of affine branch. The extra cell is disjoint from all 400 chosen by the source. Its additional weighted integral is $I=1444283345/16059278496$. By the source's summatory-totient asymptotic the change in the sufficient cutoff is $6I/(\pi^2 S)$. The unrounded original cutoff is less than $0.57329010999214$. The new cutoff exceeds $0.57329011043309$. The existing norm argument therefore proves irrationality at $113590/791$, for which $\theta\approx0.57329011034470$. The additional interval was checked in rational Python without building a Lean proof for the extension.

For these fixed sizes, on $5\le x\le543$ the integral of the maximum saving from the six published expressions equals the unrounded mass of the original 400 cells. The verification uses 762 affine cells. Above 543 the clearing exponent is already zero (\texttt{Orbit.lean}). Below 5, even making the clearing cost zero saves a weighted integral of at most 671890. Indeed the baseline $B(x)=500+v(x)-v(2x)$ is at most $500+516^2/(2x)$, since $v(x)=\sum_{j\ge1}(516-jx)_+$ is bounded by its decreasing integral. Thus the largest possible threshold using the fixed estimates, even with free clearing below 5, is less than $0.576584565$, below $\log(4)/\log(11)$. Changing the sizes or obtaining different estimates falls outside this barrier.

\section{Problem 3. Gieseking volume}
The script uses the rational nuisance coordinates and Rodrigues output lattice in corpus [045]. It takes $m=n/3$ and $K=n$, tries every three-element support with a rank-two nuisance matrix and a strictly positive kernel vector, forms the integral polynomial, divides by the greatest common divisor of all its coefficients, and scores the logarithm of its maximum normalized Bernstein weight divided by $n$. The reported minimum applies to this finite search. Degree 45 gives 4048 qualifying supports and minimum $6.75850$ (the asymptotic log budget is $1.175542931\ldots$). Arbitrary degrees and larger supports remain outside the search.

\section{Problem 4. Krzy\.{z}'s eighth coefficient}
For analytic $f$ in the open unit disk with $0<|f|\le1$, write $t=-\log|f(0)|$. When $t>0$, take its Herglotz probability measure with moments $m_j$ in the convention
\[
 \frac{f(z)}{f(0)}=\exp\!\left(-2t\sum_{j\ge1}m_jz^j\right).
\]
Set $\varepsilon=\max(|m_1|,|m_2|)$, $q=|m_3|^2+|m_4|^2+|m_5|^2$ and $B=2m_3m_5+m_4^2$. Then
\[
 0.9\le t\le1.1,\qquad \varepsilon\le q/1000
 \quad\Longrightarrow\quad |a_8|\le(2/e)(1-0.78q).
\]
The Gram matrix of $z^3,z^4,z^5$ differs from the identity by operator norm at most $2\varepsilon$. In the projections of $1$ and $z^8$ on their span, both coupling vectors have squared norm $q$. The squared norm of each projection is at least $q/(1+2\varepsilon)$. Their cross product differs from $B$ (up to simultaneous conjugation of the moment convention) by at most $2\varepsilon q/(1-2\varepsilon)$. Applying Cauchy--Schwarz to the residual vectors yields
\[
 |m_8-B|\le1-\frac{q}{1+2\varepsilon}
            +\frac{2\varepsilon q}{1-2\varepsilon},
 \qquad |B|\le q .
\]
In the exponential expansion the retained terms give $-2t m_8+2t^2B$. Every omitted composition of 8 contains an index 1 or 2. Bounding each such product by $\varepsilon$ bounds its total contribution, before multiplication by $\exp(-t)$, by $\varepsilon C(2t)$, where
\[
 C(x)=x^2+\sum_{r=2}^8\binom7{r-1}\frac{x^r}{r!}.
\]
In particular $C(2.2)<112$ and $\exp(-0.9)<0.407$. Since $q\le3$, we have $\varepsilon\le0.003$. The coefficient $1/(1+2\varepsilon)-|1-t|-2\varepsilon/(1-2\varepsilon)$ is at least $0.887$. Also $q\le1+2\varepsilon$ by the projection bound. Thus
\[
 |a_8|\le (2/e)(1-0.887q)+0.046q\le(2/e)(1-0.78q).
\]
For an explicit family, let $\mu=(1-\delta-\eta)U_8+\delta U_3+\eta\delta_1$ at $t=1$, where $U_j$ is uniform on the $j$th roots of unity, $0<\delta<1/2$ and $\eta=\delta^2/2000$. Then $\varepsilon=\eta$ and $q=(\delta+\eta)^2+2\eta^2$. These nonsymmetric examples converge to the eighth-order extremizer. When $\varepsilon=0$ the projection proof works for all $t$ and gives $|a_8|\le2t\exp(-t)\max(1,|1-t|)\le2/e$. See Instanes, Llinares and Vukoti\'c, \href{https://arxiv.org/abs/2610.23817}{arXiv:2610.23817}, for related vanishing-coefficient results.

\section{Problem 5. A necessary property of any algebraic least Mills constant}
Assume the cubic Pisot reduction and divisibility formula in corpus [010]. Under algebraicity of the least Mills constant, the late tower gaps $g_r=p_{r+1}-p_r^3$ satisfy
\[
 h_r=g_r/3=-(X+Y)(Y+Z)(Z+X),
\]
where $X=\alpha^{3^r}$, $Y=\beta^{3^r}$ and $Z=\gamma^{3^r}$ are the conjugate powers, with $\alpha>1$ and $0<|\beta|<1$. All these factors are nonzero at sufficiently large indices, and $g_r\mid g_{r+1}$. Then the number of distinct rational primes dividing $g_r$ tends to infinity.

Indeed, if the set of prime divisors of all late $g_r$ were finite, pass to the normal closure of the cubic field. Invert the primes over that finite set and over the norms of the three conjugates. Each of $X,Y,Z$ is now a unit in the ring of $S$-integers. Each of $X+Y,Y+Z,Z+X$ is an algebraic integer whose product is a unit there, so each is a unit there too. Consequently $X/(X+Y)+Y/(X+Y)=1$ is an $S$-unit equation with infinitely many different solutions, since the ratio $(\alpha/\beta)^{3^r}$ has strictly increasing absolute value. This contradicts the two-variable $S$-unit finiteness theorem of Beukers and Schlickewei (\emph{Acta Arithmetica} 78). Thus infinitely many rational primes appear, and persistence by divisibility proves the stated limit. This is a conditional property under the algebraicity hypothesis.

\section{Problem 6. Baker's question}
Take the primes $p=160169$ and $\ell=(p-1)/8=20021$. Compute the four $\beta_{\ell u}$ modulo $p$, $u=1,3,5,7$, using the Fermat quotient expression in corpus [053]. Their residues are $52672,6782,101934,159050$, respectively. The finite-field eighth root in the script is primitive. For all eight choices of eighth root and four sign characters, all 32 sums $S_\beta$ are nonzero. This suffices for condition $C_p$ of [053] and hence for the rational linear independence of $L(1,\chi)$ for all nonprincipal characters of modulus $p^2$. The computation concerns the displayed prime.

\section{Problems 7 and 9. Finite denominator exclusion}
For each of $\zeta(7)$ and Euler's $\gamma$, \texttt{denominators.py} gives an exact rational interval excluding every rational of reduced denominator at most $10^{150}$. It uses $N=512$, 120 Euler--Maclaurin corrections and the periodic Bernoulli bound $|B_{2m}(x)|\le|B_{2m}|$ on $[0,1]$. The absolute integral remainder is bounded by the absolute value of the last correction for the reciprocal powers in use. For $\gamma$, take the limit of $H_L-\log L$ after summing between $N$ and $L$. For $\zeta(7)$, sum the infinite tail starting at $N$. Enclose $\log(512)=18\operatorname{atanh}(1/3)$ by its positive series and a geometric upper tail. The nearest rational of denominator at most the cutoff to the midpoint, obtained by continued fractions with \texttt{Fraction.limit\_denominator}, lies strictly outside the closed interval, as verified by exact rational comparison. Irrationality of either constant remains outside the scope of this finite bound.

\section{Problem 8. Beukers's Conjecture 1.27}
Let $D_n$ and $\mathrm{hv}_n$ be the sums of squared multinomial coefficients on four and five parts, respectively. Put $W_d(X)=D_d(-1/8)^dP_d(X)$, with the ordinary Legendre polynomial, and $S_p=\sum_{n<p}\mathrm{hv}_n(-1/7)^n$. For every odd prime $p\ne7$ and $h=(p-1)/2$,
\[
 S_p\equiv 2\sum_{d<h}W_d(3/4)+W_h(3/4)\pmod p,\qquad
 \sum_{d<p}(2d+1)W_d(X)\equiv0 \quad\text{in }\mathbb F_p[X].
\]
Indeed $\mathrm{hv}_n=\sum_d\binom nd^2D_d$. For $y=-1/7$ use
\[
 \sum_{n\ge0}\binom nd^2y^n
 =\frac{y^d}{(1-y)^{d+1}}P_d\!\left(\frac{1+y}{1-y}\right).
\]
Modulo $p$, $\binom nd$ is $p$-periodic in $n$ for $d<p$. Multiplying by $1-y^p$ gives the finite sum. At this particular $y$ in $\mathbb F_p$, $1-y^p=1-y$, giving a sum of $W_d(3/4)$ over all $d<p$. The reflection $P_{p-1-d}=P_d$ modulo $p$ follows termwise from the binomial expansion in $(X-1)/2$. Also $D_{p-1-d}\equiv64^{-d}D_d\pmod p$. Indeed grouping the four parts in pairs gives $D_{p-1}\equiv\binom{p-1}{h}^{4}\equiv1$ by the carry in every other central binomial term. In the Domb recurrence
\[
 (n+1)^3 D_{n+1}-A(n)D_n+64n^3D_{n-1}=0,\quad
 A(n)=2(2n+1)(5n^2+5n+2),
\]
$A(n)$ changes sign under $n\mapsto-1-n$. Thus $64^nD_{p-1-n}$ satisfies the same recurrence modulo $p$, with initial values 1 and 4 (the second follows by taking $n=p-1$ in the original recurrence). Hence $W_{p-1-d}\equiv W_d\pmod p$. Pairing terms gives both assertions. The comparison $S_p\equiv a_p\pmod{p^2}$ with the newform 14.4.a.a remains open in this work. See Beukers, \href{https://arxiv.org/abs/2403.03301}{arXiv:2403.03301}, Conjecture 1.27.

\section{Problem 10. Positive Laguerre expression on the central interval}
In the Rodgers--Tao normalization
\[
 H_t(x)=\int_0^\infty\Phi(u)e^{tu^2}\cos(xu)\,du,\qquad
 \mu_j(t)=\int_0^\infty u^j\Phi(u)e^{tu^2}\,du .
\]
Fix real $T$ and let $B_T=\sqrt{2\mu_2(T)/\mu_4(T)}$. For all $t\le T$ and $|x|<B_T$, both $H_t(x)$ and $-H_t''(x)$ are strictly positive, hence $(H_t')^2-H_tH_t''>0$. Indeed $\cos y\ge1-y^2/2$ gives lower bounds $\mu_0-x^2\mu_2/2$ and $\mu_2-x^2\mu_4/2$. Cauchy--Schwarz gives $\mu_0/\mu_2\ge\mu_2/\mu_4$ and $d(\mu_j/\mu_{j+2})/dt\le0$. This provides a uniform central interval for the corpus [126] heat-flow approach. At $T=11/50$, \texttt{heat\_window.py} certifies $2\mu_2-49\mu_4>0$ by interval integration on 4096 subintervals with two analytic tail bounds. This covers the closed interval $|x|\le7$. The sign at unbounded $x$ in the low positive-time region remains uncontrolled. The kernel convention is from Rodgers and Tao, \emph{The de Bruijn--Newman constant is non-negative}.

\unitstart{Appendix B}{supp}{Proofs and verification material}{Lambert determinant estimates, programs, and interval data}{}
\section{The determinant with changed dimensions}
\begingroup\setlength{\emergencystretch}{3em}
Let $F(q)=\sum_{n\ge1}(q^n-1)^{-1}$. All transfers below refer to commit
\texttt{008e3136\allowbreak{}a95e1e1e\allowbreak{}a605837c\allowbreak{}a9742b86\allowbreak{}13b2034a} of
\texttt{clambro/\allowbreak{}lambert-nonquadraticity}. The new bounds and cells are proved
in ordinary mathematics and checked in exact rational arithmetic. The
specialized Lean theorem and the upstream Lean library have not been rerun
for the new dimensions.

\paragraph{Statement.}
For coprime positive integers $a>b$ with
$\log b/\log a<0.57329026631440$, $F(a/b)$ is irrational. This includes
$20814/299$ and all of its positive integer powers. The old dimensions, with
the same six coefficientwise majorants and both old endpoint estimates held
fixed, cannot reach $20814/299$ even on integrating their full envelope.
No claim of priority over other constructions is made.

\paragraph{Archimedean size and normalization.}
Fix $h=3000$, $k=161$, $d=97$, $s=245$ and $L=h+d=3097$. Each size is
multiplied by an integer $N$. Use the source's rectangular pole Hankel
determinant $\Delta$ with rank $hN$, initial pole $kN$, pole length $LN$ and
moment shift $sN$. Its signed infinity normalizer is $U^{-g_N}$, where
\[
 g_N=(s-d)N\sum_{i=0}^{hN-1}\bigl((k+d)N+i\bigr),\qquad
 \frac{g_N}{N^3}\longrightarrow g=(s-d)h(k+d+h/2).
\]
See the generic theorem
\texttt{lambert\allowbreak{}Translated\allowbreak{}Infinity\allowbreak{}Exponent\_\allowbreak{}of\_\allowbreak{}le} in
\texttt{Refinement/\allowbreak{}Infinity.lean}. The generic \texttt{Decay.lean} theorem
\texttt{lambert\allowbreak{}Pole\allowbreak{}Determinant\_\allowbreak{}log\_\allowbreak{}bounds}, by comparing positive moment
determinants against the explicit Cauchy determinant, proves
$Q_N(q,F(q))>0$ and
\[
 \begin{aligned}
 \frac{\log Q_N(q,F(q))}{N^3}&\longrightarrow -S\log q,\qquad
 Q_N=U^{-g_N}\Delta,\\
 S&=h^3/3+sh^2/2+h(Lk+L^2/2)+g=26766016500 .\end{aligned}
\]
The normalized determinant has no coefficient pole at infinity by the
generic rearrangement bound. The partial-fraction construction uses only
$U$ and differences of its integer powers. All other poles are cyclotomic.
Monic clearing therefore produces integer coefficient polynomials with
degree in $U$ at most the clearing degree, and degree in $X$ at most $hN$.

\paragraph{Zero-place cost.}
The original proof of its zero valuation via finite-part moment blocks
applies at general $h,k,L,s$ (\texttt{Determinants/\allowbreak{}Zero\allowbreak{}Block},
\texttt{ZeroPole}, \texttt{ZeroCost}). Scale the ordered slot index and its
node by $N$, and write $r\in[0,h]$ for the scaled slot. Up to $O(1/N)$, the
scaled node satisfies $y\ge u$ with $u=r$ for $r<k$ and $u=(r+k)/2$ for
$r\ge k$. The scaled cost is linear with slope $2h+s-L-2r$ before $y=k$.
On $[k,k+L]$ it is
\[
 w(r,y)=\bigl(y^2+2(2h+s-k-L-2r)y+k^2\bigr)/2,
\]
and it is nondecreasing beyond $k+L$. Finite and limiting costs differ by
$O(N)$ per slot after restoring the $N^2$ scaling, uniformly since one may
minimize on $y\le k+L$. The quadratic is strictly convex and joins the
first branch with matching derivative. Because $2h+s-L-2k>0$, for
$0\le r\le k$ the minimum is at $y=r$. For $k\le r\le r_0$ it is at
$y=(r+k)/2$, where $r_0=(4h+2s-k-2L)/3=2045$. For $r_0\le r\le h$ it is at
the vertex $y=2r-(2h+s-k-L)$. The checker verifies the ordering and all
endpoint feasibility conditions. Hence
\[
 A=g-\int_0^h\min_{y\ge u} w(r,y)\,dr=2054711745
\]
is a sufficient asymptotic zero-clearing cost for $Q_N$. The elementary
Riemann sums converge. Rounding the finite clearing power upwards costs
$o(N^3)$. The verifier integrates each piece by the exact three-point formula
for quadratics.

\paragraph{Cyclotomic costs.}
The six generic numerator-weighted presentations use
\[
\begin{aligned}
 (A,b,d',s')\ \in\ \bigl\{&
 (0,k,d,s),\ (d,s,0,k),\ (s,k+d,s-d,0),\\
 & (0,s-d,d,k+d),\ (d,k+d,0,s-d),\ (s-d,k,s,d)\bigr\}.
\end{aligned}
\]
They follow by applying the two polynomial identities
\texttt{lambert\allowbreak{}Refinement\allowbreak{}Determinant\_\allowbreak{}transpose} and
\texttt{lambert\allowbreak{}Refinement\allowbreak{}Determinant\_\allowbreak{}prefix\_\allowbreak{}swap} in the same order as
\texttt{Refinement/\allowbreak{}Orbit.lean}. The required inequalities hold because
$0\le d\le s\le k+d$. Each state's coefficientwise bound uses
\texttt{Refinement/\allowbreak{}Cyclotomic.lean}, which is generic in all four
parameters, as are the transpose and swap transports in
\texttt{Refinement/\allowbreak{}Transport.lean}. The expression-tree generator gives
their exact homogeneous values at $x=n/N$ for each cyclotomic index $n$.
Equivalently, substitute these dimensions in the six displayed general
expressions in the accompanying code.

For $x\ge k+L$ the original state has $v_L(x)=0$, tail zero, and zero
residual. The latter vanishes since $s+h\le k+L$ and $d\ge0$. Thus the
clearing exponent is zero there. Elsewhere it is bounded by the baseline
\[
 B(x)=h+v_L(x)-v_L(2x),\qquad v_L(x)=\sum_{j\ge1}(L-jx)_+ .
\]
Let $c(x)$ be the maximum of zero and the minimum of the six signed local
upper bounds. Each selected cell $[l,u]$ certifies
$B(x)-c(x)\ge a+bx\ge0$ there. Cells meet only at endpoints. The TSV file
contains 860 such disjoint cells. The program \texttt{retuned.py}
independently checks every floor coefficient and max branch at both
endpoints. The list's lower integral $I$ is the exact sum of
$\int_l^u x(a+bx)\,dx$, evaluated rationally by the script.

Partial summation of $\sum_{n\le x}\varphi(n)\sim3x^2/\pi^2$ gives limiting
mass
\[
 (6/\pi^2)\int xB(x)\,dx=L^3/8+3h(k+L)^2/\pi^2 .
\]
Indeed $v_L(x)-v_L(2x)$ is the odd-index tent sum and
$\sum_{j\ {\rm odd}}j^{-2}=\pi^2/8$. Domination by the convergent sum of
$j^{-2}$ permits passage through its infinite series. The 860 disjoint
cell masses subtract $6I/\pi^2$. Hence the integer clearing-polynomial
degree divided by $N^3$ has limsup at most
\[
 D=A+L^3/8+\bigl(3h(k+L)^2-6I\bigr)/\pi^2 .
\]
The cyclotomic logarithms at any fixed $q>1$ differ from
$\varphi(n)\log q$ by $O_q(1)$ per index, and total exponents are
$O(N^2\log N)=o(N^3)$. With clearing polynomial $J_N$,
$\deg_U(J_NQ_N)\le\deg J_N$. Under $F(a/b)=u/v$, a nonzero integer is
obtained by multiplying $(J_NQ_N)(a/b,u/v)$ by
$b^{\deg J_N}v^{hN}$. Its cubic log limsup is at most
$D\log a-S\log(a/b)$, negative under the displayed threshold, a
contradiction. This proves the irrationality criterion. The cutoffs and
the logarithm comparison are certified by rational bounds for $\pi$ and
logarithms. The artifact contains no replacement of the source estimates
by sampled finite-$N$ inequalities.

\paragraph{Earlier six profiles.}
The previous verification, \texttt{keio\_\allowbreak{}ch\_\allowbreak{}six\_\allowbreak{}envelope.py}, established
on 762 exact affine cells that the maximum saving of the six old expressions
integrated on $[5,543]$ equals the 400 source-cell integral exactly. All old
costs vanish beyond 543. On $0<x\le5$ all source filtration terms
$v_{h+d'+A}(2x)$ exceed $h$ (already the first two summands do), so the
comparison correction reduces to a signed sum of ordinary and double-spaced
tents, truncated if necessary by the nonnegative-part rule. For
$D_{K,c}(x)=v_K(cx)$, write $t=\{K/(cx)\}$. Its exact expansion is
\[
 D_{K,c}(x)=K^2/(2cx)-K/2+(cx/2)t(1-t) .
\]
In all six signed corrections the inverse and the constant terms cancel
individually. Discarding negative contributions to the last terms shows
that the six savings are at most $0,0,5x/8,x/2,x/2,x/2$, respectively.
The program \texttt{OLD\_BOUND.py} expands and checks the coefficients.
Thus the entire unexamined integral is at most $625/24$. Together the old
full envelope has cutoff $<0.573290237682$, strictly below
$\log(299)/\log(20814)$. The claim concerns the exact published upper
estimates at the fixed old parameters. Extra divisibility not captured
there could supersede them.

\paragraph{Reproducing the archive checks.}
The accompanying README, dated 13 November 2026, describes the graph and
Lambert papers as research drafts for independent review. The graph
results use analytic proofs without a finite certificate. The Lambert
specialization uses the upstream parameter-independent estimates and the
supplemental exact checks. To reproduce those checks, use Python 3.9 or
later with assertions enabled from the root of the extracted archive.
\begin{verbatim}
git clone https://github.com/clambro/lambert-nonquadraticity upstream
git -C upstream checkout 008e3136a95e1e1ea605837ca9742b8613b2034a
python checks/retuned.py upstream
python checks/keio_ch_six_envelope.py upstream
python checks/OLD_BOUND.py upstream
\end{verbatim}
The first checker covers each used cell for the changed dimensions. The
second covers the old full envelope for $x\ge5$, and the third bounds
$0<x<5$ and combines the two parts. All three use exact rational arithmetic.
Keep \texttt{cells.tsv} beside \texttt{retuned.py}. Keep the previous
one-cell checker beside the old envelope checker, which invokes it
automatically.
\par\endgroup

\section{Reproducing the archive checks}
The two original archives share two programs. Each program appears once below. The data file is also preserved in the single source. To reconstruct the files from this appendix, save the following listing as \texttt{extract.py} beside the volume source, then run \texttt{python3 extract.py Collected\_Research\_Papers.tex}. The files are placed in \texttt{research-checks/}. The program checks the SHA-256 digest of each file. The scripts within the individual papers follow their original in-paper instructions, scoped to that paper's two source boundary comments.

\begin{Verbatim}[fontsize=\small]
from hashlib import sha256
from pathlib import Path
from collections import defaultdict
import re, sys
s = Path(sys.argv[1]).read_text(encoding="utf-8")
files = defaultdict(list)
pat = (r"(?m)^% VOLFILE ([\w.]+) ([a-f0-9]{64})\n"
       r"\\begin\{Verbatim\}[^\n]*\n([\s\S]*?)^\\end\{Verbatim\}")
for name, digest, chunk in re.findall(pat, s):
    files[name].append((digest, chunk))
assert len(files) == 11
out = Path("research-checks")
out.mkdir(exist_ok=True)
for name, pieces in files.items():
    raw = "".join(text for digest, text in pieces).encode("utf-8")
    assert sha256(raw).hexdigest() == pieces[0][0]
    (out / name).write_bytes(raw)
\end{Verbatim}
\subsection*{\texttt{baker.py}}
% VOLFILE baker.py ba7a8324cae9bce3c47b32da94b7f9515517481c017a306955ad46f5d424a02c
\begin{Verbatim}[fontsize=\footnotesize,breakanywhere=false]
from sympy import isprime, primitive_root
p=next(p for p in range(160001,170000,8) if isprime(p) and isprime((p-1)//8))
ell=(p-1)//8
U=(1,3,5,7)
B=[0]*4
G=[0]*4
# b^(ell u+1) = b*(b^ell)^u
for b in range(1,p):
 q=(pow(b,p-1,p*p)-1)//p
 a=pow(b,ell,p); pair=(b*q)%p; z=a; step=a*a%p
 for i in range(4):
  B[i]=(B[i]-pair*z)%p
  G[i]=(G[i]+pair*q*z)%p
  z=z*step%p
G=[x*pow(2,-1,p)%p for x in G]
eps=pow(int(primitive_root(p)),ell,p)
rows=[]
for j in range(8):
 xi=pow(eps,j,p)
 for a in (1,-1):
  for c in (1,-1):
   phases=[s*pow(xi,u,p)%p for s,u in zip((1,a,c,a*c),U)]
   rows.append(tuple(sum(s*v for s,v in zip(phases,Z))%p for Z in (B,G)))
assert len(rows) == 32 and all(first != 0 for first,second in rows)
print(p,ell,p*p,eps,B,G, len(rows),sum(q==0 for q,r in rows),sum(q==r==0 for q,r in rows))
\end{Verbatim}
\subsection*{\texttt{beukers.py}}
% VOLFILE beukers.py 1852e32012a7fe64448c7dfdc11cecd8b7749b40445544313bfc2167e9d179b0
\begin{Verbatim}[fontsize=\footnotesize,breakanywhere=false]
"""Finite confirmations of the two polynomial/number congruences in NOTES.
Standard library only. The proof for every odd p other than 7 is in NOTES.
"""
from math import comb, isqrt

for p in range(3, 102, 2):
    if p == 7 or any(p % d == 0 for d in range(2, isqrt(p)+1)):
        continue
    # Work with X = 1+2Y. This substitution is invertible in F_p.
    values = [1]*p
    for level in range(2,6):
        values = [sum(comb(n,j)**2 * values[j] for j in range(n+1)) % p
                  for n in range(p)]
        if level == 4:
            domb = values
    powers = [[comb(n,j)*comb(n+j,j)*domb[n]*pow(-8,-n,p) % p
               for j in range(n+1)] for n in range(p)]
    for j in range(p):
        assert sum((2*n+1)*powers[n][j] for n in range(j,p)) % p == 0
    numerics = [sum(c*pow(-8,-j,p) for j,c in enumerate(row)) % p for row in powers]
    h = (p-1)//2
    assert (sum(v*pow(-7,-n,p) for n,v in enumerate(values))
            -2*sum(numerics[:h])-numerics[h]) % p == 0
print('All primes tested through 101 passed both identities.')
\end{Verbatim}
\subsection*{\texttt{denominators.py}}
% VOLFILE denominators.py 365fdad52c6c0db0eb4714ac9525103f2c5d6171e9b2e8ce0e93372989f3d259
\begin{Verbatim}[fontsize=\footnotesize,breakanywhere=false]
"""Exact finite denominator exclusions for zeta(7) and Euler's constant.

Python 3 standard library. Euler--Maclaurin is used with the periodic
Bernoulli remainder. Its absolute value is at most the absolute size of
the last correction for the two decreasing reciprocal-power tails.
"""
from fractions import Fraction as F
from math import comb, factorial, prod

N = 512
M = 120
B = 10**150

bernoulli = [F(1)]
for m in range(1, 2*M+1):
    bernoulli.append(-sum(comb(m+1,j)*val for j,val in enumerate(bernoulli))/(m+1))

def exclude(label, lo, hi):
    midpoint = (lo+hi)/2
    best = midpoint.limit_denominator(B)
    assert lo < hi and abs(best-midpoint) > (hi-lo)/2
    print(label, 'contains no rational with positive denominator <= 10^150')

# log(512)=9*log(2), with log(2)=2*atanh(1/3).
z = F(1,3)
J = 430
log_lo = 18*sum(z**(2*j+1)/(2*j+1) for j in range(J))
log_hi = log_lo + 18*z**(2*J+1)/((2*J+1)*(1-z*z))
gamma_main = sum(F(1,j) for j in range(1,N+1))-F(1,2*N)
for k in range(1, M+1):
    term = bernoulli[2*k]/(2*k*N**(2*k))
    gamma_main += term
exclude('Euler gamma', gamma_main-log_hi-abs(term), gamma_main-log_lo+abs(term))

s = 7
value = sum(F(1,j**s) for j in range(1,N))+F(1,(s-1)*N**(s-1))+F(1,2*N**s)
for k in range(1,M+1):
    term = bernoulli[2*k]*prod(range(s,s+2*k-1))/(factorial(2*k)*N**(s+2*k-1))
    value += term
exclude('zeta(7)', value-abs(term), value+abs(term))
\end{Verbatim}
\subsection*{\texttt{gieseking.py}}
% VOLFILE gieseking.py 3875e0b8dd428890b7ebfb87feeefd56f2052b109ffefb1eb28f7be49503e214
\begin{Verbatim}[fontsize=\footnotesize,breakanywhere=false]
from fractions import Fraction as F
from math import comb,gcd,lcm,log
from itertools import combinations
def get(n):
 m=n//3; K=n
 J=[(F(0),F(1,3)),(F(-2),F(0))]; I={0:(F(0),F(0))}
 for j in range(2,2*n+1):
  J.append(tuple(F(4*(j-1),j)*x-c for x,c in zip(J[-2],(F(2,j),0))))
  if j%2==0: I[j]=tuple(F(4*(j-1),j)*a-F(4,j)*b+v/j for a,b,v in zip(I[j-2],J[j-2],J[j]))
 moments=[]
 for j in range(2*n+1):
  M=[F(0),F(0)]
  for i in range(0,2*m+1,2):
   v=I[i] if i==j else [(a-b)/(i-j) for a,b in zip(J[j],J[i])]
   for k in range(2): M[k]+=(-1)**(i//2)*comb(m,i//2)*v[k]
  moments.append(M)
 cols=[]; polys=[]
 for d in range(K+1):
  P=[0]*(2*n+1)
  for j in range(n+K-d+1): P[d+j]=(-1)**(n+j)*comb(n+K-d,j)*comb(n+d+j,n)
  polys.append(P)
  cols.append([sum(v*x[j] for v,x in zip(P,moments)) for j in range(2)])
 return polys,cols
def test(n):
 polys,C=get(n)
 def det(i,j): return C[i][0]*C[j][1]-C[j][0]*C[i][1]
 vals={}
 for i,j,k in combinations(range(n+1),3):
  v=(det(j,k),det(k,i),det(i,j))
  if all(x<0 for x in v): v=tuple(-x for x in v)
  if min(v)<=0: continue
  k0=lcm(*(v0.denominator for v0 in v)); w=[int(x*k0) for x in v]
  P=[sum(x*y for x,y in zip(w,z)) for z in zip(polys[i],polys[j],polys[k])]
  div=gcd(*P)
  cost=F(max(w),div)
  vals[(i,j,k)]=cost
 if vals:
  best=min(vals,key=vals.get); return n,len(vals),best,str(vals[best]) if n<=9 else '', round(log(vals[best])/n,5)
 return n,0
if __name__=='__main__':
 for n in range(3,46,3): print(test(n),flush=True)
\end{Verbatim}
\subsection*{\texttt{heat\_window.py}}
% VOLFILE heat_window.py d36336e5b1aee1948aa727624eda37c01c8c5fdccdee7be0b501d9737077291a
\begin{Verbatim}[fontsize=\footnotesize,breakanywhere=false]
"""Interval certification that B_(11/50)>7 in the Rodgers--Tao convention."""
from mpmath import iv

iv.dps = 28
T = iv.mpf(11)/50
size = 4096
total = iv.mpf(0)
for j in range(size):
    u = iv.mpf([j,j+1])/size
    v = iv.exp(4*u)
    terms = sum((2*iv.pi**2*n**4*iv.exp(9*u)-3*iv.pi*n**2*iv.exp(5*u))
                *iv.exp(-iv.pi*n*n*v) for n in (1,2,3))
    total += (u**2*(2-49*u**2)*iv.exp(T*u**2)*terms)/size

# For 0<=u<=1, n>=4, v^(9/4)*exp(-pi*n^2*(v-1))<=1.
# The summands of n^4 exp(-3 n^2) decrease with ratio <= 3 exp(-27).
small_tail = 49*2*20*256*iv.exp(-48)/(1-3*iv.exp(-27))
# For u>=1, u^4<=exp(4u), T*u^2<=exp(4u)/4,
# 13u<=exp(4u)/4, and sum n^4 exp(-pi*n^2*v)<=2 exp(-3v).
# Use e^(4u)>= e^4 + 4*(u-1) for the remaining exponential integral.
large_tail = 1960*iv.exp(-iv.mpf(5)*iv.exp(4)/2)/10
assert total.a > small_tail.b+large_tail.b
print('2 mu_2(11/50)-49 mu_4(11/50)>0; finite interval enclosure:',total)
\end{Verbatim}
\subsection*{\texttt{keio\_ch\_onecell\_verify.py}}
% VOLFILE keio_ch_onecell_verify.py 50a9551fa8c52fc703f728a29ed6a44b393db84ab199f17ac38b2e1bb9689a6a
\begin{Verbatim}[fontsize=\footnotesize,breakanywhere=false]
"""Exact rational verification of a disjoint additional Lambert profile cell.

Optional command argument is a checkout of clambro/lambert-nonquadraticity.
"""
from fractions import Fraction as F
from functools import cache
from pathlib import Path
import sys
source=Path(sys.argv[1] if len(sys.argv)>1 else '/tmp/keio_clambro')
sys.path.insert(0,str(source/'formalization/certificate'))
import lambert_profile as lp
from decimal import Decimal,localcontext
left,right=F(86,19),F(527,116)
mid=(left+right)/2
def aff(e):
    tag=e[0]
    if tag=='lin':return e[1:]
    if tag=='def':
        k,c=e[1:]; m=int(F(k,c)/mid)
        assert m*c*right<=k<=(m+1)*c*left
        return m*k,-F(m*(m+1)*c,2)
    if tag=='res':
        M,d=e[1:]
        assert right<=d
        return 0,0
    p,q=aff(e[1]),aff(e[2])
    if tag=='add':return p[0]+q[0],p[1]+q[1]
    if tag=='sub':return p[0]-q[0],p[1]-q[1]
    assert tag=='max'
    if all(p[0]+p[1]*t>=q[0]+q[1]*t for t in (left,right)): return p
    assert all(p[0]+p[1]*t<=q[0]+q[1]*t for t in (left,right))
    return q
expr=lp.expressions(500,27,516,41)[0][5]
assert aff(expr)==(500,-110)
assert 500-110*right>0
old_data=(source/'formalization/Lambert/Refinement/ProfileData.lean').read_text()
import re
intervals=re.findall(r'^\s+⟨(\(\d+/\d+\)|\d+), (\(\d+/\d+\)|\d+), [0-5],',old_data,re.M)
assert len(intervals)==400
assert right<min(F(l.strip('()')) for l,_ in intervals)
I=500*(right**2-left**2)/2-110*(right**3-left**3)/3
assert I==F(1444283345,16059278496)
# Bound pi via the alternating Machin identity at the n-th partial sum.
def atn(z,n=26):
    partial=sum((-1)**i*z**(2*i+1)/(2*i+1) for i in range(n))
    next_term=(-1)**n*z**(2*n+1)/(2*n+1)
    return min(partial,partial+next_term),max(partial,partial+next_term)
a,b=atn(F(1,5));c,d=atn(F(1,239))
pil,pih=16*a-4*d,16*b-4*c
S=F(371952500,3)
def theta(pi,I):
    return 1-(F(80135650,3)+(F(517036731039,2000)-6*I)/pi**2)/S
# log z = 2 arctanh((z-1)/(z+1)), for 1<=z<=2, incl geometric tail.
def lg(z,n=32):
    u=(z-1)/(z+1)
    l=2*sum(u**(2*i+1)/(2*i+1) for i in range(n))
    return l,l+2*u**(2*n+1)/((2*n+1)*(1-u**2))
log2=lg(F(2))
def ln_nat(k):
    e=k.bit_length()-1
    p,q=lg(F(k,2**e))
    return e*log2[0]+p,e*log2[1]+q
al,au=ln_nat(113590); bl,bu=ln_nat(791)
old=theta(pih,0);new=theta(pil,I)
assert old < bl/au <= bu/al < new
print('Certified old cutoff < theta(113590/791) < enlarged cutoff.')
print('Exact added integral:',I)
print('Certified added cutoff >', '0.57329011043309', 'old cutoff <', '0.57329010999213')
assert F(57329011043309,10**14)<new
assert old<F(57329010999213,10**14)
\end{Verbatim}
\subsection*{\texttt{keio\_ch\_six\_envelope.py}}
% VOLFILE keio_ch_six_envelope.py 49bd4cf9f98217fc81ea1fe05432e33ebaffff0d8bc99f9220342fb325cbbcb1
\begin{Verbatim}[fontsize=\footnotesize,breakanywhere=false]
"""Exact Riemann-profile optimality on 5 <= n/N <= 543 for the fixed six costs."""
from fractions import Fraction as Q
from pathlib import Path
from functools import cache
import re,sys
source=Path(sys.argv[1] if len(sys.argv)>1 else '/tmp/keio_clambro')
sys.path.insert(0,str(source/'formalization/certificate'))
import lambert_profile as lp
ex,_=lp.expressions(500,27,516,41)
nodes=set()
def walk(e):
    if e in nodes:return
    nodes.add(e)
    if e[0] in ('add','sub','max'):
        walk(e[1]);walk(e[2])
for e in ex:walk(e)
lo=Q(5);hi=Q(543)
cuts={lo,hi}
def include(x):
    if lo<x<hi:cuts.add(x)
def grid(k,c=1):
    for j in range(1,int(Q(k,c)/lo)+1):include(Q(k,c*j))
for e in nodes:
    if e[0]=='def':grid(*e[1:])
    if e[0]=='res':
        m,d=e[1:];include(Q(d));grid(m);grid(m-d)
# Find all changes of branch, including presentation-vs-presentation (absent in source generator)
while True:
    n=len(cuts)
    for left,right in zip(sorted(cuts),sorted(cuts)[1:]):
        mid=(left+right)/2
        fs=[lp.affine(e,mid,left,right,cuts) for e in ex]+[(Q(0),Q(0))]
        for j,(a,b) in enumerate(fs):
            for c,d in fs[:j]:
                if b!=d and left<(c-a)/(b-d)<right: cuts.add((c-a)/(b-d))
    if len(cuts)==n: break
# Certify that for every cell each leaf and each max is valid on both closed endpoints;
# for a pole/deficit vertex both adjacent formulae have the same value.
def coefficients(left,right):
    mid=(left+right)/2
    @cache
    def aff(e):
        op=e[0]
        if op=='lin':return e[1:]
        if op=='def':
            k,c=e[1:];m=int(Q(k,c)/mid)
            assert m*c*right<=k<=(m+1)*c*left
            return m*k,-Q(m*(m+1)*c,2)
        if op=='res':
            M,d=e[1:]
            if right<=d or M<=left: return Q(0),Q(0)
            m=int(Q(M)/mid)
            assert left>=d and m>0 and m*right<=M<=(m+1)*left
            if M-d>=m*right:return Q(M-m*d),Q(-1)
            assert M-d<=m*left
            return Q(-(m-1)*d),Q(m-1)
        p,q=aff(e[1]),aff(e[2])
        if op=='add':return p[0]+q[0],p[1]+q[1]
        if op=='sub':return p[0]-q[0],p[1]-q[1]
        return maxline(p,q)
    def maxline(p,q):
        if all(p[0]+p[1]*t>=q[0]+q[1]*t for t in (left,right)):return p
        assert all(q[0]+q[1]*t>=p[0]+p[1]*t for t in (left,right))
        return q
    fs=[aff(e) for e in ex]
    w=max(fs+[(Q(0),Q(0))],key=lambda f:f[0]+f[1]*mid)
    assert all(w[0]+w[1]*t>=f[0]+f[1]*t for t in (left,right) for f in fs+[(Q(0),Q(0))])
    return w
integral=0
for l,r in zip(sorted(cuts),sorted(cuts)[1:]):
    a,b=coefficients(l,r)
    integral+=a*(r**2-l**2)/2+b*(r**3-l**3)/3
old_data=(source/'formalization/Lambert/Refinement/ProfileData.lean').read_text()
rows=re.findall(r'^\s+⟨(\(\d+/\d+\)|\d+), (\(\d+/\d+\)|\d+), [0-5], \((-?\d+), (-?\d+)\), (\d+)⟩',old_data,re.M)
assert len(rows)==400
old=sum(Q(a)*(Q(v.strip('()'))**2-Q(u.strip('()'))**2)/2 +Q(b)*(Q(v.strip('()'))**3-Q(u.strip('()'))**3)/3 for u,v,a,b,_ in rows)
assert integral==old
print('Exact mass of the full best-of-six profile over [5,543] equals source 400-cell mass.',len(cuts)-1,'affine intervals verified.')
# A universal upper bound for any further savings at 0<x<5, even hypothetical zero cost.
# For x=n/N, within (0,5) the original baseline is
#   B(x)=500+v_516(x)-v_516(2x),
#   v_K(x)=sum_{j>=1}(K-j*x)_+ <= integral_0^infty (K-u*x)_+ du.
# The six clearing costs have nonnegative part; even making every such cost
# zero saves <= integral_0^5 x*(500+516^2/(2*x)) dx = 671890.
raw_budget=Q(500*5**2,2)+Q(516**2*5,2)
assert raw_budget==671890
# Beyond x=543 the original clearing exponent is already 0
# (Orbit.lean, lambertRefinementClearingExponent_eq_zero).
# Reuse only exact rational elementary enclosures from the one-cell checker.
import runpy
helpers=runpy.run_path(str(Path(__file__).with_name('keio_ch_onecell_verify.py')))
pil,pih,theta,ln_nat=[helpers[key] for key in ('pil','pih','theta','ln_nat')]
S=Q(371952500,3)
rounding=old-Q(122503422987,4000)
assert 0<rounding<Q(18,10**7)
old_lo=theta(pil,rounding)
old_hi=theta(pih,rounding)
assert Q(57329010999212,10**14)<old_lo<=old_hi<Q(57329010999214,10**14)
barrier=old_hi + 6*raw_budget/(pil**2*S)
l4,u4=ln_nat(4)
l11,u11=ln_nat(11)
assert barrier < Q(576584565,10**9) < l4/u11 <= u4/l11
print('The exact 400-cell cutoff lies strictly between 0.57329010999212 and 0.57329010999214.')
print('Fixed-six certificate with free clearing for 0<x<5: cutoff < 0.576584565 < log(4)/log(11).')
\end{Verbatim}
\subsection*{Original reproduction instructions}
The instructions below describe the original archive layout. After extraction from this volume, all programs and the interval file share one directory. Run the listed commands there, omitting the path prefix in front of each program.
% VOLFILE README.txt 5348c3498c006391de8291c81c9dbe783326283499fbd32fe0d90e43f6888e0d
\begin{Verbatim}[fontsize=\footnotesize,breakanywhere=false]
Randić radius and rational Lambert bases
13 November 2026. Research drafts for independent review.

The two PDFs and their individual, standalone TeX sources are included.
The graph results have an analytic proof and use no finite certificate.
The Lambert proof relies on parameter-independent estimates in the upstream
work and the supplemental exact arithmetic checks below. The new
specialization has not been formally verified in Lean.

To reproduce the Lambert certificates, extract this archive, then run from
its root with Python 3.9 or later and assertions enabled:

git clone https://github.com/clambro/lambert-nonquadraticity upstream
git -C upstream checkout 008e3136a95e1e1ea605837ca9742b8613b2034a
python checks/retuned.py upstream
python checks/keio_ch_six_envelope.py upstream
python checks/OLD_BOUND.py upstream

The first checker verifies every used cell for the changed dimensions.
The second verifies the full envelope of the old dimensions for x>=5.
The last bounds the old dimensions for 0<x<5 and combines the two
parts. All three use exact rational arithmetic. The cell file belongs in
the same directory as retuned.py. The previous one-cell checker belongs
beside the old envelope checker; it is invoked automatically.
\end{Verbatim}
\subsection*{\texttt{retuned.py}}
% VOLFILE retuned.py 41dad1f6442e36219b4386eb75af8015800a3a8bd753ed349098d966cde86ee2
\begin{Verbatim}[fontsize=\footnotesize,breakanywhere=false]
#!/usr/bin/env python3
"""Exact rational certificate descended from clambro/lambert-nonquadraticity @ 008e3136.
Run python retuned.py /path/to/checkout. No non-stdlib Python dependencies."""
import sys
from pathlib import Path
from fractions import Fraction as Q
from functools import cache
from math import gcd
repo=Path(sys.argv[1]) if len(sys.argv)>1 else Path('/tmp/keio_clambro')
sys.path.insert(0,str(repo/'formalization/certificate'))
import lambert_profile as lp
h,k,d,s=(3000,161,97,245)
L=h+d
# Cost in the Riemann limit of the zero-place capacity argument.
r0=Q(4*h+2*s-k-2*L,3)
assert 0<k<r0<h and s>=d and s<=k+d and 2*h+s-L-2*k>0
def cost1(r):return (2*h+s-L-2*r)*r
def cost2(r,j):
    return (j*j+2*(2*h+s-k-L-2*r)*j+k*k)/2
def costmid(r):return cost2(r,(r+k)/2)
def costlast(r):return cost2(r,2*r-(2*h+s-k-L))
assert (h+k)/2<L+k and k<=2*h-(2*h+s-k-L)<=k+L
def quadint(f,l,u):
    l,u=Q(l),Q(u)
    return (u-l)*(f(l)+4*f((l+u)/2)+f(u))/6
pzero=-(quadint(cost1,0,k)+quadint(costmid,k,r0)+quadint(costlast,r0,h))
ainf=Q(s-d)*(h*(k+d)+Q(h*h,2))
A=pzero+ainf
S=Q(h**3,3)+Q(s*h*h,2)+h*(L*k+Q(L**2,2))+ainf
# Source cyclotomic majorant integrated before the finite orbit corrections.
base_const=Q(L**3,8)
base_pi=3*h*(k+L)**2
expressions,_=lp.expressions(h,k,L,s)
# Recompute the generator's proposed formula independently on both endpoints.
def prove_affine(e,lo,hi):
    mid=(lo+hi)/2
    def maxline(p,q):
        if all(p[0]+p[1]*t>=q[0]+q[1]*t for t in (lo,hi)):return p
        assert all(q[0]+q[1]*t>=p[0]+p[1]*t for t in (lo,hi))
        return q
    @cache
    def rec(e):
        op=e[0]
        if op=='lin': return e[1:]
        if op=='def':
            K,c=e[1:];m=int(Q(K,c)/mid)
            assert m*c*hi<=K<=(m+1)*c*lo
            return m*K,-Q(m*(m+1)*c,2)
        if op=='res':
            M,d=e[1:]
            if hi<=d or M<=lo:return Q(0),Q(0)
            m=int(Q(M)/mid)
            assert lo>=d and m>0 and m*hi<=M<=(m+1)*lo
            if M-d>=m*hi:return M-m*d,-1
            assert M-d<=m*lo
            return -(m-1)*d,m-1
        p,q=rec(e[1]),rec(e[2])
        if op=='add':return p[0]+q[0],p[1]+q[1]
        if op=='sub':return p[0]-q[0],p[1]-q[1]
        assert op=='max'
        return maxline(p,q)
    return rec(e)
cell_path=Path(__file__).with_name('cells.tsv')
if cell_path.exists():
    rows=[line.strip().split('\t') for line in cell_path.read_text().splitlines()]
    cells=[(Q(lo),Q(hi),int(j),Q(a),Q(b)) for lo,hi,j,a,b in rows]
else:
    cells=[row[:5] for row in lp.generate(h,k,L,s,minimum=Q(15))]
    cell_path.write_text(''.join('\t'.join(map(str,r))+'\n' for r in cells))
I=Q(0);prev=Q(0)
for lo,hi,j,a,b in cells:
    assert 0<lo<hi<=k+L and prev<=lo
    assert prove_affine(expressions[j],lo,hi)==(a,b)
    assert all(a+b*t>=0 for t in (lo,hi))
    I+=a*(hi**2-lo**2)/2+b*(hi**3-lo**3)/3
    prev=hi
# Machin pi identity and atanh logarithm bounds with signed/rational remainders.
def atan(z,n=26):
    t=sum((-1)**i*z**(2*i+1)/(2*i+1) for i in range(n))
    u=t+(-1)**n*z**(2*n+1)/(2*n+1)
    return min(t,u),max(t,u)
al,au=atan(Q(1,5));bl,bu=atan(Q(1,239))
pil,pih=16*al-4*bu,16*au-4*bl
def theta(p):
    assert base_pi-6*I>0 and S>0
    return 1-(A+base_const+(base_pi-6*I)/p**2)/S
def lg(z,n=32):
    u=(z-1)/(z+1)
    assert 0<=u<=Q(1,3)
    t=2*sum(u**(2*i+1)/(2*i+1) for i in range(n))
    return t,t+2*u**(2*n+1)/((2*n+1)*(1-u**2))
def ln_int(z):
    e=z.bit_length()-1
    low,up=lg(Q(z,2**e));a,b=lg(Q(2))
    return e*a+low,e*b+up
def ratio(a,b):
    u,v=ln_int(a);p,q=ln_int(b)
    return p/v,q/u
print('Parametric certificate',h,k,L,s,'cells',len(cells),'A',A,'S',S)
print('Lower cutoff ~',float(theta(pil)), 'upper expression ~',float(theta(pih)))
for a,b in [(21548,305),(20814,299),(11,4)]:
    l,u=ratio(a,b)
    print(a,b,float(l),'certified' if u<theta(pil) else 'outside')
assert gcd(20814,299)==1 and gcd(21548,305)==1
assert theta(pil)>Q(57329026631440,10**14)
assert Q(573290237682,10**12)<ratio(20814,299)[0]
assert ratio(20814,299)[1]<theta(pil)
assert ratio(21548,305)[1]<theta(pil)
assert theta(pih)<ratio(11,4)[0]
\end{Verbatim}
\subsection*{\texttt{OLD\_BOUND.py}}
% VOLFILE OLD_BOUND.py 84cb1b3ec15ee023a6facb4d9466b59e231cb79d55dbd3ed48f6f500f128ff5e
\begin{Verbatim}[fontsize=\footnotesize,breakanywhere=false]
"""Old fixed six-profile upper bound at x<5 and exact source rounding outside it."""
from fractions import Fraction as F
from collections import Counter
from pathlib import Path
import sys,re
source=Path(sys.argv[1]) if len(sys.argv)>1 else Path('/tmp/keio_clambro')
sys.path.insert(0,str(source/'formalization/certificate'))
import lambert_profile as lp
ex,st=lp.expressions(500,27,516,41)
# For x=n/N <=5, v_{h+d+A}(2x)>= (500-10)+(500-20)>500.
# Thus the filtration dominates the residue corank <=h for all six expressions.
# The positive part after transport only decreases the potential saving.
def filtration(k):
    if k[0] in ('lin','def'):return k
    if k[0]=='max':
        assert k[1][0]=='def' and k[1][2]==2
        return k[1]
    return k[0],filtration(k[1]),filtration(k[2])
def tally(k,co,c):
    if k[0] in ('lin','def'):c[k]+=co
    elif k[0]=='add':tally(k[1],co,c);tally(k[2],co,c)
    else:
        assert k[0]=='sub'
        tally(k[1],co,c);tally(k[2],-co,c)
bounds=[]
for e in ex:
    c=Counter();tally(e[1],1,c);tally(filtration(e[2][1]),-1,c)
    assert all(not co or k[0]!='lin' or k[1:]==(0,0) for k,co in c.items())
    terms=[(k[1],k[2],co) for k,co in c.items() if k[0]=='def' and k[1] and co]
    assert sum(co*F(v*v,2*d) for v,d,co in terms)==0
    assert sum(co*v for v,d,co in terms)==0
    # v_V(d x)=V²/(2dx)-V/2+(dx/2) frac(V/(dx))(1-frac(V/(dx))).
    bounds.append(sum(F(co*d,8) for v,d,co in terms if co>0))
assert max(bounds)==F(5,8)
Imax=F(5,8)*F(5**3,3)
# Existing proof: keio_ch_six_envelope.py checks dominance on [5,543]; >=543 zero.
# Read the exact 400 source-cell masses, including the original rounding gain:
t=(source/'formalization/Lambert/Refinement/ProfileData.lean').read_text()
rows=re.findall(r'^\s+⟨(\(\d+/\d+\)|\d+), (\(\d+/\d+\)|\d+), [0-5], \((-?\d+), (-?\d+)\), (\d+)⟩',t,re.M)
assert len(rows)==400
old=sum(F(a)*(F(v.strip('()'))**2-F(u.strip('()'))**2)/2 +F(b)*(F(v.strip('()'))**3-F(u.strip('()'))**3)/3 for u,v,a,b,_ in rows)
S=F(371952500,3); eps=old-F(122503422987,4000)
def rateuplim(pi):
    return 1-(F(80135650,3)+(F(517036731039,2000)-6*(eps+Imax))/pi**2)/S
pi=F(314159265358979323847,10**20) # elementary strict upper bound for pi
print('Old six-profile theoretical cutoff upper bound',float(rateuplim(pi)),bounds)
assert rateuplim(pi)<F(573290237682,10**12)
\end{Verbatim}
\clearpage\section{The 860 interval records}
The two rational endpoints delimit a certificate interval. The index $i$ selects the presentation. The two integer coefficients give the affine expression $b+mx$ on that interval. The record number follows the order of the original tab-separated file. Read down the left table, then down the right table.

\noindent{\small\sffamily\bfseries Records 1--86}\par\vspace{4pt}
\noindent\begin{minipage}[t]{.49\linewidth}\centering\footnotesize\setlength{\tabcolsep}{2pt}\renewcommand{\arraystretch}{1.09}
\begin{tabular*}{\linewidth}{@{}r@{\extracolsep{\fill}}rrrrr@{}}\toprule
No. & Left & Right & $i$ & $b$ & $m$\\\midrule
1 & 3245/216 & 3097/206 & 5 & -3245 & 216\\
2 & 3097/206 & 3161/210 & 5 & -148 & 10\\
3 & 3161/210 & 3148/209 & 5 & -148 & 10\\
4 & 3148/209 & 1703/113 & 5 & 3000 & -199\\
5 & 1703/113 & 3000/199 & 5 & 3000 & -199\\
6 & 3245/214 & 258/17 & 5 & -3245 & 214\\
7 & 258/17 & 3097/204 & 5 & -3245 & 214\\
8 & 3097/204 & 3161/208 & 5 & -148 & 10\\
9 & 3161/208 & 1703/112 & 5 & -148 & 10\\
10 & 1703/112 & 3148/207 & 5 & -148 & 10\\
11 & 3148/207 & 1629/107 & 5 & 3000 & -197\\
12 & 1629/107 & 3000/197 & 2 & 3000 & -197\\
13 & 3406/223 & 1574/103 & 3 & -3406 & 223\\
14 & 1574/103 & 1131/74 & 3 & -258 & 17\\
15 & 1131/74 & 1086/71 & 3 & -258 & 17\\
16 & 1086/71 & 750/49 & 3 & 3000 & -196\\
17 & 3245/212 & 245/16 & 5 & -3245 & 212\\
18 & 245/16 & 3097/202 & 5 & -3000 & 196\\
19 & 3097/202 & 1703/111 & 5 & 97 & -6\\
20 & 1703/111 & 3161/206 & 5 & 97 & -6\\
21 & 3161/206 & 3148/205 & 5 & 97 & -6\\
22 & 3148/205 & 1629/106 & 5 & 3245 & -211\\
23 & 1629/106 & 3245/211 & 2 & 3245 & -211\\
24 & 3161/205 & 3393/220 & 3 & -3161 & 205\\
25 & 3393/220 & 787/51 & 3 & -3161 & 205\\
26 & 787/51 & 3258/211 & 3 & -13 & 1\\
27 & 3258/211 & 649/42 & 3 & 3245 & -210\\
28 & 1500/97 & 1703/110 & 5 & -3000 & 194\\
29 & 1703/110 & 3097/200 & 5 & -3000 & 194\\
30 & 3097/200 & 3161/204 & 5 & 97 & -6\\
31 & 3161/204 & 3148/203 & 5 & 97 & -6\\
32 & 3148/203 & 543/35 & 5 & 3245 & -209\\
33 & 543/35 & 295/19 & 2 & 3245 & -209\\
34 & 109/7 & 1574/101 & 3 & -3161 & 203\\
35 & 1574/101 & 3258/209 & 3 & -13 & 1\\
36 & 3258/209 & 3245/208 & 3 & 3245 & -208\\
37 & 125/8 & 3097/198 & 5 & -3000 & 192\\
38 & 3097/198 & 3161/202 & 5 & 97 & -6\\
39 & 3161/202 & 3148/201 & 5 & 97 & -6\\
40 & 3148/201 & 1629/104 & 5 & 3245 & -207\\
41 & 1629/104 & 3245/207 & 2 & 3245 & -207\\
42 & 3161/201 & 362/23 & 3 & -3161 & 201\\
43 & 362/23 & 787/50 & 3 & 97 & -6\\
\bottomrule\end{tabular*}\end{minipage}\hfill
\begin{minipage}[t]{.49\linewidth}\centering\footnotesize\setlength{\tabcolsep}{2pt}\renewcommand{\arraystretch}{1.09}
\begin{tabular*}{\linewidth}{@{}r@{\extracolsep{\fill}}rrrrr@{}}\toprule
No. & Left & Right & $i$ & $b$ & $m$\\\midrule
44 & 787/50 & 3245/206 & 3 & 3245 & -206\\
45 & 300/19 & 3097/196 & 5 & -3000 & 190\\
46 & 3097/196 & 3161/200 & 5 & 97 & -6\\
47 & 3161/200 & 1629/103 & 5 & 97 & -6\\
48 & 1629/103 & 3148/199 & 2 & 97 & -6\\
49 & 3148/199 & 649/41 & 2 & 3245 & -205\\
50 & 3161/199 & 3258/205 & 3 & -3161 & 199\\
51 & 3258/205 & 1574/99 & 3 & 97 & -6\\
52 & 1574/99 & 3245/204 & 3 & 3245 & -204\\
53 & 750/47 & 3097/194 & 5 & -3000 & 188\\
54 & 3097/194 & 3161/198 & 5 & 97 & -6\\
55 & 3161/198 & 543/34 & 5 & 97 & -6\\
56 & 543/34 & 3148/197 & 2 & 97 & -6\\
57 & 3148/197 & 3245/203 & 2 & 3245 & -203\\
58 & 3161/197 & 3097/193 & 3 & -3161 & 197\\
59 & 3097/193 & 3258/203 & 3 & -3161 & 197\\
60 & 3258/203 & 787/49 & 3 & 97 & -6\\
61 & 787/49 & 3245/202 & 3 & 3245 & -202\\
62 & 500/31 & 3097/192 & 5 & -3000 & 186\\
63 & 3097/192 & 3406/211 & 5 & 97 & -6\\
64 & 3406/211 & 3148/195 & 5 & 97 & -6\\
65 & 3148/195 & 3245/201 & 5 & 3245 & -201\\
66 & 97/6 & 1086/67 & 3 & -97 & 6\\
67 & 1086/67 & 3161/195 & 3 & 3161 & -195\\
68 & 649/40 & 1574/97 & 2 & -3245 & 200\\
69 & 1574/97 & 1629/100 & 2 & -97 & 6\\
70 & 1629/100 & 3161/194 & 5 & -97 & 6\\
71 & 3161/194 & 3406/209 & 5 & -97 & 6\\
72 & 3406/209 & 163/10 & 5 & -97 & 6\\
73 & 163/10 & 375/23 & 5 & 3000 & -184\\
74 & 3245/199 & 3148/193 & 3 & -3245 & 199\\
75 & 3148/193 & 261/16 & 3 & -97 & 6\\
76 & 261/16 & 49/3 & 3 & -97 & 6\\
77 & 49/3 & 3258/199 & 3 & 148 & -9\\
78 & 3258/199 & 131/8 & 3 & 3406 & -208\\
79 & 1000/61 & 787/48 & 2 & -3000 & 183\\
80 & 787/48 & 148/9 & 2 & 148 & -9\\
81 & 3245/197 & 3097/188 & 5 & -3245 & 197\\
82 & 3097/188 & 3148/191 & 5 & -148 & 9\\
83 & 3148/191 & 1500/91 & 5 & 3000 & -182\\
84 & 649/39 & 3097/186 & 5 & -3245 & 195\\
85 & 3097/186 & 3148/189 & 5 & -148 & 9\\
86 & 3148/189 & 50/3 & 5 & 3000 & -180\\
\bottomrule\end{tabular*}\end{minipage}\par
\clearpage
\noindent{\small\sffamily\bfseries Records 87--172}\par\vspace{4pt}
\noindent\begin{minipage}[t]{.49\linewidth}\centering\footnotesize\setlength{\tabcolsep}{2pt}\renewcommand{\arraystretch}{1.09}
\begin{tabular*}{\linewidth}{@{}r@{\extracolsep{\fill}}rrrrr@{}}\toprule
No. & Left & Right & $i$ & $b$ & $m$\\\midrule
87 & 3245/193 & 3161/188 & 5 & -3245 & 193\\
88 & 3161/188 & 3097/184 & 5 & -3245 & 193\\
89 & 3097/184 & 3148/187 & 5 & -148 & 9\\
90 & 3148/187 & 1500/89 & 5 & 3000 & -178\\
91 & 3245/191 & 3161/186 & 5 & -3245 & 191\\
92 & 3161/186 & 3148/185 & 5 & -3245 & 191\\
93 & 3148/185 & 3097/182 & 5 & -97 & 6\\
94 & 3097/182 & 1703/100 & 5 & 3000 & -176\\
95 & 1703/100 & 375/22 & 5 & 3000 & -176\\
96 & 3245/189 & 3161/184 & 5 & -3245 & 189\\
97 & 3161/184 & 86/5 & 5 & -3245 & 189\\
98 & 86/5 & 1703/99 & 5 & -3245 & 189\\
99 & 1703/99 & 3148/183 & 5 & -3245 & 189\\
100 & 3148/183 & 3097/180 & 5 & -97 & 6\\
101 & 3097/180 & 362/21 & 5 & 3000 & -174\\
102 & 362/21 & 500/29 & 2 & 3000 & -174\\
103 & 3406/197 & 1574/91 & 3 & -3406 & 197\\
104 & 1574/91 & 3097/179 & 3 & -258 & 15\\
105 & 3097/179 & 3393/196 & 3 & -258 & 15\\
106 & 3393/196 & 1629/94 & 3 & -258 & 15\\
107 & 1629/94 & 3000/173 & 3 & 3000 & -173\\
108 & 295/17 & 3161/182 & 5 & -3245 & 187\\
109 & 3161/182 & 1703/98 & 5 & -3245 & 187\\
110 & 1703/98 & 3148/181 & 5 & -3245 & 187\\
111 & 3148/181 & 3097/178 & 5 & -97 & 6\\
112 & 3097/178 & 3258/187 & 5 & 3000 & -172\\
113 & 3258/187 & 750/43 & 2 & 3000 & -172\\
114 & 262/15 & 787/45 & 3 & -3406 & 195\\
115 & 787/45 & 3393/194 & 3 & -258 & 15\\
116 & 3393/194 & 3097/177 & 3 & -258 & 15\\
117 & 3097/177 & 35/2 & 3 & -258 & 15\\
118 & 35/2 & 543/31 & 3 & -13 & 1\\
119 & 543/31 & 649/37 & 3 & 3245 & -185\\
120 & 1000/57 & 1703/97 & 5 & -3000 & 171\\
121 & 1703/97 & 3161/180 & 5 & -3000 & 171\\
122 & 3161/180 & 3148/179 & 5 & -3000 & 171\\
123 & 3148/179 & 3097/176 & 5 & 148 & -8\\
124 & 3097/176 & 3258/185 & 5 & 3245 & -184\\
125 & 3258/185 & 3245/184 & 2 & 3245 & -184\\
126 & 3161/179 & 1131/64 & 3 & -3161 & 179\\
127 & 1131/64 & 1574/89 & 3 & -3161 & 179\\
128 & 1574/89 & 3097/175 & 3 & -13 & 1\\
129 & 3097/175 & 1629/92 & 3 & -13 & 1\\
\bottomrule\end{tabular*}\end{minipage}\hfill
\begin{minipage}[t]{.49\linewidth}\centering\footnotesize\setlength{\tabcolsep}{2pt}\renewcommand{\arraystretch}{1.09}
\begin{tabular*}{\linewidth}{@{}r@{\extracolsep{\fill}}rrrrr@{}}\toprule
No. & Left & Right & $i$ & $b$ & $m$\\\midrule
130 & 1629/92 & 3245/183 & 3 & 3245 & -183\\
131 & 3000/169 & 3161/178 & 5 & -3000 & 169\\
132 & 3161/178 & 3148/177 & 5 & -3000 & 169\\
133 & 3148/177 & 3097/174 & 5 & 148 & -8\\
134 & 3097/174 & 1086/61 & 5 & 3245 & -182\\
135 & 1086/61 & 3245/182 & 2 & 3245 & -182\\
136 & 3161/177 & 787/44 & 3 & -3161 & 177\\
137 & 787/44 & 161/9 & 3 & -13 & 1\\
138 & 161/9 & 1629/91 & 3 & 148 & -8\\
139 & 1629/91 & 3097/173 & 3 & 3406 & -190\\
140 & 3097/173 & 1703/95 & 3 & 3406 & -190\\
141 & 3000/167 & 3148/175 & 2 & -3000 & 167\\
142 & 3148/175 & 18 & 2 & 148 & -8\\
143 & 18 & 3097/172 & 5 & 148 & -8\\
144 & 3097/172 & 3406/189 & 5 & 3245 & -180\\
145 & 3406/189 & 649/36 & 5 & 3245 & -180\\
146 & 1500/83 & 1574/87 & 3 & -3000 & 166\\
147 & 1574/87 & 181/10 & 3 & 148 & -8\\
148 & 181/10 & 163/9 & 3 & 3406 & -188\\
149 & 163/9 & 1703/94 & 3 & 3406 & -188\\
150 & 200/11 & 3148/173 & 2 & -3000 & 165\\
151 & 3148/173 & 3258/179 & 2 & 148 & -8\\
152 & 3258/179 & 3406/187 & 5 & 148 & -8\\
153 & 3406/187 & 3097/170 & 5 & 148 & -8\\
154 & 3097/170 & 3245/178 & 5 & 3245 & -178\\
155 & 750/41 & 787/43 & 3 & -3000 & 164\\
156 & 787/43 & 1629/89 & 3 & 148 & -8\\
157 & 1629/89 & 1703/93 & 3 & 3406 & -186\\
158 & 3000/163 & 1086/59 & 2 & -3000 & 163\\
159 & 1086/59 & 3148/171 & 5 & -3000 & 163\\
160 & 3148/171 & 3406/185 & 5 & 148 & -8\\
161 & 3406/185 & 129/7 & 5 & 148 & -8\\
162 & 129/7 & 3097/168 & 5 & 148 & -8\\
163 & 3097/168 & 295/16 & 5 & 3245 & -176\\
164 & 37/2 & 1703/92 & 5 & -148 & 8\\
165 & 1703/92 & 1629/88 & 5 & -148 & 8\\
166 & 1629/88 & 1574/85 & 2 & -148 & 8\\
167 & 1574/85 & 500/27 & 2 & 3000 & -162\\
168 & 3406/183 & 3258/175 & 3 & -3406 & 183\\
169 & 3258/175 & 3148/169 & 3 & -148 & 8\\
170 & 3148/169 & 3000/161 & 3 & 3000 & -161\\
171 & 3245/174 & 3097/166 & 5 & -3245 & 174\\
172 & 3097/166 & 3161/169 & 5 & -148 & 8\\
\bottomrule\end{tabular*}\end{minipage}\par
\clearpage
\noindent{\small\sffamily\bfseries Records 173--258}\par\vspace{4pt}
\noindent\begin{minipage}[t]{.49\linewidth}\centering\footnotesize\setlength{\tabcolsep}{2pt}\renewcommand{\arraystretch}{1.09}
\begin{tabular*}{\linewidth}{@{}r@{\extracolsep{\fill}}rrrrr@{}}\toprule
No. & Left & Right & $i$ & $b$ & $m$\\\midrule
173 & 3161/169 & 131/7 & 5 & -148 & 8\\
174 & 131/7 & 543/29 & 5 & -148 & 8\\
175 & 543/29 & 787/42 & 2 & -148 & 8\\
176 & 787/42 & 75/4 & 2 & 3000 & -160\\
177 & 3406/181 & 3258/173 & 3 & -3406 & 181\\
178 & 3258/173 & 245/13 & 3 & -148 & 8\\
179 & 245/13 & 377/20 & 3 & 97 & -5\\
180 & 377/20 & 3148/167 & 3 & 97 & -5\\
181 & 3148/167 & 3245/172 & 3 & 3245 & -172\\
182 & 1000/53 & 3097/164 & 5 & -3000 & 159\\
183 & 3097/164 & 1703/90 & 5 & 97 & -5\\
184 & 1703/90 & 3161/167 & 5 & 97 & -5\\
185 & 3161/167 & 1629/86 & 5 & 97 & -5\\
186 & 1629/86 & 1574/83 & 2 & 97 & -5\\
187 & 1574/83 & 3245/171 & 2 & 3245 & -171\\
188 & 3161/166 & 362/19 & 3 & -3161 & 166\\
189 & 362/19 & 3393/178 & 3 & 97 & -5\\
190 & 3393/178 & 3148/165 & 3 & 97 & -5\\
191 & 3148/165 & 649/34 & 3 & 3245 & -170\\
192 & 3000/157 & 3097/162 & 5 & -3000 & 157\\
193 & 3097/162 & 1703/89 & 5 & 97 & -5\\
194 & 1703/89 & 3161/165 & 5 & 97 & -5\\
195 & 3161/165 & 1629/85 & 5 & 97 & -5\\
196 & 1629/85 & 787/41 & 2 & 97 & -5\\
197 & 787/41 & 3245/169 & 2 & 3245 & -169\\
198 & 3161/164 & 3258/169 & 3 & -3161 & 164\\
199 & 3258/169 & 3393/176 & 3 & 97 & -5\\
200 & 3393/176 & 3148/163 & 3 & 97 & -5\\
201 & 3148/163 & 3245/168 & 3 & 3245 & -168\\
202 & 600/31 & 3097/160 & 5 & -3000 & 155\\
203 & 3097/160 & 3161/163 & 5 & 97 & -5\\
204 & 3161/163 & 543/28 & 5 & 97 & -5\\
205 & 543/28 & 97/5 & 2 & 97 & -5\\
206 & 3245/166 & 3148/161 & 5 & -3245 & 166\\
207 & 3148/161 & 1703/87 & 5 & -97 & 5\\
208 & 1703/87 & 3097/158 & 5 & -97 & 5\\
209 & 3097/158 & 1000/51 & 5 & 3000 & -153\\
210 & 3245/164 & 3148/159 & 5 & -3245 & 164\\
211 & 3148/159 & 1703/86 & 5 & -97 & 5\\
212 & 1703/86 & 258/13 & 5 & -97 & 5\\
213 & 258/13 & 3097/156 & 5 & -97 & 5\\
214 & 3097/156 & 1629/82 & 5 & 3000 & -151\\
215 & 1629/82 & 3000/151 & 2 & 3000 & -151\\
\bottomrule\end{tabular*}\end{minipage}\hfill
\begin{minipage}[t]{.49\linewidth}\centering\footnotesize\setlength{\tabcolsep}{2pt}\renewcommand{\arraystretch}{1.09}
\begin{tabular*}{\linewidth}{@{}r@{\extracolsep{\fill}}rrrrr@{}}\toprule
No. & Left & Right & $i$ & $b$ & $m$\\\midrule
216 & 3406/171 & 1574/79 & 3 & -3406 & 171\\
217 & 1574/79 & 3393/170 & 3 & -258 & 13\\
218 & 3393/170 & 3097/155 & 3 & -258 & 13\\
219 & 3097/155 & 3258/163 & 3 & -258 & 13\\
220 & 3258/163 & 20 & 3 & 3000 & -150\\
221 & 3245/162 & 1703/85 & 5 & -3245 & 162\\
222 & 1703/85 & 3148/157 & 5 & -3245 & 162\\
223 & 3148/157 & 3097/154 & 5 & -97 & 5\\
224 & 3097/154 & 181/9 & 5 & 3000 & -149\\
225 & 181/9 & 161/8 & 2 & 3000 & -149\\
226 & 161/8 & 3161/157 & 5 & 3000 & -149\\
227 & 3161/157 & 3000/149 & 5 & 3000 & -149\\
228 & 3245/161 & 787/39 & 3 & -3245 & 161\\
229 & 787/39 & 1131/56 & 3 & -97 & 5\\
230 & 1131/56 & 3258/161 & 3 & -97 & 5\\
231 & 3258/161 & 3097/153 & 3 & 3161 & -156\\
232 & 3097/153 & 3161/156 & 3 & 3161 & -156\\
233 & 649/32 & 3148/155 & 2 & -3245 & 160\\
234 & 3148/155 & 1629/80 & 2 & -97 & 5\\
235 & 1629/80 & 163/8 & 5 & -97 & 5\\
236 & 163/8 & 3161/155 & 5 & 3000 & -147\\
237 & 3161/155 & 3406/167 & 5 & 3000 & -147\\
238 & 3406/167 & 1000/49 & 5 & 3000 & -147\\
239 & 3245/159 & 245/12 & 3 & -3245 & 159\\
240 & 245/12 & 3393/166 & 3 & -3000 & 147\\
241 & 3393/166 & 1574/77 & 3 & -3000 & 147\\
242 & 1574/77 & 1086/53 & 3 & 148 & -7\\
243 & 1086/53 & 3097/151 & 3 & 3406 & -166\\
244 & 3097/151 & 1703/83 & 3 & 3406 & -166\\
245 & 1500/73 & 3148/153 & 2 & -3000 & 146\\
246 & 3148/153 & 1629/79 & 2 & 148 & -7\\
247 & 1629/79 & 3406/165 & 5 & 148 & -7\\
248 & 3406/165 & 3097/150 & 5 & 148 & -7\\
249 & 3097/150 & 3161/153 & 5 & 3245 & -157\\
250 & 3161/153 & 3245/157 & 5 & 3245 & -157\\
251 & 600/29 & 787/38 & 3 & -3000 & 145\\
252 & 787/38 & 3258/157 & 3 & 148 & -7\\
253 & 3258/157 & 1703/82 & 3 & 3406 & -164\\
254 & 125/6 & 3148/151 & 2 & -3000 & 144\\
255 & 3148/151 & 543/26 & 2 & 148 & -7\\
256 & 543/26 & 3406/163 & 5 & 148 & -7\\
257 & 3406/163 & 3097/148 & 5 & 148 & -7\\
258 & 3097/148 & 3161/151 & 5 & 3245 & -155\\
\bottomrule\end{tabular*}\end{minipage}\par
\clearpage
\noindent{\small\sffamily\bfseries Records 259--344}\par\vspace{4pt}
\noindent\begin{minipage}[t]{.49\linewidth}\centering\footnotesize\setlength{\tabcolsep}{2pt}\renewcommand{\arraystretch}{1.09}
\begin{tabular*}{\linewidth}{@{}r@{\extracolsep{\fill}}rrrrr@{}}\toprule
No. & Left & Right & $i$ & $b$ & $m$\\\midrule
259 & 3161/151 & 649/31 & 5 & 3245 & -155\\
260 & 3000/143 & 1574/75 & 3 & -3000 & 143\\
261 & 1574/75 & 3258/155 & 3 & 148 & -7\\
262 & 3258/155 & 1703/81 & 3 & 3406 & -162\\
263 & 1500/71 & 3148/149 & 2 & -3000 & 142\\
264 & 3148/149 & 148/7 & 2 & 148 & -7\\
265 & 3245/153 & 3097/146 & 5 & -3245 & 153\\
266 & 3097/146 & 3161/149 & 5 & -148 & 7\\
267 & 3161/149 & 787/37 & 5 & -148 & 7\\
268 & 787/37 & 1000/47 & 5 & 3000 & -141\\
269 & 3245/151 & 43/2 & 5 & -3245 & 151\\
270 & 43/2 & 3161/147 & 5 & -3245 & 151\\
271 & 3161/147 & 3097/144 & 5 & -3245 & 151\\
272 & 3097/144 & 1703/79 & 5 & -148 & 7\\
273 & 1703/79 & 1574/73 & 5 & -148 & 7\\
274 & 1574/73 & 3258/151 & 5 & 3000 & -139\\
275 & 3258/151 & 3000/139 & 2 & 3000 & -139\\
276 & 3406/157 & 3148/145 & 3 & -3406 & 157\\
277 & 3148/145 & 543/25 & 3 & -258 & 12\\
278 & 543/25 & 500/23 & 3 & 3000 & -138\\
279 & 3245/149 & 109/5 & 5 & -3245 & 149\\
280 & 109/5 & 3097/142 & 5 & -3245 & 149\\
281 & 3097/142 & 131/6 & 5 & -148 & 7\\
282 & 131/6 & 787/36 & 5 & -148 & 7\\
283 & 787/36 & 3258/149 & 5 & 3000 & -137\\
284 & 3258/149 & 3000/137 & 2 & 3000 & -137\\
285 & 3406/155 & 1629/74 & 3 & -3406 & 155\\
286 & 1629/74 & 3148/143 & 3 & -148 & 7\\
287 & 3148/143 & 3393/154 & 3 & 3000 & -136\\
288 & 3393/154 & 375/17 & 3 & 3000 & -136\\
289 & 3245/147 & 3161/143 & 5 & -3245 & 147\\
290 & 3161/143 & 1703/77 & 5 & -3245 & 147\\
291 & 1703/77 & 3097/140 & 5 & -3245 & 147\\
292 & 3097/140 & 1086/49 & 5 & -148 & 7\\
293 & 1086/49 & 1574/71 & 2 & -148 & 7\\
294 & 1574/71 & 200/9 & 2 & 3000 & -135\\
295 & 3406/153 & 245/11 & 3 & -3406 & 153\\
296 & 245/11 & 3097/139 & 3 & -3161 & 142\\
297 & 3097/139 & 1629/73 & 3 & -3161 & 142\\
298 & 1629/73 & 3393/152 & 3 & 97 & -4\\
299 & 3393/152 & 3148/141 & 3 & 97 & -4\\
300 & 3148/141 & 649/29 & 3 & 3245 & -145\\
301 & 1500/67 & 1703/76 & 5 & -3000 & 134\\
\bottomrule\end{tabular*}\end{minipage}\hfill
\begin{minipage}[t]{.49\linewidth}\centering\footnotesize\setlength{\tabcolsep}{2pt}\renewcommand{\arraystretch}{1.09}
\begin{tabular*}{\linewidth}{@{}r@{\extracolsep{\fill}}rrrrr@{}}\toprule
No. & Left & Right & $i$ & $b$ & $m$\\\midrule
302 & 1703/76 & 3161/141 & 5 & -3000 & 134\\
303 & 3161/141 & 3097/138 & 5 & -3000 & 134\\
304 & 3097/138 & 3258/145 & 5 & 97 & -4\\
305 & 3258/145 & 787/35 & 2 & 97 & -4\\
306 & 787/35 & 3245/144 & 2 & 3245 & -144\\
307 & 3161/140 & 3097/137 & 3 & -3161 & 140\\
308 & 3097/137 & 1131/50 & 3 & -3161 & 140\\
309 & 1131/50 & 181/8 & 3 & -3161 & 140\\
310 & 181/8 & 3148/139 & 3 & 97 & -4\\
311 & 3148/139 & 295/13 & 3 & 3245 & -143\\
312 & 250/11 & 3161/139 & 5 & -3000 & 132\\
313 & 3161/139 & 3097/136 & 5 & -3000 & 132\\
314 & 3097/136 & 3258/143 & 5 & 97 & -4\\
315 & 3258/143 & 1574/69 & 2 & 97 & -4\\
316 & 1574/69 & 3245/142 & 2 & 3245 & -142\\
317 & 3161/138 & 3393/148 & 3 & -3161 & 138\\
318 & 3393/148 & 3097/135 & 3 & -3161 & 138\\
319 & 3097/135 & 1629/71 & 3 & -3161 & 138\\
320 & 1629/71 & 3148/137 & 3 & 97 & -4\\
321 & 3148/137 & 23 & 3 & 3245 & -141\\
322 & 23 & 1703/74 & 3 & 3406 & -148\\
323 & 300/13 & 1086/47 & 2 & -3000 & 130\\
324 & 1086/47 & 3097/134 & 5 & -3000 & 130\\
325 & 3097/134 & 787/34 & 5 & 97 & -4\\
326 & 787/34 & 3406/147 & 5 & 3245 & -140\\
327 & 3406/147 & 649/28 & 5 & 3245 & -140\\
328 & 1000/43 & 1629/70 & 3 & -3000 & 129\\
329 & 1629/70 & 163/7 & 3 & 258 & -11\\
330 & 163/7 & 3148/135 & 3 & 258 & -11\\
331 & 3148/135 & 1703/73 & 3 & 3406 & -146\\
332 & 375/16 & 3258/139 & 2 & -3000 & 128\\
333 & 3258/139 & 258/11 & 5 & -3000 & 128\\
334 & 258/11 & 3097/132 & 5 & -3000 & 128\\
335 & 3097/132 & 3406/145 & 5 & 97 & -4\\
336 & 3406/145 & 1574/67 & 5 & 97 & -4\\
337 & 1574/67 & 3245/138 & 5 & 3245 & -138\\
338 & 500/21 & 262/11 & 5 & -3000 & 126\\
339 & 262/11 & 3097/130 & 5 & -3000 & 126\\
340 & 3097/130 & 787/33 & 5 & 97 & -4\\
341 & 787/33 & 3245/136 & 5 & 3245 & -136\\
342 & 750/31 & 3097/128 & 5 & -3000 & 124\\
343 & 3097/128 & 1574/65 & 5 & 97 & -4\\
344 & 1574/65 & 3245/134 & 5 & 3245 & -134\\
\bottomrule\end{tabular*}\end{minipage}\par
\clearpage
\noindent{\small\sffamily\bfseries Records 345--430}\par\vspace{4pt}
\noindent\begin{minipage}[t]{.49\linewidth}\centering\footnotesize\setlength{\tabcolsep}{2pt}\renewcommand{\arraystretch}{1.09}
\begin{tabular*}{\linewidth}{@{}r@{\extracolsep{\fill}}rrrrr@{}}\toprule
No. & Left & Right & $i$ & $b$ & $m$\\\midrule
345 & 97/4 & 1629/67 & 3 & -97 & 4\\
346 & 1629/67 & 3161/130 & 3 & 3161 & -130\\
347 & 3245/133 & 3148/129 & 2 & -3245 & 133\\
348 & 3148/129 & 3258/133 & 2 & -97 & 4\\
349 & 3258/133 & 49/2 & 5 & -97 & 4\\
350 & 49/2 & 3406/139 & 5 & 148 & -6\\
351 & 3406/139 & 3161/129 & 5 & 148 & -6\\
352 & 3161/129 & 3097/126 & 5 & 148 & -6\\
353 & 3097/126 & 295/12 & 5 & 3245 & -132\\
354 & 1500/61 & 787/32 & 3 & -3000 & 122\\
355 & 787/32 & 74/3 & 3 & 148 & -6\\
356 & 649/26 & 3097/124 & 5 & -3245 & 130\\
357 & 3097/124 & 1574/63 & 5 & -148 & 6\\
358 & 1574/63 & 25 & 5 & 3000 & -120\\
359 & 3245/128 & 3097/122 & 5 & -3245 & 128\\
360 & 3097/122 & 787/31 & 5 & -148 & 6\\
361 & 787/31 & 1703/67 & 5 & 3000 & -118\\
362 & 1703/67 & 1500/59 & 5 & 3000 & -118\\
363 & 3245/126 & 129/5 & 5 & -3245 & 126\\
364 & 129/5 & 1703/66 & 5 & -3245 & 126\\
365 & 1703/66 & 1574/61 & 5 & -3245 & 126\\
366 & 1574/61 & 3097/120 & 5 & -97 & 4\\
367 & 3097/120 & 181/7 & 5 & 3000 & -116\\
368 & 181/7 & 750/29 & 2 & 3000 & -116\\
369 & 26 & 3148/121 & 3 & -3406 & 131\\
370 & 3148/121 & 3097/119 & 3 & -258 & 10\\
371 & 3097/119 & 3258/125 & 3 & -258 & 10\\
372 & 3258/125 & 600/23 & 3 & 3000 & -115\\
373 & 3245/124 & 131/5 & 5 & -3245 & 124\\
374 & 131/5 & 787/30 & 5 & -3245 & 124\\
375 & 787/30 & 3097/118 & 5 & -97 & 4\\
376 & 3097/118 & 1629/62 & 5 & 3000 & -114\\
377 & 1629/62 & 500/19 & 2 & 3000 & -114\\
378 & 3406/129 & 3148/119 & 3 & -3406 & 129\\
379 & 3148/119 & 3097/117 & 3 & -258 & 10\\
380 & 3097/117 & 1086/41 & 3 & -258 & 10\\
381 & 1086/41 & 3393/128 & 3 & 3000 & -113\\
382 & 3393/128 & 3000/113 & 3 & 3000 & -113\\
383 & 3245/122 & 1703/64 & 5 & -3245 & 122\\
384 & 1703/64 & 1574/59 & 5 & -3245 & 122\\
385 & 1574/59 & 3097/116 & 5 & -97 & 4\\
386 & 3097/116 & 1629/61 & 5 & 3000 & -112\\
387 & 1629/61 & 375/14 & 2 & 3000 & -112\\
\bottomrule\end{tabular*}\end{minipage}\hfill
\begin{minipage}[t]{.49\linewidth}\centering\footnotesize\setlength{\tabcolsep}{2pt}\renewcommand{\arraystretch}{1.09}
\begin{tabular*}{\linewidth}{@{}r@{\extracolsep{\fill}}rrrrr@{}}\toprule
No. & Left & Right & $i$ & $b$ & $m$\\\midrule
388 & 3406/127 & 161/6 & 3 & -3406 & 127\\
389 & 161/6 & 3148/117 & 3 & -3245 & 121\\
390 & 3148/117 & 3258/121 & 3 & -97 & 4\\
391 & 3258/121 & 377/14 & 3 & 3161 & -117\\
392 & 377/14 & 3097/115 & 3 & 3161 & -117\\
393 & 3097/115 & 3161/117 & 3 & 3161 & -117\\
394 & 649/24 & 787/29 & 2 & -3245 & 120\\
395 & 787/29 & 543/20 & 2 & -97 & 4\\
396 & 543/20 & 163/6 & 5 & -97 & 4\\
397 & 163/6 & 245/9 & 5 & 3000 & -110\\
398 & 245/9 & 3406/125 & 5 & 3245 & -119\\
399 & 3406/125 & 109/4 & 5 & 3245 & -119\\
400 & 109/4 & 3245/119 & 5 & 3245 & -119\\
401 & 300/11 & 3393/124 & 3 & -3000 & 110\\
402 & 3393/124 & 3148/115 & 3 & -3000 & 110\\
403 & 3148/115 & 3258/119 & 3 & 148 & -5\\
404 & 3258/119 & 3097/113 & 3 & 3406 & -124\\
405 & 3097/113 & 1703/62 & 3 & 3406 & -124\\
406 & 3000/109 & 1629/59 & 2 & -3000 & 109\\
407 & 1629/59 & 1574/57 & 5 & -3000 & 109\\
408 & 1574/57 & 3097/112 & 5 & 148 & -5\\
409 & 3097/112 & 3406/123 & 5 & 3245 & -117\\
410 & 3406/123 & 3161/114 & 5 & 3245 & -117\\
411 & 3161/114 & 3245/117 & 5 & 3245 & -117\\
412 & 250/9 & 3393/122 & 3 & -3000 & 108\\
413 & 3393/122 & 362/13 & 3 & -3000 & 108\\
414 & 362/13 & 3148/113 & 3 & 258 & -9\\
415 & 3148/113 & 3097/111 & 3 & 3406 & -122\\
416 & 3097/111 & 1703/61 & 3 & 3406 & -122\\
417 & 3000/107 & 1629/58 & 2 & -3000 & 107\\
418 & 1629/58 & 787/28 & 5 & -3000 & 107\\
419 & 787/28 & 3406/121 & 5 & 148 & -5\\
420 & 3406/121 & 3097/110 & 5 & 148 & -5\\
421 & 3097/110 & 649/23 & 5 & 3245 & -115\\
422 & 1500/53 & 3258/115 & 3 & -3000 & 106\\
423 & 3258/115 & 3148/111 & 3 & 258 & -9\\
424 & 3148/111 & 1703/60 & 3 & 3406 & -120\\
425 & 200/7 & 543/19 & 2 & -3000 & 105\\
426 & 543/19 & 1574/55 & 5 & -3000 & 105\\
427 & 1574/55 & 3406/119 & 5 & 148 & -5\\
428 & 3406/119 & 86/3 & 5 & 148 & -5\\
429 & 86/3 & 3097/108 & 5 & 148 & -5\\
430 & 3097/108 & 3245/113 & 5 & 3245 & -113\\
\bottomrule\end{tabular*}\end{minipage}\par
\clearpage
\noindent{\small\sffamily\bfseries Records 431--516}\par\vspace{4pt}
\noindent\begin{minipage}[t]{.49\linewidth}\centering\footnotesize\setlength{\tabcolsep}{2pt}\renewcommand{\arraystretch}{1.09}
\begin{tabular*}{\linewidth}{@{}r@{\extracolsep{\fill}}rrrrr@{}}\toprule
No. & Left & Right & $i$ & $b$ & $m$\\\midrule
431 & 3000/103 & 787/27 & 5 & -3000 & 103\\
432 & 787/27 & 3097/106 & 5 & 148 & -5\\
433 & 3097/106 & 3245/111 & 5 & 3245 & -111\\
434 & 3406/115 & 1629/55 & 3 & -3406 & 115\\
435 & 1629/55 & 1574/53 & 3 & -148 & 5\\
436 & 1574/53 & 3000/101 & 3 & 3000 & -101\\
437 & 3245/109 & 3097/104 & 5 & -3245 & 109\\
438 & 3097/104 & 3161/106 & 5 & -148 & 5\\
439 & 3161/106 & 1703/57 & 5 & -148 & 5\\
440 & 1703/57 & 3258/109 & 5 & -148 & 5\\
441 & 3258/109 & 3148/105 & 2 & -148 & 5\\
442 & 3148/105 & 30 & 2 & 3000 & -100\\
443 & 3406/113 & 181/6 & 3 & -3406 & 113\\
444 & 181/6 & 787/26 & 3 & -148 & 5\\
445 & 787/26 & 3393/112 & 3 & 3000 & -99\\
446 & 3393/112 & 1000/33 & 3 & 3000 & -99\\
447 & 3245/107 & 3097/102 & 5 & -3245 & 107\\
448 & 3097/102 & 3161/104 & 5 & -148 & 5\\
449 & 3161/104 & 1703/56 & 5 & -148 & 5\\
450 & 1703/56 & 3258/107 & 5 & -148 & 5\\
451 & 3258/107 & 3148/103 & 2 & -148 & 5\\
452 & 3148/103 & 1500/49 & 2 & 3000 & -98\\
453 & 3161/103 & 1629/53 & 3 & -3161 & 103\\
454 & 1629/53 & 3393/110 & 3 & 97 & -3\\
455 & 3393/110 & 1574/51 & 3 & 97 & -3\\
456 & 1574/51 & 649/21 & 3 & 3245 & -105\\
457 & 3000/97 & 1703/55 & 5 & -3000 & 97\\
458 & 1703/55 & 3097/100 & 5 & -3000 & 97\\
459 & 3097/100 & 3161/102 & 5 & 97 & -3\\
460 & 3161/102 & 1086/35 & 5 & 97 & -3\\
461 & 1086/35 & 3148/101 & 2 & 97 & -3\\
462 & 3148/101 & 3245/104 & 2 & 3245 & -104\\
463 & 3161/101 & 1629/52 & 3 & -3161 & 101\\
464 & 1629/52 & 377/12 & 3 & 97 & -3\\
465 & 377/12 & 787/25 & 3 & 97 & -3\\
466 & 787/25 & 3245/103 & 3 & 3245 & -103\\
467 & 600/19 & 3097/98 & 5 & -3000 & 95\\
468 & 3097/98 & 3161/100 & 5 & 97 & -3\\
469 & 3161/100 & 3258/103 & 5 & 97 & -3\\
470 & 3258/103 & 3148/99 & 2 & 97 & -3\\
471 & 3148/99 & 3245/102 & 2 & 3245 & -102\\
472 & 3161/99 & 543/17 & 3 & -3161 & 99\\
473 & 543/17 & 3393/106 & 3 & 97 & -3\\
\bottomrule\end{tabular*}\end{minipage}\hfill
\begin{minipage}[t]{.49\linewidth}\centering\footnotesize\setlength{\tabcolsep}{2pt}\renewcommand{\arraystretch}{1.09}
\begin{tabular*}{\linewidth}{@{}r@{\extracolsep{\fill}}rrrrr@{}}\toprule
No. & Left & Right & $i$ & $b$ & $m$\\\midrule
474 & 3393/106 & 1574/49 & 3 & 97 & -3\\
475 & 1574/49 & 3245/101 & 3 & 3245 & -101\\
476 & 1000/31 & 3097/96 & 5 & -3000 & 93\\
477 & 3097/96 & 97/3 & 5 & 97 & -3\\
478 & 649/20 & 3148/97 & 3 & -3245 & 100\\
479 & 3148/97 & 1629/50 & 3 & -97 & 3\\
480 & 1629/50 & 3161/97 & 3 & 3161 & -97\\
481 & 295/9 & 787/24 & 2 & -3245 & 99\\
482 & 787/24 & 362/11 & 2 & -97 & 3\\
483 & 362/11 & 3161/96 & 5 & -97 & 3\\
484 & 3161/96 & 3097/94 & 5 & -97 & 3\\
485 & 3097/94 & 3000/91 & 5 & 3000 & -91\\
486 & 3245/98 & 3148/95 & 3 & -3245 & 98\\
487 & 3148/95 & 1629/49 & 3 & -97 & 3\\
488 & 1629/49 & 1131/34 & 3 & 3161 & -95\\
489 & 1131/34 & 3161/95 & 3 & 3161 & -95\\
490 & 3245/97 & 1574/47 & 2 & -3245 & 97\\
491 & 1574/47 & 3258/97 & 2 & -97 & 3\\
492 & 3258/97 & 3161/94 & 5 & -97 & 3\\
493 & 3161/94 & 3097/92 & 5 & -97 & 3\\
494 & 3097/92 & 3000/89 & 5 & 3000 & -89\\
495 & 3245/96 & 3148/93 & 3 & -3245 & 96\\
496 & 3148/93 & 3393/100 & 3 & -97 & 3\\
497 & 3393/100 & 543/16 & 3 & -97 & 3\\
498 & 543/16 & 3161/93 & 3 & 3161 & -93\\
499 & 649/19 & 787/23 & 2 & -3245 & 95\\
500 & 787/23 & 3258/95 & 2 & -97 & 3\\
501 & 3258/95 & 3161/92 & 5 & -97 & 3\\
502 & 3161/92 & 3406/99 & 5 & -97 & 3\\
503 & 3406/99 & 3097/90 & 5 & -97 & 3\\
504 & 3097/90 & 1000/29 & 5 & 3000 & -87\\
505 & 3245/94 & 3148/91 & 3 & -3245 & 94\\
506 & 3148/91 & 3393/98 & 3 & -97 & 3\\
507 & 3393/98 & 1629/47 & 3 & -97 & 3\\
508 & 1629/47 & 3161/91 & 3 & 3161 & -91\\
509 & 3245/93 & 1574/45 & 2 & -3245 & 93\\
510 & 1574/45 & 35 & 2 & -97 & 3\\
511 & 35 & 1086/31 & 2 & 148 & -4\\
512 & 1086/31 & 3406/97 & 5 & 148 & -4\\
513 & 3406/97 & 3161/90 & 5 & 148 & -4\\
514 & 3161/90 & 3097/88 & 5 & 148 & -4\\
515 & 3097/88 & 3245/92 & 5 & 3245 & -92\\
516 & 600/17 & 1131/32 & 3 & -3000 & 85\\
\bottomrule\end{tabular*}\end{minipage}\par
\clearpage
\noindent{\small\sffamily\bfseries Records 517--602}\par\vspace{4pt}
\noindent\begin{minipage}[t]{.49\linewidth}\centering\footnotesize\setlength{\tabcolsep}{2pt}\renewcommand{\arraystretch}{1.09}
\begin{tabular*}{\linewidth}{@{}r@{\extracolsep{\fill}}rrrrr@{}}\toprule
No. & Left & Right & $i$ & $b$ & $m$\\\midrule
517 & 1131/32 & 3148/89 & 3 & -3000 & 85\\
518 & 3148/89 & 1629/46 & 3 & 148 & -4\\
519 & 1629/46 & 1703/48 & 3 & 3406 & -96\\
520 & 250/7 & 787/22 & 2 & -3000 & 84\\
521 & 787/22 & 3258/91 & 2 & 148 & -4\\
522 & 3258/91 & 3406/95 & 5 & 148 & -4\\
523 & 3406/95 & 3161/88 & 5 & 148 & -4\\
524 & 3161/88 & 3097/86 & 5 & 148 & -4\\
525 & 3097/86 & 649/18 & 5 & 3245 & -90\\
526 & 3000/83 & 3148/87 & 3 & -3000 & 83\\
527 & 3148/87 & 181/5 & 3 & 148 & -4\\
528 & 181/5 & 1703/47 & 3 & 3406 & -94\\
529 & 1500/41 & 1574/43 & 2 & -3000 & 82\\
530 & 1574/43 & 3258/89 & 2 & 148 & -4\\
531 & 3258/89 & 3406/93 & 5 & 148 & -4\\
532 & 3406/93 & 3161/86 & 5 & 148 & -4\\
533 & 3161/86 & 258/7 & 5 & 148 & -4\\
534 & 258/7 & 3097/84 & 5 & 148 & -4\\
535 & 3097/84 & 295/8 & 5 & 3245 & -88\\
536 & 37 & 1703/46 & 5 & -148 & 4\\
537 & 1703/46 & 1629/44 & 5 & -148 & 4\\
538 & 1629/44 & 3148/85 & 2 & -148 & 4\\
539 & 3148/85 & 1000/27 & 2 & 3000 & -81\\
540 & 262/7 & 1086/29 & 3 & -3406 & 91\\
541 & 1086/29 & 787/21 & 3 & -148 & 4\\
542 & 787/21 & 75/2 & 3 & 3000 & -80\\
543 & 3245/86 & 3097/82 & 5 & -3245 & 86\\
544 & 3097/82 & 1703/45 & 5 & -148 & 4\\
545 & 1703/45 & 1629/43 & 5 & -148 & 4\\
546 & 1629/43 & 3148/83 & 2 & -148 & 4\\
547 & 3148/83 & 3000/79 & 2 & 3000 & -79\\
548 & 3406/89 & 3258/85 & 3 & -3406 & 89\\
549 & 3258/85 & 1574/41 & 3 & -148 & 4\\
550 & 1574/41 & 500/13 & 3 & 3000 & -78\\
551 & 3245/84 & 1703/44 & 5 & -3245 & 84\\
552 & 1703/44 & 3097/80 & 5 & -3245 & 84\\
553 & 3097/80 & 543/14 & 5 & -148 & 4\\
554 & 543/14 & 3148/81 & 2 & -148 & 4\\
555 & 3148/81 & 3000/77 & 2 & 3000 & -77\\
556 & 3406/87 & 3097/79 & 3 & -3406 & 87\\
557 & 3097/79 & 3258/83 & 3 & -3406 & 87\\
558 & 3258/83 & 787/20 & 3 & -148 & 4\\
559 & 787/20 & 3393/86 & 3 & 3000 & -76\\
\bottomrule\end{tabular*}\end{minipage}\hfill
\begin{minipage}[t]{.49\linewidth}\centering\footnotesize\setlength{\tabcolsep}{2pt}\renewcommand{\arraystretch}{1.09}
\begin{tabular*}{\linewidth}{@{}r@{\extracolsep{\fill}}rrrrr@{}}\toprule
No. & Left & Right & $i$ & $b$ & $m$\\\midrule
560 & 3393/86 & 750/19 & 3 & 3000 & -76\\
561 & 3245/82 & 1703/43 & 5 & -3245 & 82\\
562 & 1703/43 & 3097/78 & 5 & -3245 & 82\\
563 & 3097/78 & 1629/41 & 5 & -148 & 4\\
564 & 1629/41 & 3148/79 & 2 & -148 & 4\\
565 & 3148/79 & 40 & 2 & 3000 & -75\\
566 & 3406/85 & 3097/77 & 3 & -3406 & 85\\
567 & 3097/77 & 362/9 & 3 & -3406 & 85\\
568 & 362/9 & 161/4 & 3 & -148 & 4\\
569 & 161/4 & 1574/39 & 3 & 13 & 0\\
570 & 1574/39 & 1131/28 & 3 & 3161 & -78\\
571 & 1131/28 & 3161/78 & 3 & 3161 & -78\\
572 & 649/16 & 1629/40 & 2 & -3245 & 80\\
573 & 1629/40 & 163/4 & 5 & -3245 & 80\\
574 & 163/4 & 245/6 & 5 & -148 & 4\\
575 & 245/6 & 3148/77 & 5 & 97 & -2\\
576 & 3148/77 & 3406/83 & 5 & 3245 & -79\\
577 & 3406/83 & 3161/77 & 5 & 3245 & -79\\
578 & 3161/77 & 3245/79 & 5 & 3245 & -79\\
579 & 3000/73 & 3258/79 & 3 & -3000 & 73\\
580 & 3258/79 & 3097/75 & 3 & 258 & -6\\
581 & 3097/75 & 3393/82 & 3 & 258 & -6\\
582 & 3393/82 & 787/19 & 3 & 258 & -6\\
583 & 787/19 & 1703/41 & 3 & 3406 & -82\\
584 & 125/3 & 543/13 & 2 & -3000 & 72\\
585 & 543/13 & 3097/74 & 5 & -3000 & 72\\
586 & 3097/74 & 3148/75 & 5 & 97 & -2\\
587 & 3148/75 & 3406/81 & 5 & 3245 & -77\\
588 & 3406/81 & 295/7 & 5 & 3245 & -77\\
589 & 3000/71 & 3258/77 & 3 & -3000 & 71\\
590 & 3258/77 & 3393/80 & 3 & 258 & -6\\
591 & 3393/80 & 3097/73 & 3 & 258 & -6\\
592 & 3097/73 & 1574/37 & 3 & 258 & -6\\
593 & 1574/37 & 1703/40 & 3 & 3406 & -80\\
594 & 300/7 & 1629/38 & 2 & -3000 & 70\\
595 & 1629/38 & 43 & 5 & -3000 & 70\\
596 & 43 & 3097/72 & 5 & -3000 & 70\\
597 & 3097/72 & 3406/79 & 5 & 97 & -2\\
598 & 3406/79 & 3148/73 & 5 & 97 & -2\\
599 & 3148/73 & 649/15 & 5 & 3245 & -75\\
600 & 750/17 & 3406/77 & 5 & -3000 & 68\\
601 & 3406/77 & 3097/70 & 5 & -3000 & 68\\
602 & 3097/70 & 3148/71 & 5 & 97 & -2\\
\bottomrule\end{tabular*}\end{minipage}\par
\clearpage
\noindent{\small\sffamily\bfseries Records 603--688}\par\vspace{4pt}
\noindent\begin{minipage}[t]{.49\linewidth}\centering\footnotesize\setlength{\tabcolsep}{2pt}\renewcommand{\arraystretch}{1.09}
\begin{tabular*}{\linewidth}{@{}r@{\extracolsep{\fill}}rrrrr@{}}\toprule
No. & Left & Right & $i$ & $b$ & $m$\\\midrule
603 & 3148/71 & 3245/73 & 5 & 3245 & -73\\
604 & 500/11 & 3097/68 & 5 & -3000 & 66\\
605 & 3097/68 & 3148/69 & 5 & 97 & -2\\
606 & 3148/69 & 3245/71 & 5 & 3245 & -71\\
607 & 375/8 & 3097/66 & 5 & -3000 & 64\\
608 & 3097/66 & 3148/67 & 5 & 97 & -2\\
609 & 3148/67 & 3245/69 & 5 & 3245 & -69\\
610 & 1500/31 & 3097/64 & 5 & -3000 & 62\\
611 & 3097/64 & 3148/65 & 5 & 97 & -2\\
612 & 3148/65 & 3245/67 & 5 & 3245 & -67\\
613 & 97/2 & 3258/67 & 3 & -97 & 2\\
614 & 3258/67 & 3161/65 & 3 & 3161 & -65\\
615 & 3000/61 & 787/16 & 2 & -3000 & 61\\
616 & 787/16 & 148/3 & 2 & 148 & -3\\
617 & 649/13 & 3097/62 & 5 & -3245 & 65\\
618 & 3097/62 & 3148/63 & 5 & -148 & 3\\
619 & 3148/63 & 50 & 5 & 3000 & -60\\
620 & 3245/63 & 258/5 & 5 & -3245 & 63\\
621 & 258/5 & 1703/33 & 5 & -3245 & 63\\
622 & 1703/33 & 3148/61 & 5 & -3245 & 63\\
623 & 3148/61 & 3097/60 & 5 & -97 & 2\\
624 & 3097/60 & 362/7 & 5 & 3000 & -58\\
625 & 362/7 & 1500/29 & 2 & 3000 & -58\\
626 & 262/5 & 787/15 & 3 & -3406 & 65\\
627 & 787/15 & 3097/59 & 3 & -258 & 5\\
628 & 3097/59 & 1629/31 & 3 & -258 & 5\\
629 & 1629/31 & 1000/19 & 3 & 3000 & -57\\
630 & 3245/61 & 1703/32 & 5 & -3245 & 61\\
631 & 1703/32 & 3148/59 & 5 & -3245 & 61\\
632 & 3148/59 & 3097/58 & 5 & -97 & 2\\
633 & 3097/58 & 3258/61 & 5 & 3000 & -56\\
634 & 3258/61 & 375/7 & 2 & 3000 & -56\\
635 & 649/12 & 1574/29 & 3 & -3245 & 60\\
636 & 1574/29 & 543/10 & 3 & -97 & 2\\
637 & 543/10 & 163/3 & 3 & 3161 & -58\\
638 & 163/3 & 109/2 & 3 & 3161 & -58\\
639 & 55 & 3258/59 & 2 & -3245 & 59\\
640 & 3258/59 & 3148/57 & 5 & -3245 & 59\\
641 & 3148/57 & 3097/56 & 5 & -97 & 2\\
642 & 3097/56 & 3161/57 & 5 & 3000 & -54\\
643 & 3161/57 & 500/9 & 5 & 3000 & -54\\
644 & 3245/58 & 1629/29 & 3 & -3245 & 58\\
645 & 1629/29 & 787/14 & 3 & 13 & 0\\
\bottomrule\end{tabular*}\end{minipage}\hfill
\begin{minipage}[t]{.49\linewidth}\centering\footnotesize\setlength{\tabcolsep}{2pt}\renewcommand{\arraystretch}{1.09}
\begin{tabular*}{\linewidth}{@{}r@{\extracolsep{\fill}}rrrrr@{}}\toprule
No. & Left & Right & $i$ & $b$ & $m$\\\midrule
646 & 787/14 & 3097/55 & 3 & 3161 & -56\\
647 & 3097/55 & 3161/56 & 3 & 3161 & -56\\
648 & 3245/57 & 1086/19 & 2 & -3245 & 57\\
649 & 1086/19 & 3148/55 & 5 & -3245 & 57\\
650 & 3148/55 & 3097/54 & 5 & -97 & 2\\
651 & 3097/54 & 3161/55 & 5 & 3000 & -52\\
652 & 3161/55 & 750/13 & 5 & 3000 & -52\\
653 & 3245/56 & 1629/28 & 3 & -3245 & 56\\
654 & 1629/28 & 1574/27 & 3 & 13 & 0\\
655 & 1574/27 & 3097/53 & 3 & 3161 & -54\\
656 & 3097/53 & 117/2 & 3 & 3161 & -54\\
657 & 117/2 & 3161/54 & 3 & 3161 & -54\\
658 & 59 & 3258/55 & 2 & -3245 & 55\\
659 & 3258/55 & 3148/53 & 5 & -3245 & 55\\
660 & 3148/53 & 3097/52 & 5 & -97 & 2\\
661 & 3097/52 & 3161/53 & 5 & 3000 & -50\\
662 & 3161/53 & 3406/57 & 5 & 3000 & -50\\
663 & 3406/57 & 60 & 5 & 3000 & -50\\
664 & 3245/54 & 181/3 & 3 & -3245 & 54\\
665 & 181/3 & 787/13 & 3 & 13 & 0\\
666 & 787/13 & 3393/56 & 3 & 3161 & -52\\
667 & 3393/56 & 3097/51 & 3 & 3161 & -52\\
668 & 3097/51 & 3161/52 & 3 & 3161 & -52\\
669 & 3245/53 & 245/4 & 2 & -3245 & 53\\
670 & 245/4 & 3258/53 & 2 & -3000 & 49\\
671 & 3258/53 & 3148/51 & 5 & -3000 & 49\\
672 & 3148/51 & 3406/55 & 5 & 148 & -2\\
673 & 3406/55 & 3097/50 & 5 & 148 & -2\\
674 & 3097/50 & 3161/51 & 5 & 3245 & -52\\
675 & 3161/51 & 3245/52 & 5 & 3245 & -52\\
676 & 125/2 & 1629/26 & 3 & -3000 & 48\\
677 & 1629/26 & 377/6 & 3 & 258 & -4\\
678 & 377/6 & 1574/25 & 3 & 258 & -4\\
679 & 1574/25 & 1703/27 & 3 & 3406 & -54\\
680 & 3000/47 & 1086/17 & 2 & -3000 & 47\\
681 & 1086/17 & 3148/49 & 5 & -3000 & 47\\
682 & 3148/49 & 3406/53 & 5 & 148 & -2\\
683 & 3406/53 & 129/2 & 5 & 148 & -2\\
684 & 129/2 & 3161/49 & 5 & 148 & -2\\
685 & 3161/49 & 3097/48 & 5 & 148 & -2\\
686 & 3097/48 & 649/10 & 5 & 3245 & -50\\
687 & 200/3 & 3406/51 & 5 & -3000 & 45\\
688 & 3406/51 & 3148/47 & 5 & -3000 & 45\\
\bottomrule\end{tabular*}\end{minipage}\par
\clearpage
\noindent{\small\sffamily\bfseries Records 689--774}\par\vspace{4pt}
\noindent\begin{minipage}[t]{.49\linewidth}\centering\footnotesize\setlength{\tabcolsep}{2pt}\renewcommand{\arraystretch}{1.09}
\begin{tabular*}{\linewidth}{@{}r@{\extracolsep{\fill}}rrrrr@{}}\toprule
No. & Left & Right & $i$ & $b$ & $m$\\\midrule
689 & 3148/47 & 3161/47 & 5 & 148 & -2\\
690 & 3161/47 & 3097/46 & 5 & 148 & -2\\
691 & 3097/46 & 3245/48 & 5 & 3245 & -48\\
692 & 3000/43 & 3148/45 & 5 & -3000 & 43\\
693 & 3148/45 & 3161/45 & 5 & 148 & -2\\
694 & 3161/45 & 3097/44 & 5 & 148 & -2\\
695 & 3097/44 & 3245/46 & 5 & 3245 & -46\\
696 & 3000/41 & 3148/43 & 5 & -3000 & 41\\
697 & 3148/43 & 3161/43 & 5 & 148 & -2\\
698 & 3161/43 & 3097/42 & 5 & 148 & -2\\
699 & 3097/42 & 295/4 & 5 & 3245 & -44\\
700 & 74 & 1703/23 & 5 & -148 & 2\\
701 & 1703/23 & 1629/22 & 5 & -148 & 2\\
702 & 1629/22 & 1574/21 & 2 & -148 & 2\\
703 & 1574/21 & 75 & 2 & 3000 & -40\\
704 & 3406/45 & 3258/43 & 3 & -3406 & 45\\
705 & 3258/43 & 3148/41 & 3 & -148 & 2\\
706 & 3148/41 & 1000/13 & 3 & 3000 & -39\\
707 & 3245/42 & 1703/22 & 5 & -3245 & 42\\
708 & 1703/22 & 3097/40 & 5 & -3245 & 42\\
709 & 3097/40 & 543/7 & 5 & -148 & 2\\
710 & 543/7 & 787/10 & 2 & -148 & 2\\
711 & 787/10 & 1500/19 & 2 & 3000 & -38\\
712 & 3406/43 & 3097/39 & 3 & -3406 & 43\\
713 & 3097/39 & 3258/41 & 3 & -3406 & 43\\
714 & 3258/41 & 161/2 & 3 & -148 & 2\\
715 & 161/2 & 3148/39 & 3 & 13 & 0\\
716 & 3148/39 & 1131/14 & 3 & 3161 & -39\\
717 & 1131/14 & 3161/39 & 3 & 3161 & -39\\
718 & 649/8 & 1629/20 & 2 & -3245 & 40\\
719 & 1629/20 & 163/2 & 5 & -3245 & 40\\
720 & 163/2 & 245/3 & 5 & -148 & 2\\
721 & 245/3 & 1574/19 & 5 & 97 & -1\\
722 & 1574/19 & 3406/41 & 5 & 3245 & -39\\
723 & 3406/41 & 3161/38 & 5 & 3245 & -39\\
724 & 3161/38 & 3245/39 & 5 & 3245 & -39\\
725 & 250/3 & 1086/13 & 3 & -3000 & 36\\
726 & 1086/13 & 3097/37 & 3 & 258 & -3\\
727 & 3097/37 & 3393/40 & 3 & 258 & -3\\
728 & 3393/40 & 3148/37 & 3 & 258 & -3\\
729 & 3148/37 & 1703/20 & 3 & 3406 & -40\\
730 & 600/7 & 1629/19 & 2 & -3000 & 35\\
731 & 1629/19 & 86 & 5 & -3000 & 35\\
\bottomrule\end{tabular*}\end{minipage}\hfill
\begin{minipage}[t]{.49\linewidth}\centering\footnotesize\setlength{\tabcolsep}{2pt}\renewcommand{\arraystretch}{1.09}
\begin{tabular*}{\linewidth}{@{}r@{\extracolsep{\fill}}rrrrr@{}}\toprule
No. & Left & Right & $i$ & $b$ & $m$\\\midrule
732 & 86 & 3097/36 & 5 & -3000 & 35\\
733 & 3097/36 & 262/3 & 5 & 97 & -1\\
734 & 262/3 & 787/9 & 5 & 97 & -1\\
735 & 787/9 & 3245/37 & 5 & 3245 & -37\\
736 & 1000/11 & 3097/34 & 5 & -3000 & 33\\
737 & 3097/34 & 3406/37 & 5 & 97 & -1\\
738 & 3406/37 & 1574/17 & 5 & 97 & -1\\
739 & 1574/17 & 649/7 & 5 & 3245 & -35\\
740 & 3000/31 & 3097/32 & 5 & -3000 & 31\\
741 & 3097/32 & 97 & 5 & 97 & -1\\
742 & 295/3 & 787/8 & 3 & -3245 & 33\\
743 & 787/8 & 1086/11 & 3 & -97 & 1\\
744 & 1086/11 & 3161/32 & 3 & 3161 & -32\\
745 & 3245/32 & 3148/31 & 2 & -3245 & 32\\
746 & 3148/31 & 1629/16 & 2 & -97 & 1\\
747 & 1629/16 & 3161/31 & 5 & -97 & 1\\
748 & 3161/31 & 3406/33 & 5 & -97 & 1\\
749 & 3406/33 & 3097/30 & 5 & -97 & 1\\
750 & 3097/30 & 3000/29 & 5 & 3000 & -29\\
751 & 3245/31 & 1574/15 & 3 & -3245 & 31\\
752 & 1574/15 & 3258/31 & 3 & -97 & 1\\
753 & 3258/31 & 3161/30 & 3 & 3161 & -30\\
754 & 649/6 & 3148/29 & 2 & -3245 & 30\\
755 & 3148/29 & 543/5 & 2 & -97 & 1\\
756 & 543/5 & 109 & 5 & -97 & 1\\
757 & 109 & 3406/31 & 5 & -97 & 1\\
758 & 3406/31 & 3097/28 & 5 & -97 & 1\\
759 & 3097/28 & 1000/9 & 5 & 3000 & -27\\
760 & 3245/29 & 3258/29 & 3 & -3245 & 29\\
761 & 3258/29 & 787/7 & 3 & 13 & 0\\
762 & 787/7 & 3161/28 & 3 & 3161 & -28\\
763 & 3245/28 & 1629/14 & 2 & -3245 & 28\\
764 & 1629/14 & 3148/27 & 5 & -3245 & 28\\
765 & 3148/27 & 3161/27 & 5 & -97 & 1\\
766 & 3161/27 & 3406/29 & 5 & -97 & 1\\
767 & 3406/29 & 3097/26 & 5 & -97 & 1\\
768 & 3097/26 & 120 & 5 & 3000 & -25\\
769 & 3245/27 & 362/3 & 3 & -3245 & 27\\
770 & 362/3 & 1574/13 & 3 & 13 & 0\\
771 & 1574/13 & 3393/28 & 3 & 3161 & -26\\
772 & 3393/28 & 3161/26 & 3 & 3161 & -26\\
773 & 125 & 1629/13 & 2 & -3000 & 24\\
774 & 1629/13 & 3148/25 & 5 & -3000 & 24\\
\bottomrule\end{tabular*}\end{minipage}\par
\clearpage
\noindent{\small\sffamily\bfseries Records 775--860}\par\vspace{4pt}
\noindent\begin{minipage}[t]{.49\linewidth}\centering\footnotesize\setlength{\tabcolsep}{2pt}\renewcommand{\arraystretch}{1.09}
\begin{tabular*}{\linewidth}{@{}r@{\extracolsep{\fill}}rrrrr@{}}\toprule
No. & Left & Right & $i$ & $b$ & $m$\\\midrule
775 & 3148/25 & 3406/27 & 5 & 148 & -1\\
776 & 3406/27 & 3161/25 & 5 & 148 & -1\\
777 & 3161/25 & 129 & 5 & 148 & -1\\
778 & 129 & 3097/24 & 5 & 148 & -1\\
779 & 3097/24 & 649/5 & 5 & 3245 & -25\\
780 & 1500/11 & 3148/23 & 5 & -3000 & 22\\
781 & 3148/23 & 3161/23 & 5 & 148 & -1\\
782 & 3161/23 & 3097/22 & 5 & 148 & -1\\
783 & 3097/22 & 3245/23 & 5 & 3245 & -23\\
784 & 3406/23 & 1629/11 & 3 & -3406 & 23\\
785 & 1629/11 & 3148/21 & 3 & -148 & 1\\
786 & 3148/21 & 150 & 3 & 3000 & -20\\
787 & 3245/21 & 1703/11 & 5 & -3245 & 21\\
788 & 1703/11 & 3097/20 & 5 & -3245 & 21\\
789 & 3097/20 & 1086/7 & 5 & -148 & 1\\
790 & 1086/7 & 787/5 & 2 & -148 & 1\\
791 & 787/5 & 3000/19 & 2 & 3000 & -19\\
792 & 649/4 & 1629/10 & 3 & -3245 & 20\\
793 & 1629/10 & 163 & 3 & 13 & 0\\
794 & 163 & 3148/19 & 3 & 13 & 0\\
795 & 3148/19 & 3161/19 & 3 & 3161 & -19\\
796 & 3245/19 & 3258/19 & 2 & -3245 & 19\\
797 & 3258/19 & 3097/18 & 5 & -3245 & 19\\
798 & 3097/18 & 1574/9 & 5 & -148 & 1\\
799 & 1574/9 & 3161/18 & 5 & 3000 & -17\\
800 & 3161/18 & 3000/17 & 5 & 3000 & -17\\
801 & 3245/18 & 181 & 3 & -3245 & 18\\
802 & 181 & 3097/17 & 3 & 13 & 0\\
803 & 3097/17 & 3148/17 & 3 & 13 & 0\\
804 & 3148/17 & 3161/17 & 3 & 3161 & -17\\
805 & 3245/17 & 3258/17 & 2 & -3245 & 17\\
806 & 3258/17 & 3097/16 & 5 & -3245 & 17\\
807 & 3097/16 & 787/4 & 5 & -148 & 1\\
808 & 787/4 & 3161/16 & 5 & 3000 & -15\\
809 & 3161/16 & 200 & 5 & 3000 & -15\\
810 & 3245/16 & 1629/8 & 3 & -3245 & 16\\
811 & 1629/8 & 3097/15 & 3 & 13 & 0\\
812 & 3097/15 & 3148/15 & 3 & 13 & 0\\
813 & 3148/15 & 3161/15 & 3 & 3161 & -15\\
814 & 649/3 & 1086/5 & 2 & -3245 & 15\\
815 & 1086/5 & 3097/14 & 5 & -3245 & 15\\
816 & 3097/14 & 1574/7 & 5 & -148 & 1\\
817 & 1574/7 & 3161/14 & 5 & 3000 & -13\\
\bottomrule\end{tabular*}\end{minipage}\hfill
\begin{minipage}[t]{.49\linewidth}\centering\footnotesize\setlength{\tabcolsep}{2pt}\renewcommand{\arraystretch}{1.09}
\begin{tabular*}{\linewidth}{@{}r@{\extracolsep{\fill}}rrrrr@{}}\toprule
No. & Left & Right & $i$ & $b$ & $m$\\\midrule
818 & 3161/14 & 3406/15 & 5 & 3000 & -13\\
819 & 3406/15 & 3000/13 & 5 & 3000 & -13\\
820 & 3245/14 & 1629/7 & 3 & -3245 & 14\\
821 & 1629/7 & 3097/13 & 3 & 13 & 0\\
822 & 3097/13 & 3148/13 & 3 & 13 & 0\\
823 & 3148/13 & 3393/14 & 3 & 3161 & -13\\
824 & 3393/14 & 3161/13 & 3 & 3161 & -13\\
825 & 250 & 3258/13 & 2 & -3000 & 12\\
826 & 3258/13 & 258 & 5 & -3000 & 12\\
827 & 258 & 3097/12 & 5 & -3000 & 12\\
828 & 3097/12 & 262 & 5 & 97 & 0\\
829 & 262 & 787/3 & 5 & 97 & 0\\
830 & 787/3 & 3161/12 & 5 & 3245 & -12\\
831 & 3161/12 & 3245/12 & 5 & 3245 & -12\\
832 & 300 & 3406/11 & 5 & -3000 & 10\\
833 & 3406/11 & 3097/10 & 5 & -3000 & 10\\
834 & 3097/10 & 1574/5 & 5 & 97 & 0\\
835 & 1574/5 & 3161/10 & 5 & 3245 & -10\\
836 & 3161/10 & 649/2 & 5 & 3245 & -10\\
837 & 375 & 3406/9 & 5 & -3000 & 8\\
838 & 3406/9 & 3097/8 & 5 & -3000 & 8\\
839 & 3097/8 & 787/2 & 5 & 97 & 0\\
840 & 787/2 & 3161/8 & 5 & 3245 & -8\\
841 & 3161/8 & 3245/8 & 5 & 3245 & -8\\
842 & 500 & 3097/6 & 5 & -3000 & 6\\
843 & 3097/6 & 1574/3 & 5 & 97 & 0\\
844 & 1574/3 & 3161/6 & 5 & 3245 & -6\\
845 & 3161/6 & 3245/6 & 5 & 3245 & -6\\
846 & 750 & 3097/4 & 5 & -3000 & 4\\
847 & 3097/4 & 787 & 5 & 97 & 0\\
848 & 787 & 3161/4 & 5 & 3245 & -4\\
849 & 3161/4 & 3245/4 & 5 & 3245 & -4\\
850 & 1500 & 3097/2 & 5 & -3000 & 2\\
851 & 3097/2 & 1574 & 5 & 97 & 0\\
852 & 1574 & 3161/2 & 5 & 3245 & -2\\
853 & 3161/2 & 3245/2 & 5 & 3245 & -2\\
854 & 1703 & 3503/2 & 0 & -3406 & 2\\
855 & 3503/2 & 3000 & 0 & -3406 & 2\\
856 & 3000 & 3097 & 0 & -3406 & 2\\
857 & 3097 & 3148 & 0 & -3406 & 2\\
858 & 3148 & 3161 & 0 & -258 & 1\\
859 & 3161 & 3245 & 0 & -258 & 1\\
860 & 3245 & 3258 & 0 & -258 & 1\\
\bottomrule\end{tabular*}\end{minipage}\par
\iffalse
% Inline original tab-separated bytes for the appendix extraction script.
% VOLFILE cells.tsv a4848ccae106f2d3e53734a44697e8e599479b4b10ee420e4f723abebb08d332
\begin{Verbatim}[fontsize=\scriptsize]
3245/216	3097/206	5	-3245	216
3097/206	3161/210	5	-148	10
3161/210	3148/209	5	-148	10
3148/209	1703/113	5	3000	-199
1703/113	3000/199	5	3000	-199
3245/214	258/17	5	-3245	214
258/17	3097/204	5	-3245	214
3097/204	3161/208	5	-148	10
3161/208	1703/112	5	-148	10
1703/112	3148/207	5	-148	10
3148/207	1629/107	5	3000	-197
1629/107	3000/197	2	3000	-197
3406/223	1574/103	3	-3406	223
1574/103	1131/74	3	-258	17
1131/74	1086/71	3	-258	17
1086/71	750/49	3	3000	-196
3245/212	245/16	5	-3245	212
245/16	3097/202	5	-3000	196
3097/202	1703/111	5	97	-6
1703/111	3161/206	5	97	-6
3161/206	3148/205	5	97	-6
3148/205	1629/106	5	3245	-211
1629/106	3245/211	2	3245	-211
3161/205	3393/220	3	-3161	205
3393/220	787/51	3	-3161	205
787/51	3258/211	3	-13	1
3258/211	649/42	3	3245	-210
1500/97	1703/110	5	-3000	194
1703/110	3097/200	5	-3000	194
3097/200	3161/204	5	97	-6
3161/204	3148/203	5	97	-6
3148/203	543/35	5	3245	-209
543/35	295/19	2	3245	-209
109/7	1574/101	3	-3161	203
1574/101	3258/209	3	-13	1
3258/209	3245/208	3	3245	-208
125/8	3097/198	5	-3000	192
3097/198	3161/202	5	97	-6
3161/202	3148/201	5	97	-6
3148/201	1629/104	5	3245	-207
1629/104	3245/207	2	3245	-207
3161/201	362/23	3	-3161	201
362/23	787/50	3	97	-6
787/50	3245/206	3	3245	-206
300/19	3097/196	5	-3000	190
3097/196	3161/200	5	97	-6
3161/200	1629/103	5	97	-6
1629/103	3148/199	2	97	-6
3148/199	649/41	2	3245	-205
3161/199	3258/205	3	-3161	199
3258/205	1574/99	3	97	-6
1574/99	3245/204	3	3245	-204
750/47	3097/194	5	-3000	188
3097/194	3161/198	5	97	-6
3161/198	543/34	5	97	-6
543/34	3148/197	2	97	-6
3148/197	3245/203	2	3245	-203
3161/197	3097/193	3	-3161	197
3097/193	3258/203	3	-3161	197
3258/203	787/49	3	97	-6
787/49	3245/202	3	3245	-202
500/31	3097/192	5	-3000	186
3097/192	3406/211	5	97	-6
3406/211	3148/195	5	97	-6
3148/195	3245/201	5	3245	-201
97/6	1086/67	3	-97	6
1086/67	3161/195	3	3161	-195
649/40	1574/97	2	-3245	200
1574/97	1629/100	2	-97	6
1629/100	3161/194	5	-97	6
3161/194	3406/209	5	-97	6
3406/209	163/10	5	-97	6
163/10	375/23	5	3000	-184
3245/199	3148/193	3	-3245	199
3148/193	261/16	3	-97	6
261/16	49/3	3	-97	6
49/3	3258/199	3	148	-9
3258/199	131/8	3	3406	-208
1000/61	787/48	2	-3000	183
787/48	148/9	2	148	-9
3245/197	3097/188	5	-3245	197
3097/188	3148/191	5	-148	9
3148/191	1500/91	5	3000	-182
649/39	3097/186	5	-3245	195
3097/186	3148/189	5	-148	9
3148/189	50/3	5	3000	-180
3245/193	3161/188	5	-3245	193
3161/188	3097/184	5	-3245	193
3097/184	3148/187	5	-148	9
3148/187	1500/89	5	3000	-178
3245/191	3161/186	5	-3245	191
3161/186	3148/185	5	-3245	191
3148/185	3097/182	5	-97	6
3097/182	1703/100	5	3000	-176
1703/100	375/22	5	3000	-176
3245/189	3161/184	5	-3245	189
3161/184	86/5	5	-3245	189
86/5	1703/99	5	-3245	189
1703/99	3148/183	5	-3245	189
3148/183	3097/180	5	-97	6
3097/180	362/21	5	3000	-174
362/21	500/29	2	3000	-174
3406/197	1574/91	3	-3406	197
1574/91	3097/179	3	-258	15
3097/179	3393/196	3	-258	15
3393/196	1629/94	3	-258	15
1629/94	3000/173	3	3000	-173
295/17	3161/182	5	-3245	187
3161/182	1703/98	5	-3245	187
1703/98	3148/181	5	-3245	187
3148/181	3097/178	5	-97	6
3097/178	3258/187	5	3000	-172
3258/187	750/43	2	3000	-172
262/15	787/45	3	-3406	195
787/45	3393/194	3	-258	15
3393/194	3097/177	3	-258	15
3097/177	35/2	3	-258	15
35/2	543/31	3	-13	1
543/31	649/37	3	3245	-185
1000/57	1703/97	5	-3000	171
1703/97	3161/180	5	-3000	171
3161/180	3148/179	5	-3000	171
3148/179	3097/176	5	148	-8
3097/176	3258/185	5	3245	-184
3258/185	3245/184	2	3245	-184
3161/179	1131/64	3	-3161	179
1131/64	1574/89	3	-3161	179
1574/89	3097/175	3	-13	1
3097/175	1629/92	3	-13	1
1629/92	3245/183	3	3245	-183
3000/169	3161/178	5	-3000	169
3161/178	3148/177	5	-3000	169
3148/177	3097/174	5	148	-8
3097/174	1086/61	5	3245	-182
1086/61	3245/182	2	3245	-182
3161/177	787/44	3	-3161	177
787/44	161/9	3	-13	1
161/9	1629/91	3	148	-8
1629/91	3097/173	3	3406	-190
3097/173	1703/95	3	3406	-190
3000/167	3148/175	2	-3000	167
3148/175	18	2	148	-8
18	3097/172	5	148	-8
3097/172	3406/189	5	3245	-180
3406/189	649/36	5	3245	-180
1500/83	1574/87	3	-3000	166
1574/87	181/10	3	148	-8
181/10	163/9	3	3406	-188
163/9	1703/94	3	3406	-188
200/11	3148/173	2	-3000	165
3148/173	3258/179	2	148	-8
3258/179	3406/187	5	148	-8
3406/187	3097/170	5	148	-8
3097/170	3245/178	5	3245	-178
750/41	787/43	3	-3000	164
787/43	1629/89	3	148	-8
1629/89	1703/93	3	3406	-186
3000/163	1086/59	2	-3000	163
1086/59	3148/171	5	-3000	163
3148/171	3406/185	5	148	-8
3406/185	129/7	5	148	-8
129/7	3097/168	5	148	-8
3097/168	295/16	5	3245	-176
37/2	1703/92	5	-148	8
1703/92	1629/88	5	-148	8
1629/88	1574/85	2	-148	8
1574/85	500/27	2	3000	-162
3406/183	3258/175	3	-3406	183
3258/175	3148/169	3	-148	8
3148/169	3000/161	3	3000	-161
3245/174	3097/166	5	-3245	174
3097/166	3161/169	5	-148	8
3161/169	131/7	5	-148	8
131/7	543/29	5	-148	8
543/29	787/42	2	-148	8
787/42	75/4	2	3000	-160
3406/181	3258/173	3	-3406	181
3258/173	245/13	3	-148	8
245/13	377/20	3	97	-5
377/20	3148/167	3	97	-5
3148/167	3245/172	3	3245	-172
1000/53	3097/164	5	-3000	159
3097/164	1703/90	5	97	-5
1703/90	3161/167	5	97	-5
3161/167	1629/86	5	97	-5
1629/86	1574/83	2	97	-5
1574/83	3245/171	2	3245	-171
3161/166	362/19	3	-3161	166
362/19	3393/178	3	97	-5
3393/178	3148/165	3	97	-5
3148/165	649/34	3	3245	-170
3000/157	3097/162	5	-3000	157
3097/162	1703/89	5	97	-5
1703/89	3161/165	5	97	-5
3161/165	1629/85	5	97	-5
1629/85	787/41	2	97	-5
787/41	3245/169	2	3245	-169
3161/164	3258/169	3	-3161	164
3258/169	3393/176	3	97	-5
3393/176	3148/163	3	97	-5
3148/163	3245/168	3	3245	-168
600/31	3097/160	5	-3000	155
3097/160	3161/163	5	97	-5
3161/163	543/28	5	97	-5
543/28	97/5	2	97	-5
3245/166	3148/161	5	-3245	166
3148/161	1703/87	5	-97	5
1703/87	3097/158	5	-97	5
3097/158	1000/51	5	3000	-153
3245/164	3148/159	5	-3245	164
3148/159	1703/86	5	-97	5
1703/86	258/13	5	-97	5
258/13	3097/156	5	-97	5
3097/156	1629/82	5	3000	-151
1629/82	3000/151	2	3000	-151
3406/171	1574/79	3	-3406	171
1574/79	3393/170	3	-258	13
3393/170	3097/155	3	-258	13
3097/155	3258/163	3	-258	13
3258/163	20	3	3000	-150
3245/162	1703/85	5	-3245	162
1703/85	3148/157	5	-3245	162
3148/157	3097/154	5	-97	5
3097/154	181/9	5	3000	-149
181/9	161/8	2	3000	-149
161/8	3161/157	5	3000	-149
3161/157	3000/149	5	3000	-149
3245/161	787/39	3	-3245	161
787/39	1131/56	3	-97	5
1131/56	3258/161	3	-97	5
3258/161	3097/153	3	3161	-156
3097/153	3161/156	3	3161	-156
649/32	3148/155	2	-3245	160
3148/155	1629/80	2	-97	5
1629/80	163/8	5	-97	5
163/8	3161/155	5	3000	-147
3161/155	3406/167	5	3000	-147
3406/167	1000/49	5	3000	-147
3245/159	245/12	3	-3245	159
245/12	3393/166	3	-3000	147
3393/166	1574/77	3	-3000	147
1574/77	1086/53	3	148	-7
1086/53	3097/151	3	3406	-166
3097/151	1703/83	3	3406	-166
1500/73	3148/153	2	-3000	146
3148/153	1629/79	2	148	-7
1629/79	3406/165	5	148	-7
3406/165	3097/150	5	148	-7
3097/150	3161/153	5	3245	-157
3161/153	3245/157	5	3245	-157
600/29	787/38	3	-3000	145
787/38	3258/157	3	148	-7
3258/157	1703/82	3	3406	-164
125/6	3148/151	2	-3000	144
3148/151	543/26	2	148	-7
543/26	3406/163	5	148	-7
3406/163	3097/148	5	148	-7
3097/148	3161/151	5	3245	-155
3161/151	649/31	5	3245	-155
3000/143	1574/75	3	-3000	143
1574/75	3258/155	3	148	-7
3258/155	1703/81	3	3406	-162
1500/71	3148/149	2	-3000	142
3148/149	148/7	2	148	-7
3245/153	3097/146	5	-3245	153
3097/146	3161/149	5	-148	7
3161/149	787/37	5	-148	7
787/37	1000/47	5	3000	-141
3245/151	43/2	5	-3245	151
43/2	3161/147	5	-3245	151
3161/147	3097/144	5	-3245	151
3097/144	1703/79	5	-148	7
1703/79	1574/73	5	-148	7
1574/73	3258/151	5	3000	-139
3258/151	3000/139	2	3000	-139
3406/157	3148/145	3	-3406	157
3148/145	543/25	3	-258	12
543/25	500/23	3	3000	-138
3245/149	109/5	5	-3245	149
109/5	3097/142	5	-3245	149
3097/142	131/6	5	-148	7
131/6	787/36	5	-148	7
787/36	3258/149	5	3000	-137
3258/149	3000/137	2	3000	-137
3406/155	1629/74	3	-3406	155
1629/74	3148/143	3	-148	7
3148/143	3393/154	3	3000	-136
3393/154	375/17	3	3000	-136
3245/147	3161/143	5	-3245	147
3161/143	1703/77	5	-3245	147
1703/77	3097/140	5	-3245	147
3097/140	1086/49	5	-148	7
1086/49	1574/71	2	-148	7
1574/71	200/9	2	3000	-135
3406/153	245/11	3	-3406	153
245/11	3097/139	3	-3161	142
3097/139	1629/73	3	-3161	142
1629/73	3393/152	3	97	-4
3393/152	3148/141	3	97	-4
3148/141	649/29	3	3245	-145
1500/67	1703/76	5	-3000	134
1703/76	3161/141	5	-3000	134
3161/141	3097/138	5	-3000	134
3097/138	3258/145	5	97	-4
3258/145	787/35	2	97	-4
787/35	3245/144	2	3245	-144
3161/140	3097/137	3	-3161	140
3097/137	1131/50	3	-3161	140
1131/50	181/8	3	-3161	140
181/8	3148/139	3	97	-4
3148/139	295/13	3	3245	-143
250/11	3161/139	5	-3000	132
3161/139	3097/136	5	-3000	132
3097/136	3258/143	5	97	-4
3258/143	1574/69	2	97	-4
1574/69	3245/142	2	3245	-142
3161/138	3393/148	3	-3161	138
3393/148	3097/135	3	-3161	138
3097/135	1629/71	3	-3161	138
1629/71	3148/137	3	97	-4
3148/137	23	3	3245	-141
23	1703/74	3	3406	-148
300/13	1086/47	2	-3000	130
1086/47	3097/134	5	-3000	130
3097/134	787/34	5	97	-4
787/34	3406/147	5	3245	-140
3406/147	649/28	5	3245	-140
1000/43	1629/70	3	-3000	129
1629/70	163/7	3	258	-11
163/7	3148/135	3	258	-11
3148/135	1703/73	3	3406	-146
375/16	3258/139	2	-3000	128
3258/139	258/11	5	-3000	128
258/11	3097/132	5	-3000	128
3097/132	3406/145	5	97	-4
3406/145	1574/67	5	97	-4
1574/67	3245/138	5	3245	-138
500/21	262/11	5	-3000	126
262/11	3097/130	5	-3000	126
3097/130	787/33	5	97	-4
787/33	3245/136	5	3245	-136
750/31	3097/128	5	-3000	124
3097/128	1574/65	5	97	-4
1574/65	3245/134	5	3245	-134
97/4	1629/67	3	-97	4
1629/67	3161/130	3	3161	-130
3245/133	3148/129	2	-3245	133
3148/129	3258/133	2	-97	4
3258/133	49/2	5	-97	4
49/2	3406/139	5	148	-6
3406/139	3161/129	5	148	-6
3161/129	3097/126	5	148	-6
3097/126	295/12	5	3245	-132
1500/61	787/32	3	-3000	122
787/32	74/3	3	148	-6
649/26	3097/124	5	-3245	130
3097/124	1574/63	5	-148	6
1574/63	25	5	3000	-120
3245/128	3097/122	5	-3245	128
3097/122	787/31	5	-148	6
787/31	1703/67	5	3000	-118
1703/67	1500/59	5	3000	-118
3245/126	129/5	5	-3245	126
129/5	1703/66	5	-3245	126
1703/66	1574/61	5	-3245	126
1574/61	3097/120	5	-97	4
3097/120	181/7	5	3000	-116
181/7	750/29	2	3000	-116
26	3148/121	3	-3406	131
3148/121	3097/119	3	-258	10
3097/119	3258/125	3	-258	10
3258/125	600/23	3	3000	-115
3245/124	131/5	5	-3245	124
131/5	787/30	5	-3245	124
787/30	3097/118	5	-97	4
3097/118	1629/62	5	3000	-114
1629/62	500/19	2	3000	-114
3406/129	3148/119	3	-3406	129
3148/119	3097/117	3	-258	10
3097/117	1086/41	3	-258	10
1086/41	3393/128	3	3000	-113
3393/128	3000/113	3	3000	-113
3245/122	1703/64	5	-3245	122
1703/64	1574/59	5	-3245	122
1574/59	3097/116	5	-97	4
3097/116	1629/61	5	3000	-112
1629/61	375/14	2	3000	-112
3406/127	161/6	3	-3406	127
161/6	3148/117	3	-3245	121
3148/117	3258/121	3	-97	4
3258/121	377/14	3	3161	-117
377/14	3097/115	3	3161	-117
3097/115	3161/117	3	3161	-117
649/24	787/29	2	-3245	120
787/29	543/20	2	-97	4
543/20	163/6	5	-97	4
163/6	245/9	5	3000	-110
245/9	3406/125	5	3245	-119
3406/125	109/4	5	3245	-119
109/4	3245/119	5	3245	-119
300/11	3393/124	3	-3000	110
3393/124	3148/115	3	-3000	110
3148/115	3258/119	3	148	-5
3258/119	3097/113	3	3406	-124
3097/113	1703/62	3	3406	-124
3000/109	1629/59	2	-3000	109
1629/59	1574/57	5	-3000	109
1574/57	3097/112	5	148	-5
3097/112	3406/123	5	3245	-117
3406/123	3161/114	5	3245	-117
3161/114	3245/117	5	3245	-117
250/9	3393/122	3	-3000	108
3393/122	362/13	3	-3000	108
362/13	3148/113	3	258	-9
3148/113	3097/111	3	3406	-122
3097/111	1703/61	3	3406	-122
3000/107	1629/58	2	-3000	107
1629/58	787/28	5	-3000	107
787/28	3406/121	5	148	-5
3406/121	3097/110	5	148	-5
3097/110	649/23	5	3245	-115
1500/53	3258/115	3	-3000	106
3258/115	3148/111	3	258	-9
3148/111	1703/60	3	3406	-120
200/7	543/19	2	-3000	105
543/19	1574/55	5	-3000	105
1574/55	3406/119	5	148	-5
3406/119	86/3	5	148	-5
86/3	3097/108	5	148	-5
3097/108	3245/113	5	3245	-113
3000/103	787/27	5	-3000	103
787/27	3097/106	5	148	-5
3097/106	3245/111	5	3245	-111
3406/115	1629/55	3	-3406	115
1629/55	1574/53	3	-148	5
1574/53	3000/101	3	3000	-101
3245/109	3097/104	5	-3245	109
3097/104	3161/106	5	-148	5
3161/106	1703/57	5	-148	5
1703/57	3258/109	5	-148	5
3258/109	3148/105	2	-148	5
3148/105	30	2	3000	-100
3406/113	181/6	3	-3406	113
181/6	787/26	3	-148	5
787/26	3393/112	3	3000	-99
3393/112	1000/33	3	3000	-99
3245/107	3097/102	5	-3245	107
3097/102	3161/104	5	-148	5
3161/104	1703/56	5	-148	5
1703/56	3258/107	5	-148	5
3258/107	3148/103	2	-148	5
3148/103	1500/49	2	3000	-98
3161/103	1629/53	3	-3161	103
1629/53	3393/110	3	97	-3
3393/110	1574/51	3	97	-3
1574/51	649/21	3	3245	-105
3000/97	1703/55	5	-3000	97
1703/55	3097/100	5	-3000	97
3097/100	3161/102	5	97	-3
3161/102	1086/35	5	97	-3
1086/35	3148/101	2	97	-3
3148/101	3245/104	2	3245	-104
3161/101	1629/52	3	-3161	101
1629/52	377/12	3	97	-3
377/12	787/25	3	97	-3
787/25	3245/103	3	3245	-103
600/19	3097/98	5	-3000	95
3097/98	3161/100	5	97	-3
3161/100	3258/103	5	97	-3
3258/103	3148/99	2	97	-3
3148/99	3245/102	2	3245	-102
3161/99	543/17	3	-3161	99
543/17	3393/106	3	97	-3
3393/106	1574/49	3	97	-3
1574/49	3245/101	3	3245	-101
1000/31	3097/96	5	-3000	93
3097/96	97/3	5	97	-3
649/20	3148/97	3	-3245	100
3148/97	1629/50	3	-97	3
1629/50	3161/97	3	3161	-97
295/9	787/24	2	-3245	99
787/24	362/11	2	-97	3
362/11	3161/96	5	-97	3
3161/96	3097/94	5	-97	3
3097/94	3000/91	5	3000	-91
3245/98	3148/95	3	-3245	98
3148/95	1629/49	3	-97	3
1629/49	1131/34	3	3161	-95
1131/34	3161/95	3	3161	-95
3245/97	1574/47	2	-3245	97
1574/47	3258/97	2	-97	3
3258/97	3161/94	5	-97	3
3161/94	3097/92	5	-97	3
3097/92	3000/89	5	3000	-89
3245/96	3148/93	3	-3245	96
3148/93	3393/100	3	-97	3
3393/100	543/16	3	-97	3
543/16	3161/93	3	3161	-93
649/19	787/23	2	-3245	95
787/23	3258/95	2	-97	3
3258/95	3161/92	5	-97	3
3161/92	3406/99	5	-97	3
3406/99	3097/90	5	-97	3
3097/90	1000/29	5	3000	-87
3245/94	3148/91	3	-3245	94
3148/91	3393/98	3	-97	3
3393/98	1629/47	3	-97	3
1629/47	3161/91	3	3161	-91
3245/93	1574/45	2	-3245	93
1574/45	35	2	-97	3
35	1086/31	2	148	-4
1086/31	3406/97	5	148	-4
3406/97	3161/90	5	148	-4
3161/90	3097/88	5	148	-4
3097/88	3245/92	5	3245	-92
600/17	1131/32	3	-3000	85
1131/32	3148/89	3	-3000	85
3148/89	1629/46	3	148	-4
1629/46	1703/48	3	3406	-96
250/7	787/22	2	-3000	84
787/22	3258/91	2	148	-4
3258/91	3406/95	5	148	-4
3406/95	3161/88	5	148	-4
3161/88	3097/86	5	148	-4
3097/86	649/18	5	3245	-90
3000/83	3148/87	3	-3000	83
3148/87	181/5	3	148	-4
181/5	1703/47	3	3406	-94
1500/41	1574/43	2	-3000	82
1574/43	3258/89	2	148	-4
3258/89	3406/93	5	148	-4
3406/93	3161/86	5	148	-4
3161/86	258/7	5	148	-4
258/7	3097/84	5	148	-4
3097/84	295/8	5	3245	-88
37	1703/46	5	-148	4
1703/46	1629/44	5	-148	4
1629/44	3148/85	2	-148	4
3148/85	1000/27	2	3000	-81
262/7	1086/29	3	-3406	91
1086/29	787/21	3	-148	4
787/21	75/2	3	3000	-80
3245/86	3097/82	5	-3245	86
3097/82	1703/45	5	-148	4
1703/45	1629/43	5	-148	4
1629/43	3148/83	2	-148	4
3148/83	3000/79	2	3000	-79
3406/89	3258/85	3	-3406	89
3258/85	1574/41	3	-148	4
1574/41	500/13	3	3000	-78
3245/84	1703/44	5	-3245	84
1703/44	3097/80	5	-3245	84
3097/80	543/14	5	-148	4
543/14	3148/81	2	-148	4
3148/81	3000/77	2	3000	-77
3406/87	3097/79	3	-3406	87
3097/79	3258/83	3	-3406	87
3258/83	787/20	3	-148	4
787/20	3393/86	3	3000	-76
3393/86	750/19	3	3000	-76
3245/82	1703/43	5	-3245	82
1703/43	3097/78	5	-3245	82
3097/78	1629/41	5	-148	4
1629/41	3148/79	2	-148	4
3148/79	40	2	3000	-75
3406/85	3097/77	3	-3406	85
3097/77	362/9	3	-3406	85
362/9	161/4	3	-148	4
161/4	1574/39	3	13	0
1574/39	1131/28	3	3161	-78
1131/28	3161/78	3	3161	-78
649/16	1629/40	2	-3245	80
1629/40	163/4	5	-3245	80
163/4	245/6	5	-148	4
245/6	3148/77	5	97	-2
3148/77	3406/83	5	3245	-79
3406/83	3161/77	5	3245	-79
3161/77	3245/79	5	3245	-79
3000/73	3258/79	3	-3000	73
3258/79	3097/75	3	258	-6
3097/75	3393/82	3	258	-6
3393/82	787/19	3	258	-6
787/19	1703/41	3	3406	-82
125/3	543/13	2	-3000	72
543/13	3097/74	5	-3000	72
3097/74	3148/75	5	97	-2
3148/75	3406/81	5	3245	-77
3406/81	295/7	5	3245	-77
3000/71	3258/77	3	-3000	71
3258/77	3393/80	3	258	-6
3393/80	3097/73	3	258	-6
3097/73	1574/37	3	258	-6
1574/37	1703/40	3	3406	-80
300/7	1629/38	2	-3000	70
1629/38	43	5	-3000	70
43	3097/72	5	-3000	70
3097/72	3406/79	5	97	-2
3406/79	3148/73	5	97	-2
3148/73	649/15	5	3245	-75
750/17	3406/77	5	-3000	68
3406/77	3097/70	5	-3000	68
3097/70	3148/71	5	97	-2
3148/71	3245/73	5	3245	-73
500/11	3097/68	5	-3000	66
3097/68	3148/69	5	97	-2
3148/69	3245/71	5	3245	-71
375/8	3097/66	5	-3000	64
3097/66	3148/67	5	97	-2
3148/67	3245/69	5	3245	-69
1500/31	3097/64	5	-3000	62
3097/64	3148/65	5	97	-2
3148/65	3245/67	5	3245	-67
97/2	3258/67	3	-97	2
3258/67	3161/65	3	3161	-65
3000/61	787/16	2	-3000	61
787/16	148/3	2	148	-3
649/13	3097/62	5	-3245	65
3097/62	3148/63	5	-148	3
3148/63	50	5	3000	-60
3245/63	258/5	5	-3245	63
258/5	1703/33	5	-3245	63
1703/33	3148/61	5	-3245	63
3148/61	3097/60	5	-97	2
3097/60	362/7	5	3000	-58
362/7	1500/29	2	3000	-58
262/5	787/15	3	-3406	65
787/15	3097/59	3	-258	5
3097/59	1629/31	3	-258	5
1629/31	1000/19	3	3000	-57
3245/61	1703/32	5	-3245	61
1703/32	3148/59	5	-3245	61
3148/59	3097/58	5	-97	2
3097/58	3258/61	5	3000	-56
3258/61	375/7	2	3000	-56
649/12	1574/29	3	-3245	60
1574/29	543/10	3	-97	2
543/10	163/3	3	3161	-58
163/3	109/2	3	3161	-58
55	3258/59	2	-3245	59
3258/59	3148/57	5	-3245	59
3148/57	3097/56	5	-97	2
3097/56	3161/57	5	3000	-54
3161/57	500/9	5	3000	-54
3245/58	1629/29	3	-3245	58
1629/29	787/14	3	13	0
787/14	3097/55	3	3161	-56
3097/55	3161/56	3	3161	-56
3245/57	1086/19	2	-3245	57
1086/19	3148/55	5	-3245	57
3148/55	3097/54	5	-97	2
3097/54	3161/55	5	3000	-52
3161/55	750/13	5	3000	-52
3245/56	1629/28	3	-3245	56
1629/28	1574/27	3	13	0
1574/27	3097/53	3	3161	-54
3097/53	117/2	3	3161	-54
117/2	3161/54	3	3161	-54
59	3258/55	2	-3245	55
3258/55	3148/53	5	-3245	55
3148/53	3097/52	5	-97	2
3097/52	3161/53	5	3000	-50
3161/53	3406/57	5	3000	-50
3406/57	60	5	3000	-50
3245/54	181/3	3	-3245	54
181/3	787/13	3	13	0
787/13	3393/56	3	3161	-52
3393/56	3097/51	3	3161	-52
3097/51	3161/52	3	3161	-52
3245/53	245/4	2	-3245	53
245/4	3258/53	2	-3000	49
3258/53	3148/51	5	-3000	49
3148/51	3406/55	5	148	-2
3406/55	3097/50	5	148	-2
3097/50	3161/51	5	3245	-52
3161/51	3245/52	5	3245	-52
125/2	1629/26	3	-3000	48
1629/26	377/6	3	258	-4
377/6	1574/25	3	258	-4
1574/25	1703/27	3	3406	-54
3000/47	1086/17	2	-3000	47
1086/17	3148/49	5	-3000	47
3148/49	3406/53	5	148	-2
3406/53	129/2	5	148	-2
129/2	3161/49	5	148	-2
3161/49	3097/48	5	148	-2
3097/48	649/10	5	3245	-50
200/3	3406/51	5	-3000	45
3406/51	3148/47	5	-3000	45
3148/47	3161/47	5	148	-2
3161/47	3097/46	5	148	-2
3097/46	3245/48	5	3245	-48
3000/43	3148/45	5	-3000	43
3148/45	3161/45	5	148	-2
3161/45	3097/44	5	148	-2
3097/44	3245/46	5	3245	-46
3000/41	3148/43	5	-3000	41
3148/43	3161/43	5	148	-2
3161/43	3097/42	5	148	-2
3097/42	295/4	5	3245	-44
74	1703/23	5	-148	2
1703/23	1629/22	5	-148	2
1629/22	1574/21	2	-148	2
1574/21	75	2	3000	-40
3406/45	3258/43	3	-3406	45
3258/43	3148/41	3	-148	2
3148/41	1000/13	3	3000	-39
3245/42	1703/22	5	-3245	42
1703/22	3097/40	5	-3245	42
3097/40	543/7	5	-148	2
543/7	787/10	2	-148	2
787/10	1500/19	2	3000	-38
3406/43	3097/39	3	-3406	43
3097/39	3258/41	3	-3406	43
3258/41	161/2	3	-148	2
161/2	3148/39	3	13	0
3148/39	1131/14	3	3161	-39
1131/14	3161/39	3	3161	-39
649/8	1629/20	2	-3245	40
1629/20	163/2	5	-3245	40
163/2	245/3	5	-148	2
245/3	1574/19	5	97	-1
1574/19	3406/41	5	3245	-39
3406/41	3161/38	5	3245	-39
3161/38	3245/39	5	3245	-39
250/3	1086/13	3	-3000	36
1086/13	3097/37	3	258	-3
3097/37	3393/40	3	258	-3
3393/40	3148/37	3	258	-3
3148/37	1703/20	3	3406	-40
600/7	1629/19	2	-3000	35
1629/19	86	5	-3000	35
86	3097/36	5	-3000	35
3097/36	262/3	5	97	-1
262/3	787/9	5	97	-1
787/9	3245/37	5	3245	-37
1000/11	3097/34	5	-3000	33
3097/34	3406/37	5	97	-1
3406/37	1574/17	5	97	-1
1574/17	649/7	5	3245	-35
3000/31	3097/32	5	-3000	31
3097/32	97	5	97	-1
295/3	787/8	3	-3245	33
787/8	1086/11	3	-97	1
1086/11	3161/32	3	3161	-32
3245/32	3148/31	2	-3245	32
3148/31	1629/16	2	-97	1
1629/16	3161/31	5	-97	1
3161/31	3406/33	5	-97	1
3406/33	3097/30	5	-97	1
3097/30	3000/29	5	3000	-29
3245/31	1574/15	3	-3245	31
1574/15	3258/31	3	-97	1
3258/31	3161/30	3	3161	-30
649/6	3148/29	2	-3245	30
3148/29	543/5	2	-97	1
543/5	109	5	-97	1
109	3406/31	5	-97	1
3406/31	3097/28	5	-97	1
3097/28	1000/9	5	3000	-27
3245/29	3258/29	3	-3245	29
3258/29	787/7	3	13	0
787/7	3161/28	3	3161	-28
3245/28	1629/14	2	-3245	28
1629/14	3148/27	5	-3245	28
3148/27	3161/27	5	-97	1
3161/27	3406/29	5	-97	1
3406/29	3097/26	5	-97	1
3097/26	120	5	3000	-25
3245/27	362/3	3	-3245	27
362/3	1574/13	3	13	0
1574/13	3393/28	3	3161	-26
3393/28	3161/26	3	3161	-26
125	1629/13	2	-3000	24
1629/13	3148/25	5	-3000	24
3148/25	3406/27	5	148	-1
3406/27	3161/25	5	148	-1
3161/25	129	5	148	-1
129	3097/24	5	148	-1
3097/24	649/5	5	3245	-25
1500/11	3148/23	5	-3000	22
3148/23	3161/23	5	148	-1
3161/23	3097/22	5	148	-1
3097/22	3245/23	5	3245	-23
3406/23	1629/11	3	-3406	23
1629/11	3148/21	3	-148	1
3148/21	150	3	3000	-20
3245/21	1703/11	5	-3245	21
1703/11	3097/20	5	-3245	21
3097/20	1086/7	5	-148	1
1086/7	787/5	2	-148	1
787/5	3000/19	2	3000	-19
649/4	1629/10	3	-3245	20
1629/10	163	3	13	0
163	3148/19	3	13	0
3148/19	3161/19	3	3161	-19
3245/19	3258/19	2	-3245	19
3258/19	3097/18	5	-3245	19
3097/18	1574/9	5	-148	1
1574/9	3161/18	5	3000	-17
3161/18	3000/17	5	3000	-17
3245/18	181	3	-3245	18
181	3097/17	3	13	0
3097/17	3148/17	3	13	0
3148/17	3161/17	3	3161	-17
3245/17	3258/17	2	-3245	17
3258/17	3097/16	5	-3245	17
3097/16	787/4	5	-148	1
787/4	3161/16	5	3000	-15
3161/16	200	5	3000	-15
3245/16	1629/8	3	-3245	16
1629/8	3097/15	3	13	0
3097/15	3148/15	3	13	0
3148/15	3161/15	3	3161	-15
649/3	1086/5	2	-3245	15
1086/5	3097/14	5	-3245	15
3097/14	1574/7	5	-148	1
1574/7	3161/14	5	3000	-13
3161/14	3406/15	5	3000	-13
3406/15	3000/13	5	3000	-13
3245/14	1629/7	3	-3245	14
1629/7	3097/13	3	13	0
3097/13	3148/13	3	13	0
3148/13	3393/14	3	3161	-13
3393/14	3161/13	3	3161	-13
250	3258/13	2	-3000	12
3258/13	258	5	-3000	12
258	3097/12	5	-3000	12
3097/12	262	5	97	0
262	787/3	5	97	0
787/3	3161/12	5	3245	-12
3161/12	3245/12	5	3245	-12
300	3406/11	5	-3000	10
3406/11	3097/10	5	-3000	10
3097/10	1574/5	5	97	0
1574/5	3161/10	5	3245	-10
3161/10	649/2	5	3245	-10
375	3406/9	5	-3000	8
3406/9	3097/8	5	-3000	8
3097/8	787/2	5	97	0
787/2	3161/8	5	3245	-8
3161/8	3245/8	5	3245	-8
500	3097/6	5	-3000	6
3097/6	1574/3	5	97	0
1574/3	3161/6	5	3245	-6
3161/6	3245/6	5	3245	-6
750	3097/4	5	-3000	4
3097/4	787	5	97	0
787	3161/4	5	3245	-4
3161/4	3245/4	5	3245	-4
1500	3097/2	5	-3000	2
3097/2	1574	5	97	0
1574	3161/2	5	3245	-2
3161/2	3245/2	5	3245	-2
1703	3503/2	0	-3406	2
3503/2	3000	0	-3406	2
3000	3097	0	-3406	2
3097	3148	0	-3406	2
3148	3161	0	-258	1
3161	3245	0	-258	1
3245	3258	0	-258	1
\end{Verbatim}
\fi
\end{document}
